\documentclass[11pt]{amsbook}

\usepackage[utf8]{inputenc}
\usepackage[OT2,T1]{fontenc}
\usepackage{lmodern}

\usepackage{amssymb,amsmath,amsthm,mathrsfs,mathtools}
\usepackage{dsfont}
\usepackage{enumitem}
\usepackage{multicol}
\usepackage{booktabs,longtable,array}
\usepackage{subcaption}
\usepackage{xcolor}
\usepackage{tikz}
\usetikzlibrary{arrows.meta,positioning,calc,decorations.pathreplacing,
  decorations.markings,shapes.geometric,shapes.misc,patterns,fit,
  backgrounds,matrix}
\usepackage{tikz-cd}
\usepackage{pgfplots}
\pgfplotsset{compat=1.18}
\usepgfplotslibrary{fillbetween,groupplots}
\usepackage{adjustbox}
\usepackage{float}
\usepackage{placeins}
\usepackage[hidelinks,hypertexnames=false]{hyperref}

\DeclareRobustCommand{\taylor}[1]{\textcolor{blue}{[Taylor: #1]}}
\newcommand{\todo}[1]{\textcolor{red}{[todo: #1]}}
\newcommand{\agentreply}[1]{\textcolor{gray}{[agent: #1]}}

\theoremstyle{plain}
\newtheorem{theorem}{Theorem}[section]
\newtheorem{proposition}[theorem]{Proposition}
\newtheorem{conjecture}[theorem]{Conjecture}
\theoremstyle{definition}
\newtheorem{definition}[theorem]{Definition}
\newtheorem{question}[theorem]{Question}
\newtheorem{example}[theorem]{Example}
\theoremstyle{remark}
\newtheorem{remark}[theorem]{Remark}
\renewcommand{\thesection}{\thechapter.\arabic{section}}
\numberwithin{equation}{section}
\numberwithin{figure}{chapter}
\numberwithin{table}{chapter}

\newcommand{\Q}{\mathbb{Q}}
\newcommand{\Z}{\mathbb{Z}}
\newcommand{\R}{\mathbb{R}}
\renewcommand{\C}{\mathbb{C}}
\newcommand{\F}{\mathbb{F}}
\newcommand{\eps}{\varepsilon}
\newcommand{\Ch}{\mathrm{Ch}}
\newcommand{\lev}{\mathscr{L}}
\newcommand{\Gz}{G_{\zeta}}
\newcommand{\Sz}{S_{\zeta}}
\newcommand{\Kept}{\operatorname{Kept}}
\newcommand{\wt}{\operatorname{wt}}
\newcommand{\Gal}{\operatorname{Gal}}
\newcommand{\Frob}{\operatorname{Frob}}
\newcommand{\rk}{\operatorname{rk}}
\newcommand{\GG}{\mathfrak{G}}
\newcommand{\pt}{\operatorname{pt}}
\DeclareMathOperator{\Aut}{Aut}
\DeclareMathOperator{\Spec}{Spec}
\DeclareMathOperator{\Real}{Re}
\newcommand{\GPsi}{G^{\mathrm{ar}}}
\newcommand{\PS}{\mathrm{PS}}
\newcommand{\MT}{\mathrm{MT}}
\newcommand{\DF}{\mathrm{DF}}
\newcommand{\Top}{\mathcal{T}}
\newcommand{\gap}{\operatorname{gap}}
\newcommand{\Err}{\mathrm{Err}}
\newcommand{\Main}{\operatorname{Main}}
\newcommand{\Words}{\mathcal{U}}
\newcommand{\Cg}{\mathrm{Cg}}
\newcommand{\lgv}{\operatorname{lg}}
\newcommand{\Fo}{\mathrm{Fo}}
\newcommand{\Ga}{\mathrm{Ga}}
\newcommand{\Ta}{\mathrm{Ta}}
\newcommand{\IO}{\mathrm{IO}}
\newcommand{\Lo}{\mathrm{Lo}}
\newcommand{\Fam}{\mathfrak{F}}
\newcommand{\Win}{\mathrm{Win}}
\DeclareSymbolFont{cyrletters}{OT2}{wncyr}{m}{n}
\DeclareMathSymbol{\Sha}{\mathalpha}{cyrletters}{"58}
\newcommand{\shalen}{\operatorname{len}}
\newcommand{\Lsp}{\mathcal{L}}
\newcommand{\ord}{\operatorname{ord}}
\newcommand{\degh}{\widehat{\deg}}
\newcommand{\covol}{\operatorname{covol}}
\newcommand{\Fr}{\mathrm{Fr}}
\newcommand{\ZZ}{\mathbb{Z}}
\newcommand{\calL}{\mathcal{L}}
\newcommand{\calO}{\mathcal{O}}
\newcommand{\calE}{\mathcal{E}}
\newcommand{\Fb}{\overline{\F}_q}
\newcommand{\Tr}{\operatorname{Tr}}
\newcommand{\cF}{\mathcal{F}}
\newcommand{\repopath}[1]{\nolinkurl{#1}}

\newcommand{\vtwoskillurl}{https://github.com/tdupu/%
agent-skills-public/blob/8d13efb5e66306aeae6e3b948bc463e7e2db1b65}
\newcommand{\vtwomathlib}{https://github.com/leanprover-community/%
mathlib4/commit/0df444a360eaa60ab8c11dca51a86af692955474}
\newcommand{\vtwoskill}[1]{%
\href{\vtwoskillurl/skills/#1/SKILL.md}{\texttt{#1}}}

\newcommand{\bidata}[1]{alpoge/data/#1}
\newcommand{\sbdata}[1]{stepanov-bombieri/data/#1}
\newcommand{\aqthreedata}{alpoge/data/aq3/}
\pgfplotsset{table/search path={.,alpoge}}
\input{alpoge/figs/alpoge-style.tex}
\input{alpoge/figs/v2/v2-style.tex}
\definecolor{sbPole}{HTML}{D55E00}  % vermillion: poles, pole divisors, P_0,
\definecolor{sbZero}{HTML}{0072B2}  % blue: zeros, Frobenius places, the
\definecolor{sbWeil}{HTML}{009E73}  % bluish green: what the Riemann
\definecolor{sbFail}{HTML}{CC79A7}  % reddish purple: the formal
\definecolor{sbBound}{HTML}{E69F00} % orange: what the Stepanov--Bombieri
\definecolor{sbAux}{HTML}{56B4E9}   % sky blue: auxiliary data -- the Siegel
\definecolor{sbOther}{HTML}{888888} % grey: neutral, other, slack
\definecolor{sbInk}{HTML}{000000}   % black: text, axes, structure

\colorlet{sbRes0}{black}
\colorlet{sbRes1}{black!75}
\colorlet{sbRes2}{black!55}
\colorlet{sbRes3}{black!40}
\colorlet{sbRes4}{black!25}

\tikzset{
  sb/.style={font=\small, >={Stealth[length=2.2mm]}},
  sb pole/.style={sbPole, line width=0.9pt, cross out, draw,
    minimum size=4.5pt, inner sep=0pt},             % x-mark
  sb zero/.style={circle, fill=sbZero, draw=sbZero, minimum size=4pt,
    inner sep=0pt},                                 % filled dot
  sb weil pt/.style={diamond, fill=sbWeil, draw=sbWeil, minimum size=5pt,
    inner sep=0pt},                                 % filled diamond
  sb fail pt/.style={circle, draw=sbFail, fill=white, line width=0.9pt,
    minimum size=5pt, inner sep=0pt},               % open circle
  sb aux pt/.style={rectangle, fill=sbAux, draw=sbAux, minimum size=4pt,
    inner sep=0pt},                                 % filled square
  sb weil line/.style={sbWeil, line width=0.9pt},             % solid
  sb fail line/.style={sbFail, line width=0.9pt, dashed},     % dashed
  sb bound line/.style={sbBound, line width=0.9pt, dash dot}, % dash-dot
  sb pole line/.style={sbPole, line width=0.9pt},             % solid
  sb zero line/.style={sbZero, line width=0.9pt},             % solid
  sb aux line/.style={sbAux, line width=0.9pt, densely dotted},
  sb other line/.style={sbOther, line width=0.6pt},
  sb weil fill/.style={fill=sbWeil!15},
  sb fail fill/.style={pattern={north east lines}, pattern color=sbFail},
  sb bound fill/.style={fill=sbBound!20},
  sb aux fill/.style={fill=sbAux!20},
  sb box/.style={draw=sbInk, rounded corners=2pt, align=center,
    inner sep=4pt, font=\small},
  sb aux box/.style={sb box, fill=sbAux!15},
  sb bound box/.style={sb box, fill=sbBound!15},
  sb weil box/.style={sb box, fill=sbWeil!15},
  sb fail box/.style={sb box, fill=sbFail!12, dashed},
  sb exact/.style={sbInk, line width=1.1pt},
  sb loose/.style={sbInk!60, line width=0.9pt, dashed},
  sb none/.style={sbFail, line width=0.9pt, densely dotted},
  sb schematic/.style={font=\scriptsize\itshape, text=sbOther},
}

\pgfplotsset{
  sb axis/.style={
    width=0.82\linewidth, height=0.5\linewidth,
    tick label style={font=\footnotesize},
    label style={font=\small},
    legend style={font=\footnotesize, draw=none, fill=white,
      fill opacity=0.85, text opacity=1, cells={anchor=west}},
    axis line style={sbInk},
    every axis plot/.append style={line join=round},
  },
  sb half axis/.style={sb axis, width=0.47\linewidth, height=0.42\linewidth},
  sb pole plot/.style={sbPole, mark=x, mark size=2.5pt, line width=0.9pt},
  sb zero plot/.style={sbZero, mark=*, mark size=1.6pt, line width=0.9pt},
  sb weil plot/.style={sbWeil, mark=diamond*, mark size=2pt,
    line width=0.9pt},
  sb fail plot/.style={sbFail, mark=o, mark size=2pt, line width=0.9pt,
    dashed},
  sb bound plot/.style={sbBound, mark=triangle*, mark size=2pt,
    line width=0.9pt, dash dot},
  sb aux plot/.style={sbAux, mark=square*, mark size=1.6pt,
    line width=0.9pt, densely dotted},
  sb other plot/.style={sbOther, line width=0.6pt},
}

\newcommand{\sbschematic}[1][schematic, not to scale]{%
  {\scriptsize\itshape\color{sbOther}#1}}
\newcommand{\sbfig}[2][tbp]{\par\noindent
  See Figure~\ref{fig:#2}.\begin{figure}[#1]\centering
  \input{stepanov-bombieri/figs/#2}\end{figure}\par}
\newcommand{\sbfloat}[2][tbp]{\begin{figure}[#1]\centering
  \input{stepanov-bombieri/figs/#2}\end{figure}}
\newcommand{\sbsee}[1]{\ See Figure~\ref{fig:#1}.}

\emergencystretch=3em
\setcounter{tocdepth}{1}

\title{Notes on Alp\"oge's Proof of The Riemann Hypothesis}
\author{Printed by Taylor Dupuy \\
AI GENERATED \\[2ex]
21 November 2026}
\date{21 November 2026}

\begin{document}

\frontmatter
\maketitle
\tableofcontents
\clearpage
\chapter*{Preface}


These are lecture notes on the argument for the Riemann Hypothesis
announced by L.~Alp\"oge in the note reproduced after this preface
\cite{AlpogeRemark2026}, together with a Lean archive
\cite{AlpogeArchive2026}. They are written for professional
mathematicians and graduate students. The aim is that the reader first
sees how the argument fits together and then meets the tools one
chapter at a time.

Chapter~\ref{ch:idea} explains the idea of the argument.
Chapter~\ref{ch:function-fields} works out the model case over a finite
field: the Stepanov--Bombieri proof of the Riemann hypothesis for
curves, and the conditional sign-only argument through the identity
$qN = M^k + R$.
Chapter~\ref{ch:assembly} states the main theorem and gives its whole
proof, with each tool used as a black box and a reference to the
chapter where it is treated. Chapters~\ref{ch:landau}
to~\ref{ch:families} develop the tools.
Appendix~\ref{ch:ai} describes the software and AI methods used to
produce the text.

\subsection*{AI notice}
AI agents wrote this text, and the AI and Software Discussion at
the end of the Introduction (page~\pageref{intro:ai}) and
Appendix~\ref{ch:ai} describe how.
\clearpage
\clearpage
\chapter*{A remark by Levent Alp\"oge}
\begingroup
\newcommand{\remarkdate}[1]{\textbf{#1}}
\newcommand{\remarkpart}[1]{\par\bigskip\noindent
  {\large\bfseries #1}\par\nopagebreak\medskip\noindent\ignorespaces}
\newcommand{\remarkdq}{\textquotedbl}
\newcommand{\remarkone}{\mathds{1}}
\setcounter{footnote}{0}
\renewcommand{\labelenumi}{\arabic{enumi}.}

\begin{quotation}
\noindent
The Riemann Hypothesis is the classic problem in number theory. It asks for
optimal control on the count of prime numbers of a given size. Its shadows
have guided mathematicians for $150$ years, and its many rhymes and
reflections, even reaching physics, taught us the unity of our fields.

I believe I have, working with Claude, discovered a
proof\footnote{\url{https://alpo.ge/lean.tar.gz.part01},\dots,%
\url{https://alpo.ge/lean.tar.gz.part14}.} of the Riemann
Hypothesis. In fact, the method easily extends much further, resolving the
so-called Generalized Riemann Hypothesis, and I expect, with more time, the
Extended Riemann Hypothesis too. Claude and I have also formalized the
arguments without assumptions. It is natural to feel that this is an end, in
my case a loss of a principal guiding light, but what I saw during the course
of progress so amazed me that I believe it is one dawn across science of
something far greater.

It was my secret in my student days that my plans were for this problem in
decades, which meant then that I tried to become familiar with all tools that
I could possibly recognize in the problem when it was time. Claude and I have
had a great time together with such breadth in recent months, but this time
Claude far outpaced me near the finish line. No, Claude flew, and danced,
and sparkled near that finish line. I had us on the line of combining two
quite different ideas present in a toy setting, one of which was already new
and shocking, making very simple one of the shadows of Riemann over which we
mathematicians felt most triumph. I guided us to the right situation
analogous to that toy setting, and then Claude began to leave me in awe.

Tens of such occurred, but I will give one line which will suffice to make
fellow mathematicians fall out of their chairs. In Stepanov's argument
resolving Riemann over function fields one uses a divisor involving
multiplicity. This argument is a descendant of Stepanov's, but multiplicity
is hard to find over the integers, corresponding to something like
differentiating whole numbers. The simple, new, Fourier-analytic ideas in the
toy setting led us to examine chaining theta lifts if we were to have the
same tricks recur but work in a natively automorphic manner. The line: for
the von Mangoldt sum of length $X$, Claude's suggested length of chain of
theta lifts, in view of each mimicking multiplicity, was $\sqrt{X}$. (And
yet, what else could it be?)

An astounding amount of mathematics is present in the argument, far more
than I had imagined would be necessary as a student. My intuition insists
it will be massively simplified, likely immediately. But I hope in having the
experience of seeing all of number theory's beauty brought together I am just
one of the first, and those who love the mathematics will all be able to
enjoy such transcendence in the future.

The argument is incredibly long at the moment, and while its formalization
is finished and it is routine to produce a huge, unreadable document
containing sufficient detail, in my view of my path I have just taken a step.
I intend to do my best to understand and impart understanding of this
argument and its certain simplifications, to speak about it, to write it
properly.

In a time of revolution in mathematics I hope the argument provides a
glimpse towards the miracles that lie ahead, and shows us the joy and wonder
we are about to find in the simplest things.

\remarkpart{1}
It is obvious to me that I should remark on the following.

The precise timing is horrible. If I could have, I'd have finished this all
much earlier. Hell, strictly before this age would have sure been nice. This
all happened for me late Sunday night, and I messaged Dario,
\remarkdq Hypothetically, if I were to have \dots\remarkdq, eventually
asking about my ability to announce. We went and figured it out and the fact
is that there was a constraint to respect, though I did ask about
progressively more bizarre ideas. This problem has been \emph{the} sacred
object in mathematics for me since my childhood, and I am very touchy about
the treatment of great intellectual goals, though seriously imperfect. Our
leadership and the few trusted mathematicians I consulted hold extremely
similar views on the right things to do in these circumstances (it is
occasionally possible for physicists and mathematicians to agree\dots), and
I am clearly unilaterally responsible for feeling I just could not live if I
waited even longer. There is an obvious cynical view, to do with the timing,
but it is false, and I am pleased that I am in circumstances where it is
obviously false.

This puts us in the rather extraordinary situation this age has made
possible, in which I am presenting the world with a complete argument before
having mastered a massive neighbourhood around its mathematics. It being
obviously inappropriate to attempt to post some massive unreadable travesty
of \remarkdq slop\remarkdq\footnote{I of course amused myself, and hit
$250,000+$ pages.\par
Claude cutely called it the GRH \dots\ \remarkdq book\remarkdq{} --- I looked
at Theorem 1 and it began by noting that because the reals are totally
ordered if $\sigma \neq \frac{1}{2}$ then $\sigma > \frac{1}{2}$ or
$\sigma < \frac{1}{2}$.}, I decided to only present Lean.

In a different vein, I have a lot of people I want to praise, and I will do a
proper job in a proper writeup, alas pardon the accidental omissions ---
among many others, Manjul Bhargava, Kai Gill, Hamish Kerr, Max Leiter, John
Ma, Tal Magdish, Chris Paika, Maksym Radziwi\l\l, Jascha Sohl-Dickstein, Peter
Sarnak, Jarred Sumner, Amir Tajdini, and especially David Dohan, Ben Michel,
Ralph Furman, Eric Price, and Cora Vinther, and my miraculous colleagues. I
want them to be proud.

I have benefited from the good taste of the people at OpenAI involved
with various mathematical results, so I am quite happy to praise them too.
The Ehrhart conjecture was one such tasteful choice some time ago; playing
further and working out the equality case I saw accurate handling of
solutions of the Monge-Amp\`ere equation and immediately understood
Navier-Stokes was going to fall. Similarly, their quiet inclusion of the
Siegel argument, whose result was strictly dominated by the
$(2,3,6)$-triangle group trick, again showed great taste; though I did not
comment so publicly, I immediately recognized it as Stepanov without
multiplicity, though it was only clear to me during this rather harebrained
tornado through the whole of number theory that this problem would fall
given enough all-nighters and physical toll when I had the
$qN = M^k + R$ sign-only proof of Weil and Deligne.

\remarkpart{2}
I had a stream-of-consciousness jaunt through the experience of finding all
the pieces of the argument through incredible confusion with Claude, but a
trusted mathematician advised me to just get to the point and have some sort
of timeline. I emphasize that once we saw multiplicity $2$ on $O(2,n)$ via
the \remarkdq natively-automorphic\remarkdq{} determinant that I'd been
searching for via analogy to the definite sign of the left-hand side
of\footnote{%
\[
  \sum_{h,x_1,\ldots,x_k} t(x_1)\cdots t(x_k)\cdot
  e_q\left(h\left(x_1\cdot x_2+\cdots+x_{k-1}\cdot x_k\right)\right),
\]
sum over $h$ first for the left-hand side, with e.g.\ quadric coefficients
too Bertini gives a definite sign from middle cohomology, one off the one on
$M^k$, which is the zero-mode on the right-hand side with $R$
square-root-quality up to an inch or so as Katz would say by iterating
Parseval $k$ times, opposite signs meaning no cancellation of inverse roots,
QED; vectorize to go up dimensions. Gabber's hypersurface in one of Katz's
arguments reminds one of the quadric, which comes from looking for a
simplification of averaging over the family
$\chi\mapsto\chi(\mathrm{disc})$, corresponding to $a\mapsto y(\cdot)^a$
with $y$ the normalized Gauss sum of order $\#|\mu_K|$ who lives via
$\theta$ on the $\#|\mu_K|$-fold cover of $\mathrm{GL}_2$, the reciprocity
law of $y(c\cdot c')$ vs.\ $y(c)\cdot y(c')$ showing some sort of complete
graph, simplifying to the above \remarkdq path\remarkdq{} quadric, Gauss sum
cancellation replaced by iterated Parseval.}
$qN = M^k + R$ meaning holo/automorphy over a function field, winning more
multiplicity with more theta lifts, we were going to win but I was also, so
to speak, out of my element, in that I personally have little mastery of
automorphic forms on higher rank groups (youthful foolish stubbornness). I
also emphasize that I may have baked lots of strange suboptimal choices into
the argument, and expect short arguments to follow fairly quickly, though I
have not had time to search for such myself (yet!).

So here is a rough timeline Claude provided after a thorough
search.\footnote{The parts in which my vague intuition worked out due to
Claude's technical mastery are of course less familiar to me today, and
probably I'll have to edit some of these descriptions on reflection. The
parts in which Claude is the expert and not me I am sort of paraphrasing a
\remarkdq slop\remarkdq{} summary, which might make zero sense, and I
apologize, though the Lean should suffice to clarify everything. I'm just too
inexpert in the relevant stuff at the moment, I hope the reader sees the
breadth of this argument!}

\begin{enumerate}
\item \remarkdate{10/6/26:} Stepanov determinant.

\item \remarkdate{10/9/26:} $D^p \equiv 0$ on a CM curve for $p$ inert in
$K$, $D$ the invariant derivation, and
$\frac{D^a\cdot\theta(0)}{\Omega^a}$ is divisible by
$p^{\frac{a}{p}+\frac{a}{p^3}+\cdots-O(\log a)}$.

\item \remarkdate{10/19/26+:} $qN = M^k + R$, using
$x_1\cdot x_2+\cdots+x_{k-1}\cdot x_k$, wanting
\remarkdq automorphy\remarkdq{} of $N$ since $\mu$ is automorphic from
$\theta$ on sixfold cover of $\mathrm{GL}_5$ as a period integral of thetas,
quadric has natural interpretation on corresponding CM e.g.\
$\mathbb{C}/\mathbb{Z}[\omega]$ as Weil\footnote{A lot is obviously being
lossily compressed here, but especially after gritting with Sussman through
Ba\l aban's fantastic induction I had to think to myself about how hard
fighting for some sort of $\mathbb{F}_1$ \`a la e.g.\ Borger would be.}, put
Poincar\'e bundle on the edges to get $\theta$s.

\item \remarkdate{10/6/26+:} $W\left(\remarkone_S\right)\geq 0$ via a
determinant, interpret as Stepanov with multiplicity one again; bootstrap!!

\item \remarkdate{10/20/26:} definitively switch to $\Lambda$ aka
$\sum_{n\leq X}\Lambda(n)\cdot\left(\frac{X}{n}-1\right)_{+}-X\log X+
(1+\gamma)\cdot X$; Guinand-Weil had the sign, variance test function trick
to make things look trace-formula-like (but how can I pigeonhole / Tur\'an
over $n$? $X$? $it$?), Mehta-like integral over twists or at infinity\dots

\item \remarkdate{10/8/26:}
$W\left(\remarkone_{(0,1]}+(\sim X^{\eta})\cdot\remarkone_{(0,1/X]}\right)
\geq 0$ gives $>1-\eta$.

\item \remarkdate{10/14/26:} Determinants for $W(\cdot)$'s containing the
desired periods, indeed GGP periods, with values $0,1$ meaning interlacing
weights, or up to $m$ non-interlacing ones.

\item \remarkdate{11/15/26:} Iterated theta lifts
$U(X_0)\to U(X_1)\to\cdots\to U(X_m)$ between totally definite unitary
groups, $\pi\mapsto\pi\boxplus\mu_j[t]$, wins multiplicity $m$. Time to take
$m$ like $\sqrt{X}$\dots

\item \remarkdate{11/15/26:} Norms of the lifts and level and cancellation
of order-$m$ terms with proven divisibility via Chowla-Selberg, Weyl, and
Siegel mass.

\item \remarkdate{11/15/26:} Sufficient lower bound on divisibility of the
$m$-fold iterated lift.

\item \remarkdate{10/6/26:} Iwasawa tower via analogy to pigeonhole on
$\mathbb{F}_{q^n}/\mathbb{F}_q$.

\item \remarkdate{11/15/26:} Seesaw: the period of the iterated lift is a
Bessel period independent of the Hecke characters, Rallis.

\item \remarkdate{10/6/26:} Bost slope inequality.

\item \remarkdate{11/15/26:} precise computation of the
\remarkdq cost\remarkdq{} of each lift, showing that one would stop at $2/3$
a priori, attempt to get a main term and transfer constant to a function
field or kill via twisting, need for \remarkdq division\remarkdq.

\item \remarkdate{11/15/26:} Glue together an integral structure across
Hecke characters for a \remarkdq gcd\remarkdq{} i.e.\ extra averaging.

\item \remarkdate{10/14/26:} Relative trace formula, seen in a different
half-plane argument via sixfold cover of $\mathrm{GL}_5$ automorphy of $\mu$
and reducing via Cauchy-Schwartz to a bound on the norm of a suitably
projected Poincar\'e series, approximate orthogonality of the columns of the
integral structure.

\item \remarkdate{11/15/26:} for the classes in
$H^1\left(\mathfrak{o}_E\left[\frac{1}{\ell}\right],
\mathbb{Q}_\ell/\mathbb{Z}_\ell(-e)\right)$ that we couldn't kill by
twisting, fight for a significant weakening of Vandiver.\footnote{This is
still pretty strong!! I believe vanishing of a number of $K$-groups of
$\mathbb{Z}$ follows, and cyclicity up to $100$ish maybe? I had Lichtenbaum,
so I knew we would have a lot of slack in fighting these $K$'s.}

\item \remarkdate{11/15/26:} computation of covolumes, theta series' index,
congruences between lifts and forms.

\item \remarkdate{10/19/26:} precise choice of the Iwasawa tower,
specifically using Behrend!!\footnote{This is the sort of trick that
intuitively to me is obviously going to be unnecessary.}

\item \remarkdate{11/16/26:} GRH.
\end{enumerate}
\end{quotation}

\begin{flushright}
--- Levent Alp\"oge, untitled note.\\
{\small Source: \url{https://alpo.ge/remark.pdf}, retrieved 2026-11-21
\cite{AlpogeRemark2026}; PDF creation date 2026-11-20. Reproduced
verbatim; the bold dates, the part headings and the list layout are the
transcriber's formatting.}
\end{flushright}
\endgroup
\clearpage

\mainmatter
\clearpage
\theoremstyle{plain}
\newtheorem*{maintheorem}{Main Theorem}
\chapter*{Introduction}
\label{ch:introduction}

The Riemann hypothesis states that every zero of the Riemann zeta
function $\zeta(s)$ in the strip $0<\operatorname{Re}s<1$ has real
part $1/2$ \cite{Riemann1859}. The same statement for the
$L$-function $L(s,\chi)$ of every Dirichlet character $\chi$ is the
generalized Riemann hypothesis for Dirichlet $L$-functions. The
first statement is a sharp bound on the oscillation of the primes
around their expected density. The second is the corresponding
bound for the primes in residue classes.

In November 2026 L.~Alp\"oge published a note
\cite{AlpogeRemark2026} and an archive of Lean source files
\cite{AlpogeArchive2026}. The note is reproduced before this
introduction. It states that its author believes that he has
discovered, working with Claude, a proof of the Riemann hypothesis.
It also states that the method
gives the generalized Riemann hypothesis and that the arguments are
formalized in Lean without assumptions. The note contains a
timeline of twenty steps. It gives no proof of the theorem. The
argument is published only as the Lean archive.

These notes explain that argument for professional mathematicians
and graduate students. They state the main theorem, show how the
proof is put together from its parts, and then treat the parts one
chapter at a time. They do not verify the argument or the Lean
archive (Remark~\ref{rem:status}).

\section*{The main theorem}

Dirichlet characters and their $L$-functions are defined in
Section~\ref{sec:intro}. The character modulo~$1$ is the primitive
principal character, and its $L$-function is $\zeta$. The statement
that Alp\"oge claims is Theorem~\ref{thm:main} of
Chapter~\ref{ch:assembly}.

\begin{maintheorem}[Alp\"oge; Theorem~\ref{thm:main}]
Let $\chi$ be a Dirichlet character. Then $L(s,\chi)\ne0$ for
every $s$ with $\operatorname{Re}s>1/2$, except $s=1$ when $\chi$
is principal, and every zero of $L(s,\chi)$ with
$0<\operatorname{Re}s<1$ has real part $1/2$. Every zero of
$\zeta(s)$ other than $-2,-4,-6,\dots$ has real part $1/2$.
\end{maintheorem}

The theorem concerns $\zeta$ and Dirichlet $L$-functions only
(Remark~\ref{rem:scope}). Chapter~\ref{ch:assembly} proves it
relative to twenty statements that are used as black boxes
(Table~\ref{tab:bb-status}). Fifteen of them are statements of
Alp\"oge's argument, two are cited from the literature and three
are unattributed. These notes give complete proofs of two of the
fifteen and of the three unattributed black boxes, each relative to
results that it quotes. For the other thirteen they give an outline
of the proof, a description of its inputs or an explanatory account
of a standard method. For some of these black boxes, parts have
none of the three (Table~\ref{tab:bb-status}). So these notes do
not establish the Main Theorem.

\section*{The idea of the argument}

The model is the proof of the Riemann hypothesis for curves over
finite fields by the method of Stepanov \cite{Stepanov1969}, in the
form that Bombieri gave it \cite{Bombieri1974}. Let $C$ be a
smooth, projective, geometrically irreducible curve of genus $g$
over $\F_q$, and let $N=\#C(\F_q)$. One constructs a nonzero
rational function on $C$ that vanishes to order at least $m$ at
every rational point except one, and that has a single pole of
controlled order. A nonzero function has as many zeros as poles,
so $m(N-1)$ is at most the order of the pole. For $q$ a square with
$q>(g+1)^4$ and $m=\sqrt q$ this gives
\[
  N\le q+(2g+1)\sqrt q
\]
(Proposition~\ref{prop:stepanov}). The multiplicity $m$ costs only
one vanishing condition at each point, because $m$ is a power of
the characteristic and the map $f\mapsto f^m$ is additive. In
general the error is at most the sum of two terms, $(g+1)q/m$ from
the vanishing conditions and $2gm$ from the linear algebra that
produces the function (Proposition~\ref{prop:balancing}). The first
term falls as $m$ grows and the second rises. They have the same
order at $m\approx\sqrt q$, and no choice of the parameters gives
an error $o(\sqrt q)$ (Proposition~\ref{prop:sqrtq-forced}).

Over the integers the counterpart of the comparison of zeros with
poles is the product formula. For a nonzero integer $\Delta$,
\[
  \sum_p\ord_p(\Delta)\log p=\log|\Delta| .
\]
If $\Delta$ is known to be divisible by $p^{v_p}$ for many primes
$p$, then $\sum_pv_p\log p\le\log|\Delta|$. Chebyshev's bound for
the primes between $n$ and $2n$ is this inequality for the
binomial coefficient $\Delta=\binom{2n}{n}$
(Proposition~\ref{prop:chebyshev}). It has multiplicity one, and it
gives the right order of magnitude without a square-root error.
The ring $\Z$ has no endomorphism with the properties of the map
$f\mapsto f^m$ (item (B1) of Section~\ref{sec:dictionary}).
Alp\"oge's note says that multiplicity is ``hard to find over the
integers'' \cite[p.~1]{AlpogeRemark2026}.

The quantity that the argument bounds is a smoothed Chebyshev
function. With $\Lambda$ the von Mangoldt function and $\gamma$
Euler's constant, put
\[
  \Sz(x)=\sum_{n\le x}\Lambda(n)\Bigl(\frac xn-1\Bigr),
  \qquad
  \Gz(x)=\Sz(x)-x\log x+(1+\gamma)x .
\]
A one-sided bound for $\Gz$ is sufficient. If $0<\theta<1$ and
$\Gz(x)\le x^\theta$ for all large $x$, then $\zeta(s)\ne0$ for
$\operatorname{Re}s>\theta$
(Proposition~\ref{prop:idea-landau}). The proof uses a theorem of
Landau on Mellin transforms of nonnegative functions
\cite{Landau1905}. The target is therefore an upper bound
$\Gz(x)\le x^{1/2+\eps}$ for every $\eps>0$, together with its
analogue for Dirichlet characters.

The integers $\Delta$ are replaced by periods of theta lifts of
algebraic automorphic forms on definite unitary groups. The groups
are defined over an abelian CM field that contains an imaginary
quadratic field $K$. On the archimedean side the periods have a
size that is given by a Stirling term and a power of a CM period.
On the finite side they are divisible, to an order of the type of
Legendre's formula, by every prime that is inert in $K$. A slope
inequality for the period map between two Hermitian lattices takes
the place of the comparison of zeros with poles. Fourier analysis
on the Galois group of the CM field brings in the primes that are
not inert in $K$ and the Dirichlet characters.

Multiplicity comes from a chain of $m$ theta lifts. According to
the note, each lift collects the forced divisibility again
(Section~\ref{subsec:lifts-derivatives}). After normalisation one part
of the error is of size $x/m$ and another is of size $m$, and the
two are equal at $m\approx\sqrt x$
(Proposition~\ref{prop:sqrtx}, Remark~\ref{rem:exponent-half}).
This is the chain length $\sqrt X$ for a von Mangoldt sum of
length $X$ that the note names \cite[p.~1]{AlpogeRemark2026}.

Chapter~\ref{ch:idea} gives this overview with more detail and
with proofs of its elementary steps.
Chapter~\ref{ch:function-fields} gives the function-field model in
full. It contains a complete proof of the Riemann hypothesis for
curves by the method of Stepanov and Bombieri
(Theorem~\ref{thm:rh}). It also treats the second argument that the
note mentions for finite fields, the sign-only identity
\[
  qN=M^k+R
\]
(Proposition~\ref{prop:identity}). In this identity $M$ is the
character sum to be bounded, $N$ is a weighted point count on a
quadric in $k$ variables, and $R$ is bounded by Fourier analysis.
The bound for $M$ that Chapter~\ref{ch:function-fields} derives
from the identity is conditional (Theorem~\ref{thm:sign}).

The model and the argument differ in the direction of the
inequality. An upper bound for the function-field counterpart of
$\Gz$ is a lower bound for point counts
(Proposition~\ref{prop:which-bound}). The method of Stepanov and
Bombieri gives an upper bound for point counts. The model and the
argument also differ in the auxiliary object. The argument
constructs no auxiliary function and no space of small sections.
Its slope inequality is applied to the period map between two
given lattices (Remark~\ref{rem:periodmap}).

\section*{The structure of the proof}

Chapter~\ref{ch:assembly} gives the proof of
Theorem~\ref{thm:main} with each tool as a black box.
Figure~\ref{fig:dag} shows the dependencies. Fix a primitive
Dirichlet character $\psi$. The proof has four parts (Figure~\ref{fig:chapter-map}).

The first part is the analytic end. To $\psi$ is attached a
combination
$G_{\lev_\psi}$, with positive weights, of smoothed Chebyshev
functions of $\zeta$ and of the primitive characters that induce
the powers of $\psi$ (Definition~\ref{def:level}). Suppose that
$G_{\lev_\psi}(x)\le x^{\sigma_0}$ on windows $[Y,Y^{3/2}]$ with
$Y$ arbitrarily large. Then $L(s,\psi)\ne0$ for
$\operatorname{Re}s>\sigma_0$, except at $s=1$ when $\psi$ is
trivial. This is the windowed Landau
criterion (Theorems~\ref{thm:windowed-landau}
and~\ref{thm:level-criterion}). Because the weights are positive,
an upper bound is sufficient (Remark~\ref{rem:positivity}).
Chapter~\ref{ch:landau} treats the criterion. It also contains the
step from primitive characters and the half-plane
$\operatorname{Re}s>1/2$ to all Dirichlet characters and the
critical line (Proposition~\ref{prop:imprimitive}), which uses the
functional equation (Theorem~\ref{thm:fe}).

The second part is the inequality. The Chebyshev inequality on
windows (Theorem~\ref{thm:chebyshev-windows}) concerns a family of
automorphic data that is \emph{tied} to a slope record at every
large point of its window set. For the families that
Chapter~\ref{ch:families} produces it gives
$G_{\lev_\psi}(x)\le x^{1/2+9\eps}$ on the windows of the family.

Five chapters lead to the inequality.
Chapter~\ref{ch:hosts} constructs the CM fields, which are called
host fields (Theorem~\ref{thm:fields-exist}). It bounds the
characters of their Galois groups that are not wanted by an
average over two auxiliary primes
(Theorem~\ref{thm:good-pairs}). Chapter~\ref{ch:balance} proves
the balance identity (Proposition~\ref{prop:balance}). The
identity writes a prime sum of the type of Legendre's formula as a
main term plus a combination of twisted Chebyshev functions.
Chapter~\ref{ch:balance} also contains the core bound
(Theorem~\ref{thm:core-bound}), which turns an upper bound for the
prime sum into an upper bound for $G_{\lev_\psi}$.

Chapter~\ref{ch:theta} defines the tower of unitary groups, the
chain of theta lifts, Hermitian lattices and slope records. Its
main statement is the frame theorem (Theorem~\ref{thm:frame}), a
slope comparison between many source lattices and one target
lattice. Chapter~\ref{ch:covolume} states the covolume bound for
lattices of integral automorphic forms
(Theorem~\ref{thm:covolume}). This theorem is the one place where
arithmetic geometry enters the argument.
Chapter~\ref{ch:chebyshev} combines these results in the gap bound
(Theorem~\ref{thm:gap}) and deduces the inequality. The exponent
$1/2$ comes from Proposition~\ref{prop:sqrtx}.

The third part consists of two arithmetic inputs. The lattices of
the chain differ from the ideal ones by congruence defects. The
defects are bounded
by Tate--Shafarevich groups of Tate twists over the host fields.
Chapter~\ref{ch:sha} states the bound that is needed, the Sha
bound on windows (Theorem~\ref{thm:sha-windows}). The bound is
obtained by a pigeonhole argument that gives it only on windows
(Remark~\ref{rem:windows-origin}). For this reason the
Landau criterion is needed on windows
(Remark~\ref{rem:landau-compare}). The Sha bound uses Galois
representations attached to automorphic representations over CM
fields, with a truncated form of local--global compatibility.
Chapter~\ref{ch:galois} states this input
(Theorem~\ref{thm:galois-reps}).

The fourth part is existence. Chapter~\ref{ch:families} states
that hosts are supplied on windows
(Theorem~\ref{thm:host-supply}), that packages
of automorphic data exist (Theorem~\ref{thm:packages}) and that
tied families exist (Theorem~\ref{thm:tied-families}). It uses
Green's parametrisation of the cuspidal characters of
$\mathrm{GL}_n(\F_p)$ \cite{Green1955}. From these statements, the
Chebyshev inequality on windows and the Landau criterion it then
deduces that $L(s,\psi)\ne0$ for $\operatorname{Re}s>1/2$, except
at $s=1$ when $\psi$ is trivial
(Theorem~\ref{thm:zero-free-primitive}).

The proof of Theorem~\ref{thm:main} in
Section~\ref{sec:proof-main} is a chain of six implications in
this order: Galois representations, the Sha bound on windows, the
supply of hosts, the existence of tied families, zero-freeness for
primitive characters, and the step to all characters and the
critical line. The deduction contains no substantial argument of
its own (Remark~\ref{rem:content}).
{\renewcommand{\thefigure}{I.\arabic{figure}}%
\begin{figure}[tbp]
\centering
\input{book/figs/ch00-chapter-map.tex}
\end{figure}}

\section*{What is proved in these notes}

For each result of the argument these notes give the statement
and the place where it enters the proof. They then give one of the
following: a complete proof, an outline of Alp\"oge's route, a
description of the inputs of the proof, an explanatory account of a
standard method, or a citation. The attribution of each
statement and the kind of each proof are marked where the statement
occurs (Section~\ref{sec:conventions}). Table~\ref{tab:bb-status}
shows what the notes give for each black box, and
Remark~\ref{rem:bb-open} lists the points that the later chapters
leave open. Two open points are the main ones. The first is the
step from the windows $[Y,Y^2]$ of the Sha bound to the windows
$[W,W^{3/2}]$ of the supply of hosts: the smaller window lies in
the larger one and Theorem~\ref{thm:zero-free-primitive} accepts
every window exponent $\lambda>1$, so the open question is only
whether the scale also changes between the two statements
(Question~\ref{q:window-exponent},
Remark~\ref{rem:window-exponent}). The second is
Hypothesis~\ref{H2} of the sign-only argument, a vanishing
statement in \'etale cohomology that the bound of
Theorem~\ref{thm:sign} assumes and that is not proved
(Question~\ref{q:sign-h2}).

\section*{Prerequisites and conventions}

Chapters~\ref{ch:idea} and~\ref{ch:function-fields} use the
algebraic geometry of curves: divisors, the Riemann--Roch theorem
and function fields. The sections on the sign-only argument also
use $\ell$-adic sheaves on varieties over finite fields, the
Grothendieck--Lefschetz trace formula, Artin's vanishing theorem
and Poincar\'e duality.

Chapters~\ref{ch:landau}, \ref{ch:hosts} and~\ref{ch:balance}
assume the theory of Dirichlet characters and their $L$-functions
at the level of a first course in analytic number theory
\cite{Davenport2000,MontgomeryVaughan2007}, which includes the
functional equation and the explicit formula.
Remark~\ref{rem:classical-inputs} lists the further classical
results that Chapter~\ref{ch:hosts} cites.
Chapters~\ref{ch:theta}, \ref{ch:chebyshev}
and~\ref{ch:families} use Hermitian lattices and their slopes
\cite{Bost1996}, algebraic automorphic forms on definite unitary
groups \cite{Gross1999} and the theta correspondence for unitary
dual pairs \cite{Howe1989,Rallis1984}. Chapter~\ref{ch:covolume}
cites results on integral models and compactifications of Shimura
varieties and on Arakelov geometry. Chapter~\ref{ch:sha} uses
Galois cohomology of number fields, which includes Poitou--Tate
duality \cite{Tate1963,Poitou1967}. Chapter~\ref{ch:galois} cites
the construction of Galois representations attached to
automorphic representations of $\mathrm{GL}_n$ over CM fields
\cite{HarrisTaylor2001,Shin2011,ChenevierHarris2013}. Results
from these subjects are cited where they are used.

Theorems, propositions, definitions, remarks and questions share
one counter within each section: Theorem~\ref{thm:main} is the
first numbered statement of Section~\ref{sec:intro}. Equations are
numbered within sections. The conventions for attributions, for
the four kinds of proof block and for schematic definitions are in
Section~\ref{sec:conventions}. They hold in
Chapters~\ref{ch:assembly}--\ref{ch:families}.
Section~\ref{sec:bb-vocabulary} recalls the terms in which the
black boxes are stated: levels, windows, scales, hosts, packages,
families and ties. From
Section~\ref{sec:notation} to the end of Section~\ref{sec:role} the
letters $M$, $N$, $R$, $k$ and $t$ have the meaning that they have
in the identity $qN=M^k+R$. In the first part of
Chapter~\ref{ch:function-fields} they have a different meaning.

\section*{Organization of the notes}

The front matter contains a preface and Alp\"oge's note, which is
reproduced in full. Chapter~\ref{ch:idea} is the informal overview.
Table~\ref{tab:idea-roadmap} gives, for each step of the overview,
the chapter that develops it. Chapter~\ref{ch:function-fields} is
the function-field model. Its Section~\ref{sec:dictionary} compares
each step of the model with the part of the argument that has the
same role, and rates each comparison. Chapter~\ref{ch:assembly}
states the main theorem and all black boxes and gives the proof.
Chapters~\ref{ch:landau}--\ref{ch:families} treat the black boxes
in the order of the section on the structure of the proof, above.
Appendix~\ref{ch:ai} describes the software and the AI methods that
produced the text.

Chapters~\ref{ch:idea} and~\ref{ch:function-fields} are not logical
inputs to the later chapters. A reader who wants only the
structure of the argument can start with
Chapter~\ref{ch:assembly}. Chapters~\ref{ch:landau}--\ref{ch:families} use from
Chapter~\ref{ch:assembly} only the statement of
Theorem~\ref{thm:main}, and each of them opens with a paragraph
that says where its results enter the proof. The lattices and
automorphic forms of Chapter~\ref{ch:theta} enter the analytic
argument of Chapters~\ref{ch:balance} and~\ref{ch:chebyshev} only
through the slope record (Definitions~\ref{def:record}
and~\ref{def:size-data}).

\section*{AI and Software Discussion}
\label{intro:ai}

AI agents assembled and rewrote the text of these notes on
2026-11-21 from drafts that AI agents wrote on 2026-11-20 and
2026-11-21. The models for the assembly were Claude Fable~5.2
(\texttt{claude-fable-5-2}) \cite{ClaudeFable52} and Claude
Opus~5.5 (\texttt{claude-opus-5-5}) \cite{ClaudeOpus55}, in Claude
Code \cite{ClaudeCode}. The model \texttt{claude-fable-5-2} drafted
this introduction. Appendix~\ref{ch:ai} gives the methods, the
citations for all software that was used, which includes the
software used for the drafts, and the token accounting. No human
has reviewed the text.
\clearpage
\clearpage
\chapter{The basic idea}
\label{ch:idea}

This chapter is an overview of the whole book. It describes the shape
of the argument for the Riemann Hypothesis announced by L.~Alp\"oge
\cite{AlpogeRemark2026}, and it does so by following one classical
proof. Over a finite field the Riemann hypothesis for curves has an
elementary proof, due to Stepanov \cite{Stepanov1969} and put in
general form by Bombieri \cite{Bombieri1974}. One builds a nonzero
function that vanishes to order at least $m$ at every rational point
and whose number of poles is controlled. Since a nonzero function has
as many zeros as poles, one gets
\[
  (\text{number of rational points}) \cdot m
  \;\le\; (\text{number of poles}),
\]
and the choice $m \approx \sqrt q$ gives the Weil bound.

Alp\"oge describes his argument as a descendant of Stepanov's. He
names the main difficulty: over the integers multiplicity is ``hard
to find'', and corresponds to something like differentiating whole
numbers \cite[p.~1]{AlpogeRemark2026}. The purpose of the chapter is
to make that sentence precise enough to read the rest of the book.
Nothing here is used later as a proved input: the chapter states what
each later chapter is for, and every step ends with a pointer to the
place where it is developed. Table~\ref{tab:idea-roadmap} collects
these pointers.

Three kinds of statement occur, and they should be kept apart.
Statements proved here are propositions with proofs. Statements
proved elsewhere in the book are quoted with a reference to the
chapter that proves them. Statements about Alp\"oge's argument that
are only reports, of the note \cite{AlpogeRemark2026} or of
statements read in the Lean archive \cite{AlpogeArchive2026}, are
marked as reports, and the interpretation that connects them is
marked as an interpretation.

\begin{table}[htbp]
\centering
\begin{tabular}{@{}p{0.50\linewidth}p{0.44\linewidth}@{}}
\toprule
Step of the overview & Where it is developed \\
\midrule
Stepanov's method over $\F_q$
  (Section~\ref{sec:stepanov})
  & Chapter~\ref{ch:function-fields},
    Sections~\ref{sec:upper}--\ref{sec:rh} \\
The sign-only argument and $qN = M^k + R$
  & Chapter~\ref{ch:function-fields},
    Sections~\ref{sec:identity}--\ref{sec:role} \\
The dictionary (Section~\ref{sec:idea-dictionary})
  & Chapter~\ref{ch:function-fields},
    Section~\ref{sec:dictionary};
    Chapter~\ref{ch:assembly} \\
The quantity to bound, a one-sided criterion
  (Section~\ref{sec:idea-landau})
  & Chapter~\ref{ch:landau} \\
Seeing every prime: characters and host fields
  & Chapters~\ref{ch:hosts} and~\ref{ch:balance} \\
Divisibility is at most size (Section~\ref{sec:slope})
  & Chapters~\ref{ch:balance}, \ref{ch:theta}
    and~\ref{ch:covolume} \\
Multiplicity and the exponent $1/2$
  (Section~\ref{sec:multiplicity})
  & Chapter~\ref{ch:chebyshev} \\
Theta lifts (Section~\ref{sec:theta})
  & Chapter~\ref{ch:theta} \\
The costs the outline omits (Section~\ref{sec:idea-roadmap})
  & Chapters~\ref{ch:sha}, \ref{ch:galois}
    and~\ref{ch:families} \\
The argument in one place
  & Chapter~\ref{ch:assembly} \\
\bottomrule
\end{tabular}
\caption{The steps of this chapter and the chapters that develop
them.}
\label{tab:idea-roadmap}
\end{table}

\section{Stepanov's method over a finite field}\label{sec:stepanov}

Throughout this section $C$ is a smooth, projective, geometrically
irreducible curve of genus $g$ over a finite field $\F_q$ of
characteristic $p$, and $N = \#C(\F_q)$. For a divisor $D$ on $C$
defined over $\F_q$, let $\Lsp(D)$ be the $\F_q$-vector space of
rational functions $f$ on $C$ defined over $\F_q$ with $f = 0$ or
$\operatorname{div}(f) + D \ge 0$, and let $\ell(D) = \dim \Lsp(D)$.

Two facts from the theory of curves are used; for both see
\cite[Exercises~16.3.E and~18.4.B, \S18.5.3]{Vakil2025}. The
first is Riemann's inequality $\ell(D) \ge \deg D + 1 - g$ for every
$D$, with equality when $\deg D \ge 2g - 1$; both statements are
part of the Riemann--Roch theorem. The second is that a nonzero
rational function has as many zeros as poles, counted with
multiplicity and degree.

\begin{proposition}[Stepanov--Bombieri bound]\label{prop:stepanov}
Let $C$ be a smooth, projective, geometrically irreducible curve of
genus $g$ over $\F_q$. Suppose that $q$ is a square and that
$q > (g+1)^4$. Then
\[
  \#C(\F_q) \;\le\; q + (2g+1)\sqrt q.
\]
\end{proposition}

This is the case $\sigma = \mathrm{id}$ of
Theorem~\ref{thm:bombieri}, and the complete proof is given there.
What follows is a sketch. It is meant to show where each of the
three ideas enters, because each of them has a counterpart in the
rest of the book.

\subsection*{Sketch of the proof}

Fix a rational point $x_0$, write $m = \sqrt q$, a power of $p$, and
let $A$ and $B$ be nonnegative integers with $mA < q$, to be chosen
at the end. The whole sketch rests on the additivity of
$f \mapsto f^m$, which holds because $m$ is a power of the
characteristic. This additivity replaces differentiation.

The first idea is independence. Choose a basis $s_1, \dots, s_r$ of
$\Lsp(Ax_0)$ whose elements have pairwise distinct pole orders
$n_1 < \dots < n_r \le A$ at $x_0$. If $w_1, \dots, w_r$ are
rational functions with $\sum_i s_i^m w_i^q = 0$, then every $w_i$
is zero. Indeed a nonzero term has valuation
$-mn_i + q\ord_{x_0}(w_i)$ at $x_0$, which is congruent to $-mn_i$
modulo $q$, and the integers $mn_i$ are distinct and lie in
$[0, q)$. So the nonzero terms have pairwise distinct valuations,
and they cannot cancel. In Chapter~\ref{ch:function-fields} this is
Propositions~\ref{prop:adapted-basis} and~\ref{prop:independence}.

The second idea is that linear algebra produces an auxiliary
function. Consider the $\F_q$-linear map
\[
  \Phi \colon \Lsp(Bx_0)^r \to \Lsp\bigl((mA + B)x_0\bigr), \qquad
  (v_1, \dots, v_r) \mapsto \sum_{i=1}^r s_i^m v_i .
\]
By Riemann--Roch its source has dimension at least
$(A + 1 - g)(B + 1 - g)$ and its target has dimension
$mA + B + 1 - g$ when $mA + B \ge 2g - 1$. So if
\begin{equation}\label{eq:kernel}
  (A + 1 - g)(B + 1 - g) > mA + B + 1 - g ,
  \qquad A + 1 - g > 0 ,
\end{equation}
then $\Phi$ has a nonzero kernel element $(v_1, \dots, v_r)$. Put
\[
  f = \sum_{i=1}^r s_i^m v_i^q .
\]
By the first idea $f \ne 0$. Its only pole is at $x_0$, of order at
most $mA + qB$.

The third idea is to count zeros against poles. Let $P \ne x_0$ be a
rational point. The values $v_i(P)$ lie in $\F_q$, so $v_i(P)^q = v_i(P)$
and $f(P) = \Phi(v_1, \dots, v_r)(P) = 0$. But
$f = \bigl(\sum_i s_i v_i^{q/m}\bigr)^m$ is an $m$-th power of a
function that is regular at $P$, so one vanishing condition gives
vanishing to order $m$. Counting zeros against poles,
\begin{equation}\label{eq:zeros-poles}
  m\,(N - 1) \;\le\; mA + qB .
\end{equation}
With $A = m - 1$ and $B = m + 2g$, condition \eqref{eq:kernel}
becomes $m > (g+1)^2$, which is the hypothesis, and
\eqref{eq:zeros-poles} becomes $N \le q + (2g+1)\sqrt q$.
(Remark~\ref{rem:idea-pair} checks this choice against the general
statement of Chapter~\ref{ch:function-fields}.)

Figure~\ref{fig:idea-zero-pole-budget} shows the count behind
\eqref{eq:zeros-poles} for small numbers.
\begin{figure}[tbp]
\centering
\adjustbox{max width=\linewidth}{%
\input{alpoge/figs/zero-pole-budget.tex}}
\caption{Zeros against poles in the sketch of
Proposition~\ref{prop:stepanov}. The auxiliary function $f$ vanishes
to order at least $m$ at every rational point other than $x_0$ and
has a single pole, of order at most $mA+qB$, at $x_0$. The ledger on
the right is drawn to scale for $q=64$, $g=1$, $m=8$, $A=7$, $B=10$,
which satisfy the hypotheses and condition \eqref{eq:kernel}: at
most $87$ blocks of length $8$ fit under $696$, so
$N\le88=q+(2g+1)\sqrt q$. The curve is schematic.}
\label{fig:idea-zero-pole-budget}
\end{figure}

\subsection*{The balance}

The same balance of parameters reappears in
Section~\ref{sec:multiplicity}, so we isolate it here. Inequality
\eqref{eq:zeros-poles} says $N \le 1 + A + qB/m$. The constraint
$mA < q$ forces $A < q/m$, and the kernel condition needs $B$ of
size about $m$ plus a constant times $g$. A short calculation with
$A = q/m - 1$ and $B = m + 2g$ shows that \eqref{eq:kernel} holds
only when $m^2 \lesssim q(g+1)/g$, and then
\[
  N \;\le\; q + (2g + 1)\,\frac{q}{m} + O(1).
\]
So the error term is $q/m$: more multiplicity helps. With these
values of $A$ and $B$ the kernel condition fails when $m$ exceeds
about $\sqrt q$. A larger $B$ restores the condition and adds to
the order of the pole. In general the error is at most
$(g+1)q/m + 2gm$ (Proposition~\ref{prop:balancing}): the first
term is the price of the vanishing conditions and falls as $m$
grows, and the second is the price of the kernel condition and
rises. The two terms have the same order at $m \approx \sqrt q$,
and no choice of the parameters makes the method give an error
$o(\sqrt q)$ (Proposition~\ref{prop:sqrtq-forced}).
Figure~\ref{fig:stepanov-balance} computes this balance exactly in
one case.
\begin{figure}[tbp]
\centering
\adjustbox{max width=\linewidth}{%
\input{alpoge/figs/stepanov-balance.tex}}
\caption{The bound on $N-q$ given by \eqref{eq:zeros-poles}, divided
by $\sqrt q$, as a function of the multiplicity $m=2^j$, for $p=2$,
$q=2^{24}$, $g=3$ and $A=q/m-1$, in exact integer arithmetic. With
$B=m+2g$ (triangles) the bound is $(2g+1)q/m$, but condition
\eqref{eq:kernel} fails for $m>\sqrt q$ (shaded). With the least $B$
satisfying \eqref{eq:kernel} (dots) the bound is about
$(g+1)q/m+gm$ (grey lines). Both are smallest at $m=\sqrt q$, where
the bound is $q+(2g+1)\sqrt q$.}
\label{fig:stepanov-balance}
\end{figure}

\subsection*{From the point count to the zeta function}

\begin{remark}[From the upper bound to the Riemann hypothesis]
\label{rem:rh-curves}
Let $\alpha_1, \dots, \alpha_{2g}$ be the reciprocal roots of the
numerator of the zeta function of $C$, so that
$\#C(\F_{q^n}) = q^n + 1 - \sum_i \alpha_i^n$ for all $n \ge 1$
(Theorem~\ref{thm:rationality}). Proposition~\ref{prop:stepanov},
applied over $\F_{q^n}$ for even $n$ large, gives an \emph{upper}
bound $\#C(\F_{q^n}) \le q^n + (2g+1) q^{n/2}$. Bombieri obtains a
matching lower bound $\#C(\F_{q^n}) \ge q^n - c\,q^{n/2}$, with $c$
independent of $n$, by applying the upper bound to the twists of a
Galois cover of $C$ \cite{Bombieri1974}. This is
Section~\ref{sec:route1}. With both bounds,
$\sum_i (\alpha_i^2)^k = O(q^k)$ for $k$ large. The power series
$\sum_k \bigl(\sum_i \alpha_i^{2k}\bigr) t^k
= \sum_i \alpha_i^2 t / (1 - \alpha_i^2 t)$ then converges for
$|t| < 1/q$, so it has no pole there, and $|\alpha_i| \le \sqrt q$
for every $i$. The complete argument is Theorem~\ref{thm:rh}.
\end{remark}

The argument that works for $\zeta$ in
Proposition~\ref{prop:idea-landau} below does not derive
$|\alpha_i| \le \sqrt q$ from the upper bound alone. The reason is
the first essential difference between the finite-field model and
the integers.

\begin{remark}[The Pringsheim argument does not suffice for curves]
\label{rem:one-sided-curves}
Let $\alpha_1, \dots, \alpha_{2g}$ be the reciprocal roots of the
numerator of the zeta function of $C$. For even $n$ large, the upper
bound $\#C(\F_{q^n}) \le q^n + (2g+1)q^{n/2}$ of
Remark~\ref{rem:rh-curves} says only that
$\sum_i \alpha_i^n \ge 1 - (2g+1)q^{n/2}$. If one $\alpha_i$ were
real with $|\alpha_i| > \sqrt q$, then for even $n$ the sum would be
large and \emph{positive}, so the upper bound would not be
contradicted. Pringsheim's theorem on power series with nonnegative
coefficients, applied to
$\sum_k \bigl(\sum_i\alpha_i^{2k} + (2g+1)q^k - 1\bigr)t^k$ after
changing finitely many coefficients, shows only this: if some
$|\alpha_i| > \sqrt q$, then some $\alpha_i$ of maximal modulus is
real. It cannot rule that case out. For $\zeta$ a one-sided bound
suffices (Proposition~\ref{prop:idea-landau}). The difference lies
in the direction of the inequality. The proof for $\zeta$ does not
use the absence of zeros of $\zeta$ on the real segment $(0,1)$. In
the function-field model an upper bound for the counterpart of
$\Gz$ is a \emph{lower} bound for point counts
(Proposition~\ref{prop:which-bound}). For a bound in that direction
a real root of maximal modulus enters with a sign that the bound
forbids, and the one-sided bound suffices
(Theorem~\ref{thm:route2prime}). The proof of
Proposition~\ref{prop:idea-landau} excludes real zeros of $\zeta$
in the same way, by the sign of a residue. The Stepanov--Bombieri
bound is the bound in the other direction, and Bombieri uses a
second argument for curves. The remark shows that this method
fails. It does not show that no other argument derives the
conclusion from the upper bound. Section~\ref{sec:route2} makes all
of this exact. There the one-sided bound proves ``the Riemann
hypothesis or a real zero of
the $L$-polynomial in one explicit interval''
(Theorems~\ref{thm:dichotomy} and~\ref{thm:route2-input}), and a
formal example shows that the second case cannot be excluded by the
formal properties of zeta functions
(Proposition~\ref{prop:formal-cex}).
\end{remark}

\section{The dictionary}\label{sec:idea-dictionary}

The strategy is to translate each line of the sketch of
Proposition~\ref{prop:stepanov} into arithmetic.

\medskip
\begin{center}
\begin{tabular}{@{}ll@{}}
\toprule
Curve over $\F_q$ & Integers \\
\midrule
rational point $P$ & prime $p$ \\
vanishing to order $m$ at $P$ & divisibility by $p^{m(\cdots)}$ \\
number of poles of $f$ & log of the size of an algebraic number \\
zeros $=$ poles for $f \ne 0$ & product formula; slope inequality \\
upper bound for $N$ & upper bound for a prime sum $\Gz(x)$ \\
$v^q = v$ on $\F_q$-points & reduction modulo $p$ \\
$f = (\cdots)^m$ with $m$ a power of $p$
  & a chain of $m$ theta lifts \\
\bottomrule
\end{tabular}
\end{center}
\medskip

Each row is explained below. The guiding principle is the rank-one
case of the third and fourth rows. Let $\Delta$ be a nonzero
integer. Then
\begin{equation}\label{eq:product-formula}
  \sum_{p} \ord_p(\Delta) \log p \;=\; \log |\Delta| .
\end{equation}
Read the left side as the number of zeros of $\Delta$, each prime
counted with its multiplicity $\ord_p(\Delta)$ and weight $\log p$.
Read the right side as the number of poles. If we know in advance
that $\Delta$ is divisible by $p^{v_p}$ for many primes, and we can
bound $|\Delta|$ from above, then
\begin{equation}\label{eq:div-size}
  \sum_p v_p \log p \;\le\; \log |\Delta| .
\end{equation}
This is ``divisibility is at most size'', the arithmetic analogue of
\eqref{eq:zeros-poles}. The whole method consists of choosing
$\Delta$, or a lattice of such numbers, so that the left side is a
prime sum we want to bound. Figure~\ref{fig:product-formula} draws
\eqref{eq:product-formula} next to its function-field counterpart.
\begin{figure}[tbp]
\centering
\adjustbox{max width=\linewidth}{%
\input{alpoge/figs/product-formula.tex}}
\caption{Zeros equal poles, on a curve and on
$\operatorname{Spec}\Z$. Top (schematic): a nonzero function has as
many zeros as poles, counted with order. Bottom: for
$\Delta=\binom{20}{10}=2^2\cdot11\cdot13\cdot17\cdot19$, the stems
$\ord_p(\Delta)\log p$ stacked at the archimedean place $\infty$ add
up to $\log|\Delta|\approx12.13$, which lies below Chebyshev's bound
$20\log2$ (Proposition~\ref{prop:chebyshev}).}
\label{fig:product-formula}
\end{figure}

\begin{proposition}[Chebyshev]\label{prop:chebyshev}
Let $n \ge 1$ be an integer. Then
$\sum_{n < p \le 2n} \log p < 2n \log 2$, the sum running over
primes.
\end{proposition}

\begin{proof}
Every prime $p$ with $n < p \le 2n$ divides $(2n)!$ and not $n!^2$,
so it divides $\Delta = \binom{2n}{n}$. Apply \eqref{eq:div-size}
with $v_p = 1$ for those primes, and $\Delta < (1+1)^{2n}$.
\end{proof}

Chebyshev's argument is Stepanov's method with multiplicity one and
without a good main term. It gives upper bounds of the right order
of magnitude for prime sums \cite{Chebyshev1852}, but no asymptotic.
Two things must be added: an exact cancellation of main terms, and
multiplicity. The four steps below supply them.

The same comparison organises Chapter~\ref{ch:assembly}: its
Table~\ref{tab:analogy} sets Chebyshev's factorial argument beside
the parts of the argument that replace it, and
Remark~\ref{rem:factorial} states in what sense the argument is
Chebyshev's, twisted by characters and sharpened by a chain of
lifts. A second dictionary, from the function-field side and with
each row rated, is in Section~\ref{sec:dictionary}.

\section{The quantity to bound}\label{sec:idea-landau}

\begin{definition}\label{def:sz-gz}
Let $\Lambda$ be the von Mangoldt function and $\gamma$ Euler's
constant. For real $x \ge 1$ put
\begin{enumerate}[label=(\roman*)]
\item $\Sz(x) = \sum_{n \le x} \Lambda(n)\,(x/n - 1)$;
\item $\Gz(x) = \Sz(x) - x\log x + (1 + \gamma)\,x$.
\end{enumerate}
\end{definition}

These are the smoothed Chebyshev function of $\zeta$ and its
normalised form. Definition~\ref{def:chebyshev} repeats them
together with their twists by Dirichlet characters.
Figure~\ref{fig:gz-size} shows $\Gz$ and the size the Riemann
Hypothesis predicts for it.
\begin{figure}[tbp]
\centering
\adjustbox{max width=\linewidth}{\input{book/figs/gz-size.tex}}
\caption{The quantity to bound. (a) $\Gz(x)$ for $x\le2000$, computed from
the von Mangoldt function, inside the band
$\log2\pi\pm c_0\sqrt x$, where
$c_0=2+\gamma-\log4\pi\approx0.046$ equals $\sum_\rho|\rho|^{-2}$
under the Riemann Hypothesis. (b) $(\Gz(x)-\log2\pi)/\sqrt x$ for
$10\le x\le10^7$ (logarithmic axis), with the explicit formula
evaluated from the first $500$ zeros (triangles). The plot samples
$2500$ values of $x$; it illustrates the prediction and proves
nothing.}
\label{fig:gz-size}
\end{figure}

Alp\"oge's note calls the switch to this sum one of the decisive
steps: ``Guinand--Weil had the sign''
\cite[p.~3, step~5]{AlpogeRemark2026}. Since
$\sum_{n \le x}\Lambda(n)(x/n - 1) = x\sum_{n\le x}\Lambda(n)
\int_n^x t^{-2}\,dt$, we have
$\Sz(x) = x\int_1^x \Ch(t)\,t^{-2}\,dt$, where
$\Ch(t) = \sum_{n \le t}\Lambda(n)$ is Chebyshev's function. So
$\Sz$ is a smoothed form of $\Ch$.

\begin{proposition}[Mellin transform]\label{prop:mellin}
Let $s$ be a complex number with $\operatorname{Re} s > 1$. Then
\[
  \int_1^\infty \Sz(x)\, x^{-s-1}\,dx
  \;=\; -\frac{\zeta'(s)}{\zeta(s)} \cdot \frac{1}{s(s-1)} ,
\]
and
\[
  \int_1^\infty \Gz(x)\, x^{-s-1}\,dx
  \;=\; -\frac{\zeta'(s)}{\zeta(s)} \cdot \frac{1}{s(s-1)}
  - \frac{1}{(s-1)^2} + \frac{1 + \gamma}{s - 1} .
\]
\end{proposition}

\begin{proof}
For each $n$, $\int_n^\infty (x/n - 1)\,x^{-s-1}\,dx
= n^{-s}\bigl(\frac{1}{s-1} - \frac{1}{s}\bigr)
= \frac{n^{-s}}{s(s-1)}$. Summing against $\Lambda(n)$, which is
allowed by absolute convergence, gives the first identity, since
$\sum_n \Lambda(n) n^{-s} = -\zeta'(s)/\zeta(s)$. The second follows
from $\int_1^\infty x\log x\cdot x^{-s-1}\,dx = (s-1)^{-2}$ and
$\int_1^\infty x\cdot x^{-s-1}\,dx = (s-1)^{-1}$.
\end{proof}

The right side of the second identity is meromorphic in
$\operatorname{Re} s > 0$. Near $s = 1$ we have
$-\zeta'/\zeta(s) = (s-1)^{-1} - \gamma + O(s-1)$ and
$1/(s(s-1)) = (s-1)^{-1} - 1 + O(s-1)$, so the first term is
$(s-1)^{-2} - (1+\gamma)(s-1)^{-1} + O(1)$. This is exactly what is
subtracted: the right side is regular at $s = 1$. Its only poles in
$\operatorname{Re} s > 0$ are at the zeros $\rho$ of $\zeta$, each
with residue $-\operatorname{mult}(\rho)/(\rho(\rho-1)) \ne 0$,
where $\operatorname{mult}(\rho)$ is the multiplicity of the zero.
This explains the choice of $\Gz$. The terms
$x\log x - (1+\gamma)x$ remove the pole at $s = 1$, and the kernel
$1/(s(s-1))$ makes the sum over zeros in the explicit formula
converge absolutely. Under the Riemann Hypothesis,
$\Gz(x) = O(\sqrt x)$, by the explicit formula
$\Gz(x) = -\sum_\rho x^\rho/(\rho(\rho-1)) + O(1)$ and
$\sum_\rho |\rho|^{-2} < \infty$; see \cite{Ingham1932} and
Remark~\ref{rem:explicit}. The converse needs only a one-sided
bound, by the following theorem of Landau.

\begin{theorem}[Landau \cite{Landau1905}; see \cite{Ingham1932}]
\label{thm:idea-landau}
Let $x_0 \ge 1$, and let $h\colon [x_0, \infty) \to [0, \infty)$ be
locally integrable. Suppose
$F(s) = \int_{x_0}^\infty h(x)\,x^{-s-1}\,dx$ converges for some
real $s$ and diverges for some real $s$, and let $\sigma_c$ be the
infimum of the real $s$ for which it converges. Then $F$, which is
analytic in $\operatorname{Re} s > \sigma_c$, has no analytic
continuation to any neighbourhood of $s = \sigma_c$.
\end{theorem}

The substitution $x = e^u$ turns this into the statement for Laplace
transforms of nondecreasing functions, which is the form quoted as
Theorem~\ref{thm:landau} in Chapter~\ref{ch:landau}.

\begin{proposition}[One-sided criterion]\label{prop:idea-landau}
Let $0 < \theta < 1$ and $x_0 \ge 1$. Suppose $\Gz(x) \le x^\theta$
for all $x \ge x_0$. Then $\zeta(s) \ne 0$ for
$\operatorname{Re} s > \theta$.
\end{proposition}

\begin{proof}
Put $h(x) = x^\theta - \Gz(x) \ge 0$ on $[x_0, \infty)$. Since
$\Gz(x) = O(x\log x)$, the integral $F(s)$ of
Theorem~\ref{thm:idea-landau} converges for
$\operatorname{Re} s > 1$. By Proposition~\ref{prop:mellin}, $F$
equals $x_0^{\theta - s}/(s - \theta)$ minus the right side of the
second identity, plus an entire function (the integral over
$[1, x_0]$). This expression is meromorphic in
$\operatorname{Re} s > 0$, with poles only at $s = \theta$ and at
the zeros of $\zeta$.

The residue of the expression at $s = \theta$ is
$1 + m/(\theta(\theta - 1))$, where $m \ge 0$ is the order of
vanishing of $\zeta$ at $\theta$. It is not zero: if $m \ge 1$ then
$m/(\theta(1-\theta)) \ge 4$. So the expression has a pole at
$s = \theta$. At a real zero $\beta \in (0,1)$ of $\zeta$ of
multiplicity $m$, other than $\theta$, its residue is
\[
  \frac{m}{\beta(\beta - 1)} \;=\; -\frac{m}{\beta(1-\beta)} \;<\; 0 .
\]

Let $\sigma_c$ be the infimum of the real $s$ for which the integral
converges. The integral cannot converge for every real $s$: then
$F$ would be analytic in $\operatorname{Re} s > 0$, and it has a
pole at $s = \theta$. So Theorem~\ref{thm:idea-landau} applies, and
$\sigma_c$ is a real pole of the meromorphic expression. Suppose
$\sigma_c > \theta$. The point $s = 1$ is not a pole, so
$\sigma_c = \beta$ is a zero of $\zeta$ in $(\theta, 1)$. For real
$s > \beta$ we have $F(s) \ge 0$, because $h \ge 0$. But the residue
of $F$ at $\beta$ is negative, so $F(s) \to -\infty$ as $s$
decreases to $\beta$. This is a contradiction. Hence
$\sigma_c \le \theta$, so the integral converges and is analytic in
$\operatorname{Re} s > \theta$. It follows that the meromorphic
expression has no poles there, and so $\zeta$ has no zeros there.
\end{proof}

The proof does not use the classical fact that $\zeta$ has no zero
on the real segment $(0,1)$, which follows from
$(1 - 2^{1-\sigma})\zeta(\sigma)
= \sum_{n \ge 1}(-1)^{n-1}n^{-\sigma} > 0$ for $0 < \sigma < 1$. A
real zero in $(\theta, 1)$ is excluded by the sign of its residue
alone. The same mechanism serves for
$L$-functions, whose real zeros are not excluded in advance
(Remark~\ref{rem:positivity}), and in the function-field model
(Theorem~\ref{thm:route2prime}).

Figure~\ref{fig:landau-plane} shows the argument in the $s$-plane.
\begin{figure}[tbp]
\centering
\adjustbox{max width=\linewidth}{%
\input{alpoge/figs/landau-plane.tex}}
\caption{The proof of Proposition~\ref{prop:idea-landau}. In
$\operatorname{Re}s>0$ the Mellin transform $F$ of
$x^\theta-\Gz(x)$ has poles only at $s=\theta$ and at the zeros of
$\zeta$ (crosses, at their true ordinates); $s=1$ is not a pole.
A real zero in $(\theta,1)$ would give $F$ a pole with negative
residue, which $F\ge0$ on the real axis forbids; so Landau's theorem
makes $F$ analytic in $\operatorname{Re}s>\theta$ (shaded), and a
zero there (open circle) cannot exist. The value $\theta=0.7$ and the
open circle are illustrative.}
\label{fig:landau-plane}
\end{figure}

Combining Proposition~\ref{prop:idea-landau} with the functional
equation: if for every $\varepsilon > 0$ one has
$\Gz(x) \le x^{1/2 + \varepsilon}$ for all large $x$, then every
nontrivial zero of $\zeta$ lies on $\operatorname{Re} s = 1/2$. So
the target is an upper bound $\Gz(x) \le x^{1/2+\varepsilon}$. In
the dictionary of Section~\ref{sec:idea-dictionary}, $\Gz$ plays the
role of the error $\#C(\F_{q^n}) - q^n - 1 = -\sum_i \alpha_i^n$ of
a point count, with the sign reversed: its exact counterpart is a
tail-smoothed sum of the numbers
$q^n + 1 - \#C(\F_{q^n}) = \sum_i\alpha_i^n$, computed in
Section~\ref{sec:route2}
(Definition~\ref{def:smoothed-sum},
Proposition~\ref{prop:which-bound}). Because of the reversal, an
upper bound for the counterpart of $\Gz$ is a
\emph{lower} bound for point counts, so the one-sided bound of
this section is not the transcription of the Stepanov--Bombieri
upper bound (Figure~\ref{fig:sign-direction}). Proposition~\ref{prop:idea-landau} plays the role of
Theorem~\ref{thm:route2prime}, and for a bound in this direction no
second bound is needed (Remark~\ref{rem:one-sided-curves}).
\begin{figure}[tbp]
\centering
\input{book/figs/ch01-sign-direction.tex}
\end{figure}

\begin{remark}[What the argument uses in place of this criterion]
\label{rem:archive-landau}
The criterion that the argument uses differs from
Proposition~\ref{prop:idea-landau} in two ways, and
Chapter~\ref{ch:landau} develops both.
\begin{enumerate}[label=(\arabic*)]
\item The twisted sum of one Dirichlet character need not be
real, and real zeros of $L(s,\chi)$ with $\chi$ real are not
excluded in advance. The argument works with a \emph{positive}
combination of $\Gz$ and the real parts of twisted sums, over a
finite set of primitive characters closed under inversion
(Definition~\ref{def:level}). With positive weights every real zero
contributes a residue of one sign. The residue argument in the
proof of Proposition~\ref{prop:idea-landau} then applies, so a
one-sided bound still suffices (Remark~\ref{rem:positivity}).
\item The inequality is needed, and is proved, only on windows
$[Y, Y^{1+\delta}]$ with $Y$ arbitrarily large. The criterion in
that form is Theorem~\ref{thm:windowed-landau}. The classical form
of Landau's theorem is not enough, because the bounds that the
argument proves hold only on windows
(Remark~\ref{rem:landau-compare}).
\end{enumerate}
\end{remark}

\section{Divisibility is at most size}\label{sec:slope}

\subsection{The brief version}

One builds a lattice of algebraic numbers attached to a theta
function on an abelian variety with complex multiplication, and maps
it to another lattice. Each image is forced to be divisible by many
primes, to a known order. Its archimedean size is also known. The
arithmetic analogue of ``zeros $=$ poles'' then gives
\[
  \underbrace{\textstyle\sum_p (\text{forced order at } p)\,\log p}
  _{\text{arithmetic side}}
  \;\le\;
  \underbrace{\log(\text{size})}_{\text{archimedean side}} .
\]
Figure~\ref{fig:outline-status} maps the steps and their status. The arithmetic side is, up to error, $b\,\Sz(a/b)$ for parameters
$a, b$. The archimedean side is computed from periods. The main
terms of the two sides cancel, by a period formula of
Chowla--Selberg type, and what remains is an upper bound for
$\Gz(a/b)$.
\begin{figure}[tbp]
\centering
\input{book/figs/ch01-outline-status.tex}
\end{figure}

\subsection{What the divisibility says about the smoothed sum}

The arithmetic side has the shape of the prime factorization of a
factorial. By Legendre's formula \cite{Legendre1808}, for an integer
$a \ge 1$,
\[
  \ord_p(a!) = \sum_{i \ge 1} \lfloor a/p^i \rfloor,
  \qquad\text{so}\qquad
  \log a! = \sum_{n \le a} \Lambda(n)\,\lfloor a/n \rfloor .
\]
Drop the floors and subtract a constant $b > 0$ from each term,
keeping only positive parts. The result is
\begin{equation}\label{eq:smoothed-factorial}
  \sum_{n \le a/b} \Lambda(n)\Bigl(\frac{a}{n} - b\Bigr)
  \;=\; b\,\Sz(a/b)
  \;=\; a\Bigl(\log\frac ab - 1 - \gamma\Bigr) + b\,\Gz(a/b),
\end{equation}
where the second equality is Definition~\ref{def:sz-gz}. So
$b\,\Sz(a/b)$ is the logarithm of a ``smoothed factorial''.
Figure~\ref{fig:smoothed-factorial} compares the factorial with its
smoothing.
\begin{figure}[tbp]
\centering
\adjustbox{max width=\linewidth}{%
\input{alpoge/figs/smoothed-factorial.tex}}
\caption{What the divisibility says about $\Sz$. (a) For $a=100$ and
$b=1$, the exponent $\ord_p(100!)$ (dots) and its smoothing
$\sum_i(100/p^i-1)_+$ (triangles); weighted by $\log p$ and summed
over $p$ they give $\log 100!$ and $\Sz(100)$. (b) With the main
terms removed ($b=1$), $\log a!$ leaves about $\frac12\log2\pi a$
(grey), while $\Sz(a)$ leaves $\Gz(a)$ (blue), which fluctuates on
the scale $\sqrt a$ of the band $\log2\pi\pm c_0\sqrt a$. The
arithmetic content is in this error.}
\label{fig:smoothed-factorial}
\end{figure}

Suppose an algebraic number is known to be divisible, at each prime
$p$, to order at least $\sum_i (a/p^i - b)_+$. Then
\eqref{eq:div-size} bounds $b\,\Sz(a/b)$ from above by the logarithm
of its size. By \eqref{eq:smoothed-factorial}, if the size is
$a(\log(a/b) - 1 - \gamma) + \mathfrak E$, the inequality becomes
\[
  \Gz(a/b) \;\le\; \mathfrak E/b .
\]
So the content of the bound is the error $\mathfrak E$. Chebyshev's
binomial coefficients are the oldest example of such numbers, but
their sizes have the wrong main term.

In the book this bookkeeping is done once and exactly. The balance
identity, Proposition~\ref{prop:balance}, relates a prime sum of the
smoothed Legendre shape to an archimedean main term and a
normalised Chebyshev function, for the characters of the Galois
group of a host field; Remark~\ref{rem:legendre} explains the
relation to Legendre's formula.

\subsection{The numbers used}

The following description is a report of the note
\cite[p.~3, step~2]{AlpogeRemark2026} and of statements in the
archive that are readable but not on its proof path
\cite[\nolinkurl{Formalism/StatementCM.lean},
\nolinkurl{Formalism/ProductFormula/StatementOCCM.lean}]
{AlpogeArchive2026}. It is not proved here. It describes the first,
abelian form of the construction. The end of the subsection
compares it with the form that Chapters~\ref{ch:balance}
and~\ref{ch:theta} develop.

Let $X$ be an elliptic curve with complex multiplication by an
imaginary quadratic field $K$, let $\vartheta$ be a theta function
on $X$, let $\partial$ be the invariant derivation (the derivation
dual to an invariant differential), and let $\Omega$ be a period.
The numbers are the normalized Taylor coefficients
$\partial^a\vartheta(0)/\Omega^a$ at the origin. They are algebraic,
and their size is about $a\log a - a - a\log\Omega$ in logarithmic
terms. At a prime $p$ inert in $K$ at which $X$ has good reduction,
the reduction is supersingular, and $\partial^p \equiv 0$ modulo
$p$. Each block of $p$ derivatives therefore contributes a factor
$p$. The note states the resulting divisibility as
\[
  p^{\,a/p + a/p^3 + \cdots - O(\log a)}
  \;\Big|\; \frac{\partial^a\vartheta(0)}{\Omega^a},
\]
an exponent of Legendre type in which only the odd powers of $p$
occur; compare Remark~\ref{rem:legendre}. A single CM curve sees
only the inert primes. The archive's abelian statements use abelian
varieties with complex multiplication by $\Q(\zeta_\ell)$ for a
prime $\ell$. A CM type of $\Q(\zeta_\ell)$ is a set of embeddings
containing one from each complex-conjugate pair. It is encoded
there as a subset $\Phi \subset (\Z/\ell)^\times$ containing exactly
one of $g$ and $-g$ for each $g$, together with a dual type. Which
primes $p$ gain is decided by the class of $p^i$ modulo $\ell$
relative to $\Phi$, and $\Phi$ together with its dual type sees
every prime except those congruent to $\pm 1$, up to an error of
about $a/\ell$. This is how the full sum $b\,\Sz(a/b)$ appears.

The archimedean main term involves $\log\Omega$. The Chowla--Selberg
formula \cite{ChowlaSelberg1967} expresses the periods of CM
elliptic curves through values of the Gamma function, equivalently
through $L'/L(0,\chi)$ for the quadratic character $\chi$ of $K$.
Colmez conjectured a generalization to CM abelian varieties, in
terms of logarithmic derivatives at $s = 0$ of Artin $L$-functions,
and proved it for abelian CM fields up to a rational multiple of
$\log 2$ \cite{Colmez1993}; Obus removed that ambiguity
\cite{Obus2013}. By the functional equation these are combinations
of $\gamma$, $\log\pi$ and $L'/L(1,\chi)$, the constants in the main
terms of $\Gz$ and of its twisted analogues. This is the mechanism
by which the main terms can cancel.

One point of these notes is unresolved. The readable statements do
not show how the parameter $b$ enters the archimedean side: the size
quoted above has main term $a\log a$, while
\eqref{eq:smoothed-factorial} has $a\log(a/b)$. The archive's
abelian estimate takes the needed relation between the periods as a
hypothesis,
$\gamma + \log(\pi/v) - \ell' \le (\log\Omega_1 + \log\Omega_2)/2$
for auxiliary parameters $v$ and $\ell'$
\cite[\nolinkurl{Formalism/ProductFormula/StatementOCCM.lean}]
{AlpogeArchive2026}.

The argument as developed in this book keeps the mechanism and
changes the cast. The CM field is an imaginary quadratic field
$K = \Q(\sqrt{-d})$ inside a \emph{host field} $E$, chosen with
auxiliary primes. The primes that gain are those inert in $K$.
Fourier analysis on $\Gal(E/\Q)$ reaches the other primes
(Chapters~\ref{ch:hosts} and~\ref{ch:balance}). The report above
suggests a sum of $L'/L(1,\chi)$ over many characters of Colmez
type in the main term, but only
$L'(1,\chi_{-d})/L(1,\chi_{-d})$ enters it, through a period of
Chowla--Selberg type (Remarks~\ref{rem:stirling}
and~\ref{rem:no-colmez}).

\subsection{Bost's slope inequality}

To use many numbers at once one replaces \eqref{eq:div-size} by an
inequality for lattices.

\begin{definition}\label{def:slope}
A \emph{Hermitian lattice} is a pair $\overline L = (L, \|\cdot\|)$,
where $L$ is a free $\Z$-module of finite rank and $\|\cdot\|$ is a
Euclidean norm on $L \otimes \R$. Let $\overline L$ be a Hermitian
lattice with $L \ne 0$.
\begin{enumerate}[label=(\roman*)]
\item The \emph{Arakelov degree} is
$\degh\,\overline L = -\log\covol(L)$, where $\covol(L)$ is the
covolume of $L$ in $L\otimes\R$.
\item The \emph{slope} is
$\mu(\overline L) = \degh\,\overline L/\rk L$.
\item The \emph{maximal slope} $\mu_{\max}(\overline L)$ is the
maximum of $\mu(\overline{L'})$ over nonzero subgroups
$L' \subseteq L$ with the restricted norm.
\end{enumerate}
\end{definition}

The maximum exists by reduction theory; see \cite{Bost1996}. For
$L = \Z$ with norm $|n|\cdot c$ the degree is $-\log c$. A nonzero
integer $\Delta$ generates a sublattice of $\Z$ (standard norm) of
degree $-\log|\Delta|$.

\begin{proposition}[Slope inequality; cf.\ \cite{Bost1996}]
\label{prop:slope}
Let $\overline L$ and $\overline M$ be Hermitian lattices with
$L \ne 0$. Let $\varphi\colon L \to M$ be an injective homomorphism,
and let $\kappa > 0$ satisfy $\|\varphi(x)\| \le \kappa\|x\|$ for
all $x \in L\otimes\R$.
\begin{enumerate}[label=(\alph*), ref=\alph*]
\item\label{item:slope} Then
$\mu(\overline L) \le \mu_{\max}(\overline M) + \log \kappa$.
\item\label{item:slope-div} Let $d \ge 1$ be an integer with
$\varphi(L) \subseteq dM$. Then
$\log d \le \mu_{\max}(\overline M) - \mu(\overline L)
+ \log\kappa$.
\end{enumerate}
\end{proposition}

\begin{proof}[Proof of \eqref{item:slope}]
Let $r = \rk L$, and give $\varphi(L)$ the norm restricted from $M$.
The map $\varphi\otimes\R$ multiplies $r$-dimensional volumes by the
product of its singular values, each at most $\kappa$. So
$\covol\varphi(L) \le \kappa^r\covol(L)$, that is,
$\degh\,\overline{\varphi(L)} \ge \degh\,\overline L - r\log\kappa$.
By definition of $\mu_{\max}$,
$\degh\,\overline{\varphi(L)} \le r\,\mu_{\max}(\overline M)$.
Combine the two inequalities and divide by $r$.
\end{proof}

\begin{proof}[Proof of \eqref{item:slope-div}]
Give $dM$ the norm restricted from $M$. Multiplication by $d$ scales
covolumes of rank-$k$ sublattices by $d^k$, so
$\mu_{\max}(\overline{dM}) = \mu_{\max}(\overline M) - \log d$.
Apply \eqref{item:slope} to $\varphi$ viewed as a map $L \to dM$.
\end{proof}

Part \eqref{item:slope-div} is \eqref{eq:div-size} for lattices. The
forced divisibility of the theta coefficients enters as $d$, with
$\log d$ about $b\,\Sz(a/b)$, and the sizes enter through
$\mu_{\max}(\overline M)$ and $\log\kappa$. The archive states the
inequality of part~\eqref{item:slope} in
\cite[\nolinkurl{Formalism/LatticeSlopes/Statement.lean}]
{AlpogeArchive2026}. Bost formulated the slope method as the
Arakelov-geometric form of the auxiliary-function method of
transcendence theory \cite{Bost1996}. That is the place to look for
an arithmetic version of ``a nonzero function has as many zeros as
poles''. Figure~\ref{fig:slope-lattice} works part
\eqref{item:slope-div} out in rank two.
\begin{figure}[tbp]
\centering
\adjustbox{max width=\linewidth}{%
\input{book/figs/slope-lattice.tex}}
\caption{Proposition~\ref{prop:slope}\eqref{item:slope-div} in rank
two. Here
$L=M=\Z^2$ with the standard norm and
$\varphi=2\left(\begin{smallmatrix}2&1\\0&1
\end{smallmatrix}\right)$, so $\varphi(L)\subseteq dM$ with $d=2$.
The covolume grows from $1$ to $8\le\kappa^2\approx20.9$, where
$\kappa=\|\varphi\|=2\sqrt{3+\sqrt5}$, and
$\log2\le\mu_{\max}(\overline M)-\mu(\overline L)+\log\kappa
=\log\kappa\approx1.52$.}
\label{fig:slope-lattice}
\end{figure}

Chapter~\ref{ch:theta} sets up Hermitian lattices over the ring of
integers of a number field (Definition~\ref{def:hlat}) and quotes
Bost's inequality in that generality (Theorem~\ref{thm:bost}). The
slope input of the argument is a comparison of an \emph{average}
slope with a slope, the frame theorem (Theorem~\ref{thm:frame}).
Bost's inequality is one step of its proof, and is not applied to
each period map. Remark~\ref{rem:not-bost} spells out the
difference. The argument does not compute the size of the target
lattice. A covolume bound for lattices of integral forms bounds it,
and this is where Arakelov geometry enters
(Theorem~\ref{thm:covolume} in Chapter~\ref{ch:covolume}).

\section{Multiplicity is the bottleneck}\label{sec:multiplicity}

The following is a report of statements in the archive that are
readable but not on its proof path
\cite[\nolinkurl{Formalism/ProductFormula/StatementOCCM.lean},
\texttt{OCCMEstimate}, \texttt{ZetaQuarterPowerCorollary}]
{AlpogeArchive2026}. It is not proved here. After the main terms
cancel, the abelian construction gives a bound of the shape
\[
  \Gz(x) \;\ll\; x^{o(1)}\Bigl(
    \frac{x}{m} + \frac{xm}{k} + \frac{x}{\ell} + k\ell \Bigr)
\]
plus terms that are negligible for suitable parameters. Here $m$ is
the number of times the forced divisibility is collected, $k$ bounds
the order of the jets, and $\ell$ is the prime with CM field
$\Q(\zeta_\ell)$. The four terms are read here as follows.

\begin{enumerate}[label=(\roman*)]
\item $x/m$: the fixed loss in the divisibility at each prime is
shared among $m$ copies. This is the term $q/m$ of
Section~\ref{sec:stepanov}.
\item $xm/k$: each copy costs archimedean size. This is the
constraint that $m$ cannot exceed about $\sqrt q$ in
Section~\ref{sec:stepanov}.
\item $x/\ell$: a CM type sees a prime only through its class
modulo $\ell$, and prime powers near $\ell$ are miscounted.
\item $k\ell$: a larger CM field costs size.
\end{enumerate}

Balancing with $m \approx x^{1/4}$, $k \approx x^{1/2}$ and
$\ell \approx x^{1/4}$ makes all four terms $x^{3/4}$, and no choice
does better, since the product of the four terms is $x^3$. By
Proposition~\ref{prop:idea-landau} such a bound would give
$\zeta(s) \ne 0$ for $\operatorname{Re} s > 3/4$. The archive's
statement of this consequence is the definition of a proposition,
not a theorem: its hypothesis is that suitable parameters and CM
data exist for all large $x$, and it is off the proof path. Its
conclusion alone, a zero-free half-plane
$\operatorname{Re} s > 3/4$, would be far beyond any known zero-free
region, so it should be read as a description of the method, not as
an established result.

To reach $1/2$, the first two terms alone require
$m \approx \sqrt x$ and $k \approx x$. The coupling $k\ell$ between
the order of the jets and the size of the CM field must then
disappear. A further off-path statement
\cite[\nolinkurl{Formalism/HalfLinks/Statement.lean}]
{AlpogeArchive2026} includes, among others, the constraints
$x^{1-\theta} \le m$ and $xm \ll x^\theta k$. At $\theta = 1/2$ they
say $m \ge \sqrt x$. This is the chain length named in the note
\cite[p.~1]{AlpogeRemark2026}. Figure~\ref{fig:exponent-budget}
carries out this balancing exactly.
\begin{figure}[tbp]
\centering
\adjustbox{max width=\linewidth}{%
\input{alpoge/figs/exponent-budget.tex}}
\caption{The exponent budget of Section~\ref{sec:multiplicity}, with
$m=x^\mu$, $k=x^\kappa$ and $\ell=x^\lambda$. (a) The smallest
possible largest exponent of the four terms is
$\max(1-\mu,(2+\mu)/3)$, computed by exact linear programming; its
minimum is $3/4$, at $\mu=1/4$. (b) Making $x/m$ and $xm/k$ at most
$x^{1/2}$ forces $\mu\ge\frac12$ and $\kappa\ge\mu+\frac12\ge1$,
which is incompatible with $k\ell\le x^{1/2}$ (hatched). The four
terms are a reading of an off-path statement of the archive.}
\label{fig:exponent-budget}
\end{figure}

In the argument as developed in this book the first two terms
survive and the last two are gone. Proposition~\ref{prop:sqrtx}
bounds the normalised error by terms of sizes $x/m$, $xm/k$ and $m$
under the constraints $xm \le k$ and
$x^{1/2-\eps} \le m \le x^{1/2+\eps/2}$, and
Remark~\ref{rem:exponent-half} explains why their balance at
$m \approx \sqrt x$ is the source of the exponent $1/2$. The
inequality that results is the Chebyshev inequality on windows,
Theorem~\ref{thm:chebyshev-windows}, which is reached through the
gap bound, Theorem~\ref{thm:gap}. Both are in
Chapter~\ref{ch:chebyshev}.

\section{Theta lifts instead of differentiating integers}
\label{sec:theta}

\subsection{Where multiplicity came from over a finite field}

In the sketch of Proposition~\ref{prop:stepanov}, multiplicity came
from one identity: $f = (\sum_i s_i v_i^{q/m})^m$ is an $m$-th
power, because $m$ is a power of the characteristic. In particular,
in a local parameter all Hasse derivatives of $f$ of orders
$1, \dots, m - 1$ vanish. There is no such operation on $\Z$. One
would want a derivation ``in the direction of $p$'' whose repeated
vanishing at $\Delta$ means divisibility of $\Delta$ by a high power
of $p$. Buium's $p$-derivations, such as
$\delta_p(n) = (n - n^p)/p$, are one formal answer. Alp\"oge's
argument does not use them. The Lean archive contains
$\delta$-ring structures only in $p$-adic Hodge theory, far from
this step.

Stepanov's original argument did use derivatives, and that form of
the method, with a chain of about $\sqrt p$ derivative conditions,
is the closer model for what follows. It is proved in the simplest
case in Section~\ref{sec:variantJ}, and the comparison is made
there.

\subsection{What a theta lift is}

For theta correspondences see \cite{Howe1989,Rallis1984}, and for
algebraic modular forms see \cite{Gross1999}. Let $K$ be a CM field.
Let $V$ be a Hermitian space over $K$ that is \emph{totally
definite}, meaning positive definite at every infinite place. Then
the unitary group $U(V)$ is compact at infinity, and an automorphic
form on $U(V)$ of a given level is a function on a \emph{finite}
set: the classes of lattices in a genus. These are Gross's algebraic
modular forms, the higher-dimensional analogue of functions on the
finite set of ideal classes of a definite quaternion algebra. A
theta correspondence pairs $U(V)$ with $U(W)$ for a skew-Hermitian
space $W$ (a dual pair inside a symplectic group), and sends
automorphic forms on one group to the other by integration against
a theta kernel.

This section does not claim that the global theta lift between two
definite groups is a matrix of representation numbers. That is not
the standard description. What follows is the archive's finite-form
version, read from
\cite[\nolinkurl{Formalism/ThetaLiftIndependence/Statement.lean}]
{AlpogeArchive2026}.
\begin{enumerate}[label=(\roman*)]
\item The lift sends a function $f$ on class representatives $L$ to
$\sum_L w(L)\,f(L)\,\theta_L(\tau)$, a weighted sum of theta series
of the lattices, with weights $w(L) = 1/n_L$ for positive integers
$n_L$ and a common denominator.
\item A chain composes $m$ kernel operators
$f \mapsto \bigl(j \mapsto \sum_i w_i f(i) K(i,j)\bigr)$ between
finite class sets. If the scaled weights and the kernel entries are
integral, the composite times the product of the denominators is
integral.
\item There is an adjunction identity of seesaw type for such
kernels.
\item If a linear map $\theta$ satisfies
$\langle\theta x, \theta y\rangle = c\cdot\langle x, y\rangle$ on a
subspace $S$ stable under $\theta^*\theta$, with $c$ a product of an
$L$-value and local factors as in the Rallis inner product formula,
then $\theta^*\theta = c$ on $S$. So the operator norm of $\theta$
on $S$ is $\sqrt{|c|}$, and a norm identity for individual forms
becomes an operator-norm bound.
\end{enumerate}
How the kernels $K$ of \cite{AlpogeArchive2026} are built from theta
correspondences is not described here.

\subsection{Why lifts substitute for derivatives}
\label{subsec:lifts-derivatives}

The following summarizes the note
\cite[pp.~3--4, steps 8--12]{AlpogeRemark2026} together with the
statements above. The connecting argument is an interpretation; it
is not proved here.

\begin{enumerate}[label=(\roman*)]
\item \emph{Integral and linear.} After clearing denominators, each
kernel operator in the chain is an integral linear map of Hermitian
lattices. So Proposition~\ref{prop:slope} applies to each lift and
to their composite, as it applied to the jets of a theta function.
\item \emph{Each lift collects the divisibility again.} Iterate
lifts $U(X_0) \to U(X_1) \to \cdots \to U(X_m)$. According to the
note, by seesaw the period of the iterated lift is a Bessel period
independent of the Hecke characters, a period of the same CM kind
as in Section~\ref{sec:slope}. So the forced Legendre-type
divisibility at each prime is collected once per lift, $m$ times in
all. This is the analogue of vanishing to order $m$.
\item \emph{Each lift has a computable cost.} The Rallis formula,
the Siegel mass formula (the total weight of a class set) and
Weyl's dimension formula control the norms and ranks along the
chain. These give the archimedean side of
Proposition~\ref{prop:slope}. Unlike $m$-fold differentiation, the
cost does not couple to the size of the CM field.
\item \emph{The length of the chain.} Balancing $x/m$ against
$xm/k$ as in Section~\ref{sec:multiplicity} gives
$m \approx \sqrt x$, as $m \approx \sqrt q$ in
Section~\ref{sec:stepanov}.
\end{enumerate}
Figure~\ref{fig:theta-chain} draws the chain.
\begin{figure}[tbp]
\centering
\adjustbox{max width=\linewidth}{%
\input{book/figs/theta-chain.tex}}
\caption{Where multiplicity comes from (schematic). Over $\F_q$ one
Frobenius power gives order $m$ at every point. Over $\Z$, a chain
of $m$ integral kernel operators between finite class sets collects
the forced divisibility once per lift, at a cost controlled by the
Rallis formula, the Siegel mass and Weyl's dimension formula.
Balancing gives $m\approx\sqrt x$. The picture follows the note
\cite{AlpogeRemark2026} and the finite-form statements of
\cite{AlpogeArchive2026}, with the interpretation of
Section~\ref{subsec:lifts-derivatives}; class-set sizes are
illustrative.}
\label{fig:theta-chain}
\end{figure}

The analogy is summarized in the following table.
\begin{center}
\begin{tabular}{@{}ll@{}}
\toprule
Curve over $\F_q$ & Integers \\
\midrule
$f = g^m$ with $m = p^\mu$ & composite of $m$ theta lifts \\
derivations of order $< m$ kill $f$
  & divisibility gained at each lift \\
$\deg f \le mA + qB$ & Rallis, Siegel mass and Weyl bounds \\
$m \approx \sqrt q$ & chain length $m \approx \sqrt x$ \\
\bottomrule
\end{tabular}
\end{center}

Chapter~\ref{ch:theta} replaces this interpretation by definitions.
The chain is the theta tower of Definition~\ref{def:spaces}
(Figure~\ref{fig:tower}). The periods are measured in the Hermitian
lattices of Definitions~\ref{def:hlat}
and~\ref{def:period-lattice}. The conditions that the lattices
along the chain must satisfy, the
divisibility at inert primes and the archimedean costs among them,
are collected in a slope record (Definitions~\ref{def:record}
and~\ref{def:size-data}). Remark~\ref{rem:factorial} states the
comparison with Stepanov's method in the terms of that chapter: the
role of vanishing to high order is played by $p$-adic divisibility
of lifted periods at all inert $p$ at once.

\section{What the outline omits, and where to read on}
\label{sec:idea-roadmap}

The outline above leaves out everything that makes the bounds
uniform, and most of the book is about those costs. The note lists
them \cite[steps 11, 14--19]{AlpogeRemark2026}. Twisting over a
$2$-power cyclotomic tower kills unwanted contributions. A weakened
form of Vandiver's conjecture bounds the classes in
$H^1(\mathcal O_F[1/\ell], \Q_\ell/\Z_\ell(-e))$, for certain CM
fields $F$, that twisting cannot kill. The archive claims a proof
of the latter, from Galois representations attached to automorphic
forms. In the book these become the following chapters.

\begin{enumerate}[label=(\arabic*)]
\item \emph{Hosts.} To isolate a given primitive character $\psi$
one places it among the characters of the Galois group of a host
field, and one needs such fields, with the right splitting
behaviour, on windows. This is Chapter~\ref{ch:hosts}. The Fourier
analysis that extracts a positive combination containing $\psi$ is
Chapter~\ref{ch:balance}.
\item \emph{Congruence defects.} The lattices of the chain differ
from the ideal ones by congruence defects, which are bounded through
Selmer groups by Tate--Shafarevich groups of Tate twists. The bound
is uniform only on windows (Theorem~\ref{thm:sha-windows} in
Chapter~\ref{ch:sha}; Remark~\ref{rem:vandiver} discusses the
relation to Vandiver's conjecture). This is the reason the
criterion of Section~\ref{sec:idea-landau} is needed on windows.
\item \emph{Galois representations.} The cases of the
Tate--Shafarevich bound that need automorphy use Galois
representations and a truncated local--global compatibility
statement (Chapter~\ref{ch:galois}).
\item \emph{Existence.} The families of arithmetic data to which
the Chebyshev inequality applies have to be shown to exist
(Chapter~\ref{ch:families}).
\end{enumerate}
Chapter~\ref{ch:assembly} states the main theorem,
Theorem~\ref{thm:main}, and gives its whole proof with each of
these tools as a black box. Figure~\ref{fig:dag} there displays the
dependencies. Before that, Chapter~\ref{ch:function-fields} carries
out the finite-field model completely, since everything in this
chapter was phrased by analogy with it.

\begin{remark}[Status]\label{rem:status}
This chapter does not verify Alp\"oge's argument, and the Lean
archive was read as text and not built. Whether it proves the
Riemann Hypothesis is not established here. The readable statements
cited in Sections~\ref{sec:slope}--\ref{sec:theta} are not on the
archive's proof path. They show what the authors had in mind, and
the on-path argument may differ. Where the later chapters of this
book differ from them, this chapter says so. The account of the
error terms and of the role of the theta lifts is an
interpretation.
\end{remark}
\clearpage
\clearpage
\newcommand{\fffloat}[2][tbp]{\begin{figure}[#1]\centering
  \let\fftikz\tikzpicture \let\endfftikz\endtikzpicture
  \renewenvironment{tikzpicture}%
    {\begin{adjustbox}{max width=\linewidth}\fftikz}%
    {\endfftikz\end{adjustbox}}%
  \input{stepanov-bombieri/figs/#2}\end{figure}}
\newcommand{\ffbookfloat}[2][tbp]{\begin{figure}[#1]\centering
  \let\fftikz\tikzpicture \let\endfftikz\endtikzpicture
  \renewenvironment{tikzpicture}%
    {\begin{adjustbox}{max width=\linewidth}\fftikz}%
    {\endfftikz\end{adjustbox}}%
  \input{book/figs/#2}\end{figure}}
\newcommand{\ffabregion}{\begingroup
  \let\ffref\ref
  \def\ffkey{def:admissible}%
  \protected\def\ref##1{\def\fftmp{##1}%
    \ifx\fftmp\ffkey\ffref{def:ff-admissible}\else\ffref{##1}\fi}%
  \fffloat{ab-region}\endgroup}
\chapter[The function-field model]{The function-field model:
  Stepanov--Bombieri and the sign-only proof}
\label{ch:function-fields}

This chapter treats the two arguments for the Riemann hypothesis
for curves over finite fields that Alp\"oge's note sets side by
side \cite{AlpogeRemark2026}. The first is a complete proof. It is
the method of Stepanov
\cite{Stepanov1969} in the form given by Bombieri
\cite{Bombieri1974}: an auxiliary function, forced to vanish to high
order at the points to be counted, and a comparison of zeros with
poles. It uses the Riemann--Roch theorem and nothing deeper, and it
is given here in full, from the definition of the zeta function to
the Riemann hypothesis
(Sections~\ref{sec:setting}--\ref{sec:rh}). The second is the
``sign-only'' proof, which the note records in one footnote as the
identity $qN = M^k + R$. It uses no auxiliary function and no
multiplicity; it uses elementary Fourier analysis and the parity of
a cohomological degree. It is reconstructed here as a complete
deduction from two cohomological hypotheses, which are stated and
not proved (Sections~\ref{sec:identity}--\ref{sec:sign}). It is
therefore a conditional derivation, and not a second proof.

Nothing in this chapter is a logical input to the later chapters.
Its output is the analogy on which they lean. Chapter~\ref{ch:idea}
was phrased by analogy with the first proof, and
Chapter~\ref{ch:assembly} assembles an argument whose parts were
found by transporting the two arguments to the integers. The dictionary
of Section~\ref{sec:dictionary} matches each step proved here with
the part of Chapters~\ref{ch:assembly}--\ref{ch:families} that
plays its role, and says where the match is exact, where it is
loose, and where it fails.

Four results of the chapter are negative.
\begin{enumerate}
\item The upper bound for point counts does not by itself give the
Riemann hypothesis through a positivity argument of
Landau--Pringsheim type. The missing input is identified exactly,
and it cannot be supplied formally (Section~\ref{sec:route2}).
\item The one-sided bound on a smoothed prime sum that the
number-field argument proves corresponds, over a function field, to
a \emph{lower} bound on point counts. It does not correspond to the
Stepanov--Bombieri upper bound (Proposition~\ref{prop:which-bound}).
\item The method forces the square root in the Stepanov--Bombieri
error. A different choice of parameters does not remove it
(Proposition~\ref{prop:sqrtq-forced}).
\item The second hypothesis of the sign-only proof fails for the
quadric as the note writes it, for some values of $k$ in the
simplest cases (Remark~\ref{rem:h2-fails}). So the sign-only proof
is conditional, and the first proof is not.
\end{enumerate}

\section{Setting}
\label{sec:setting}

Let $p$ be a prime and $q=p^e$. Sections~\ref{sec:setting}
to~\ref{sec:variantJ} are written in the language of function fields
and valuations, in which the proof needs no geometry beyond the
Riemann--Roch theorem; the translation to curves follows the
definition.

\begin{definition}\label{def:function-field}
\leavevmode
\begin{enumerate}
\item A \emph{function field} over $\F_q$ is a finitely generated
extension $K/\F_q$ of transcendence degree one in which $\F_q$ is
algebraically closed. We also say that $K$ has \emph{full constant
field} $\F_q$.
\item A \emph{place} of $K$ is the maximal ideal $P$ of a discrete
valuation ring $\calO_P$ with $\F_q\subset\calO_P\subset K$ and
fraction field $K$; $v_P$ is its valuation,
$\kappa(P)=\calO_P/P$, $\deg P=[\kappa(P):\F_q]$, and $z(P)$ is the
class of $z\in\calO_P$ in $\kappa(P)$.
\item A \emph{divisor} is an element of the free abelian group on
the places; $(z)$ is the divisor of $z\in K^\times$, and
$(z)_0$, $(z)_\infty$ are its divisors of zeros and of poles. For a
divisor $D$,
$\calL(D)=\{z\in K^\times:(z)+D\ge0\}\cup\{0\}$ and
$\ell(D)=\dim_{\F_q}\calL(D)$.
\item $B_d$ is the number of places of degree $d$, and
$N_n=\sum_{d\mid n}dB_d$.
\end{enumerate}
\end{definition}

Geometrically $K=\F_q(C)$ for a smooth projective geometrically
irreducible curve $C/\F_q$ of genus $g$, the places are the closed
points of $C$, $\calL(D)=H^0(C,\calO_C(D))$ is the space of sections
of a line bundle, and $N_n=\#C(\F_{q^n})$; see
\cite[Theorem~16.3.3]{Vakil2025} for
the equivalence between curves and function fields and
\cite[App.~B]{Stichtenoth2009} for the comparison of the two
languages. Figure~\ref{fig:places-divisors} illustrates the
definition. The notation $\calL(D)$, $\ell(D)$ is that of
Chapter~\ref{ch:idea}.
\fffloat{places-divisors}

The theorem to be proved is the following.

\begin{theorem}[Weil]\label{thm:weil}
Let $K/\F_q$ be a function field of genus $g$. There are complex
numbers $\alpha_1,\dots,\alpha_{2g}$ such that
\begin{enumerate}
\item\label{item:weil-count} $N_n=q^n+1-\sum_{i=1}^{2g}\alpha_i^n$
for all $n\ge1$;
\item\label{item:weil-fe} $\alpha\mapsto q/\alpha$ permutes the
multiset $\{\alpha_i\}$;
\item\label{item:weil-rh} $\lvert\alpha_i\rvert=\sqrt q$ for all
$i$.
\end{enumerate}
In particular $\lvert N_n-q^n-1\rvert\le2g\,q^{n/2}$ for all
$n\ge1$.
\end{theorem}

Parts \eqref{item:weil-count} and \eqref{item:weil-fe} are
Theorem~\ref{thm:rationality}; part \eqref{item:weil-rh} and the
bound are Theorem~\ref{thm:rh}. Figure~\ref{fig:proof-map} is a map
of the proof and of the two routes from the upper bound to the
Riemann hypothesis, only one of which arrives.
\fffloat{proof-map}

\subsection{Standard facts}

The proof is complete relative to the following facts. Those that
are algebraic geometry of curves are in \cite{Vakil2025}, with
locators in the numbering of its public version of October 21,
2025; the locators into \cite{Stichtenoth2009}, where the same facts
are proved in the language of valuations, are at the level of
sections.
\begin{enumerate}[label=(S\arabic*),ref=S\arabic*,itemsep=0pt]
\item\label{S1} Weak approximation: for distinct places
$P_1,\dots,P_s$, elements $z_1,\dots,z_s\in K$ and an integer $n$
there is $z\in K$ with $v_{P_j}(z-z_j)\ge n$ for all $j$
\cite[\S1.3]{Stichtenoth2009}.
\item\label{S2} $\deg(z)_0=\deg(z)_\infty=[K:\F_q(z)]$ for
$z\notin\F_q$ \cite[Exercises~16.3.E and~18.4.E]{Vakil2025},
\cite[\S1.4]{Stichtenoth2009}; so a
principal divisor has degree $0$, an effective divisor has
nonnegative degree, a function without zeros or poles is constant,
and $\calL(D)=0$ if $\deg D<0$.
\item\label{S3} Riemann--Roch:
$\ell(D)=\deg D+1-g+\ell(W-D)$ for a canonical divisor $W$, with
$\deg W=2g-2$; so $\ell(D)\ge\deg D+1-g$, with equality if
$\deg D\ge2g-1$
\cite[Exercise~18.4.B, \S18.5.3, Exercise~18.5.A]{Vakil2025},
\cite[\S1.5]{Stichtenoth2009}.
\item\label{S4} For $r\ge1$ the compositum $K_r=K\F_{q^r}$ is a
function field with full constant field $\F_{q^r}$, and over a place
of $K$ of degree $m$ lie $\gcd(m,r)$ places of $K_r$, each of degree
$m/\gcd(m,r)$ over $\F_{q^r}$; so $K_r$ has $N_r$ places of degree
one \cite[\S3.6, \S5.1]{Stichtenoth2009}.
\item\label{S5} $\F_q(x)$ has $q+1$ places of degree one and
finitely many of each degree \cite[\S1.2]{Stichtenoth2009}.
\item\label{S6} $K/\F_q(x)$ is finite and separable for some
$x\in K$ \cite[Theorem~10.5.17]{Vakil2025},
\cite[\S3.10]{Stichtenoth2009}.
\item\label{S7} In a finite extension $E/F$ of function fields each
place of $F$ has between $1$ and $[E:F]$ places above it, and
$\Gal(E/F)$ permutes them transitively if $E/F$ is Galois
\cite[\S3.1, \S3.7]{Stichtenoth2009}.
\item\label{S8} $\Gal(\F_{q^f}/\F_q)$ is cyclic of order $f$,
generated by $\Fr_q\colon c\mapsto c^q$; the group
$\F_{q^f}^\times$ is cyclic; and $\F_q=\{c:c^q=c\}$.
\item\label{S9} If $E/F$ is a finite Galois extension of fields and
$F'\supseteq F$, then $EF'/F'$ is Galois and restriction embeds
$\Gal(EF'/F')$ in $\Gal(E/F)$.
\end{enumerate}
Facts \ref{S2}, \ref{S3}, \ref{S5}, \ref{S6} and the first half of
\ref{S7} are statements about curves, line bundles and finite
morphisms of the kind assumed throughout this book. Weak
approximation, the count of places in a constant field extension,
the transitivity in \ref{S7}, and \ref{S8}, \ref{S9} are commutative
algebra and field theory, which is why the references to
\cite{Stichtenoth2009} are kept.

One more fact about valuation rings is used, and it has a
three-line proof. If $\calO\subseteq\calO'\subsetneq K$ are
valuation rings of $K$ containing $\F_q$ and $\calO$ is discrete,
then $\calO=\calO'$. Indeed, let $v$ be the valuation of $\calO$ and
$t$ a uniformizer. If some $z\in\calO'$ had $v(z)=-m<0$, then
$zt^m$ is a unit of $\calO$, so
$t^{-m}=z(zt^m)^{-1}\in\calO'$; hence $t^{-1}=t^{m-1}t^{-m}$ lies in
$\calO'$, and then $\calO'\supseteq\calO[t^{-1}]=K$.

\subsection{Frobenius and automorphisms}

Characteristic $p$ enters the whole proof through the next
proposition.

\begin{proposition}\label{prop:frobenius}
Let $K/\F_q$ be a function field, $j\ge0$ an integer and
$\pi=p^j$.
\begin{enumerate}
\item\label{item:frob-additive} $z\mapsto z^\pi$ is an injective
ring endomorphism of $K$.
\item\label{item:frob-valuation} $v_P(z^\pi)=\pi\,v_P(z)$ for every
place $P$ and $z\in K^\times$.
\item\label{item:frob-rational} If $\deg P=1$ and $z\in\calO_P$,
then $z^q-z\in P$.
\end{enumerate}
\end{proposition}

\begin{proof}
Proof of \eqref{item:frob-additive}: the map is multiplicative, and
it is additive because $\binom pi\equiv0\pmod p$ for $0<i<p$; a
ring endomorphism of a field is injective. Proof of
\eqref{item:frob-valuation}: valuations are multiplicative. Proof
of \eqref{item:frob-rational}: $z(P)\in\kappa(P)=\F_q$, and
$c^q=c$ for $c\in\F_q$ by \ref{S8}.
\end{proof}

Let $\Aut(K/\F_q)$ be the group of field automorphisms of $K$ that
fix $\F_q$. For $\sigma\in\Aut(K/\F_q)$ and a place $P$, the image
$\sigma(P)$ is a place, with $\calO_{\sigma(P)}=\sigma(\calO_P)$,
$v_{\sigma(P)}(z)=v_P(\sigma^{-1}z)$ and $\deg\sigma(P)=\deg P$; and
$\sigma(\calL(D))=\calL(\sigma D)$ for every divisor $D$. All of
this is transport of structure along $\sigma$.

\section{The zeta function}
\label{sec:zeta}

This section proves parts \eqref{item:weil-count} and
\eqref{item:weil-fe} of Theorem~\ref{thm:weil}. The argument is
F.~K.~Schmidt's \cite{FKSchmidt1931}; see also
\cite[\S5.1]{Stichtenoth2009}. All identities are between formal
power series or rational functions of $t$.

\begin{proposition}\label{prop:finiteness}
Let $K/\F_q$ be a function field of genus $g$. For each $n\ge0$ the
number $A_n$ of effective divisors of degree $n$ is finite, and the
number $h$ of divisor classes of degree $0$ is finite.
\end{proposition}

\begin{proof}
Fix $x\in K\setminus\F_q$. A place $P$ of $K$ lies over a place of
$\F_q(x)$ of degree at most $\deg P$, so $B_d<\infty$ by \ref{S5}
and \ref{S7}, and hence $A_n<\infty$. Let $D_0$ be a divisor of
degree at least $g$. For $D$ of degree $0$, \ref{S3} gives
$\ell(D+D_0)\ge1$, so $D+D_0$ is equivalent to an effective divisor
of degree $\deg D_0$. Distinct classes give distinct effective
divisors, so $h\le A_{\deg D_0}$.
\end{proof}

\begin{definition}\label{def:zeta}
Let $K/\F_q$ be a function field. Its \emph{zeta function} is
\[
Z(t)=Z_K(t)=\sum_{n\ge0}A_nt^n\in\Z[[t]],
\]
where $A_n$ is the number of effective divisors of $K$ of degree
$n$; and $\delta\ge1$ denotes the positive generator of the image of
the degree map on divisors.
\end{definition}

Unique factorisation of effective divisors into places gives the
Euler product $Z(t)=\prod_P(1-t^{\deg P})^{-1}$, and applying
$t\frac{d}{dt}\log$,
\begin{equation}\label{eq:log-deriv}
t\,\frac{Z'(t)}{Z(t)}=\sum_P\deg P\sum_{j\ge1}t^{j\deg P}
=\sum_{n\ge1}N_nt^n.
\end{equation}
So the point counts $N_n$ are the coefficients of the logarithmic
derivative of $Z$, exactly as $\Lambda(n)$ gives the coefficients of
$-\zeta'/\zeta$.

\begin{proposition}\label{prop:class-count}
Let $K/\F_q$ be a function field and $c$ a divisor class. Then $c$
contains exactly $(q^{\ell(c)}-1)/(q-1)$ effective divisors.
\end{proposition}

\begin{proof}
Fix $D\in c$. The map $f\mapsto D+(f)$ from
$\calL(D)\setminus\{0\}$ to the effective divisors equivalent to $D$
is surjective, and $(f)=(f')$ exactly when $f/f'\in\F_q^\times$, by
\ref{S2}.
\end{proof}

\begin{proposition}\label{prop:rationality-delta}\label{def:n0-F}
Let $K/\F_q$ be a function field of genus $g$ with $h$ divisor
classes of degree $0$, let $\delta$ be the positive generator of the
image of the degree map on divisors, let $n_0$ be the least multiple
of $\delta$ with $n_0\ge\max(0,2g-1)$, and let
$F(t)=\sum_{0\le\deg c\le2g-2}q^{\ell(c)}t^{\deg c}$, the sum
running over divisor classes $c$ of $K$. Then
\begin{equation}\label{eq:rationality}
(q-1)Z(t)=F(t)+h\,q^{1-g}\frac{(qt)^{n_0}}{1-(qt)^\delta}
-\frac{h}{1-t^\delta}.
\end{equation}
In particular $Z$ is a rational function of $t^\delta$ with a simple
pole at $t=1$.
\end{proposition}

\begin{proof}
There are $h$ classes of each degree in $\delta\Z$ and none of other
degrees. By Proposition~\ref{prop:class-count},
$(q-1)Z(t)=\sum_{\deg c\ge0}(q^{\ell(c)}-1)t^{\deg c}$. For
$\deg c\ge2g-1$, \ref{S3} gives $\ell(c)=\deg c+1-g$. Split the sum
at degree $2g-2$ and sum the two geometric series that remain: the
terms $q^{\ell(c)}t^{\deg c}$ with $\deg c\ge n_0$ give the middle
term of \eqref{eq:rationality}, and the terms $-t^{\deg c}$ over all
classes of degree at least $0$ give the last. The polynomial $F$ is
a polynomial in $t^\delta$; the middle term is regular at $t=1$
because $q^\delta\ne1$; the last term has a simple pole there.
\end{proof}

\begin{proposition}[F.~K.~Schmidt]\label{prop:delta-one}
Let $K/\F_q$ be a function field. Then $\delta=1$: there is a
divisor of degree one.
\end{proposition}

\begin{proof}
The plan is to compare the order of the pole at $t=1$ of both sides
of a product formula. Let $r\ge1$ and let $Z_r$ be the zeta function
of $K_r/\F_{q^r}$ (fact \ref{S4}). For a place $P$ of $K$ of degree
$m$ put $d=\gcd(m,r)$. By \ref{S4} the factors of
$Z_r(t^r)^{-1}$ at the places above $P$ multiply to
$(1-t^{rm/d})^d$. Also
$\prod_{\zeta^r=1}(1-(\zeta t)^m)=(1-t^{mr/d})^d$, because
$\zeta^m$ runs $d$ times over the $(r/d)$-th roots of unity. Hence
\begin{equation}\label{eq:zeta-tower}
Z_r(t^r)=\prod_{\zeta^r=1}Z(\zeta t).
\end{equation}
Take $r=\delta$. As $Z$ is a function of $t^\delta$, the right side
is $Z(t)^\delta$, which has a pole of order $\delta$ at $t=1$. By
Proposition~\ref{prop:rationality-delta} for $K_\delta$,
$Z_\delta(u)$ has a simple pole at $u=1$, so $Z_\delta(t^\delta)$
has a simple pole at $t=1$. Hence $\delta=1$.
\end{proof}

\begin{theorem}[F.~K.~Schmidt: rationality and functional equation]
\label{thm:rationality}
Let $K/\F_q$ be a function field of genus $g$ with $h$ divisor
classes of degree $0$.
\begin{enumerate}
\item\label{item:rat-form} $Z(t)=L(t)/((1-t)(1-qt))$ with
$L=L_K\in\Z[t]$ of degree $2g$, $L(0)=1$, leading coefficient $q^g$
and $L(1)=h$.
\item\label{item:rat-fe} $Z(1/(qt))=q^{1-g}t^{2-2g}Z(t)$,
equivalently $L(t)=q^gt^{2g}L(1/(qt))$.
\item\label{item:rat-roots} The map $\alpha\mapsto q/\alpha$
permutes the multiset $\{\alpha_1,\dots,\alpha_{2g}\}$ of
reciprocals of the roots of $L$, no $\alpha_i$ equals $1$ or
$q$, and $N_n=q^n+1-\sum_i\alpha_i^n$ for all $n\ge1$.
\end{enumerate}
\end{theorem}

The numbers $\alpha_i$, for which
$L(t)=\prod_{i=1}^{2g}(1-\alpha_it)$, are the \emph{reciprocal
roots} of $K$, and $L_K$ is its \emph{$L$-polynomial}.

\begin{proof}
Since $\delta=1$ (Proposition~\ref{prop:delta-one}),
Proposition~\ref{prop:rationality-delta} gives
\eqref{eq:rationality} with $n_0=\max(0,2g-1)$ and
$F(t)=\sum_{0\le\deg c\le2g-2}q^{\ell(c)}t^{\deg c}$. If $g=0$,
then $F=0$, $n_0=0$, and
$(q-1)Z=\frac{hq}{1-qt}-\frac h{1-t}=\frac{h(q-1)}{(1-t)(1-qt)}$;
since $Z(0)=1$, $h=1$ and $L=1$.

Let $g\ge1$, so $n_0=2g-1$. Write $(q-1)Z=F+R$ with
$R(t)=hq^{1-g}\frac{(qt)^{2g-1}}{1-qt}-\frac h{1-t}$.

\emph{Proof of \eqref{item:rat-form}.} We have
\[
(q-1)(1-t)(1-qt)Z(t)=(1-t)(1-qt)F(t)+hq^gt^{2g-1}(1-t)-h(1-qt),
\]
a polynomial of degree at most $2g$ because $\deg F\le2g-2$. So
$L:=(1-t)(1-qt)Z$ is a polynomial of degree at most $2g$ lying in
$\Z[[t]]$, hence in $\Z[t]$, and $L(0)=Z(0)=1$. At $t=1$ the
identity gives $(q-1)L(1)=h(q-1)$.

\emph{Proof of \eqref{item:rat-fe}.} The map $c\mapsto W-c$ is an
involution of the classes of degree in $[0,2g-2]$, and
$\ell(W-c)=\ell(c)-\deg c-1+g$ by \ref{S3}. Hence
\begin{align*}
F\Bigl(\frac1{qt}\Bigr)
&=\sum_cq^{\ell(W-c)}(qt)^{-(2g-2-\deg c)}\\
&=\sum_cq^{\ell(c)-\deg c-1+g}(qt)^{\deg c-2g+2}
=q^{1-g}t^{2-2g}F(t).
\end{align*}
For $R$ a direct computation gives
\[
R\Bigl(\frac1{qt}\Bigr)
=-\frac{hq^{1-g}t^{2-2g}}{1-t}+\frac{hqt}{1-qt}
=q^{1-g}t^{2-2g}R(t).
\]
This proves \eqref{item:rat-fe} for $Z$. For $L$,
$L(1/(qt))=\frac{(qt-1)(t-1)}{qt^2}\,q^{1-g}t^{2-2g}Z(t)
=q^{-g}t^{-2g}L(t)$.

\emph{Proof of \eqref{item:rat-roots}.} With $L=\sum_ja_jt^j$ the
functional equation reads $a_{2g-j}=q^{g-j}a_j$. As $a_0=1$, we get
$a_{2g}=q^g$; so $\deg L=2g$ and $\prod_i\alpha_i=q^g$. Then
\[
L(t)=q^gt^{2g}\prod_i\Bigl(1-\frac{\alpha_i}{qt}\Bigr)
=q^g\prod_i\Bigl(-\frac{\alpha_i}q\Bigr)
\prod_i\Bigl(1-\frac q{\alpha_i}t\Bigr)
=\prod_i\Bigl(1-\frac q{\alpha_i}t\Bigr),
\]
so $\{q/\alpha_i\}=\{\alpha_i\}$. Since $L(1)=h\ge1$, no $\alpha_i$
equals $1$, and hence none equals $q$. Finally \eqref{eq:log-deriv}
and \eqref{item:rat-form} give
$\sum_nN_nt^n=\frac t{1-t}+\frac{qt}{1-qt}
-\sum_i\frac{\alpha_it}{1-\alpha_it}$; compare coefficients.
\end{proof}

Figure~\ref{fig:zeta-plane} shows the zeros and poles of $Z$ and the
symmetry of the functional equation.
\ffbookfloat{zeta-plane}

\subsection{The tower of constant field extensions}

The tower $K_r=K\F_{q^r}$ is used below to pass from one point count
to all of them. Fact~\ref{S4} describes how places split in $K_r$,
but it says nothing about the genus, and the Riemann--Roch theorem
for $K_r$ involves the genus of $K_r$, not that of $K$.
Geometrically the genus is the dimension of $H^1(C,\calO_C)$, which
does not change under extension of the base field
\cite[Exercise~18.2.H]{Vakil2025}. That argument uses cohomology,
which this proof otherwise avoids. The next proposition derives the
invariance from the zeta function instead.

\begin{proposition}\label{prop:tower}
Let $K/\F_q$ be a function field of genus $g$ with reciprocal roots
$\alpha_1,\dots,\alpha_{2g}$, and let $r\ge1$. Then
$L_{K_r}(t)=\prod_i(1-\alpha_i^rt)$. In particular $K_r$ has genus
$g$.
\end{proposition}

\begin{proof}
Let $g_r$ be the genus of $K_r$. By \eqref{eq:zeta-tower} and
Theorem~\ref{thm:rationality}\eqref{item:rat-form} for $K$ and
$K_r$,
$L_{K_r}(t^r)=\prod_{\zeta^r=1}L_K(\zeta t)
=\prod_i(1-\alpha_i^rt^r)$,
using $\prod_\zeta(1-\zeta t)=1-t^r$ and
$\prod_\zeta(1-q\zeta t)=1-q^rt^r$. Comparing degrees in $t$ gives
$2g_rr=2gr$.
\end{proof}

\begin{remark}\label{rem:genus}
The invariance of the genus under constant field extension need not
be assumed alongside fact~\ref{S4}. Proposition~\ref{prop:tower}
derives it from the zeta function. The degree of $L_{K_r}$ is twice
the genus of $K_r$ by Theorem~\ref{thm:rationality} applied to
$K_r$, and the product formula \eqref{eq:zeta-tower} fixes that
degree as $2g$. The proof uses only \ref{S4} and Riemann--Roch for
each field separately.
\end{remark}

\section{The upper bound}
\label{sec:upper}

This section proves the Stepanov--Bombieri upper bound for point
counts. The proof has four named steps, shown in
Figure~\ref{fig:four-steps}: a Siegel lemma, a non-vanishing
statement, forced multiplicity, and a slope inequality. The names
are chosen so that the steps can be compared with the number-field
argument in Section~\ref{sec:dictionary}. The bound is proved in a
form twisted by an automorphism $\sigma$, which costs nothing and is
what Section~\ref{sec:route1} needs.
\fffloat{four-steps}

\begin{definition}\label{def:frob-place}
Let $K/\F_q$ be a function field and $\sigma\in\Aut(K/\F_q)$.
\begin{enumerate}
\item A place $P$ is \emph{$\sigma$-Frobenius} if
$\sigma(z)-z^q\in P$ for all $z\in\calO_P$.
\item $N(K,\sigma)=\sum\deg P$, the sum over the
$\sigma$-Frobenius places $P$.
\end{enumerate}
\end{definition}

By Proposition~\ref{prop:frobenius}\eqref{item:frob-rational} and
\ref{S8}, the $\mathrm{id}$-Frobenius places are those of degree
one, so $N(K,\mathrm{id})=N_1$. A $\sigma$-Frobenius place has
$\sigma(\calO_P)\subseteq\calO_P$, hence $\sigma(\calO_P)=\calO_P$
by the maximality of discrete valuation rings proved in
Section~\ref{sec:setting}, and therefore $\sigma(P)=P$.
Geometrically, the $\sigma$-Frobenius places are the points of $C$
at which the $q$-Frobenius agrees with the automorphism induced by
$\sigma$: the rational points of a twist of $C$.

\subsection{An adapted basis and independence}

In the chain $\calL(jP_0)$, $0\le j\le a$, each step adds at most
one dimension, so the pole orders at $P_0$ of the elements of
$\calL(aP_0)$ form a set of exactly $\ell(aP_0)$ integers. Some
basis has pairwise distinct pole orders. Not every basis does
(Remark~\ref{rem:basis-cex}).

\begin{proposition}\label{prop:adapted-basis}
Let $K/\F_q$ be a function field, $P_0$ a place of degree one and
$a\ge0$. Then $\calL(aP_0)$ has a basis $u_1,\dots,u_r$ with
$0\le-v_{P_0}(u_1)<\dots<-v_{P_0}(u_r)\le a$.
\end{proposition}

\begin{proof}
For a uniformizer $t$ at $P_0$ and $0\le j\le a$, the map
$f\mapsto(t^jf)(P_0)$ sends $\calL(jP_0)$ to $\kappa(P_0)=\F_q$,
with kernel $\calL((j-1)P_0)$. Choose one element of pole order $j$
for each $j$ for which the quotient is nonzero. These elements are
independent, because their pole orders are distinct, and their
number is $\ell(aP_0)-\ell(-P_0)=\ell(aP_0)$ by \ref{S2}.
\end{proof}

We call such a basis \emph{adapted}. Figure~\ref{fig:pole-filtration}
shows the filtration by pole order.
\fffloat{pole-filtration}

\begin{remark}\label{rem:basis-cex}
If $\ell(aP_0)\ge2$, then some basis of $\calL(aP_0)$ has two
elements with the same pole order at $P_0$. If
$-v_{P_0}(u_1)<-v_{P_0}(u_2)$, then $u_1+u_2$ and $u_2$ span the
same space and have the same pole order. So the argument below
needs an adapted basis, though its conclusion does not depend on the
basis (Remark~\ref{rem:basis-free}).
\end{remark}

\begin{proposition}\label{prop:independence}
Let $v$ be a discrete valuation on a field $K$, let $M'\ge1$ be an
integer, and let $u_1,\dots,u_r\in K^\times$ have values $v(u_i)$
pairwise incongruent modulo $M'$. If $v_1,\dots,v_r\in K$ are not
all zero, then $\sum_iu_iv_i^{M'}\ne0$.
\end{proposition}

\begin{proof}
For $v_i\ne0$,
$v(u_iv_i^{M'})=v(u_i)+M'v(v_i)\equiv v(u_i)\pmod{M'}$. So the
nonzero terms have pairwise distinct values, and the value of the
sum is the least of them.
\end{proof}

Figure~\ref{fig:residues} pictures the proof. This is the only
non-vanishing argument in the proof of the upper bound, and it is
the step for which Section~\ref{sec:dictionary} finds no
counterpart in the number-field argument.
\ffbookfloat{residues}

\begin{remark}\label{rem:basis-free}
Let $K/\F_q$ be a function field, $P_0$ a place of degree one,
$a\ge0$ an integer, $M'$ a power of $p$, and let
$(u_i)_{i\le r}$ and $(u'_j)_{j\le r}$ be two bases of
$\calL(aP_0)$, with $u_i=\sum_jc_{ij}u'_j$
and $c_{ij}\in\F_q$. Since $c\mapsto c^{M'}$ is a bijection of
$\F_q$, for all $v_1,\dots,v_r\in K$
\[
\sum_iu_iv_i^{M'}
=\sum_ju'_j\Bigl(\sum_ic_{ij}^{1/M'}v_i\Bigr)^{M'},
\]
and the matrix $(c_{ij}^{1/M'})$ is invertible. So the statement
that $\sum_iu_iv_i^{M'}=0$ forces all $v_i=0$ holds for one basis
if and only if it holds for all.
\end{remark}

\begin{remark}\label{rem:bombieri-injectivity}
Bombieri proves the same injectivity with the roles of the two
spaces exchanged \cite[pp.~236--237]{Bombieri1974}. Let $K/\F_q$
be a function field, $P_0$ a place of degree one, $q=MM'$ with $M$
a power of $p$, and $a,b\ge0$ integers. Bombieri takes a basis
$s_1,\dots,s_r$ of $\calL(bP_0)$ with
$v_{P_0}(s_1)<\dots<v_{P_0}(s_r)$ and coefficients $\sigma_i$ in
$\calL(aP_0)$. If $\sum_{i\ge\rho}\sigma_i^Ms_i^q=0$ with
$\sigma_\rho\ne0$, then, writing $v=v_{P_0}$,
\[
Mv(\sigma_\rho)+qv(s_\rho)
\ge\min_{i>\rho}\bigl(Mv(\sigma_i)+qv(s_i)\bigr)
\ge-Ma+qv(s_{\rho+1}),
\]
so $Mv(\sigma_\rho)\ge q-Ma>0$ when $a<q/M$. Then $\sigma_\rho$ has
a zero at $P_0$ and no poles, so $\sigma_\rho=0$ by \ref{S2}, a
contradiction. The hypothesis $a<q/M$ is condition (N) below.
\end{remark}

\subsection{The auxiliary section}

\begin{definition}\label{def:ff-admissible}
Let $K/\F_q$ be a function field, $P_0$ a place of degree one,
$\sigma\in\Aut(K/\F_q)$, and $q=MM'$ with $M=p^\mu$ a power of $p$.
A pair of integers $a,b\ge0$ is \emph{admissible} if it satisfies
the two conditions
\begin{enumerate}[label=(\Alph*)]
\item[(N)] $a<M'$;
\item[(S)] $\ell(aP_0)\,\ell(bP_0)
>\ell\bigl(MaP_0+b\,\sigma(P_0)\bigr)$.
\end{enumerate}
\end{definition}

Condition (N) will give non-vanishing and condition (S) is the
hypothesis of a Siegel lemma, whence the letters. The integer $M$ is
the multiplicity; it corresponds to $m$ in
Section~\ref{sec:stepanov}, and $a$, $b$ to $A$, $B$ there.

\begin{theorem}[Stepanov--Bombieri inequality, twisted]
\label{thm:aux-section}
Let $K/\F_q$ be a function field, $P_0$ a place of degree one,
$\sigma\in\Aut(K/\F_q)$, $q=MM'$ with $M=p^\mu$, and $(a,b)$ an
admissible pair. Then there is a nonzero
$w\in\calL((a+M'b)P_0)$ vanishing at every $\sigma$-Frobenius place
$P\ne P_0$. Consequently
\begin{equation}\label{eq:sb-bound}
N(K,\sigma)\le1+a+M'b.
\end{equation}
\end{theorem}

Geometrically, if $K=\F_q(C)$, then $w$ is a nonzero section of
$\calL((a+M'b)P_0)=H^0\bigl(C,\calO_C((a+M'b)P_0)\bigr)$.

\begin{proof}
Fix an adapted basis $u_1,\dots,u_r$ of $\calL(aP_0)$,
$r=\ell(aP_0)$ (Proposition~\ref{prop:adapted-basis}).

\emph{Step 1: a Siegel lemma, by pigeonhole.} This step produces
coefficients with a linear relation. Consider the map
\[
\Lambda\colon\calL(bP_0)^r
\to\calL\bigl(MaP_0+b\,\sigma(P_0)\bigr),\qquad
(v_i)\mapsto\sum_iu_i^M\sigma(v_i).
\]
It is well defined, since $u_i^M\in\calL(MaP_0)$ by
Proposition~\ref{prop:frobenius}\eqref{item:frob-valuation} and
$\sigma(v_i)\in\calL(b\,\sigma(P_0))$ by transport of structure. It
is $\F_q$-linear because $\sigma$ fixes $\F_q$. The source has
$q^{r\ell(bP_0)}$ elements and the target
$q^{\ell(MaP_0+b\sigma(P_0))}$. By (S) two distinct tuples have the
same image, and their difference is a nonzero tuple $(v_i)$ with
\begin{equation}\label{eq:siegel}
\sum_iu_i^M\sigma(v_i)=0.
\end{equation}
No height estimate is needed: the size of a solution is its
membership in $\calL(bP_0)$.

\emph{Step 2: non-vanishing.} This step defines the section. Put
$w=\sum_iu_iv_i^{M'}$. By (N) the pole orders $-v_{P_0}(u_i)$ are
distinct integers in $[0,M')$, so
Proposition~\ref{prop:independence} gives $w\ne0$. By
Proposition~\ref{prop:frobenius}\eqref{item:frob-valuation},
$w\in\calL((a+M'b)P_0)$.

\emph{Step 3: forced multiplicity.} This step shows that $w$
vanishes at the points to be counted. Since $M$ is a power of $p$
and $MM'=q$,
Proposition~\ref{prop:frobenius}\eqref{item:frob-additive} gives
\begin{equation}\label{eq:wM}
w^M=\sum_iu_i^Mv_i^q.
\end{equation}
Let $P\ne P_0$ be $\sigma$-Frobenius. All $u_i,v_i$ lie in
$\calO_P$, because their only pole is at $P_0$, and
$v_i^q-\sigma(v_i)\in P$ by the definition of a
$\sigma$-Frobenius place. By \eqref{eq:siegel} and \eqref{eq:wM},
\[
w^M=\sum_iu_i^M\bigl(v_i^q-\sigma(v_i)\bigr)\in P.
\]
As $P$ is prime, $w\in P$. Thus $v_P(w)\ge1$, and the section $w^M$
of $\calO_C((Ma+qb)P_0)$ vanishes to order at least $M$ at $P$. No
derivative was computed: in a local parameter,
$(\sum_{j\ge1}c_jt^j)^M=\sum_jc_j^Mt^{jM}$, so all jets of order
less than $M$ of an $M$-th power vanish once its value does
(Figure~\ref{fig:jets}).

\emph{Step 4: the slope inequality.} Let
$Z_\sigma$ be the sum of the $\sigma$-Frobenius places $P\ne P_0$;
it is finite because each is a zero of $w\ne0$. Then
$(w)+(a+M'b)P_0-Z_\sigma\ge0$. An effective divisor has nonnegative
degree (\ref{S2}), so $\deg Z_\sigma\le a+M'b$. Equivalently,
multiplication by $w$ is a nonzero map of line bundles
$\calO_C(Z_\sigma)\to\calO_C((a+M'b)P_0)$, and the degree of the
source is at most that of the target. Finally
$N(K,\sigma)\le\deg P_0+\deg Z_\sigma$.
\end{proof}

For $\sigma=\mathrm{id}$, \eqref{eq:sb-bound} is Bombieri's
inequality (7) \cite[p.~238]{Bombieri1974}.
Figure~\ref{fig:zero-pole-budget} draws the count of Step~4. The
function $f$ of Section~\ref{sec:stepanov} is $w^M$: the inequality
$m(N-1)\le mA+qB$ there is \eqref{eq:sb-bound} multiplied by $M$.
\ffbookfloat{jets}
\ffbookfloat{zero-pole-budget}

\subsection{Balancing}

It remains to find admissible pairs for which the bound
\eqref{eq:sb-bound} is small.

\begin{proposition}\label{prop:balancing}
Let $K/\F_q$ be a function field of genus $g$, $P_0$ a place of
degree one and $\sigma\in\Aut(K/\F_q)$. Let $M,M'\ge1$ be integers
with $MM'=q$, $M$ a power of $p$ and $M'\ge g+2$. Then the pair
$a=M'-1$, $b=M+g+\lfloor Mg/(M'-1-g)\rfloor$ is admissible, and
\begin{equation}\label{eq:balancing}
1+a+M'b=MM'+(g+1)M'+M'\Bigl\lfloor\frac{Mg}{M'-1-g}\Bigr\rfloor.
\end{equation}
If $M'\ge2(g+1)$, the right side of \eqref{eq:balancing} is at most
$MM'+(g+1)M'+2gM$.
\end{proposition}

\begin{proof}
Condition (N) is clear. As $a\ge g+1$, $b\ge g$ and
$Ma+b\ge2g+2$, \ref{S3} gives
$\ell(aP_0)\ell(bP_0)\ge(a+1-g)(b+1-g)$ and
$\ell(MaP_0+b\sigma(P_0))=Ma+b+1-g$. Now
\begin{equation}\label{eq:siegel-margin}
(a+1-g)(b+1-g)-(Ma+b+1-g)=(a-g)(b+1-g)-Ma.
\end{equation}
With $a=M'-1$,
$\frac{Ma}{a-g}=M+\frac{Mg}{M'-1-g}<b+1-g$, so
\eqref{eq:siegel-margin} is positive, which is (S). The identity
\eqref{eq:balancing} is a substitution. Finally, writing $x$ for
$Mg/(M'-1-g)$,
$M'\lfloor x\rfloor\le\frac{M'Mg}{M'-1-g}\le2gM$ when
$M'\le2(M'-1-g)$, that is when $M'\ge2(g+1)$.
\end{proof}

Figure~\ref{fig:ab-region} shows the admissible pairs in one case.
So the error in \eqref{eq:sb-bound} is at most $(g+1)q/M+2gM$: the
price of (N) falls with $M$, that of (S) rises. Bombieri takes
$M=M'=\sqrt q$, assuming $q$ a square. The two spaces then play
symmetric roles. The square root in the error does not come from
this choice (Proposition~\ref{prop:sqrtq-forced}).
\ffabregion

\begin{proposition}\label{prop:sqrtq-forced}
Let $K/\F_q$ be a function field of genus $g\ge1$, $P_0$ a place of
degree one, $\sigma\in\Aut(K/\F_q)$, and $q=MM'$ with $M$ a power
of $p$. Let $(a,b)$ be an admissible pair for which the three
Riemann--Roch dimensions take their generic values:
$\ell(aP_0)=a+1-g$, $\ell(bP_0)=b+1-g$ and
$\ell(MaP_0+b\sigma(P_0))=Ma+b+1-g$. Then
\begin{equation}\label{eq:sqrtq-forced}
1+a+M'b-q>(g-1)\frac qM+gM\qquad\text{and}\qquad
1+a+M'b-q>g+2\sqrt{gq}.
\end{equation}
\end{proposition}

\begin{proof}
By \eqref{eq:siegel-margin}, (S) reads $(a-g)(b+1-g)>Ma$. Here
$a\ge g$ because $\ell(aP_0)\ge1$, and $a=g$ is impossible because
$Ma>0$; so $a>g$ and (S) reads $b+1-g>M+\frac{Mg}{a-g}$.
Multiplying by $M'$,
\[
M'b>q+(g-1)M'+\frac{qg}{a-g}.
\]
For the first inequality use $a-g<M'$, which holds by (N), so that
$\frac{qg}{a-g}>gM$. For the second drop $(g-1)M'\ge0$ and use
$a+\frac{qg}{a-g}=g+(a-g)+\frac{qg}{a-g}\ge g+2\sqrt{gq}$.
\end{proof}

\begin{remark}\label{rem:sqrtq}
The square root does not enter the error of
Theorem~\ref{thm:aux-section} through the choice $M=M'=\sqrt q$.
By Proposition~\ref{prop:sqrtq-forced}, no choice of $M$ and $(a,b)$
with generic Riemann--Roch values gives an error $o(\sqrt q)$, and
for $g\ge2$ an error $O(\sqrt q)$ requires $M\asymp_g\sqrt q$. The
second inequality of \eqref{eq:sqrtq-forced} is the
arithmetic--geometric mean inequality. The product of the two
competing costs is $gq$ whatever the split
(Figure~\ref{fig:sqrtq-valley}). Chapter~\ref{ch:idea} follows this
balance to the integers. There the two costs are $x/m$ and $m$, and
the chain length is $m\approx\sqrt x$
(Section~\ref{sec:multiplicity}, Proposition~\ref{prop:sqrtx}).
\end{remark}
\ffbookfloat{sqrtq-valley}

\subsection{Bombieri's theorem}

\begin{theorem}[Bombieri, twisted]\label{thm:bombieri}
Let $g\ge0$ be an integer and $q=q_0^2$ with $q_0$ a power of $p$
and $q_0>(g+1)^2$. Let $K/\F_q$ be a function field of genus $g$
with a place of degree one. Then for every $\sigma\in\Aut(K/\F_q)$
\begin{equation}\label{eq:bombieri}
N(K,\sigma)\le q+2gq_0+\Bigl\lfloor\frac{gq_0}{g+1}\Bigr\rfloor+g+2
\le q+(2g+1)q_0.
\end{equation}
In particular every function field of genus $g$ over $\F_q$ has
$N_1<q+1+(2g+1)\sqrt q$.
\end{theorem}

\begin{proof}
Take $M=M'=q_0$, $b=q_0+2g$ and
$a=\lfloor gq_0/(g+1)\rfloor+g+1$ \cite[p.~239]{Bombieri1974}.
Condition (N): $a\le q_0-1$ is equivalent to
$\frac{gq_0}{g+1}<q_0-g-1$, that is to $q_0>(g+1)^2$. Condition
(S): as $a\ge g+1$, $b\ge g$ and $Ma+b\ge2g-1$, fact \ref{S3} and
\eqref{eq:siegel-margin} show that it suffices that
$(a-g)(q_0+g+1)>q_0a$, that is $a>\frac{gq_0}{g+1}+g$, which holds.
Then $1+a+q_0b$ is the middle term of \eqref{eq:bombieri}, and the
last inequality of \eqref{eq:bombieri} is $a+1\le q_0$, which is
(N). If $K$ has no place of degree one, $N_1=0$.
\end{proof}

This is Bombieri's Theorem~1 \cite[p.~236]{Bombieri1974}, and for
$\sigma=\mathrm{id}$ it is Proposition~\ref{prop:stepanov}.

\begin{remark}[The pair used in Chapter~\ref{ch:idea}]
\label{rem:idea-pair}
The sketch in Section~\ref{sec:stepanov} uses, with $M=M'=q_0$, the
pair $a=q_0-1$, $b=q_0+2g$ in place of Bombieri's. It is admissible
under the same hypothesis: (N) is clear, and by
\eqref{eq:siegel-margin} the margin in (S) is at least
$(q_0-1-g)(q_0+g+1)-q_0(q_0-1)=q_0-(g+1)^2>0$. It gives
$1+a+q_0b=q+(2g+1)q_0$, the last term of \eqref{eq:bombieri}, where
Bombieri's pair gives the slightly smaller middle term. Condition
\eqref{eq:kernel} of Chapter~\ref{ch:idea} is (S) with the
Riemann--Roch lower bounds inserted.
\end{remark}

Both (N) and (S) in the proof of Theorem~\ref{thm:bombieri} reduce
to chains of equivalences ending in $q_0>(g+1)^2$. The chain for
(S) used the lower bound $\ell(aP_0)\ge a+1-g$, which is an
equality only when $a\ge2g-1$. So the chain does not show that
$q_0>(g+1)^2$ is necessary.

\begin{remark}\label{rem:special-admissible}
Let $b=q_0+2g$. For generic Riemann--Roch values an admissible $a$
exists if and only if $q_0>(g+1)^2$. At points with special gap
sequences an admissible $a$ can exist with $q_0\le(g+1)^2$. Let
$q=9$, $q_0=3$, $K$ the function field of
$y^2=f(x)$ with $f\in\F_9[x]$ squarefree of degree $7$, so $g=3$,
$P_0$ the place at infinity and $\sigma=\mathrm{id}$. Then
$1,x\in\calL(2P_0)$, so $\ell(2P_0)=2$, while $\ell(9P_0)=7$ and
$\ell(15P_0)=13$. So $(a,b)=(2,9)$ satisfies (N) and (S), since
$2\cdot7>13$, although $q_0=3\le16$, and it gives $N_1\le30$. This
cannot happen if $q_0\ge g+1$. Indeed, let $(a,b)$ be admissible
with $b=q_0+2g$ and $M=M'=q_0$. By \ref{S3},
$\ell(bP_0)=q_0+g+1$ and the right side of (S) is
$q_0a+q_0+g+1$. If $a\ge2g-1$, then $\ell(aP_0)=a+1-g$ and the
chain of equivalences in the proof of Theorem~\ref{thm:bombieri} is
exact, so $q_0>(g+1)^2$. If $a\le2g-2$, then Clifford's theorem,
which is not proved here, gives $\ell(aP_0)\le a/2+1$, and (S)
gives $(a/2)(q_0+g+1)>q_0a$, so $a>0$ and $q_0<g+1$.
\end{remark}

The exposition \cite{Tao2014} carries out Bombieri's choice with
the same parameters, and states it for every $q\ge(g+1)^4$.

\begin{remark}\label{rem:tao}
The parameter calculation of \cite{Tao2014} does not reach every
$q\ge(g+1)^4$, for two reasons. At $q=(g+1)^4$ the
parameter $\ell=\lfloor\frac g{g+1}\sqrt q\rfloor+g+1$ of
\cite{Tao2014} equals $\sqrt q$, against its condition
$\ell<\sqrt q$. Its inequality (11) has right side
$\ell\sqrt q+m+1$ instead of $\ell\sqrt q+m+1-g$, and then holds
exactly when $g+1$ divides $g\sqrt q$; for $g=1$, $q=25$ it reads
$28>28$. With the exact value it is correct
(Theorem~\ref{thm:bombieri}).
\end{remark}

For $q=p^e$ with $e$ odd no power of $p$ equals $\sqrt q$, and the
balanced choice is unavailable. The upper bound still has an error
of order $\sqrt q$.

\begin{proposition}\label{prop:all-q}
Let $g\ge0$ be an integer and $q=p^e$ with $\sqrt q\ge2(g+1)$, and
let $K/\F_q$ be a function
field of genus $g$ with a place of degree one. Then
$N(K,\sigma)\le q+((g+1)\sqrt p+2g)\sqrt q$ for every
$\sigma\in\Aut(K/\F_q)$. Hence for every function field $K/\F_q$ of
genus $g$ and every $n$ with $q^{n/2}\ge2(g+1)$,
\begin{equation}\label{eq:all-q}
N_n\le q^n+\bigl((g+1)\sqrt p+2g\bigr)q^{n/2}.
\end{equation}
\end{proposition}

\begin{proof}
Take $M=p^{\lfloor e/2\rfloor}$ and $M'=p^{\lceil e/2\rceil}$, so
that $M\le\sqrt q\le M'\le\sqrt{pq}$ and $M'\ge2(g+1)$, and apply
Proposition~\ref{prop:balancing} and
Theorem~\ref{thm:aux-section}. For \eqref{eq:all-q} apply this to
$K_n$, which has full constant field $\F_{q^n}$, genus $g$
(Proposition~\ref{prop:tower}) and $N_n$ places of degree one
(\ref{S4}); if $N_n=0$ there is nothing to prove.
\end{proof}

\begin{remark}\label{rem:square}
The upper bound does not need $q$ to be a square to give an error
of order $\sqrt q$. By Proposition~\ref{prop:all-q}, the square
case only improves the constant, from $(g+1)\sqrt p+2g$ to $2g+1$.
Figure~\ref{fig:constants} compares the constants.
\end{remark}
\ffbookfloat{constants}

\section{Route 1: the lower bound by Galois averaging}
\label{sec:route1}

An upper bound for point counts is half of
Theorem~\ref{thm:weil}. Following \cite[p.~240]{Bombieri1974}, the
other half comes from the same upper bound, applied to all the
twists of a Galois cover of the projective line at once: the total
of their point counts is known exactly, so if none is too large,
none is too small (Figure~\ref{fig:cover-twists}). This is the
reason for the automorphism $\sigma$ in
Theorem~\ref{thm:aux-section}.
\fffloat{cover-twists}

\begin{definition}\label{def:decomposition}
Let $E/F$ be a finite Galois extension of function fields with
group $G$, both with full constant field $\F_q$, and let $P$ be a
place of $E$.
\begin{enumerate}
\item $D(P)=\{\sigma\in G:\sigma(P)=P\}$ is the \emph{decomposition
group}; each $\sigma\in D(P)$ induces an automorphism $\bar\sigma$
of $\kappa(P)$ over $\kappa(P\cap F)$.
\item $I(P)$, the \emph{inertia group}, is the kernel of
$\sigma\mapsto\bar\sigma$ on $D(P)$.
\item $\Phi(P)=\{\sigma\in G:P\text{ is }\sigma\text{-Frobenius}\}$.
\end{enumerate}
\end{definition}

\begin{proposition}\label{prop:frob-place-base}
Let $E/F$ be a finite Galois extension of function fields with
group $G$, both with full constant field $\F_q$, let $P$ be a place
of $E$, $Q=P\cap F$ and $\sigma\in G$. If $P$ is
$\sigma$-Frobenius, then $\deg Q=1$ and $\sigma\in D(P)$.
\end{proposition}

\begin{proof}
For $z\in\calO_Q$ we have $\sigma(z)=z$, so $z^q-z\in P\cap F=Q$.
Every element of $\kappa(Q)$ is then a root of $X^q-X$, so
$\kappa(Q)=\F_q$. That $\sigma(P)=P$ was shown after
Definition~\ref{def:frob-place}.
\end{proof}

\begin{proposition}\label{prop:frob-coset}
Let $E/F$ be a finite Galois extension of function fields with
group $G$, both with full constant field $\F_q$, let $Q$ be a place
of $F$ of degree one and $P$ a place of $E$ above $Q$. Then
$\Phi(P)$ is a coset of $I(P)$ in $D(P)$, in particular nonempty,
and
\begin{equation}\label{eq:frob-coset}
\lvert\Phi(P)\rvert\cdot\deg P=\lvert D(P)\rvert.
\end{equation}
\end{proposition}

\begin{proof}
Here $\kappa(Q)=\F_q$, and $\Gal(\kappa(P)/\F_q)$ is cyclic of
order $f=\deg P$, generated by $\Fr_q$ (\ref{S8}). By
Proposition~\ref{prop:frob-place-base} and the definition of a
$\sigma$-Frobenius place,
$\Phi(P)=\{\sigma\in D(P):\bar\sigma=\Fr_q\}$, which is empty or a
coset of $I(P)$.

We show first that $\Phi(P)$ is nonempty. Let $P=P_1,\dots,P_s$ be the places
above $Q$
(\ref{S7}). Let $\bar\theta$ generate the group $\kappa(P)^\times$
(\ref{S8}); so $\bar\theta\ne0$ and
$\kappa(P)=\F_q(\bar\theta)$. By \ref{S1} there is $\theta\in E$
with $\theta\in\calO_P$, $\theta(P)=\bar\theta$ and
$v_{P_j}(\theta)\ge1$ for $j\ge2$. For $\tau\in G$,
$v_P(\tau\theta)=v_{\tau^{-1}P}(\theta)\ge0$. Put
$\chi(X)=\prod_{\tau\in G}(X-\tau\theta)$. Its coefficients are
$G$-invariant and lie in $\calO_P$, hence in
$F\cap\calO_P=\calO_Q$. Its reduction
$\bar\chi=\prod_\tau(X-\overline{\tau\theta})$ lies in $\F_q[X]$
and has the root $\bar\theta$. So
$\bar\chi(\bar\theta^q)=\bar\chi(\bar\theta)^q=0$, and
$\bar\theta^q=\overline{\tau\theta}$ for some $\tau$. As
$\bar\theta^q\ne0$, $v_{\tau^{-1}P}(\theta)=0$, which forces
$\tau^{-1}P=P$. So $\tau\in D(P)$ and
$\bar\tau(\bar\theta)=\bar\theta^q$; since $\bar\theta$ generates
$\kappa(P)$, $\bar\tau=\Fr_q$.

It remains to count. The image of $D(P)$ in the cyclic group
$\langle\Fr_q\rangle$ contains a generator, so it is everything.
Hence $\lvert D(P)\rvert=\lvert I(P)\rvert f$ and
$\lvert\Phi(P)\rvert=\lvert I(P)\rvert$.
\end{proof}

Figure~\ref{fig:frob-coset} shows the coset $\Phi(P)$ at an
unramified and at a ramified place.
\fffloat{frob-coset}

The sum of $N(E,\sigma)$ over $\sigma\in G$ counts each rational
place of the base once for every way it can be a Frobenius fixed
point upstairs. At an unramified place this is the familiar
statement that Frobenius is a well-defined conjugacy class. At a
branch point it is less clear what should be counted: the residue
field extension no longer determines one element of $G$, and
Bombieri's identity \cite[p.~240, eq.~(9)]{Bombieri1974}
accordingly carries an error term $O(1)$ for the branch points,
which are finite in number. For the count of
Definition~\ref{def:frob-place} the identity is exact.

\begin{proposition}[class-sum identity]\label{prop:class-sum}
Let $E/F$ be a finite Galois extension of function fields with
group $G$, both with full constant field $\F_q$, and let $N_1(F)$
be the number of places of $F$ of degree one. Then
\begin{equation}\label{eq:class-sum}
\sum_{\sigma\in G}N(E,\sigma)=\lvert G\rvert\cdot N_1(F).
\end{equation}
\end{proposition}

\begin{proof}
The left side is $\sum_P\lvert\Phi(P)\rvert\deg P$, the sum over
the places of $E$. By Propositions~\ref{prop:frob-place-base}
and~\ref{prop:frob-coset} only the $P$ over places $Q$ of degree
one contribute, and each contributes $\lvert D(P)\rvert$. The
places $P$ over a given $Q$ form one $G$-orbit (\ref{S7}), with
stabilisers $D(P)$, so together they contribute $\lvert G\rvert$.
\end{proof}

\begin{remark}\label{rem:class-sum-error}
For the count of Definition~\ref{def:frob-place}, the averaging
identity over all twists of a Galois cover has no error term from
the ramified places. A ramified place over a rational place is
$\sigma$-Frobenius for $\lvert I(P)\rvert$ elements $\sigma$, and
this compensates exactly. Bombieri counts unramified points only,
and his identity has the error term for that reason.
\end{remark}

\begin{theorem}\label{thm:averaging}
Let $F_0=\F_q(x)$, let $E/F_0$ be a finite Galois extension with
group $G$, and let $F$ be an intermediate field, all three with
full constant field $\F_q$, and let $N_1(F)$ be the number of
places of $F$ of degree one. Let $\Theta\ge0$ and suppose
$N(E,\sigma)\le q+1+\Theta$ for all $\sigma\in G$. Then
$N(E,\tau)\ge q+1-(\lvert G\rvert-1)\Theta$ for every $\tau\in G$,
and
\[
q+1-(\lvert G\rvert-1)\Theta\le N_1(F)\le q+1+\Theta.
\]
\end{theorem}

\begin{proof}
By \eqref{eq:class-sum} for $E/F_0$ and \ref{S5},
$\sum_\sigma N(E,\sigma)=\lvert G\rvert(q+1)$. Subtracting the
upper bounds for the $\sigma\ne\tau$ gives the first bound. By
\eqref{eq:class-sum} for $E/F$, $N_1(F)$ is the average of
$N(E,\tau)$ over $\tau\in\Gal(E/F)$, and each of these lies between
the two bounds.
\end{proof}

The base $F_0$ has trivial $L$-polynomial, so its point count $q+1$
is known exactly. Figure~\ref{fig:twist-bars} shows the counts of
the twists in an example. There is no base of
this kind over the integers, which is one of the places where the
analogy breaks (Section~\ref{sec:dictionary}, item (B3)).
\ffbookfloat{twist-bars}

\subsection{The Galois closure and the tower}
\label{subsec:closure}

Let $K/\F_q$ be a function field of genus $g$. To apply
Theorem~\ref{thm:averaging} to $K$ we fix the following data.
\begin{enumerate}
\item Choose $x\in K$ with $K/\F_q(x)$ finite separable (\ref{S6}).
Let $E$ be the Galois closure of $K/\F_q(x)$ and
$\gamma=[E:\F_q(x)]$.
\item Let $\F_{q^d}$ be the algebraic closure of $\F_q$ in $E$. It
is finite over $\F_q$, because
$[\F_{q^d}(x):\F_q(x)]=d\le\gamma$. Then $E/\F_{q^d}$ is a function
field with full constant field $\F_{q^d}$; let $g_E$ be its genus.
Choose a place $P_1$ of $E$ and let $d_1$ be its degree over
$\F_{q^d}$.
\item For $n$ divisible by $d$ put $E_n=E\F_{q^n}$,
$K_n=K\F_{q^n}$ and $F_{0,n}=\F_{q^n}(x)$. By \ref{S4} and
Proposition~\ref{prop:tower} applied to $E/\F_{q^d}$, $E_n$ has full
constant field $\F_{q^n}$ and genus $g_E$; so $K_n$, lying between
$F_{0,n}$ and $E_n$, has full constant field $\F_{q^n}$. By
\ref{S9}, $E_n/F_{0,n}$ is Galois with group $G_n$ of order at most
$\gamma$, and every element of $G_n$ fixes $\F_{q^n}$.
\item If also $dd_1\mid n$, then $E_n$ has a place of degree one,
by \ref{S4} applied to $P_1$.
\end{enumerate}
Let $n_1$ be the least even multiple of $dd_1$ with
$q^{n_1/2}>(g_E+1)^2$.

\begin{theorem}\label{thm:two-sided}
Let $K/\F_q$ be a function field, let $x\in K$ be such that
$K/\F_q(x)$ is finite and separable, let $E$ be the Galois closure
of $K/\F_q(x)$ and $\gamma=[E:\F_q(x)]$, let $\F_{q^d}$ be the
algebraic closure of $\F_q$ in $E$ and $g_E$ the genus of
$E/\F_{q^d}$, let $d_1$ be the degree over $\F_{q^d}$ of a place of
$E$, and let $n_1$ be the least even multiple of $dd_1$ with
$q^{n_1/2}>(g_E+1)^2$. Then
\begin{equation}\label{eq:two-sided}
\lvert N_n-q^n-1\rvert\le\max(1,\gamma-1)(2g_E+1)\,q^{n/2}
\qquad\text{for all }n\in n_1\Z_{>0}.
\end{equation}
\end{theorem}

\begin{proof}
Let $n\in n_1\Z_{>0}$. Then $q^n$ is a square,
$q^{n/2}>(g_E+1)^2$, and $E_n$ has a place of degree one.
Theorem~\ref{thm:bombieri} for $E_n/\F_{q^n}$ gives
$N(E_n,\sigma)\le q^n+(2g_E+1)q^{n/2}$ for all $\sigma\in G_n$.
Theorem~\ref{thm:averaging} for
$E_n\supseteq K_n\supseteq F_{0,n}$ over $\F_{q^n}$, with
$\Theta=(2g_E+1)q^{n/2}-1$, gives
\[
q^n+1-(\gamma-1)(2g_E+1)q^{n/2}\le N_1(K_n)
\le q^n+(2g_E+1)q^{n/2}.
\]
By \ref{S4}, $N_1(K_n)=N_n$.
\end{proof}

The next section uses only that $\gamma$ and $g_E$ do not depend on
$n$. The constants in \eqref{eq:two-sided} are poor, and the bound
holds only along one arithmetic progression; both defects disappear
at the next step.

\section{The Riemann hypothesis}
\label{sec:rh}

\begin{proposition}\label{prop:mean-values}
Let $\phi$ be a real number and $I$ a finite nonempty interval of
integers. If $\phi\notin2\pi\Z$, then
\[
\Bigl\lvert\frac1{\lvert I\rvert}\sum_{k\in I}e^{ik\phi}\Bigr\rvert
\le\frac{2}{\lvert I\rvert\,\lvert1-e^{i\phi}\rvert};
\]
if $\phi\in2\pi\Z$ the mean is $1$.
\end{proposition}

\begin{proof}
Sum the geometric series.
\end{proof}

A set $W$ of positive integers \emph{contains arbitrarily long
intervals} if for every $n$ it contains $n$ consecutive integers.
In such a set one can choose intervals $I_\nu\subseteq W$ with
$\lvert I_\nu\rvert\to\infty$ and $\min I_\nu\to\infty$, for
instance the upper halves of longer and longer intervals in $W$.

\begin{proposition}\label{prop:power-sums}
Let $\beta_1,\dots,\beta_r\in\C$, let $R,A>0$, and let $W$ be a set
of positive integers containing arbitrarily long intervals. If
$\lvert\sum_i\beta_i^k\rvert\le AR^k$ for all $k\in W$, then
$\lvert\beta_i\rvert\le R$ for all $i$.
\end{proposition}

\begin{proof}
Suppose $\rho=\max_i\lvert\beta_i\rvert>R$. Let
$\rho e^{i\theta_1},\dots,\rho e^{i\theta_s}$ be the distinct
values of modulus $\rho$ among the $\beta_i$, with multiplicities
$\mu_j\ge1$. Then
$T_k:=\rho^{-k}\sum_i\beta_i^k
=\sum_j\mu_je^{ik\theta_j}+\epsilon_k$
with $\epsilon_k\to0$, and
$\lvert T_k\rvert\le A(R/\rho)^k\to0$ on $W$. Choose intervals
$I_\nu\subseteq W$ with $\lvert I_\nu\rvert\to\infty$ and
$\min I_\nu\to\infty$. By Proposition~\ref{prop:mean-values} the
mean of $T_ke^{-ik\theta_1}$ over $I_\nu$ tends to $\mu_1\ge1$. It
also tends to $0$, a contradiction.
\end{proof}

No sign is used here: the power sums are bounded in absolute value,
and no cancellation among the $\beta_i$ of largest modulus can hide
them all. Compare Proposition~\ref{lem:growth}, the same statement
for weighted sums and all $k$, which is where the sign-only proof
needs its sign.

\begin{theorem}[the Riemann hypothesis for $K$]\label{thm:rh}
Let $K/\F_q$ be a function field of genus $g$ with reciprocal roots
$\alpha_1,\dots,\alpha_{2g}$. Then
$\lvert\alpha_i\rvert=\sqrt q$ for all $i$, and
$\lvert N_n-q^n-1\rvert\le2g\,q^{n/2}$ for all $n\ge1$.
\end{theorem}

\begin{proof}
By Theorem~\ref{thm:rationality}\eqref{item:rat-roots} and
\eqref{eq:two-sided},
\[
\Bigl\lvert\sum_i(\alpha_i^{n_1})^k\Bigr\rvert
=\lvert q^{n_1k}+1-N_{n_1k}\rvert\le A\,(q^{n_1/2})^k
\qquad\text{for all }k\ge1,
\]
with $A=\max(1,\gamma-1)(2g_E+1)$.
Proposition~\ref{prop:power-sums} gives
$\lvert\alpha_i\rvert^{n_1}\le q^{n_1/2}$, so
$\lvert\alpha_i\rvert\le\sqrt q$. By
Theorem~\ref{thm:rationality}\eqref{item:rat-roots},
$q/\alpha_i$ is also a reciprocal root, so
$q/\lvert\alpha_i\rvert\le\sqrt q$. Hence
$\lvert\alpha_i\rvert=\sqrt q$, and
$\lvert N_n-q^n-1\rvert=\lvert\sum_i\alpha_i^n\rvert
\le2gq^{n/2}$.
\end{proof}

This proves Theorem~\ref{thm:weil}. The bounds were needed only for
$n$ in one arithmetic progression, with constants independent of
$n$; the sharp constant $2g$ and all $n$ come out at the end, from
the tensor-power step of Proposition~\ref{prop:power-sums}.
Figure~\ref{fig:roots-circle} shows the reciprocal roots of some
curves.
\ffbookfloat{roots-circle}

\section{Route 2: a one-sided bound and a Landau--Pringsheim lemma}
\label{sec:route2}

This section asks whether the upper bound of
Section~\ref{sec:upper}, without the averaging of
Section~\ref{sec:route1}, gives the Riemann hypothesis through a
positivity lemma of Landau--Pringsheim type. That is how the
analytic step of the number-field argument works: a one-sided bound
and Landau's theorem (Section~\ref{sec:idea-landau},
Theorem~\ref{thm:windowed-landau}). Over a function field it does
not work with the bound that Section~\ref{sec:upper} provides, and
we identify the missing input exactly. The section then finds the
function-field statement to which the number-field criterion does
correspond.

\subsection{The lemma}

\begin{proposition}[Landau--Pringsheim for exponential sums]
\label{prop:landau}
Let $\beta_1,\dots,\beta_s\in\C^\times$ be distinct, let
$c_1,\dots,c_s\in\C^\times$, and let $W$ be a set of positive
integers containing arbitrarily long intervals. Suppose that
$\sum_jc_j\beta_j^k$ is real and nonnegative for all $k\in W$. Then
the largest of the moduli $\lvert\beta_j\rvert$ is itself one of
the $\beta_j$; the coefficient of this $\beta_j$ is real and
positive; and it is at least $\lvert c_i\rvert$ for every $i$ for
which $\lvert\beta_i\rvert$ is largest.
\end{proposition}

\begin{proof}
Put $a_k=\sum_jc_j\beta_j^k$ and
$\rho=\max_j\lvert\beta_j\rvert$. Let $J=\{j:\lvert\beta_j\rvert=\rho\}$ and
$\beta_j=\rho e^{i\theta_j}$ for $j\in J$, the $\theta_j$ distinct
modulo $2\pi$. Put
$b_k=\rho^{-k}a_k=\sum_{j\in J}c_je^{ik\theta_j}+\epsilon_k$ with
$\epsilon_k\to0$. Choose intervals $I_\nu\subseteq W$ with
$\lvert I_\nu\rvert\to\infty$ and $\min I_\nu\to\infty$. For real
$\psi$ let
$m(\psi)=\lim_\nu\frac1{\lvert I_\nu\rvert}
\sum_{k\in I_\nu}b_ke^{-ik\psi}$.
By Proposition~\ref{prop:mean-values}, $m(\theta_j)=c_j$ for
$j\in J$, and $m(\psi)=0$ if $\psi$ is congruent to no $\theta_j$.
Fix $j_1\in J$ and a complex number $\epsilon$ of modulus one.
Since $b_k$ is real and nonnegative on $W$ and
$1+\Real(\epsilon e^{-ik\theta_{j_1}})\ge0$,
\[
0\le\lim_\nu\frac1{\lvert I_\nu\rvert}\sum_{k\in I_\nu}
b_k\bigl(1+\Real(\epsilon e^{-ik\theta_{j_1}})\bigr)
=m(0)+\Real\bigl(\epsilon\,c_{j_1}\bigr).
\]
Choosing $\epsilon=-\overline{c_{j_1}}/\lvert c_{j_1}\rvert$ gives
$m(0)\ge\lvert c_{j_1}\rvert>0$. So some $\theta_{j_0}\equiv0$,
that is $\beta_{j_0}=\rho$, and $c_{j_0}=m(0)$ is real, positive
and at least $\lvert c_{j_1}\rvert$.
\end{proof}

The generating function
$\sum_ka_kt^k=\sum_j\frac{c_j\beta_jt}{1-\beta_jt}$ has its poles
at $t=1/\beta_j$. The proposition says that if the coefficients are
nonnegative, the pole nearest to $0$ lies on the positive real axis
and its coefficient is positive
(Figure~\ref{fig:landau-poles}): Pringsheim's theorem for this
class, of which Landau's theorem is the Dirichlet-series form
\cite{Titchmarsh1939}. Restricting to a set $W$ of windows costs
nothing here, because there are finitely many frequencies; over the
integers the windowed form is a theorem with a proof of its own
(Theorem~\ref{thm:windowed-landau}).
\fffloat{landau-poles}

\subsection{What a one-sided bound gives}

Section~\ref{sec:upper} gives $N_n\le q^n+1+Aq^{n/2}$, so the
numbers $Aq^{n/2}+\sum_i\alpha_i^n$ are nonnegative. Positivity
confines the first singularity to the real axis, the functional
equation pairs roots outside $\lvert\alpha\rvert=\sqrt q$ with
roots inside, and the bound holds along the whole tower. These
inputs do not prove the Riemann hypothesis
(Remark~\ref{rem:route2}).

\begin{theorem}[dichotomy]\label{thm:dichotomy}
Let $K/\F_q$ be a function field of genus $g\ge1$ with reciprocal
roots $\alpha_i$, let $n_1\ge1$ and $A>0$, let $W$ be a set of
positive integers containing arbitrarily long intervals, and
suppose $N_n\le q^n+1+Aq^{n/2}$ for all $n\in n_1W$. Then exactly
one of the following holds.
\begin{enumerate}
\item\label{item:dich-rh} $\lvert\alpha_i\rvert=\sqrt q$ for all
$i$.
\item\label{item:dich-real} The largest modulus of an $\alpha_i$
exceeds $\sqrt q$; for some $\alpha_i$ of largest modulus the power
$\alpha_i^{n_1}$ is real and positive; and this value has the
largest multiplicity among the $\alpha_j^{n_1}$ with $\alpha_j$ of
largest modulus.
\end{enumerate}
\end{theorem}

\begin{proof}
Let $\rho=\max_i\lvert\alpha_i\rvert$; in case
\eqref{item:dich-real} the power in question is $\rho^{n_1}$. The
two cases exclude each other. For $k\in W$ put
\[
a_k:=Aq^{n_1k/2}+\sum_i\alpha_i^{n_1k}
=Aq^{n_1k/2}+q^{n_1k}+1-N_{n_1k}\ge0.
\]
Collecting equal values of $\alpha_i^{n_1}$, $a_k$ is an
exponential sum whose frequencies are the distinct
$\beta=\alpha_i^{n_1}$, with coefficients the multiplicities
$\mu_\beta\ge1$, together with the frequency $q^{n_1/2}$ with a
positive coefficient. If $\rho^{n_1}>q^{n_1/2}$, the frequencies of
largest modulus are among the $\beta$, and
Proposition~\ref{prop:landau} gives \eqref{item:dich-real}. If
$\rho\le\sqrt q$, the functional equation
(Theorem~\ref{thm:rationality}\eqref{item:rat-roots}) gives
$\lvert\alpha_i\rvert\ge\sqrt q$ for all $i$, which is
\eqref{item:dich-rh}.
\end{proof}

Theorem~\ref{thm:dichotomy} leaves one alternative to the Riemann
hypothesis, and it has the shape of a failure of positivity: a real
reciprocal root of maximal modulus that the upper bound for $N_n$
does not see. The next theorem identifies the missing input
exactly.

\begin{theorem}[the exact input Route 2 needs]
\label{thm:route2-input}
Let $K/\F_q$ be a function field with reciprocal roots $\alpha_i$
and $L$-polynomial $L_K$. The following are equivalent.
\begin{enumerate}
\item\label{item:r2-rh} $\lvert\alpha_i\rvert=\sqrt q$ for all $i$.
\item\label{item:r2-real} $L_K(t)\ne0$ for real
$t\in(q^{-1},q^{-1/2})$.
\end{enumerate}
\end{theorem}

The proof that \eqref{item:r2-real} implies \eqref{item:r2-rh} uses
only Sections~\ref{sec:zeta} and~\ref{sec:upper} and
Proposition~\ref{prop:landau}, not Section~\ref{sec:route1}.

\begin{proof}
That \eqref{item:r2-rh} implies \eqref{item:r2-real} is clear.
Assume \eqref{item:r2-real}; we may assume $g\ge1$. By
\eqref{eq:all-q}, $N_n\le q^n+1+Aq^{n/2}$ for all $n\ge n_0$, with
$A=(g+1)\sqrt p+2g$ and $n_0$ least with $q^{n_0/2}\ge2(g+1)$.
Apply Theorem~\ref{thm:dichotomy} with $n_1=1$ and
$W=\{n\ge n_0\}$. In case \eqref{item:dich-real}, $\rho>\sqrt q$ is
itself a reciprocal root. Also $\rho\le q$: otherwise apply
Proposition~\ref{prop:landau} to
$N_n=q^n+1-\sum_i\alpha_i^n\ge0$, whose frequencies of largest
modulus would be reciprocal roots, with negative coefficients. And
$\rho\ne q$ by
Theorem~\ref{thm:rationality}\eqref{item:rat-roots}. So $t=1/\rho$
is a real zero of $L_K$ in $(q^{-1},q^{-1/2})$, contradicting
\eqref{item:r2-real}.
\end{proof}

\begin{remark}\label{rem:missing-input}
By Theorem~\ref{thm:route2-input}, what Route~2 lacks is the
absence of real zeros of $L_K$ in $(q^{-1},q^{-1/2})$. This is a
statement about one real interval, and with the upper bound it is
equivalent to the Riemann hypothesis. So the one-sided bound
excludes every violation of the Riemann hypothesis except a real
one. The proof of Proposition~\ref{prop:idea-landau} has no step
that uses $\zeta(\sigma)\ne0$ for $0<\sigma<1$. There the bound is
in the other direction, and the sign of its residue excludes a real
zero. The counterpart of Proposition~\ref{prop:idea-landau} is
Theorem~\ref{thm:route2prime} below, which needs no input on real
zeros. Theorem~\ref{thm:route2-input} is the counterpart of
Landau's argument applied to an upper bound for $\Ch(x)-x$. For
that bound the residue at a real zero $\beta$ of $\zeta$, which is
$m/\beta$ for a zero of multiplicity $m$, has the sign that the
bound permits. So the absence of real zeros would have to be
supplied separately. Over a function field the corresponding
statement is the whole difficulty. The next proposition shows that
formal properties do not give it.
\end{remark}

\begin{proposition}[a formal counterexample]\label{prop:formal-cex}
Let $q=9$ and $L^*(t)=1-7t+9t^2=(1-\rho t)(1-\rho't)$, with
$\rho=(7+\sqrt{13})/2$ and $\rho'=9/\rho$, and put
$Z^*=L^*/((1-t)(1-9t))$ and $N_n^*=9^n+1-\rho^n-\rho'^n$. Then:
\begin{enumerate}
\item\label{item:cex-euler} $Z^*=\prod_d(1-t^d)^{-B_d^*}$ with
integers $B_d^*$, and $N_n^*=\sum_{d\mid n}dB_d^*$;
\item\label{item:cex-fe} $L^*$ satisfies the functional equation of
genus one, and $L^*(1)=3=N_1^*$;
\item\label{item:cex-pos} $N_n^*\ge0$ and $B_d^*\ge0$ for all
$n,d\ge1$;
\item\label{item:cex-upper} $N_n^*\le q^n+1$ for all $n\ge1$;
\item\label{item:cex-not-rh} the reciprocal roots
$\rho=5.30\ldots$, $\rho'=1.69\ldots$ are real, not of modulus $3$,
and $L^*(1/\rho)=0$ with $1/\rho\in(1/9,1/3)$.
\end{enumerate}
\end{proposition}

\begin{proof}
Proof of \eqref{item:cex-euler}: $L^*\in\Z[t]$ with $L^*(0)=1$, and
$(1-t)^{-1}(1-9t)^{-1}\in1+t\Z[[t]]$; an element of $1+t\Z[[t]]$
has a unique expansion $\prod_d(1-t^d)^{-B_d}$ with $B_d\in\Z$, and
taking $t\frac d{dt}\log$ gives the formula for $N_n^*$. Proof of
\eqref{item:cex-fe}: the coefficients $(1,-7,9)$ satisfy
$a_2=q\,a_0$, and $L^*(1)=3$, $N_1^*=10-7=3$. Proof of
\eqref{item:cex-pos}: $N_1^*=3$, and for $n\ge2$,
$N_n^*\ge9^n-2\rho^n>0$ since $\rho<5.31$; also
$B_1^*=N_1^*=3$. Let $d\ge2$ and assume $B_e^*\ge0$ for $e<d$, so
that $eB_e^*\le N_e^*$. Then
\[
dB_d^*=N_d^*-\sum_{e\mid d,\,e<d}eB_e^*
\ge N_d^*-\sum_{e\le d/2}N_e^*
\ge9^d-2(5.31)^d-2\cdot3^d>0,
\]
using $N_e^*\le9^e+1$; this is an induction on $d$. Proof of
\eqref{item:cex-upper}: $\rho,\rho'>0$. Part
\eqref{item:cex-not-rh} is arithmetic.
\end{proof}

No curve has this zeta function: by Theorem~\ref{thm:rh} a curve of
genus one over $\F_9$ has at least $4$ points. The proposition
shows that no argument using only rationality, the functional
equation, integrality, nonnegativity and one-sided upper bounds for
point counts can prove the Riemann hypothesis
(Figure~\ref{fig:counterexample-roots}). Bombieri's example
$\omega_1=q$, $\omega_2=1$, with all point counts $0$
\cite[p.~239]{Bombieri1974}, makes the same point; it has class
number $L(1)=0$, which Proposition~\ref{prop:formal-cex} avoids.
\ffbookfloat{counterexample-roots}

\begin{remark}\label{rem:route2}
The one-sided bound $N_n\le q^n+1+Aq^{n/2}$,
Proposition~\ref{prop:landau}, the functional equation and the
tower of constant field extensions do not together prove the
Riemann hypothesis without the averaging of
Section~\ref{sec:route1}. These four inputs prove only ``the
Riemann hypothesis or a real reciprocal root in $(\sqrt q,q)$''
(Theorems~\ref{thm:dichotomy} and~\ref{thm:route2-input}). The
polynomial $L^*$ of Proposition~\ref{prop:formal-cex} has all of
the formal properties those theorems use, and it realises the
second case. The averaging of Section~\ref{sec:route1}, or some
other input specific to curves, is needed.
\end{remark}

\subsection{The centred smoothed sum}

The number-field argument does not bound the Chebyshev function. It
bounds the normalised smoothed sum
$\Gz(x)=\sum_{n\le x}\Lambda(n)(x/n-1)-x\log x+(1+\gamma)x$ of
Definition~\ref{def:sz-gz} from above, and feeds the bound into a
Landau theorem. Here is the exact function-field counterpart, with
$x=q^T$ and with $N_n$ in the role of the sum of $\Lambda$ over the
elements of degree $n$.

\begin{definition}\label{def:smoothed-sum}
Let $K/\F_q$ be a function field of genus $g$ with reciprocal roots
$\alpha_i$. For integers $n,T\ge1$ and a complex number
$\alpha\ne1,q$ put
\begin{enumerate}
\item $E_n=N_n-q^n-1$ and $\calE_n=\sum_{j\le n}E_j$;
\item $S(T)=\sum_{n=1}^TN_n(q^{T-n}-1)$;
\item $\gamma_K=\sum_i\frac1{\alpha_i-1}$ and
$\delta_K=\sum_i\frac{\alpha_i}{\alpha_i-1}$;
\item $c(\alpha)=\frac{(q-1)\alpha}{(q-\alpha)(\alpha-1)}$;
\item $G(T)=S(T)-(T-1)(q^T-1)+\gamma_Kq^T+\delta_K$, the
\emph{centred smoothed sum} of $K$.
\end{enumerate}
\end{definition}

Here $\delta_K=\gamma_K+2g=L_K'(1)/L_K(1)$, the number $\gamma_K$
plays the role of Euler's constant, and $(T-1)(q^T-1)$ plays the
role of $x\log x-x$. The numbers $\gamma_K$ and $\delta_K$ are real,
because the multiset of reciprocal roots is stable under complex
conjugation.

\begin{proposition}[explicit formula]\label{prop:explicit-G}
Let $K/\F_q$ be a function field of genus $g$ with reciprocal roots
$\alpha_i$ and centred smoothed sum $G(T)$, and for complex
$\alpha\ne1,q$ let
$c(\alpha)=\frac{(q-1)\alpha}{(q-\alpha)(\alpha-1)}$. For all
$T\ge1$,
\begin{equation}\label{eq:explicit-G}
G(T)=\sum_{i=1}^{2g}c(\alpha_i)\,\alpha_i^T.
\end{equation}
Also $c(\alpha)=c(q/\alpha)$, $c(\alpha)>0$ for real
$1<\alpha<q$, and $c(\alpha)=(q-1)/\lvert1-\alpha\rvert^2$ if
$\lvert\alpha\rvert=\sqrt q$.
\end{proposition}

\begin{proof}
Sum geometric series:
$\sum_{n=1}^T(q^n+1)(q^{T-n}-1)=(T-1)(q^T-1)$ and
\[
\sum_{n=1}^T\alpha^n(q^{T-n}-1)=\frac{\alpha}{q-\alpha}\,q^T
+\frac{\alpha}{\alpha-1}-c(\alpha)\,\alpha^T.
\]
Sum the second identity over the $\alpha_i$ and use
$\sum_i\frac{\alpha_i}{q-\alpha_i}=\gamma_K$, which holds because
$\alpha\mapsto q/\alpha$ permutes the roots. The properties of $c$
follow from $c(\alpha)=(q-1)/(q+1-\alpha-q/\alpha)$.
\end{proof}

\begin{theorem}[a one-sided bound on the smoothed sum suffices]
\label{thm:route2prime}
Let $K/\F_q$ be a function field with reciprocal roots $\alpha_i$
and centred smoothed sum $G(T)$, let $A>0$, and let $W$ be a set of
positive integers containing arbitrarily long intervals. If
$G(T)\le Aq^{T/2}$ for all $T\in W$,
then $\lvert\alpha_i\rvert=\sqrt q$ for all $i$.
\end{theorem}

\begin{proof}
Suppose $\rho=\max_i\lvert\alpha_i\rvert>\sqrt q$. Then
$a_T=Aq^{T/2}-\sum_ic(\alpha_i)\alpha_i^T\ge0$ on $W$. The
frequencies of modulus $\rho$ are reciprocal roots $\beta$, with
coefficients $-\mu_\beta c(\beta)\ne0$. If $\sqrt q$ is a reciprocal
root, the coefficient $A-\mu_{\sqrt q}\,c(\sqrt q)$ of the frequency
$\sqrt q$ may vanish; in that case the frequency is omitted, which
does not affect the frequencies of modulus $\rho>\sqrt q$. By
Proposition~\ref{prop:landau}, $\rho$ is a reciprocal root and
$-\mu_\rho c(\rho)>0$, so $c(\rho)<0$. Also $\rho\le q$: otherwise
apply Proposition~\ref{prop:landau} to
$N_n=q^n+1-\sum_i\alpha_i^n\ge0$, whose frequencies of largest
modulus would be reciprocal roots, with negative coefficients. And
$\rho\ne q$ by
Theorem~\ref{thm:rationality}\eqref{item:rat-roots}. So
$\sqrt q<\rho<q$, and $c(\rho)>0$ by
Proposition~\ref{prop:explicit-G}, a contradiction. Hence
$\rho\le\sqrt q$, and the functional equation gives equality.
\end{proof}

Applied to $K_{n_1}$, whose reciprocal roots are the
$\alpha_i^{n_1}$ (Proposition~\ref{prop:tower}), the theorem works
along the tower. So the analytic step has an exact and valid
function-field form: for the centred smoothed sum, real reciprocal
roots count with a positive sign, and an upper bound alone
suffices. It remains to say what that upper bound is in terms of
point counts.

Section~\ref{sec:upper} also proves a one-sided upper bound. The
sum $S(T)$ is a combination of the $N_n$ with nonnegative weights
$q^{T-n}-1$, so an upper bound for each $N_n$ gives an upper bound
for $S(T)$, and $G(T)$ differs from $S(T)$ by terms that do not
depend on the point counts beyond the constants
$\gamma_K,\delta_K$. The two upper bounds are nevertheless different
statements, with opposite signs (Remark~\ref{rem:smoothed-sign}).

\begin{proposition}[which bound is asked for]
\label{prop:which-bound}
Let $K/\F_q$ be a function field with reciprocal roots $\alpha_i$
and centred smoothed sum $G(T)$, and put $E_n=N_n-q^n-1$,
$\calE_n=\sum_{j\le n}E_j$ and
$\delta_K=\sum_i\frac{\alpha_i}{\alpha_i-1}$. Assume
$\lvert\alpha_i\rvert<q$ for all $i$. Then for all $T\ge1$
\begin{equation}\label{eq:which-bound}
G(T)=\delta_K-\sum_{n\le T}E_n-\sum_{n>T}E_n\,q^{T-n}
=\delta_K-(1-q^{-1})\sum_{n\ge T}\calE_n\,q^{T-n}.
\end{equation}
\end{proposition}

The hypothesis $\lvert\alpha_i\rvert<q$ holds for every function
field by Theorem~\ref{thm:rh}, and it holds for the reciprocal roots
$\rho,\rho'$ of the polynomial $L^*$ of
Proposition~\ref{prop:formal-cex}.

\begin{proof}
For $\lvert\alpha\rvert<q$,
\[
\sum_{n\le T}\alpha^n+\sum_{n>T}\alpha^nq^{T-n}
=\frac{\alpha(\alpha^T-1)}{\alpha-1}
+\frac{\alpha^{T+1}}{q-\alpha}
=c(\alpha)\alpha^T-\frac\alpha{\alpha-1}.
\]
Summing over $i$ and using $E_n=-\sum_i\alpha_i^n$ and
\eqref{eq:explicit-G} gives the first form. For the second, Abel
summation gives
$\sum_{n>T}(\calE_n-\calE_{n-1})q^{T-n}
=-q^{-1}\calE_T+(1-q^{-1})\sum_{n>T}\calE_nq^{T-n}$; adding
$\sum_{n\le T}E_n=\calE_T$ gives
$(1-q^{-1})\calE_T+(1-q^{-1})\sum_{n>T}\calE_nq^{T-n}$.
\end{proof}

\begin{remark}\label{rem:smoothed-sign}
A one-sided upper bound for $G(T)$ is the counterpart of the bound
$\Gz(x)\le x^{1/2+\varepsilon}$ of Chapter~\ref{ch:idea}. It
corresponds to a lower bound for the point counts $N_n$, and the
Stepanov--Bombieri upper bound has the opposite sign. The upper
bound for $S(T)$ obtained from upper bounds for the $N_n$ is of
order $q^T$. If $\lvert\alpha_i\rvert<q$ for all reciprocal roots
$\alpha_i$, then
$\gamma_K=-\sum_{n\ge1}E_nq^{-n}$, so the centring term
$\gamma_Kq^T$ cancels the whole series
$\sum_{n\ge1}E_nq^{T-n}$ that carries this order. What is left
enters with the sign of $-E_n$, by the first form of
\eqref{eq:which-bound}. Let $A>0$ and put
$C_q=A\frac{\sqrt q+1}{\sqrt q-1}$.
Then the following hold.
\begin{enumerate}
\item If $E_n\le Aq^{n/2}$ for all $n$, which is the direction
Section~\ref{sec:upper} proves, then
$G(T)\ge\delta_K-C_qq^{T/2}$. This is a lower bound on $G$, and
Proposition~\ref{prop:landau} cannot use it.
\item If $E_n\ge-Aq^{n/2}$ for all $n$, a \emph{lower} bound on
point counts, then $G(T)\le\delta_K+C_qq^{T/2}$. This is the
hypothesis of Theorem~\ref{thm:route2prime}.
\item A real reciprocal root $\alpha\in(\sqrt q,q)$ contributes
$-\alpha^n<0$ to $E_n$ and $c(\alpha)\alpha^T>0$ to $G(T)$. An
upper bound on $E_n$ does not see it; an upper bound on $G(T)$
does. For the polynomial $L^*$ of
Proposition~\ref{prop:formal-cex}, the right side of
\eqref{eq:explicit-G} with the reciprocal roots $\rho,\rho'$ is
$G^*(T)=\frac83(\rho^T+\rho'^T)$.
\end{enumerate}
Figure~\ref{fig:smoothed-sign} shows the reversal.
\end{remark}
\fffloat{smoothed-sign}

\begin{remark}\label{rem:G-over-Q}
The same holds over $\Q$. Partial summation gives, with
$\Ch(u)=\sum_{n\le u}\Lambda(n)$,
\[
\sum_{n\le x}\Lambda(n)\Bigl(\frac xn-1\Bigr)
=x\int_1^x\frac{\Ch(u)}{u^2}\,du,
\]
so that
\[
\Gz(x)=x\Bigl(\int_1^x\frac{\Ch(u)-u}{u^2}\,du+1+\gamma\Bigr).
\]
The classical estimate
$\sum_{n\le x}\Lambda(n)/n=\log x-\gamma+o(1)$ gives
$\int_1^\infty(\Ch(u)-u)u^{-2}\,du=-1-\gamma$, hence
\[
\Gz(x)=-x\int_x^\infty\frac{\Ch(u)-u}{u^2}\,du.
\]
An upper bound on $\Gz$ is therefore a tail-averaged lower bound on
$\Ch(u)-u$.
\end{remark}

\subsection{Summary of Route 2}

With the bound that Section~\ref{sec:upper} provides, an upper
bound on $N_n$, Route~2 proves ``the Riemann hypothesis or a real
zero of $L_K$ in $(q^{-1},q^{-1/2})$'' and no more
(Theorems~\ref{thm:dichotomy} and~\ref{thm:route2-input},
Proposition~\ref{prop:formal-cex}). With an upper bound on the
centred smoothed sum it proves the Riemann hypothesis
(Theorem~\ref{thm:route2prime}), but that bound is a lower bound
for point counts (Proposition~\ref{prop:which-bound}). Route~1
supplies the lower bound, and with it both the condition of
Theorem~\ref{thm:route2-input}\eqref{item:r2-real} and the
hypothesis of Theorem~\ref{thm:route2prime}. No Stepanov-type
argument giving either directly is known to us. This leaves two
questions.

\begin{question}\label{q:real-zero}
Let $K/\F_q$ be a function field. Is there a Stepanov-type proof,
without a Galois closure, that $L_K$ has no real zero in
$(q^{-1},q^{-1/2})$? It would complete Route~2.
\end{question}

\begin{question}\label{q:smoothed}
Let $K/\F_q$ be a function field with reciprocal roots $\alpha_i$
and centred smoothed sum $G(T)$, and let
$\gamma_K=\sum_i\frac1{\alpha_i-1}$. Is there a Stepanov-type upper
bound for $G(T)$ with $\gamma_K$ built in? It would be a
function-field model of the number-field argument.
\end{question}

\section{Stepanov's derivative form over a prime field}
\label{sec:variantJ}

Bombieri notes that the constructions of Stepanov and Schmidt
obtain vanishing of order at least $m$ through derivatives or
hyperderivatives of order up to $m-1$, with $m$ about $q^{1/2}$,
whereas his own argument does not use derivations
\cite[p.~235]{Bombieri1974}. The derivative form is closer in
structure to the chain of theta lifts of Chapter~\ref{ch:theta}, so
we prove it in the simplest case. The argument is a reconstruction
of the standard one; agreement in detail with \cite{Stepanov1969}
is not claimed.

In Bombieri's argument the multiplicity comes from one application
of Frobenius, Step~3 of Theorem~\ref{thm:aux-section}. The part of
the number-field argument that produces multiplicity is the forced
$p$-adic divisibility of periods along a chain of lifts
(Section~\ref{sec:theta}). Both make high-order jets vanish without
paying for each jet separately. In structure, the derivative form
of this section matches the chain more closely
(Remark~\ref{rem:frobenius-chain}).

Throughout this section $p$ is an odd prime, $f\in\F_p[X]$ is a
polynomial of odd degree $d$ with $3\le d<p$, and
\[
c=\frac{p-1}2,\qquad D=\frac{p-d}2,
\qquad\text{so that}\qquad cd=\frac{d-1}2\,p+D.
\]
In this section the letter $\ell$ denotes a positive integer, the
multiplicity, and not the dimension of a Riemann--Roch space.

\begin{definition}\label{def:Sa}
Let $p$ be an odd prime, $f\in\F_p[X]$ and $c=(p-1)/2$.
\begin{enumerate}
\item for $a=\pm1$, $S_a=\{x\in\F_p:f(x)^c=a\}$, and
$S_0=\{x\in\F_p:f(x)=0\}$;
\item $N^{\mathrm{aff}}$ is the number of points of the affine
curve $y^2=f(x)$ over $\F_p$.
\end{enumerate}
\end{definition}

By Euler's criterion
$N^{\mathrm{aff}}=2\lvert S_1\rvert+\lvert S_0\rvert$, and
$\lvert S_1\rvert+\lvert S_{-1}\rvert+\lvert S_0\rvert=p$.

\begin{definition}\label{def:aux-poly}
Let $p$ be an odd prime, $f\in\F_p[X]$ of odd degree $d$ with
$3\le d<p$, $c=(p-1)/2$ and $D=(p-d)/2$, and let $J,\ell\ge1$ be
integers. An \emph{auxiliary polynomial} is
\begin{equation}\label{eq:aux-poly}
r(X)=f(X)^\ell\sum_{j=0}^{J-1}
\bigl(r_j(X)+s_j(X)f(X)^c\bigr)X^{jp},
\end{equation}
with $r_j,s_j\in\F_p[X]$ of degree at most $D-1$. Its \emph{jet
coefficients} are defined by $r_j^{[0]}=r_j$, $s_j^{[0]}=s_j$ and
\[
r_j^{[k+1]}=(\ell-k)f'r_j^{[k]}+f\,(r_j^{[k]})',\qquad
s_j^{[k+1]}=(\ell-k+c)f's_j^{[k]}+f\,(s_j^{[k]})'.
\]
\end{definition}

\begin{proposition}\label{prop:J-nonvanishing}
Let $J,\ell\ge1$ be integers and let $r$ be an auxiliary polynomial
\eqref{eq:aux-poly} for $f$, $J$ and $\ell$, with coefficients
$r_j,s_j$. If the $r_j,s_j$ are not all zero, then $r\ne0$.
\end{proposition}

\begin{proof}
A nonzero term $r_jX^{jp}$ has degree $jp+\deg r_j$, with residue
modulo $p$ in $[0,D-1]$ and quotient $j$. A nonzero term
$s_jf^cX^{jp}$ has degree $(j+\frac{d-1}2)p+D+\deg s_j$, with
residue in $[D,2D-1]$ (note $2D-1<p$) and quotient
$j+\frac{d-1}2$. Two distinct nonzero terms of the same kind differ
in the quotient, and two of different kinds differ in the residue.
So all nonzero terms have distinct degrees, their sum is nonzero,
and so is its product with $f^\ell$.
\end{proof}

This is the counterpart of Proposition~\ref{prop:independence}:
degrees, that is pole orders at infinity, are separated modulo $p$
(Figure~\ref{fig:J-degrees}).
\fffloat{J-degrees}

\begin{proposition}\label{prop:J-jets}
Let $J,\ell\ge1$ be integers and let $r$ be an auxiliary polynomial
\eqref{eq:aux-poly} for $f$, $J$ and $\ell$, with coefficients
$r_j,s_j$ and jet coefficients $r_j^{[k]},s_j^{[k]}$. For
$0\le k\le\ell$ the $k$-th
derivative of $r$ is
\[
r^{(k)}=f^{\ell-k}\sum_j\bigl(r_j^{[k]}+s_j^{[k]}f^c\bigr)X^{jp}.
\]
The polynomials $r_j^{[k]},s_j^{[k]}$ have degree at most
$D-1+k(d-1)$ and depend linearly on $(r_j,s_j)_j$.
\end{proposition}

\begin{proof}
Differentiate $f^{\ell-k}(u+vf^c)X^{jp}$ and use $(X^{jp})'=0$:
the derivative of $f^{\ell-k}u$ is
$f^{\ell-k-1}((\ell-k)f'u+fu')$, and that of $f^{\ell-k+c}v$ is
$f^{\ell-k-1}f^c((\ell-k+c)f'v+fv')$. Each step raises the degree
of a jet coefficient by at most $d-1$.
\end{proof}

\begin{theorem}\label{thm:variantJ}\label{def:tJB}
Let $p$ be an odd prime, let $f\in\F_p[X]$ be a polynomial of odd
degree $d$ with $3\le d<p$, and put $c=(p-1)/2$ and $D=(p-d)/2$. Let
$\ell$ be an even integer with $2\le\ell\le p-d-1$, and put
\[
t=\Bigl\lfloor\frac{\ell^2d}{2(p-d-\ell)}\Bigr\rfloor+1,\qquad
B=d+\frac{(d+1)p-2d-2}{2\ell}+\frac{(t-1)p}{\ell}.
\]
For $a=\pm1$ let $S_a=\{x\in\F_p:f(x)^c=a\}$, and let
$N^{\mathrm{aff}}$ be the number
of points of the affine curve $y^2=f(x)$ over $\F_p$. Then for
both values of $a$
\begin{equation}\label{eq:variantJ}
\lvert S_a\rvert\le\frac p2+B
\le\frac p2+d+\frac{(d+1)p}{2\ell}+\frac{\ell d}2\cdot
\frac p{p-d-\ell},
\end{equation}
and $\lvert N^{\mathrm{aff}}-p\rvert\le2B+d$.
\end{theorem}

The sets $S_a$ and the count $N^{\mathrm{aff}}$ are those of
Definition~\ref{def:Sa}.

\begin{proof}
\emph{Step 1: a Siegel lemma.} This step chooses the coefficients
$r_j,s_j$ of an auxiliary polynomial $r$ for $f$, $J=\ell/2+t$ and
$\ell$. Put
$\sigma_k(X)=\sum_j(r_j^{[k]}(X)+a\,s_j^{[k]}(X))X^j$, of degree at
most $D+J-2+k(d-1)$. Require $\sigma_k=0$ as a polynomial for
$0\le k<\ell$. These are at most
$\ell(D+J-1)+\frac{\ell(\ell-1)}2(d-1)$ linear equations in $2JD$
unknowns. With $J=\ell/2+t$ the number of unknowns minus the number
of equations is
\[
t(p-d-\ell)-\frac{\ell^2d-\ell d-\ell}2
>t(p-d-\ell)-\frac{\ell^2d}2>0
\]
by the choice of $t$. So a nonzero solution exists.

\emph{Step 2: non-vanishing.} $r\ne0$ by
Proposition~\ref{prop:J-nonvanishing}.

\emph{Step 3: multiplicity.} This step shows that $r$ vanishes to
order $\ell$ on $S_a$. Let $x\in S_a$. Then $x^{jp}=x^j$ and
$f(x)^c=a$, so by Proposition~\ref{prop:J-jets},
$r^{(k)}(x)=f(x)^{\ell-k}\sigma_k(x)=0$ for $0\le k<\ell$. Since
$\ell<p$, $x$ is a zero of $r$ of order at least $\ell$: if
$(X-x)^e$ divided $r$ exactly with $e<\ell$, then
$r^{(e)}(x)=e!\cdot h(x)\ne0$ for the cofactor $h$.

\emph{Step 4: degree count.} We get
$\ell\lvert S_a\rvert\le\deg r\le\ell d+(D-1)+cd+(J-1)p$.
Substituting $D$, $c$ and $J$ gives the first inequality of
\eqref{eq:variantJ}; the second uses
$t-1\le\frac{\ell^2d}{2(p-d-\ell)}$.

For the point count,
$N^{\mathrm{aff}}=2\lvert S_1\rvert+\lvert S_0\rvert\le p+2B+d$,
and
$N^{\mathrm{aff}}=2p-2\lvert S_{-1}\rvert-\lvert S_0\rvert
\ge p-2B-d$.
\end{proof}

The lower bound for $N^{\mathrm{aff}}$ came from the upper bound
for $\lvert S_{-1}\rvert$, and
$2\lvert S_{-1}\rvert+\lvert S_0\rvert$ is the affine point count
of the quadratic twist $\epsilon y^2=f(x)$, with $\epsilon$ a
non-square. Together with the exact identity
$\lvert S_1\rvert+\lvert S_{-1}\rvert+\lvert S_0\rvert=p$, this is
Theorem~\ref{thm:averaging} with $G=\Z/2$.

\begin{remark}\label{rem:J-balance}
In Step~4 the degree of $r$ splits into three parts:
\begin{enumerate}
\item an overhead $D-1+cd\approx\frac{(d+1)p}2$, independent of
the chain length $\ell$;
\item a main term $(J-1)p\approx\frac\ell2p$, proportional to
$\ell$, which divided by $\ell$ is the true main term $p/2$;
\item a correction $(t-1)p$, forced by the quadratic cost
$\frac{\ell(\ell-1)}2(d-1)$ of the chain of $\ell$ jet conditions
in Step~1; it needs $t\gtrsim\ell^2d/(2p)$.
\end{enumerate}
Per unit of multiplicity the error is
$\frac{(d+1)p}{2\ell}+\frac{\ell d}2(1+o(1))$, balanced at
$\ell\asymp\sqrt p$, which gives
$\lvert S_a\rvert\le p/2+O(d\sqrt p)$
(Figure~\ref{fig:J-balance}). The constant is not Weil's. Minimising
$B$ over even $\ell$, the ratio of $2B+d$ to Weil's $(d-1)\sqrt p$ is
$2.78$ for $(p,d)=(1009,3)$, at $\ell=24$; $2.45$ for $(1009,5)$;
$2.54$ for $(10007,3)$; $2.21$ for $(10007,7)$; and $2.26$ for
$(100003,5)$.
\end{remark}
\ffbookfloat{J-balance}

\begin{remark}\label{rem:frobenius-chain}
Bombieri's $q_0$-th power step, a single application of Frobenius,
is not the structural counterpart of the multiplicity chain of
Alp\"oge's argument. Theorem~\ref{thm:variantJ} matches that chain
better. It has a chain of $\ell\approx\sqrt p$ jet conditions, an
overhead independent of $\ell$, a main term proportional to $\ell$
and a quadratic cost. In the number-field argument the errors of
the first kind do not grow with the number $m$ of lifts, the main
term is a sum of $m$ copies of an archimedean term, and the balance
falls at $m\approx\sqrt x$ (Remark~\ref{rem:exponent-half}).
Bombieri's argument has no chain length to balance. This remark
compares shapes and is not a theorem. Over general $\F_q$ the
derivative form needs Hasse derivatives \cite{Schmidt1976}.
\end{remark}

\section{The sign-only identity
  \texorpdfstring{$qN = M^k + R$}{qN = M\textasciicircum k + R}}
\label{sec:identity}

We now turn to the second argument. Alp\"oge's note mentions a proof
of the Weil and Deligne bounds over finite fields that ``uses only
a sign'', written as $qN = M^k + R$ \cite{AlpogeRemark2026}, and
says that the plan for the integers was to combine ``two quite
different ideas present in a toy setting''. We read the two ideas
as Stepanov's multiplicity, which is
Sections~\ref{sec:upper}--\ref{sec:variantJ}, and the sign-only
proof. The note calls the second one ``new and shocking'' and its
ideas ``Fourier-analytic''.

\subsection{Where the identity appears}\label{sec:source}

The identity appears twice in the note. In the introduction it is
called the sign-only proof of Weil and Deligne. In the timeline,
item~3 reads, in part, ``$qN = M^k + R$, using
$x_1 \cdot x_2 + \cdots + x_{k-1}\cdot x_k$''. A footnote, attached
to the phrase ``the definite sign of the left-hand side'', sketches
the proof in one paragraph:
\begin{quote}
sum over $h$ first for the left-hand side, with e.g.\ quadric
coefficients too Bertini gives a definite sign from middle
cohomology, one off the one on $M^k$, which is the zero-mode on the
right-hand side with $R$ square-root-quality up to an inch or so as
Katz would say by iterating Parseval $k$ times, opposite signs
meaning no cancellation of inverse roots, QED; vectorize to go up
dimensions.
\end{quote}
That paragraph is the source for
Sections~\ref{sec:identity}--\ref{sec:sign}. The note defines none
of $t$, $M$, $N$, and item~3 of the timeline has not been located
in the Lean archive \cite{AlpogeArchive2026}. So the statements and
proofs below are a reconstruction: the definitions are a reading of
the footnote, the proofs are supplied here, and neither has been
confirmed by the note's author. Where the reconstruction leaves
something unproved, this is said
(Section~\ref{subsec:not-checked}).

\subsection{Two ways to bound a point count}\label{sec:two-ways}

Both arguments aim at the Riemann hypothesis for curves over
$\F_q$, in the form
$\bigl|\#C(\F_q) - q - 1\bigr| \le 2g\sqrt q$, or equivalently that
every eigenvalue $\alpha$ of Frobenius on $H^1(C)$ has
$|\alpha| = \sqrt q$. The first proves it. The second derives the
corresponding bound for character sums from two hypotheses, of
which the second is not proved (Theorem~\ref{thm:sign}). They
proceed by different methods.

Stepanov's method is elementary: it uses the Riemann--Roch theorem
and nothing else. Write $m$ for the multiplicity, as in
Chapter~\ref{ch:idea}; it is the integer called $M$ in
Theorem~\ref{thm:aux-section}. The steps are these.
\begin{enumerate}[label=(\arabic*)]
\item Choose $m$ and build, by linear algebra in Riemann--Roch
  spaces, a nonzero function that vanishes to order at least $m$
  at every rational point. The construction uses the Frobenius
  $z \mapsto z^q$, which is the identity on residue fields of
  degree one, and the additivity of $z \mapsto z^{p^j}$.
  Vanishing to order $m$ comes from an $m$-th power, in place of
  $m$ derivatives.
\item A nonzero function has as many zeros as poles. Hence $m$
  times the number of rational points is at most the degree of the
  poles.
\item The degree of the poles grows with $m$, but slower than $m$
  times $q$, as long as $m$ is at most about $\sqrt q$. Dividing by
  $m$ and taking $m \approx \sqrt q$ gives
  $\#C(\F_q) \le q + O(g\sqrt q)$.
\end{enumerate}
The lower bound follows by applying the upper bound to twists
(Section~\ref{sec:route1}), and the bounds over all $\F_{q^n}$ give
$|\alpha| \le \sqrt q$ (Section~\ref{sec:rh}). The method rests on
\emph{multiplicity}: one function is forced to vanish to high order
at many points. More multiplicity gives a better bound, but costs
degree. The balance point is $m \approx \sqrt q$.

The sign-only proof uses no auxiliary function and no multiplicity.
It uses two inputs.
\begin{enumerate}[label=(\alph*)]
\item An identity $qN = M^k + R$, where $M$ is the sum to be
  bounded, $N$ is a weighted point count on a quadric, and $R$ is
  bounded by elementary Fourier analysis
  (Section~\ref{sec:parseval}).
\item A \emph{sign}: the Frobenius eigenvalues in $qN$ and in
  $M^k$ enter the identity with opposite signs, so they cannot
  cancel against each other (Section~\ref{sec:sign}).
\end{enumerate}
The sign is a statement about cohomology, which is why the proof is
not elementary in the way Stepanov's is.

\subsection{Notation}\label{sec:notation}

From here to the end of Section~\ref{sec:role} the letters $M$,
$N$, $R$, $k$ and $t$ have the meaning they have in the note's
identity. They no longer denote the multiplicity, a point count of
a function field, a radius and an exponent as in
Sections~\ref{sec:setting}--\ref{sec:variantJ}; in particular $M$
is now a complex number, and $k$ is the number of variables of a
quadratic form.

Let $q = p^f$. For $y \in \F_q$ set
\[
  e_q(y) \;=\; \exp\!\Bigl(\frac{2\pi i}{p}\,
  \Tr_{\F_q/\F_p}(y)\Bigr),
\]
where the trace $\Tr_{\F_q/\F_p}(y) \in \F_p = \Z/p\Z$ is lifted to
$\{0, \dots, p-1\}$. Then $e_q \colon \F_q \to \C^\times$ is a
nontrivial homomorphism from the additive group, and every additive
character of $\F_q$ is $y \mapsto e_q(ay)$ for a unique
$a \in \F_q$. The only fact used about $e_q$ is orthogonality:
\begin{equation}\label{eq:orth}
  \sum_{h \in \F_q} e_q(hy) \;=\;
  \begin{cases} q, & y = 0,\\ 0, & y \ne 0.\end{cases}
\end{equation}
For $q = p$ this is $e_p(y) = e^{2\pi i y/p}$, the classical case.

For $g \colon \F_q \to \C$ define the Fourier transform
$\widehat g(\xi) = \sum_{x\in\F_q} g(x)\,e_q(x\xi)$ and the norm
$\|g\|_2^2 = \sum_x |g(x)|^2$. By \eqref{eq:orth},
\begin{equation}\label{eq:parseval}
  \|\widehat g\|_2^2 \;=\; q\,\|g\|_2^2 \qquad\text{(Parseval).}
\end{equation}

\subsection{The function $t$}\label{subsec:t}

The note never says what $t$ is. The reading adopted here is that
$t$ is a \emph{trace function}: the function
$x \mapsto \Tr(\Fr_x \mid \cF_{\bar x})$ of a lisse $\ell$-adic
sheaf $\cF$ on an open subset of the affine line, extended by
zero, where $\Fr_x$ is the geometric Frobenius at $x$. There are
two reasons.
\begin{enumerate}[label=(\roman*)]
\item The note says the proof gives ``Weil and Deligne''. Weil's
  bound is the case $t(x) = \chi(f(x))$ or $t(x) = e_q(f(x))$, for
  a multiplicative character $\chi$ and a polynomial $f$. Deligne's
  bound is the case of a general pure sheaf.
\item The footnote speaks of ``inverse roots'' and ``middle
  cohomology''. Both are about the cohomology of a sheaf whose
  trace function is $t$.
\end{enumerate}
Then $M = \sum_x t(x)$ is the character sum to be bounded. For
example, for $t(x) = \chi(f(x))$ with $\chi$ quadratic, $M$ is
$\#C(\F_q) - q$ up to the points at infinity, where $C$ is the
curve $y^2 = f(x)$; with the sets $S_{\pm1}$ and the count
$N^{\mathrm{aff}}$ of Definition~\ref{def:Sa},
$M=\lvert S_1\rvert-\lvert S_{-1}\rvert=N^{\mathrm{aff}}-p$ when
$q=p$. In what follows only these
properties of $t$ are used: $|t| \le r$ for a constant $r$ (the
rank of $\cF$), and $t$ is the trace function of a sheaf to which
the trace formula applies.

\subsection{The identity}

\begin{definition}\label{def:sign-sums}
Let $k \ge 2$ be an integer. The \emph{path form} is
\[
  Q(x) \;=\; x_1x_2 + x_2x_3 + \cdots + x_{k-1}x_k ,
\]
the quadratic form of the path graph with vertices $1, \dots, k$
and edges $\{i, i+1\}$, and $Z \subset \mathbb{A}^k$ is the affine
quadric $Q = 0$. For a function $t\colon\F_q\to\C$ put
\begin{gather*}
  S = \sum_{h\in\F_q}\;\sum_{x\in\F_q^k}
  t(x_1)\cdots t(x_k)\; e_q\bigl(h\,Q(x)\bigr),\qquad
  M = \sum_{x\in\F_q} t(x),\\
  N = \sum_{x\in Z(\F_q)} t(x_1)\cdots t(x_k),\qquad
  R = \sum_{h \ne 0}\;\sum_{x\in\F_q^k}
  t(x_1)\cdots t(x_k)\, e_q\bigl(hQ(x)\bigr).
\end{gather*}
\end{definition}

The note writes the form as
$x_1 \cdot x_2 + \cdots + x_{k-1}\cdot x_k$, which can be read in
two ways. The path form is the reading adopted here; the other
reading is treated, and set aside, in Remark~\ref{rem:hyperbolic}.

\begin{proposition}[the identity]\label{prop:identity}
Let $t\colon\F_q\to\C$ be a function, let $k\ge2$, let $Q$ be the
path form in $k$ variables, and let $Z\subset\mathbb A^k$ be the
quadric $Q=0$. Then
\begin{align*}
  &\sum_{h\in\F_q}\;\sum_{x\in\F_q^k}
  t(x_1)\cdots t(x_k)\,e_q\bigl(hQ(x)\bigr)
  = q\sum_{x\in Z(\F_q)}t(x_1)\cdots t(x_k)\\
  &\qquad= \Bigl(\sum_{x\in\F_q}t(x)\Bigr)^k
  + \sum_{h\ne0}\;\sum_{x\in\F_q^k}
  t(x_1)\cdots t(x_k)\,e_q\bigl(hQ(x)\bigr).
\end{align*}
\end{proposition}

With the sums of Definition~\ref{def:sign-sums}, the proposition
reads $S = qN = M^k + R$.

\begin{proof}
Summing over $h$ first and using \eqref{eq:orth} with $y = Q(x)$
keeps only the $x$ with $Q(x) = 0$, each with weight $q$. So
$S = qN$. Splitting $S$ by $h$, the term $h = 0$ is
$\sum_x t(x_1)\cdots t(x_k) = M^k$ and the rest is $R$.
\end{proof}

So the five symbols are:
\begin{itemize}
\item $q$, the size of the field. It enters as the value of the
  $h$-sum in \eqref{eq:orth}.
\item $N$, the point count of the quadric $Z$, with each point
  weighted by $t(x_1)\cdots t(x_k)$.
\item $M$, the sum whose bound is wanted. $M^k$ is the \emph{zero
  mode}: the frequency $h = 0$.
\item $k$, the number of variables. It is a free parameter, and the
  proof sends it to infinity.
\item $R$, the frequencies $h \ne 0$.
\end{itemize}
The identity itself is orthogonality of characters and nothing
else; it holds for every function $t$. Its content is that the
three terms can be controlled by three different means: $N$ by the
cohomology of a quadric, $M^k$ by the cohomology of a product, and
$R$ by Fourier analysis.

\section{Iterated Parseval: the bound on
  \texorpdfstring{$R$}{R}}
\label{sec:parseval}

The note says that $R$ is of ``square-root quality'' by ``iterating
Parseval $k$ times''. The following is one way to carry this out.

The Fourier transform is the one on the additive group of $\F_q$
defined in Section~\ref{sec:notation}. It will be applied to a
sequence of functions $g_1, g_2, \dots$ on $\F_q$, built from $t$
one variable at a time, starting with $g_1 = t$. This is possible
because the phase $e_q(hQ(x))$ factors along the edges of the path:
\[
  e_q\bigl(hQ(x)\bigr) = e_q(hx_1x_2)\,e_q(hx_2x_3)\cdots
  e_q(hx_{k-1}x_k).
\]
The variable $x_1$ occurs only in the first factor. So for fixed
$x_2$,
\[
  \sum_{x_1} t(x_1)\,e_q(hx_1x_2) = \widehat t(hx_2),
\]
which is the Fourier transform of $t$ at the point $hx_2$.

\subsection{The case $k = 3$}

Here $Q = x_1x_2 + x_2x_3$, and $x_1$ and $x_3$ each meet only
$x_2$. Summing over both of them,
\[
  \sum_{x_1,x_2,x_3} t(x_1)t(x_2)t(x_3)\,
  e_q\bigl(h(x_1x_2+x_2x_3)\bigr)
  = \sum_{x_2} t(x_2)\,\widehat t(hx_2)^2.
\]
If $|t| \le c$ and $h\ne0$, then the triangle inequality and
Parseval give
\[
  \Bigl|\sum_{x_2} t(x_2)\,\widehat t(hx_2)^2\Bigr|
  \le c\,\sum_{y}|\widehat t(y)|^2 = c\,q\,\|t\|_2^2 \le c^3 q^2,
\]
using that $y = hx_2$ runs over $\F_q$ once. The trivial bound is
$c^3q^3$. So one power of $q$ is saved, and no Gauss sum was
evaluated. For general $k$ one sums from one end and uses Parseval
at each step.

\subsection{The general case}

\begin{proposition}[iterated Parseval]\label{lem:parseval}
Let $t\colon\F_q\to\C$ be a function with $|t(x)| \le c$ for all
$x$, let $k\ge2$, let $Q$ be the path form in $k$ variables, and let
$R=\sum_{h\ne0}\sum_{x\in\F_q^k}t(x_1)\cdots t(x_k)\,
e_q\bigl(hQ(x)\bigr)$. For each $h \in\F_q$ with $h\ne 0$,
\[
  \Bigl|\sum_{x\in\F_q^k} t(x_1)\cdots t(x_k)\,
  e_q\bigl(hQ(x)\bigr)\Bigr| \;\le\; c^k\, q^{(k+1)/2}.
\]
Hence $|R| \le c^k q^{(k+3)/2}$.
\end{proposition}

\begin{proof}
Because $Q$ is a path, the variables can be summed one at a time
from one end. Put $g_1 = t$ and, for $2 \le j \le k$,
\[
  g_j(y) \;=\; t(y)\sum_{x\in\F_q} g_{j-1}(x)\,e_q(hxy)
  \;=\; t(y)\,\widehat{g_{j-1}}(hy).
\]
Summing over $x_1$, then $x_2$, and so on up to $x_{k-1}$ shows, by
induction on $j$, that
\[
  \sum_{x_1,\dots,x_{j-1}} t(x_1)\cdots t(x_j)\,
  e_q\bigl(h(x_1x_2+\cdots+x_{j-1}x_j)\bigr) = g_j(x_j),
\]
so the sum in the proposition is $\sum_{y} g_k(y)$. Since
$h \ne 0$, the map $y \mapsto hy$ is a bijection of $\F_q$, so by
\eqref{eq:parseval}
\[
  \|g_j\|_2 \;\le\; c\,\|\widehat{g_{j-1}}\|_2
  \;=\; c\sqrt q\,\|g_{j-1}\|_2 .
\]
Iterating from $\|g_1\|_2 \le c\sqrt q$ gives
$\|g_{k-1}\|_2 \le c^{k-1} q^{(k-1)/2}$. Last, by Cauchy--Schwarz
and Parseval once more,
\begin{align*}
  \Bigl|\sum_y g_k(y)\Bigr|
  =\Bigl|\sum_y t(y)\,\widehat{g_{k-1}}(hy)\Bigr|
  &\le \|t\|_2\,\|\widehat{g_{k-1}}\|_2\\
  &\le c\sqrt q\cdot \sqrt q\, c^{k-1} q^{(k-1)/2}
  = c^k q^{(k+1)/2}.
\end{align*}
Summing over the $q - 1$ values of $h$ gives the bound for $R$.
\end{proof}

The trivial bound for $R$ is $c^kq^{k+1}$, so the proposition saves
a factor $q^{(k-1)/2}$. Compare it with $M^k$.
If $M$ has the square-root size $c\sqrt q$ that is to be proved,
then $|M^k|\le c^kq^{k/2}$, and the bound on $R$ exceeds this by the
fixed factor $q^{3/2}$, whatever $k$ is. In the reading adopted
here, this fixed loss is the note's ``up to an inch''. The last step
of the proof of Theorem~\ref{thm:sign} removes it by letting
$k \to \infty$.

The argument depends on the shape of the path. For a general
quadratic form one would diagonalise and evaluate Gauss sums. The
note says ``Gauss sum cancellation replaced by iterated Parseval'':
for the path, each variable meets only the next one, and Parseval suffices (Figure~\ref{fig:path-parseval}). The
same proof works for the form
$Q_a=\sum_ia_ix_ix_{i+1}$ with all $a_i\in\F_q^\times$, since it
only needs $ha_i\ne0$.

The note's form $x_1 \cdot x_2 + \cdots + x_{k-1}\cdot x_k$ can
also be read as the hyperbolic form
$x_1x_2+x_3x_4+\cdots+x_{k-1}x_k$ with $k$ even, in which every
variable occurs once. That form does not give a remainder of
square-root quality.
\begin{figure}[tbp]
\centering
\input{book/figs/ch02-path-parseval.tex}
\end{figure}

\begin{remark}[the hyperbolic reading]\label{rem:hyperbolic}
Let $t\colon\F_q\to\C$ be a function, let $k$ be even, let
$Q'=x_1x_2+x_3x_4+\cdots+x_{k-1}x_k$, and let
$R'=\sum_{h\ne0}\sum_{x\in\F_q^k}t(x_1)\cdots t(x_k)\,
e_q\bigl(hQ'(x)\bigr)$, the sum $R$ of
Definition~\ref{def:sign-sums} with $Q'$ in place of $Q$. For
$h\ne0$ the sum over $x$ is a product of $k/2$ equal factors:
\begin{gather*}
  \sum_{x\in\F_q^k} t(x_1)\cdots t(x_k)\,e_q\bigl(hQ'(x)\bigr)
  = B(h)^{k/2},\\
  B(h)=\sum_{x,y\in\F_q}t(x)\,t(y)\,e_q(hxy)
  =\sum_{y\in\F_q}t(y)\,\widehat t(hy).
\end{gather*}
If $|t|\le c$, the Cauchy--Schwarz inequality and
\eqref{eq:parseval} give
$|B(h)|\le\|t\|_2\,\|\widehat t\|_2\le c^2q^{3/2}$, so
\[
  |R'|\;\le\;c^k\,q^{3k/4+1}.
\]
The exponent cannot be improved. Let $t=\chi$ be a nontrivial
multiplicative character of $\F_q$, extended by $\chi(0)=0$, with
Gauss sum $\tau(\chi)=\sum_x\chi(x)e_q(x)$, of absolute value
$\sqrt q$. Then $\widehat\chi(\xi)=\overline\chi(\xi)\,\tau(\chi)$,
so $B(h)=\overline\chi(h)\,\tau(\chi)\,(q-1)$. If $\chi^{k/2}$ is
trivial, for instance if $\chi$ is quadratic and $4$ divides $k$,
then
\[
  R'=\tau(\chi)^{k/2}\,(q-1)^{k/2+1},\qquad
  |R'|=q^{k/4}\,(q-1)^{k/2+1}.
\]
So the remainder for $Q'$ exceeds the size $c^kq^{k/2}$ of $M^k$ by
a factor of order $q^{k/4+1}$, which grows with $k$, where the
path form loses the fixed factor $q^{3/2}$. If the proof of
Theorem~\ref{thm:sign} is run with $Q'$, under the hypothesis
corresponding to \ref{H2}, its Step~2 has only the bound
$r^kq^{n(3k/4+1)}$ over $\F_{q^n}$, and Step~5 gives
$|\alpha_j|\le q^{3/4+1/k}$, in the limit only
$|\alpha_j|\le q^{3/4}$. The path form is therefore
the reading that agrees with the footnote's ``square-root-quality''
and ``iterating Parseval $k$ times'': in it every variable except
the two ends lies on two edges, and each summation is a Fourier
transform.
\end{remark}

\section{The sign-only theorem}
\label{sec:sign}

This section states the sign-only bound in full and proves it. The
statement is a reading of the note's footnote, and the two
hypotheses under which it is proved are assumptions, not results:
the second of them is the content of the note's word ``Bertini''
and is not proved here.

\subsection{Trace functions and Frobenius eigenvalues}

Let $U \subseteq \mathbb{A}^1_{\F_q}$ be a nonempty open subset
with $0 \notin U$, and let $j \colon U \to \mathbb{A}^1$ be the
inclusion. Let $\cF$ be a lisse $\overline{\Q}_\ell$-sheaf on $U$
of rank $r$, for a prime $\ell\ne p$, which is pointwise pure of
weight $0$; and let $t$ be its trace function, extended by zero to
$\F_q$. So $t$ is the trace function of $j_!\cF$, and
$|t(x)| \le r$ for all $x$, because the eigenvalues of $\Fr_x$ on
the stalk have absolute value $1$. For $n\ge1$ let $t_n$ be the
trace function over $\F_{q^n}$, and let $M_n$, $N_n$, $R_n$ be the
sums of Definition~\ref{def:sign-sums} over $\F_{q^n}$, with
$t_n$ in place of $t$ and $e_{q^n}$ in place of $e_q$. Again
$|t_n|\le r$.

Write $\Fb$ for an algebraic closure of $\F_q$ and $\Fr$ for the
geometric Frobenius of $\F_q$ acting on cohomology with compact
supports. The sheaf $(j_!\cF)^{\boxtimes k}$ on $\mathbb A^k$ has
trace function $x\mapsto t(x_1)\cdots t(x_k)$. By the
Grothendieck--Lefschetz trace formula,
\begin{equation}\label{eq:sign-trace}
\begin{aligned}
  M_n &= \sum_i (-1)^i \Tr\bigl(\Fr^n \mid
  H^i_c(\mathbb{A}^1_{\Fb}, j_!\cF)\bigr),\\
  N_n &= \sum_i (-1)^i \Tr\bigl(\Fr^n \mid
  H^i_c(Z_{\Fb}, (j_!\cF)^{\boxtimes k}|_Z)\bigr).
\end{aligned}
\end{equation}
The eigenvalues of $\Fr$ on these groups are the ``inverse roots''
of the note: the reciprocals of the zeros and poles of the
corresponding $L$-functions. They are the counterparts of the
reciprocal roots $\alpha_i$ of Theorem~\ref{thm:rationality}. In
the rest of this chapter
$d=\dim H^1_c(\mathbb{A}^1_{\Fb}, j_!\cF)$.

\subsection{The two hypotheses}

The proof assumes that the cohomology on each side lives in a
single degree:
\begin{enumerate}[label=(H\arabic*)]
\item\label{H1} $H^i_c(\mathbb{A}^1_{\Fb}, j_!\cF) = 0$ for
  $i \ne 1$;
\item\label{H2} for infinitely many $k$,
  $H^i_c\bigl(Z_{\Fb}, (j_!\cF)^{\boxtimes k}|_Z\bigr) = 0$ for
  $i \ne k - 1$.
\end{enumerate}
Hypothesis~\ref{H1} holds when $\cF$ is geometrically irreducible
and not geometrically constant: $H^0_c$ vanishes because $U$ is an
affine curve and $\cF$ is lisse, and $H^2_c$ is, up to a twist, the
coinvariants of the geometric monodromy. Hypothesis~\ref{H2} is the
``definite sign from middle cohomology'' of the note: $Z$ has
dimension $k-1$, and $H^{k-1}_c$ is its middle degree. The note
says that Bertini gives it, possibly after giving the quadric
generic coefficients, $Q_a = \sum_i a_i x_ix_{i+1}$. Neither a
proof of \ref{H2} nor a reference for it is given here. Half of it
is automatic (Remark~\ref{rem:h2-half}), and the other half fails
for some $k$ in simple cases (Remark~\ref{rem:h2-fails}); see also
Remark~\ref{rem:bertini} and Section~\ref{subsec:not-checked}.

\subsection{A growth lemma}

\begin{proposition}\label{lem:growth}
Let $\gamma_1, \dots, \gamma_s \in \C$ be distinct, let
$m_1, \dots, m_s \in \C$ be nonzero, and let $A,C>0$. If
$\bigl|\sum_i m_i\gamma_i^n\bigr| \le A\,C^n$ for all $n \ge 1$,
then $|\gamma_i| \le C$ for all $i$.
\end{proposition}

\begin{proof}
For $|z|$ small,
\[
  \sum_{n\ge1}\Bigl(\sum_i m_i\gamma_i^n\Bigr) z^n
  = \sum_i \frac{m_i\gamma_i z}{1 - \gamma_i z}.
\]
By the hypothesis the left side converges for $|z| < 1/C$, so it
defines an analytic function there, which agrees with the rational
function on the right. The right side has a pole at each
$1/\gamma_i$ with $\gamma_i \ne 0$, since the $\gamma_i$ are
distinct and $m_i \ne 0$. So every pole lies in $|z| \ge 1/C$,
that is, $|\gamma_i| \le C$.
\end{proof}

The proposition needs the $m_i$ to be nonzero after the equal
$\gamma_i$ have been collected. A sum
$\sum_i m_i\gamma_i^n$ in which large $\gamma$ occur with
coefficients of both signs can be small for every $n$, because the
coefficients of one value of $\gamma$ can add up to zero.

\subsection{The theorem}

\begin{theorem}[Sign-only bound]\label{thm:sign}
Let $U \subseteq \mathbb{A}^1_{\F_q}$ be a nonempty open subset
with $0 \notin U$, with inclusion $j\colon U\to\mathbb A^1$. Let
$\cF$ be a lisse $\overline{\Q}_\ell$-sheaf on $U$ of rank $r$, for
a prime $\ell\ne p$, pointwise pure of weight $0$, and let $t$ be its
trace function, extended by zero to $\F_q$. For $k\ge2$ let $Z$ be the affine
quadric $x_1x_2 + \cdots + x_{k-1}x_k = 0$ in $\mathbb{A}^k$.
Assume \ref{H1} and \ref{H2}. Then every eigenvalue $\alpha$ of
$\Fr$ on $H^1_c(\mathbb{A}^1_{\Fb}, j_!\cF)$ satisfies
$|\alpha| \le \sqrt q$. Consequently
\[
  \Bigl|\sum_{x\in\F_q} t(x)\Bigr| \;\le\;
  \sqrt q\,\dim H^1_c(\mathbb{A}^1_{\Fb}, j_!\cF).
\]
\end{theorem}

\begin{proof}
Let $\alpha_1, \dots, \alpha_d$ be the eigenvalues of $\Fr$ on
$H^1_c(\mathbb{A}^1_{\Fb},j_!\cF)$, and, for $k$ allowed by \ref{H2},
let $\beta_1, \dots, \beta_e$ be those on
$H^{k-1}_c(Z_{\Fb}, (j_!\cF)^{\boxtimes k}|_Z)$, both with
multiplicity. Fix such a $k$. The proof has five steps.

\emph{Step 1: the identity.} Over $\F_{q^n}$,
Proposition~\ref{prop:identity} gives $q^nN_n = M_n^k + R_n$, for
every $n\ge1$.

\emph{Step 2: the bound on $R$.}
Proposition~\ref{lem:parseval} over $\F_{q^n}$, with $c = r$, gives
$|R_n| \le r^k q^{n(k+3)/2}$ for all $n \ge 1$.

\emph{Step 3: eigenvalues.} By \eqref{eq:sign-trace} and \ref{H1},
$M_n = -\sum_j \alpha_j^n$, so
\[
  M_n^k = (-1)^k \sum_{I} \alpha_I^n,\qquad
  \alpha_I = \alpha_{i_1}\cdots\alpha_{i_k},
\]
the sum over all $I = (i_1,\dots,i_k) \in \{1,\dots,d\}^k$. By
\eqref{eq:sign-trace} and \ref{H2},
$N_n = (-1)^{k-1}\sum_j \beta_j^n$, so
$q^nN_n = (-1)^{k-1}\sum_j (q\beta_j)^n$.

\emph{Step 4: the sign.} The degrees $k$ and $k-1$ differ by one.
This is the note's ``one off''. Hence
\[
  R_n \;=\; q^nN_n - M_n^k
  \;=\; (-1)^{k-1}\Bigl(\sum_j (q\beta_j)^n
  + \sum_I \alpha_I^n\Bigr).
\]
Every eigenvalue enters with the same sign. After grouping equal
values, $R_n = (-1)^{k-1}\sum_\gamma m_\gamma \gamma^n$ with
distinct complex numbers $\gamma$ and integers $m_\gamma \ge 1$,
the same for all $n$; and every $\alpha_I$ is one of the $\gamma$.
No term cancels: this is the note's ``opposite signs meaning no
cancellation of inverse roots''.

\emph{Step 5: comparison and $k \to \infty$.} By Step~2 and
Proposition~\ref{lem:growth}, applied with $A=r^k$ and
$C=q^{(k+3)/2}$, we get $|\gamma| \le q^{(k+3)/2}$ for every
$\gamma$ that occurs in $R_n$. By Step~4, every $\alpha_I$ occurs.
Take $I = (j, \dots, j)$:
\[
  |\alpha_j|^k \le q^{(k+3)/2},\qquad\text{so}\qquad
  |\alpha_j| \le q^{1/2 + 3/(2k)} .
\]
Let $k \to \infty$ through the values allowed by \ref{H2}. This
gives $|\alpha_j| \le \sqrt q$. Last,
$\sum_x t(x) = M_1 = -\sum_j \alpha_j$ by \ref{H1}, so
$|M_1| \le d\sqrt q$.
\end{proof}

\begin{remark}[what the sign does]\label{rem:sign-positivity}
The positivity of the $m_\gamma$ is what is used. In general a sum
of eigenvalues of size $q^{k}$ could cancel down to the size of
$R$. The sign forbids this. Purity of $\cF$ is used only to bound
$|t_n|$ by $r$; no bound on the size of an eigenvalue on cohomology
is assumed.
\end{remark}

\begin{remark}[Weil and Deligne]\label{rem:weil-deligne}
For $t(x)=\chi(f(x))$ or $t(x)=e_q(f(x))$, with the corresponding
Kummer or Artin--Schreier sheaf on the open set where it is lisse,
the conclusion of Theorem~\ref{thm:sign} is the Weil bound for the
character sum; for a general lisse sheaf, pure of weight $0$, it is
Deligne's bound $|\alpha|\le\sqrt q$ for the eigenvalues on
$H^1_c$. This is the reading of ``Weil and Deligne'' in the note.
In each case the theorem gives the bound only if \ref{H1} and
\ref{H2} hold for the sheaf in question, and for the plain path
quadric \ref{H2} can fail for particular $k$ in exactly these
cases (Remark~\ref{rem:h2-fails}).
\end{remark}

\begin{remark}[the condition $0\notin U$]\label{rem:origin}
The condition $0 \notin U$ is not used in the deduction of
Theorem~\ref{thm:sign} from \ref{H1} and \ref{H2}. It is there for
the sake of \ref{H2}: the quadric $Z$ is singular, and its singular
points lie in the hyperplane $x_2=0$, so the condition keeps them
out of the support of the sheaf (Remark~\ref{rem:h2-half}). It
costs nothing. Removing the point $0$ from $U$ changes
$\sum_xt(x)$ by one term of size at most $r$. It also does not
lose eigenvalues: with $U'=U\smallsetminus\{0\}$ and
$j'\colon U'\to\mathbb A^1$, the exact sequence
$0\to j'_!(\cF|_{U'})\to j_!\cF\to\cF_0\to0$, in which $\cF_0$ is
a skyscraper sheaf at the origin, shows that
$H^1_c(\mathbb A^1_{\Fb},j_!\cF)$ is a quotient of
$H^1_c(\mathbb A^1_{\Fb},j'_!(\cF|_{U'}))$ and that \ref{H1} for
$U$ implies \ref{H1} for $U'$.
\end{remark}

\begin{remark}[half of \ref{H2} is automatic]\label{rem:h2-half}
The gradient of $Q$ at $x$ has first coordinate $x_2$, so every
singular point of $Z$ has $x_2=0$. Since $0\notin U$, the variety
$Z\cap U^k$ is therefore smooth, of dimension $k-1$ if it is not
empty, and it is affine. The sheaf $(j_!\cF)^{\boxtimes k}|_Z$ is
the extension by zero of a lisse sheaf on $Z\cap U^k$. By Artin's
vanishing theorem and Poincar\'e duality, the cohomology with
compact supports of a lisse sheaf on a smooth affine variety of
dimension $k-1$ vanishes in degrees below $k-1$. So the part of
\ref{H2} with $i<k-1$ always holds, and \ref{H2} is the statement
that the cohomology vanishes in the degrees $k,\dots,2k-2$ above
the middle.
\end{remark}

The note says ``Bertini gives a definite sign from middle
cohomology''. For the path quadric as it stands, \ref{H2} can fail
even when $\cF$ satisfies \ref{H1}.

\begin{remark}[\ref{H2} fails for the plain path quadric at $k=3$]
\label{rem:h2-fails}
For $k=3$ the form is $Q=x_2(x_1+x_3)$, and $t(0)=0$, so only the
points with $x_3=-x_1$ contribute to $N$:
\[
  N=M\cdot\sum_{x\in\F_q}t(x)\,t(-x).
\]
Take for $\cF$ the Artin--Schreier sheaf on
$U=\mathbb A^1\smallsetminus\{0\}$ with trace function
$t(x)=e_q(x)$ for $x\ne0$. It has rank one and weight $0$, and it
satisfies \ref{H1} with $d=1$ and $\alpha_1=1$, since
$M_n=\sum_{x\ne0}e_{q^n}(x)=-1$. Here $t_n(x)t_n(-x)=1$ for
$x\ne0$, so
\[
  N_n=-(q^n-1)=1^n-q^n\qquad\text{for all }n\ge1.
\]
The eigenvalues $1$ and $q$ occur with opposite signs, so the
cohomology of $Z$ is not concentrated in one degree, and \ref{H2}
fails for $k=3$.

The general statement at $k=3$ is as follows. Since $0\notin U$,
the map $(x_1,x_2)\mapsto(x_1,x_2,-x_1)$ identifies
$(U\cap-U)\times U$ with $Z\cap U^3$, and the sheaf
$(j_!\cF)^{\boxtimes3}|_Z$ with the external product of
$\cF\otimes[-1]^*\cF$ on $U\cap-U$ and $\cF$ on $U$, where $[-1]$
is the map $x\mapsto-x$. By the K\"unneth formula and \ref{H1},
\[
  H^i_c\bigl(Z_{\Fb},(j_!\cF)^{\boxtimes3}|_Z\bigr)
  = H^{i-1}_c\bigl((U\cap-U)_{\Fb},\cF\otimes[-1]^*\cF\bigr)
  \otimes H^1_c\bigl(\mathbb A^1_{\Fb},j_!\cF\bigr).
\]
The first factor vanishes outside the degrees $1$ and $2$. So
\ref{H2} fails at $k=3$ exactly when $d\ge1$ and
$H^2_c((U\cap-U)_{\Fb},\cF\otimes[-1]^*\cF)\ne0$, that is, when
$\cF\otimes[-1]^*\cF$ has nonzero coinvariants under the geometric
monodromy group. For $\cF$ geometrically irreducible with $d\ge1$
this is the condition that the pullback of $\cF$ by $x\mapsto-x$ be
geometrically isomorphic to the dual of $\cF$. It holds for the
Artin--Schreier sheaf above, and for $t=\chi\circ f$ with $\chi$ a
quadratic character and $f$ an even polynomial, where
$\cF\otimes[-1]^*\cF$ is geometrically constant.

For the same Artin--Schreier sheaf, direct computation of $N$ over
the fields $\F_q$ with $q$ odd and $q\le27$, which include $\F_9$,
$\F_{25}$ and $\F_{27}$, gives $N=q-2$ for $k=4$,
$N=3$ for $k=5$, $N=q^2-2q-4$ for $k=6$ and $N=-2q^2+5q+5$ for
$k=7$. If these formulas hold over all finite fields of odd
characteristic, then
\ref{H2} also fails for $k=4,6,7$, where coefficients of both signs
occur, while $k=5$ passes this test. The same formulas do not hold
in characteristic $2$. For $q=2,4,8,16$ the computation gives
$N=-1,-13,-61,-253$ at $k=5$, which is $3-q^2$, and
$N=0,-12,-276,-3108$ at $k=6$, where the formula for odd $q$ would
give $-4,4,44,220$. So in characteristic $2$ the test also fails
at $k=5$, on evidence of the same kind. With the coefficients
$a=(1,2,3,5)$ the same computation gives $N=1$, $-2-p$ and $3+3p$
for $k=3,4,5$ at $p=11,13,17,19$, of the single sign $(-1)^{k-1}$
that \ref{H2} requires. These computations are evidence and prove
nothing. Theorem~\ref{thm:sign} needs \ref{H2} only for infinitely
many $k$, so the failures do not make it vacuous; they show that
the choice of $k$, or of the coefficients, cannot be left
unexamined.
\end{remark}

\begin{remark}[Which Bertini]\label{rem:bertini}
The note says only ``Bertini''. It cannot mean Poonen's Bertini
theorem over finite fields \cite{Poonen2004}. That theorem produces
smooth hypersurface sections of large degree defined over $\F_q$,
and here the hypersurface is a quadric of fixed degree. The reading
adopted here is that it is a Bertini theorem over $\Fb$ with sheaf
coefficients, which would give a dense open set of coefficients $a$
for which \ref{H2} holds for $Q_a$. Such a set has points over
$\F_{q^s}$ for large $s$, and passing from $\F_q$ to $\F_{q^s}$
costs nothing, since it replaces $\alpha_j$ by $\alpha_j^s$. We do
not know a theorem of this kind that applies.

The geometry of $Z$ is not the obstruction. The quadric $Z$ is a
cone, singular at the origin, and the form $Q$ is nondegenerate
only for even $k$, since the adjacency matrix of the path has
determinant $0$ for odd $k$ and $(-1)^{k/2}$ for even $k$. Neither
fact matters. The singular locus of $Z$ lies in
$x_2=0$, outside the support of the sheaf, and $Z\cap U^k$ is
smooth for every $k$ (Remark~\ref{rem:h2-half}). What can fail is
the upper half of \ref{H2}. By Poincar\'e duality on $Z\cap U^k$ it
says that the cohomology without supports of the dual lisse sheaf
vanishes in degrees at most $k-2$, and $Z\cap U^k$ is not proper:
this is a condition on the behaviour of the sheaf along the
boundary. For the plain path
quadric \ref{H2} fails in cases (Remark~\ref{rem:h2-fails}).
Whether it holds for infinitely many $k$ for a given $\cF$, or for
generic coefficients, is not known to us; what generic coefficients
can and cannot do is the subject of Remark~\ref{rem:torus-orbit}.
\end{remark}

\begin{remark}[the forms $Q_a$ are one torus orbit]
\label{rem:torus-orbit}
The forms $Q_a=\sum_ia_ix_ix_{i+1}$ make up the linear system
spanned by the monomials $x_ix_{i+1}$. Its members with all
$a_i\ne0$ are a single orbit of the diagonal torus: the
substitution $x_i=\lambda_iy_i$ with $\lambda_1=1$ and
$\lambda_{i+1}=(a_i\lambda_i)^{-1}$ turns $Q_a(x)$ into $Q(y)$. So
the weighted count for $Q_a$ and $t$ is the weighted count for the
plain path form $Q$ in which $t$ is replaced, in the $i$-th
variable, by $x\mapsto t(\lambda_ix)$; in terms of sheaves, $\cF$ is
replaced in the $i$-th factor by its pullback under
$x\mapsto\lambda_ix$. Changing the coefficients therefore moves the
sheaf, not the quadric, and no member of the family is more
general than another as a variety.

Generic coefficients do not help every sheaf. They change nothing
for a trace function that is homogeneous, $t(\lambda x)=\chi(\lambda)\,t(x)$
for a
multiplicative character $\chi$: then the count for $Q_a$ is
$\chi(\lambda_1\cdots\lambda_k)$ times the count for $Q$, over
$\F_q$ and, with $\chi(\lambda_1\cdots\lambda_k)^n$, over
$\F_{q^n}$, so \ref{H2} holds for $Q_a$ exactly when it holds for
$Q$. Such a $t$ is a multiple of $\chi$, the trace function of a
Kummer sheaf on $\mathbb A^1\smallsetminus\{0\}$, and for
nontrivial $\chi$ its sum $M$ is zero, so
Theorem~\ref{thm:sign} has no content for it. The Artin--Schreier
sheaf is not of this kind: its pullbacks under distinct dilations
are not geometrically isomorphic, and the counts do change with
$a$ (Remark~\ref{rem:h2-fails}).
\end{remark}

\subsection{What is not checked}\label{subsec:not-checked}

The proof of Theorem~\ref{thm:sign} is complete as a deduction from
\ref{H1} and \ref{H2}. The following is not established.
\begin{itemize}
\item Hypothesis \ref{H2} is not proved here, and no reference for
  it is known to us. Its lower half is automatic
  (Remark~\ref{rem:h2-half}); its upper half fails for the plain
  path quadric at $k=3$ in simple cases, and on numerical evidence
  at other small $k$ (Remark~\ref{rem:h2-fails}). The singularity
  of $Z$ and the degeneracy of the path form for odd $k$ are not
  the obstruction (Remark~\ref{rem:bertini}). The note suggests
  generic coefficients,
  $Q_a = \sum_i a_ix_ix_{i+1}$ with $a_i \ne 0$, which amounts to
  moving the sheaf by dilations (Remark~\ref{rem:torus-orbit}). All
  of the proof
  above works for $Q_a$:
  Proposition~\ref{prop:identity} holds for any polynomial in place
  of $Q$, and Proposition~\ref{lem:parseval} only needs
  $ha_i \ne 0$. If \ref{H2} fails for every $Q_a$, Step~4 fails and
  the proof gives nothing.
\item The note's ``vectorize to go up dimensions'' suggests that
  the same argument, with $x_i$ in $\F_q^m$, gives Deligne's
  theorem in higher dimension. This has not been worked out.
\item That this is the argument the note intends has not been
  confirmed.
\end{itemize}
The first item leaves the following question.

\begin{question}\label{q:sign-h2}
For which sheaves $\cF$ satisfying \ref{H1}, and which $k$, does
Hypothesis~\ref{H2} hold for the path quadric, or for the quadric
$Q_a$ with suitable coefficients? By Remark~\ref{rem:h2-fails} it
does not hold for all $k$. An answer for infinitely many $k$ would
complete the sign-only proof.
\end{question}

\section{Multiplicity versus sign}
\label{sec:discussion}

The two arguments of this chapter aim at the same bound from
opposite ends, the second conditionally, and the number-field
argument was assembled from ideas of both. This section compares them.

\subsection{What the sign-only proof uses}

The proof of Theorem~\ref{thm:sign} uses four inputs.
\begin{enumerate}[label=(\arabic*)]
\item The identity $qN = M^k + R$. This is orthogonality of $e_q$
  and nothing else.
\item The bound on $R$. This is Parseval on $\F_q$ and the bound
  $|t| \le r$. For $t = \chi\circ f$ or $t = e_q\circ f$ the bound
  is $|t| \le 1$, which is trivial.
\item The Grothendieck--Lefschetz trace formula, to write $M_n$ and
  $N_n$ as alternating sums of $n$-th powers of eigenvalues.
\item Hypotheses \ref{H1} and \ref{H2}: the cohomology of each side
  lives in a single degree.
\end{enumerate}
It does not use any bound on the size of an eigenvalue. Hypotheses
\ref{H1} and \ref{H2} say in which degrees cohomology can be
nonzero. They say nothing about weights, and statements of this
kind are proved by geometric means: Artin vanishing, Poincar\'e
duality, and Lefschetz or Bertini arguments. So the proof is not
circular.

\subsection{Where the sign is}

The ``sign'' is not the sign of a number, since $N$ is complex. It
is the parity of a cohomological degree (Figure~\ref{fig:sign-degrees}):
\begin{itemize}
\item every eigenvalue in $M_n^k$ comes from degree $k$ of the
  $k$-fold product, with sign $(-1)^k$;
\item every eigenvalue in $q^nN_n$ comes from degree $k-1$, with
  sign $(-1)^{k-1}$.
\end{itemize}
In $R_n = q^nN_n - M_n^k$ the minus sign flips the second, so all
terms carry the same sign $(-1)^{k-1}$. Without this, the identity
would say nothing. A sum of eigenvalues of size $q^{k}$ could
cancel against another sum down to the size of $R_n$. The
single-degree hypothesis rules this out. Then the large
eigenvalues $\alpha_I$ must appear in $R_n$, and $R_n$ is small.
\begin{figure}[tbp]
\centering
\input{book/figs/ch02-sign-degrees.tex}
\end{figure}

Compare Section~\ref{sec:route2}. There too everything turned on a
sign: the upper bound for point
counts controls the sums $\sum_i\alpha_i^n$ from one side only, and
a real root of maximal modulus escapes it
(Theorem~\ref{thm:dichotomy}). In the sign-only proof the
eigenvalues are controlled from one side as well, but
Step~4 arranges that every eigenvalue enters with a positive
integer coefficient, so that Proposition~\ref{lem:growth} applies
and nothing can escape.

\subsection{The role of $k$}

The bound on $R$ loses a fixed factor $q^{3/2}$, and the eigenvalue
$\alpha_j$ appears as $\alpha_j^k$. Taking $k$-th roots turns the
fixed loss into $q^{3/(2k)}$, which tends to $1$ (Figure~\ref{fig:k-exponents}). This is the
tensor power trick, the same device as
Proposition~\ref{prop:power-sums} in the first proof, where the
fixed loss was the constant $A$ and the parameter sent to infinity
was the degree $n$ of the constant field extension. Deligne's first
proof of the Weil conjectures \cite{Deligne1974} also gains by
passing to large tensor powers, and also rests on a positivity that
prevents cancellation. There the positivity is that of the
coefficients of an $L$-function of $\cF \otimes \overline{\cF}$.
Here it is the parity of a degree. This comparison is ours and is
not in the note.
\begin{figure}[p]
\centering
\input{book/figs/ch02-k-exponents.tex}
\end{figure}

\subsection{The comparison}

\begin{center}
\begin{tabular}{@{}p{0.22\linewidth}p{0.35\linewidth}
p{0.35\linewidth}@{}}
\toprule
 & Stepanov--Bombieri & Sign-only \\
\midrule
Resource & multiplicity: one function vanishing to order $m$ at
every point & a sign: all eigenvalues enter $R_n$ with one sign \\
Where $\sqrt q$ comes from & the balance between $m$ and the degree
it costs (Proposition~\ref{prop:sqrtq-forced}) & the size of one
Fourier sum, by Parseval \\
Fixed loss, and its removal & the constant $(2g+1)$, later $A$;
removed by $n\to\infty$ & the factor $q^{3/2}$; removed by
$k\to\infty$ \\
What is elementary & everything: Riemann--Roch only & the identity
and the bound on $R$ \\
What is not & nothing & the trace formula and
\ref{H1}, \ref{H2} \\
The other sign & from twists of a Galois cover
(Section~\ref{sec:route1}) & not needed \\
\bottomrule
\end{tabular}
\end{center}

\begin{itemize}
\item \emph{Where $\sqrt q$ comes from.} In Stepanov's method,
  $\sqrt q$ is the balance point between the multiplicity and the
  degree it costs. In the sign-only proof, $\sqrt q$ is the size of
  one Fourier sum, by Parseval.
\item \emph{What is elementary.} Stepanov's method needs only
  Riemann--Roch. The sign-only proof needs $\ell$-adic cohomology
  for the trace formula and the vanishing hypotheses. Its Fourier
  analysis is elementary.
\item \emph{What carries over to $\Q$.} Neither proof carries over
  directly. There is no Frobenius on $\Z$ for Stepanov's $m$-th
  powers, and there is no cohomology for the sign. This is why the
  note looks for an automorphic source of the sign
  (Section~\ref{sec:role}), and why
  Chapter~\ref{ch:idea} replaces the $m$-th power by a chain of
  theta lifts (Section~\ref{sec:theta}).
\end{itemize}

\section{What the identity is for}\label{sec:role}

This section reports what the note says about the role of the
identity in the number-field argument, and what can be inferred
from it. The inferences are marked. No fuller account is known to
us, and item~3 of the timeline was not found in the Lean archive.

\subsection{What ``automorphy of $N$'' means}

The footnote is attached to this phrase: the definite sign of the
left side of $qN = M^k + R$, ``meaning holo/automorphy over a
function field''. Our reading is as follows. Over $\F_q$ the sign
of $qN$ comes from cohomology: Hypothesis~\ref{H2} puts all of $N$
in one degree. Over $\Q$ there is no such cohomology. The idea was
to find another source for the sign: a way to write $N$ as an
\emph{automorphic} object, for instance a period integral of theta
functions. Its sign would then come from holomorphy or automorphy,
and not from the degree of a cohomology group.

Item~3 of the timeline sketches how. It says that $\mu$ is
automorphic, coming from a theta function on a sixfold cover of
$\mathrm{GL}_5$, as a period integral of thetas. The note does not
say what $\mu$ is. Item~3 also says that the quadric has a natural
meaning on a CM elliptic curve such as $E = \C/\Z[\omega]$, and
that one should ``put Poincar\'e bundle on the edges to get
$\theta$s''. Our reading of this sentence:
\begin{itemize}
\item In the finite-field sum, the edge $\{i, i+1\}$ contributes
  the bilinear character $e_q(h x_ix_{i+1})$.
\item On $E \times E$ the analogue of a bilinear pairing is the
  Poincar\'e line bundle $P$. Its Riemann form is a bilinear form.
\item Pulling $P$ back along $(x_i, x_{i+1}) \colon E^k \to E^2$
  for each edge and tensoring gives a line bundle on $E^k$ attached
  to the path. Its sections are theta functions, which are
  automorphic.
\item $\#\mu_K = 6$ for $K = \Q(\omega)$, which would explain
  ``sixfold''. Why the group is $\mathrm{GL}_5$ is not explained.
\end{itemize}

\subsection{What $O(2,n)$ is}

The note says the decisive observation was ``multiplicity $2$ on
$O(2,n)$'' via a ``natively-automorphic determinant'', found by
analogy with the sign of $qN$. The note does not define $O(2,n)$.
In standard usage it is the orthogonal group of a real quadratic
space of signature $(2,n)$, or a rational form of it. It is the
orthogonal group whose symmetric space is Hermitian (a type~IV
domain), so it has holomorphic modular forms. It is the target of
the classical theta lifts from $\mathrm{SL}_2$, the Borcherds and
Kudla--Millson lifts. Which quadratic space, which $n$, and which
determinant are meant is not said.

Our reading: ``multiplicity'' is used as in Stepanov's method, as
an order of vanishing or of divisibility. ``Multiplicity $2$'' on
$O(2,n)$ was the first case beyond multiplicity one. The note says
that more theta lifts then gave more multiplicity.

\subsection{Is there an $N$, $M$, $k$ and $R$ in the final
argument?}

We have found no evidence that the identity $qN=M^k+R$ has a
counterpart in the argument of
Chapters~\ref{ch:assembly}--\ref{ch:families}. The note does not mention the
identity
after item~3. What it gives instead is the following.
\begin{itemize}
\item The quantity to bound is a smoothed prime sum. In the
  dictionary of Section~\ref{sec:idea-dictionary}, the function
  $\Gz$ of Definition~\ref{def:sz-gz} plays the role of a
  tail-smoothed sum of the numbers $q^n + 1 - \#C(\F_{q^n})$, with
  the sign reversed against the error of a point count
  (Proposition~\ref{prop:which-bound}). So the point count of the
  curve is replaced by a weighted count of primes.
\item The sign is that of the Guinand--Weil explicit formula.
  Item~4 of the timeline is ``$W(1_S) \ge 0$ via a determinant'',
  read there as Stepanov with multiplicity one. In the book the
  sign appears as the positivity of the weights in the windowed
  Landau criterion (Remark~\ref{rem:positivity}).
\item Multiplicity comes from a chain of $m$ theta lifts between
  definite unitary groups, with $m$ of order $\sqrt X$ for a prime
  sum of length $X$. This is the analogue of $m \approx \sqrt q$ in
  Stepanov's method.
\end{itemize}
So the sign-only identity seems to have survived only as an
analogy. Its definite sign became the target, a ``natively
automorphic determinant''. One could match $k \to \infty$ with the
length of the chain of lifts, since in both cases a parameter is
sent to infinity to remove a fixed loss. Nothing in the note
supports this match.

\subsection{Where the quadric came from}

The footnote says the quadric arose from simplifying an average
over the family $\chi \mapsto \chi(\mathrm{disc})$. It names a
normalized Gauss sum $y$ of order $\#\mu_K$, which lives via
$\theta$ on the $\#\mu_K$-fold cover of $\mathrm{GL}_2$. Here is
our reading, from what is known about these Gauss sums; the
account of the literature follows the survey in
\cite[\S1]{BrokerHoffstein2013}.
\begin{itemize}
\item Let $K$ be a number field containing the group $\mu_K$ of
  roots of unity, and let $n = \#\mu_K$. For a modulus $c$ there is
  an $n$-th order Gauss sum $g(c)$, built from the $n$-th power
  residue symbol modulo $c$. Let $y(c) = g(c)/|g(c)|$.
\item The Fourier coefficients of the Eisenstein series on the
  $n$-fold cover of $\mathrm{GL}_2$ are Dirichlet series built from
  Gauss sums of order $n$, and the theta function on that cover,
  introduced by Kubota, is a residue of this Eisenstein series
  \cite[\S1, pp.~3--4]{BrokerHoffstein2013}. The footnote's wording
  suggests that the $y(c)$ are the Fourier coefficients of the
  theta function for every $n$. This is known only for $n=3$: Patterson
  determined the
  coefficients of the cubic theta function, and at squarefree $c$
  they are the normalized cubic Gauss sums \cite{Patterson1977};
  see \cite[\S1, p.~4]{BrokerHoffstein2013}. Kazhdan and Patterson
  \cite{KazhdanPatterson1984} proved a periodicity property and
  Hecke relations for every $n$, and these determine all the
  coefficients for $n=3$ but only some of them for $n\ge4$
  \cite[\S1, pp.~5--6]{BrokerHoffstein2013}. For $n=4$ the squares
  of the coefficients are predicted by a conjecture of Patterson.
  For $n=6$, which is the case $\#\mu_K=6$ of the note, a weaker
  conjecture of Chinta, Friedberg and Hoffstein
  \cite{ChintaFriedbergHoffstein2012} relates the coefficients at
  squares to cubic Gauss sums, and does not determine all of them
  \cite[\S1, pp.~6--8]{BrokerHoffstein2013}; see also
  \cite{Patterson2014}.
\item $y$ is not multiplicative. For coprime $c, c'$ one has
  $y(cc') = \varepsilon(c, c')\,y(c)\,y(c')$, where
  $\varepsilon(c, c')$ is a power residue symbol. By the
  reciprocity law, $\varepsilon$ behaves like a bilinear pairing.
\item For a product of $k$ coprime moduli this gives
  \[
    y(c_1\cdots c_k) \;=\; \prod_i y(c_i)\,
    \prod_{i<j}\varepsilon(c_i, c_j),
  \]
  with one factor for each pair $\{i, j\}$. That is the complete
  graph on $k$ vertices.
\item In a finite-field model a bilinear pairing becomes
  $e_q(x_ix_j)$. The complete graph gives
  $e_q\bigl(\sum_{i<j} x_ix_j\bigr)$. Keeping only the edges
  $\{i, i+1\}$ gives the path quadric.
\end{itemize}
The phrases ``$\chi \mapsto \chi(\mathrm{disc})$'' and
``$a \mapsto y(\cdot)^a$'' could not be reconstructed more
precisely.

The footnote also says that Gabber's hypersurface in one of Katz's
arguments ``reminds one of the quadric''. Katz \cite{Katz2015} uses
the hypersurface $X_1^d + \sum_i X_iX_{i+1}^{d-1} = 0$ with
$p \mid d$, due to Gabber, which he treats by Delsarte's method and
Gauss sums. Its monomials $X_iX_{i+1}^{d-1}$ link each variable to
the next, as the path does. This link is an inference.

\section{Dictionary with Alp\"oge's argument}
\label{sec:dictionary}

This section matches the steps of the Stepanov--Bombieri proof with
the parts of the number-field argument. The argument is taken in
five components, each with the place where the book develops it.
\begin{enumerate}[label=(\roman*)]
\item A one-sided bound $G(x)\le x^{1/2+O(\varepsilon)}$ on windows
for a positive combination $G$ of smoothed Chebyshev functions
(Theorem~\ref{thm:chebyshev-windows}), fed into a windowed Landau
theorem (Theorem~\ref{thm:windowed-landau}).
\item A set of characters closed under inversion
(Definition~\ref{def:level}), so that every zero counts with the
same sign (Remark~\ref{rem:positivity}).
\item Hermitian lattices of integral automorphic forms, with a
slope inequality for the period maps between them in the role of
``zeros at most poles'' (Theorems~\ref{thm:bost}
and~\ref{thm:frame}; \cite{Bost1996}). No small sections are
constructed (Remark~\ref{rem:periodmap}).
\item Forced $p$-adic divisibility, then a chain of
$m\approx\sqrt x$ theta lifts, with an error independent of $m$
against a main term proportional to $m$
(Chapter~\ref{ch:theta}, Remark~\ref{rem:exponent-half}).
\item Cyclotomic towers, which the note's outline places in the
role of $\F_{q^n}/\F_q$. In the book cyclotomic fields enter in the
construction of hosts (Chapter~\ref{ch:hosts}) and layers of a
cyclotomic tower in the Tate--Shafarevich bound
(Chapter~\ref{ch:sha}); that they play the role of the tower of
constant field extensions is not established here.
\end{enumerate}
The ratings in the table are judgments: ``exact'' is the same
statement under the usual dictionary, ``loose'' the same role by
another mechanism, ``none'' no counterpart or a failing one. Each
row names the statement of
Chapters~\ref{ch:assembly}--\ref{ch:families} against which it was
rated, and the table agrees with Table~\ref{tab:bb-model} of
Chapter~\ref{ch:assembly}. Figure~\ref{fig:dictionary-diagram}
draws the table as a graph.

{\footnotesize
\begin{longtable}{>{\raggedright}p{0.47\textwidth}
>{\raggedright}p{0.34\textwidth}l}
\toprule
Step in this chapter & Component & Rating\\
\midrule
\endhead
Section $w$ (Theorem~\ref{thm:aux-section}); pole bound, degree
& (iii) no constructed section; the period map between two given
lattices, its height and the slopes (Remark~\ref{rem:periodmap})
& loose\\
Step~1: Riemann--Roch and pigeonhole
& (iii) no Siegel lemma; arithmetic Hilbert--Samuel enters only
inside the covolume bound (Theorem~\ref{thm:covolume}) & none\\
Step~4: effective divisors have nonnegative degree
& (iii) slope inequality (Theorem~\ref{thm:bost}), one step of
the frame theorem and not used in rank one
(Remark~\ref{rem:not-bost}) & loose\\
Step~3: multiplicity from $w^M=\sum u_i^Mv_i^q$, a ring
endomorphism
& (iv) forced $p$-adic divisibility along the chain of lifts
(Proposition~\ref{prop:sqrtx}) & loose\\
Section~\ref{sec:variantJ}: $\ell\approx\sqrt p$ derivative
conditions
& (iv) chain of $m\approx\sqrt x$ lifts
(Remark~\ref{rem:exponent-half}) & loose\\
Balance $(g+1)q/M+2gM$; Remark~\ref{rem:J-balance}
& (iv) errors $x/m$ and $m$ (Proposition~\ref{prop:sqrtx})
& loose\\
Non-vanishing (Propositions~\ref{prop:independence},
\ref{prop:J-nonvanishing}) & not identified; no section is
constructed (Remark~\ref{rem:periodmap}) & none\\
Proposition~\ref{prop:landau}, a finite case & (i) windowed Landau
theorem (Theorem~\ref{thm:windowed-landau}) & exact\\
Theorem~\ref{thm:dichotomy}; fails by
Proposition~\ref{prop:formal-cex}
& (i) ``the upper bound alone gives RH''; the bound of
Theorem~\ref{thm:chebyshev-windows} is in the other direction
& none\\
$G(T)$, $\gamma_K$, $\delta_K$; Theorem~\ref{thm:route2prime}
& (i) $\Gz$, $G_\chi$, $1+\gamma$, $L'/L(1,\chi)$
(Definition~\ref{def:chebyshev}) & exact\\
Route~1, which also needs a base with $L=1$
& (ii) characters closed under inversion
(Definition~\ref{def:level}); Fourier analysis on a Galois group
(Theorem~\ref{thm:core-bound}) & loose\\
All roots of $L_E$ enter $N(E,\mathrm{id})$ with one sign
& (ii) every zero counts with one sign
(Remark~\ref{rem:positivity}) & exact\\
The tower $K_n$ (Propositions~\ref{prop:tower},
\ref{prop:power-sums}) & (v) cyclotomic tower, role not established
(Chapters~\ref{ch:hosts}, \ref{ch:sha}) & loose\\
Functional equation, Theorem~\ref{thm:rh} & functional equation
(Theorem~\ref{thm:fe}) & exact\\
\bottomrule
\end{longtable}
}
\ffbookfloat{dictionary-diagram}

The sign-only proof has no row in this table. As
Section~\ref{sec:role} explains, the identity $qN=M^k+R$ is not
known to have a counterpart in the final argument; what survives
of it is the demand for a definite sign, which is component (ii).

\subsection{Where the analogy breaks}
\begin{enumerate}[label=(B\arabic*)]
\item \emph{Frobenius.} Step~3 of
Theorem~\ref{thm:aux-section} uses that $z\mapsto z^M$ is
additive; $\Z$ has no such endomorphism. Vanishing to order $M$ at
$N$ points costs $\ell(MaP_0+b\sigma(P_0))$ conditions, not about
$MN$ as in characteristic zero. Section~\ref{sec:variantJ} still
uses $(X^p)'=0$.
\item \emph{Sign.} The bound that the number-field argument proves
is, in function-field terms, a lower bound on point counts
(Proposition~\ref{prop:which-bound}), not the upper bound proved
in Section~\ref{sec:upper}. Over $\Q$ the corresponding statement
is Remark~\ref{rem:G-over-Q}.
\item \emph{The base.} Route~1 needs a base with $L$-polynomial
$1$; $\zeta$ has zeros. This is consistent with the use of a
Landau criterion over $\Q$ in place of an averaging argument.
\item \emph{Finiteness.} Here there are $2g$ roots and
\eqref{eq:explicit-G} is exact; not so over $\Q$, where the
explicit formula is an infinite sum (Remark~\ref{rem:explicit}).
\item \emph{The centring constant.}
Theorem~\ref{thm:route2prime} needs $\gamma_K$ exactly, which
Section~\ref{sec:upper} never produces. In the number-field
argument the corresponding constants enter through a period of
Chowla--Selberg type (Remark~\ref{rem:stirling}).
\end{enumerate}
A faithful function-field model of the number-field argument would
add to Theorem~\ref{thm:route2prime} a Stepanov-type upper bound
for $G(T)$ (Question~\ref{q:smoothed}). Nothing in this chapter
confirms or refutes Alp\"oge's claim.

\section{Remarks and questions}
\label{sec:remarks}

\subsection{The positivity that Weil uses and this proof avoids}

Weil's proofs use positivity on a surface or on the Jacobian: the
Castelnuovo--Severi inequality, a form of the Hodge index theorem,
on $C\times C$ applied to the diagonal and the graph of Frobenius;
or the positivity of the Rosati involution on
$\mathrm{End}(J)\otimes\Q$ \cite{Weil1948}. Either gives both signs
of $\lvert N-q-1\rvert\le2g\sqrt q$ at once.
Figure~\ref{fig:weil-vs-sb} contrasts the two pictures.
\ffbookfloat{weil-vs-sb}

The Stepanov--Bombieri proof uses, on the curve itself, only
Riemann--Roch and the nonnegativity of the degree of an effective
divisor. This positivity is one-sided. The lower bound comes from
averaging (Section~\ref{sec:route1}), not from a second positivity.
In explicit-formula terms the proof establishes an inequality for
the one-parameter family of counts $N_n$ only; it does not
construct a positive form. The sharp constant $2g$ is not produced
by the auxiliary section: Theorem~\ref{thm:bombieri} gives $2g+1$,
and Proposition~\ref{prop:sqrtq-forced} shows that the method
cannot do much better. The constant $2g$ appears only after
Proposition~\ref{prop:power-sums}.

\subsection{Where the characteristic is used}

Characteristic $p$ enters the first proof in three places:
Proposition~\ref{prop:frobenius}\eqref{item:frob-additive} in
\eqref{eq:wM}, the only use in the upper bound;
Proposition~\ref{prop:frobenius}\eqref{item:frob-rational}, which
defines the points counted; and, in Section~\ref{sec:variantJ},
$(X^{jp})'=0$ and $x^{jp}=x^j$. Step~1 and Step~4 of
Theorem~\ref{thm:aux-section},
Proposition~\ref{prop:independence}, Section~\ref{sec:zeta},
Section~\ref{sec:route1} apart from the definition of Frobenius
places, and Propositions~\ref{prop:power-sums}
and~\ref{prop:landau} do not use the characteristic. Stepanov
treated $y^2=f(x)$ over prime fields and W.~Schmidt the general
case \cite[p.~234]{Bombieri1974}, \cite{Schmidt1973}.

\subsection{Number fields}

The following are heuristic remarks, not theorems.
\begin{enumerate}[label=(\alph*)]
\item No endomorphism of $\Z$ is known that plays the role of
Frobenius, and no object $\Spec\Z\times\Spec\Z$ over a base on
which a graph of Frobenius could live. Item (B1) of
Section~\ref{sec:dictionary} is the concrete form of this.
\item Route~1 has no base with trivial $L$-function.
\item The analytic step does transpose: a one-sided bound on a
centred smoothed prime sum, with Landau's theorem, is a sound
criterion (Theorem~\ref{thm:route2prime},
Remark~\ref{rem:G-over-Q}, Proposition~\ref{prop:idea-landau}). The
open matter over $\Q$ is the bound itself.
\item Stepanov's method has the form of the auxiliary-function
arguments of Diophantine approximation (Thue, Siegel). This is a
resemblance of form, not a transfer.
\end{enumerate}

\subsection{Questions}

The questions that this chapter leaves open are stated where they
arise. Questions~\ref{q:real-zero} and~\ref{q:smoothed} close the
summary of Route~2 in Section~\ref{sec:route2}, and
Question~\ref{q:sign-h2} follows the list of
Section~\ref{subsec:not-checked}.
\clearpage
\clearpage
\chapter{The argument assembled}
\label{ch:assembly}

This chapter states the main theorem and gives its whole proof. Each
tool that the proof uses is taken as a black box: a result proved in
one of Chapters~\ref{ch:landau}--\ref{ch:families}, restated here
with the same hypotheses and the same conclusion, and followed by a
few sentences on what it says, why the proof needs it and where it is
proved. The aim is that the reader sees the flow of the argument and
how its parts fit together before any of the parts is opened.
Chapter~\ref{ch:idea} describes the idea informally, and
Chapter~\ref{ch:function-fields} works out the model over a finite
field. The present chapter is the precise assembly. The later
chapters use from it only the statement of Theorem~\ref{thm:main}.

The order is as follows. Section~\ref{sec:intro} states the theorem.
Section~\ref{sec:bb-plan} gives the logical skeleton of the proof and
its dependency diagram, and Section~\ref{sec:bb-vocabulary} recalls the
terms in which the black boxes are stated.
Sections~\ref{sec:bb-galois}--\ref{sec:bb-line} state the black boxes
in the order in which the proof calls on them, and
Section~\ref{sec:proof-main} is the proof.
Section~\ref{sec:bb-analogy} compares the assembled argument with
Chebyshev's and with the function-field model of
Chapter~\ref{ch:function-fields}, and Section~\ref{sec:bb-status}
records the status of each input: which statements are proved in this
book, which are outlined, and which are cited.

\section{The main theorem}
\label{sec:intro}

A \emph{Dirichlet character} modulo $N\ge1$ is a homomorphism
$(\Z/N\Z)^\times\to\C^\times$, extended to a function $\chi$ on $\Z$
that vanishes on the integers not prime to $N$. Its $L$-function is
$L(s,\chi)=\sum_{n\ge1}\chi(n)n^{-s}$ for $\operatorname{Re}s>1$,
continued meromorphically to $\C$; the continuation is entire when
$\chi$ is nonprincipal, and has a simple pole at $s=1$ and no other
pole when $\chi$ is principal. A character is \emph{primitive} if it
is not induced from a character of smaller modulus. The character
modulo~$1$ is the primitive principal character; its $L$-function is
$\zeta$.

\begin{theorem}[Alp\"oge]\label{thm:main}
Let $\chi$ be a Dirichlet character. Then $L(s,\chi)\ne0$ for
every $s$ with $\operatorname{Re}s>1/2$, except $s=1$ when $\chi$ is
principal, and every zero of $L(s,\chi)$ with
$0<\operatorname{Re}s<1$ has real part $1/2$. Every zero of
$\zeta(s)$ other than $-2,-4,-6,\dots$ has real part $1/2$.
\end{theorem}

\begin{remark}[scope of the conclusion]\label{rem:scope}
For the principal character the point $s=1$ is excluded because
its $L$-function has a pole there. The zeros of $\zeta$ at the
negative even integers are its trivial zeros. For an imprimitive
character, the extra Euler factors can have zeros on
$\operatorname{Re}s=0$; they do not affect the critical-strip
statement of Theorem~\ref{thm:main}. The theorem concerns
$\zeta$ and Dirichlet $L$-functions.
\end{remark}

\begin{remark}[the final statements of the archive]\label{rem:ai-final}
The half-plane statement of Alp\"oge's Lean archive
\cite{AlpogeArchive2026} is not produced for every character
directly. The final statements use Mathlib's own definitions
of the Riemann zeta function and of Dirichlet $L$-functions, with no
versions of the archive's own. The upstream output concerns primitive
characters only. At the character of modulus~$1$ and $s=1$ it is a
statement about the finite value $(\gamma-\log4\pi)/2$ that Mathlib
assigns to $\zeta$ at its pole. Theorem~\ref{thm:main} excludes that
point. The archive was read as text and not built
(Section~\ref{sec:ai-formalization}).
\end{remark}

The proof is given in Section~\ref{sec:proof-main}. Its last
deduction is short: a bound on windows gives nonvanishing for
primitive characters, the Euler factors extend it to imprimitive
characters, and the functional equation reflects zeros across the
critical line. The work of the argument is the bound on windows and
the construction of the arithmetic data to which it applies.

\section{The skeleton of the proof}
\label{sec:bb-plan}

The Riemann hypothesis asks for a sharp bound on the oscillation of
primes around their expected density, and for a Dirichlet character
the same question concerns primes in residue classes. The argument
approaches these questions through an upper bound for smoothed,
character-weighted Chebyshev functions. The bound is needed only on
long multiplicative windows that occur arbitrarily far out, and a
windowed converse to the explicit formula then excludes zeros to the
right of the critical line. The difficulty lies in producing the
windows and in obtaining the square-root exponent in the bound.

Fix a primitive Dirichlet character $\psi$. The terms used in the
next four paragraphs are recalled in Section~\ref{sec:bb-vocabulary}.
\begin{enumerate}[label=(\Roman*)]
\item\label{item:bb-sk-landau}
\emph{The analytic end.} To $\psi$ is attached a positive
combination $G_{\lev_\psi}$ of smoothed Chebyshev functions, one for
each primitive character inducing a power of $\psi$. If, for some
$\sigma_0\in(0,1)$, the inequality $G_{\lev_\psi}(x)\le x^{\sigma_0}$
holds on windows $[Y,Y^{3/2}]$ with $Y$ arbitrarily large, then
$L(s,\psi)\ne0$ for $\operatorname{Re}s>\sigma_0$. This is the
windowed Landau criterion (Section~\ref{sec:bb-landau}). Positivity
of the weights makes a one-sided upper bound sufficient, and the
bound is not required at every large $x$.
\item\label{item:bb-sk-cheb}
\emph{The inequality.} For every small $\eps>0$ one needs
$G_{\lev_\psi}(x)\le x^{1/2+9\eps}$ on such windows. This is the
conclusion of the Chebyshev inequality on windows
(Section~\ref{sec:bb-chebyshev}). Its hypothesis is the existence of
a \emph{family} of automorphic data at $\eps$, indexed by the points
$x$ of a set $\Win$ with windows, that is \emph{tied} at every large
$x\in\Win$ and satisfies an explicit-formula condition for
$\lev_\psi$.
\item\label{item:bb-sk-fam}
\emph{Existence.} Tied families exist: this is the top existence
statement, which is proved from the existence of packages and of
admissible-family data (Section~\ref{sec:bb-families}). It needs two
inputs. The first is a supply of host fields for $\psi$ on windows,
meeting the explicit-formula condition and carrying a small Sha
budget (Section~\ref{sec:bb-hosts}). The second is the Sha bound on
windows, which provides that budget (Section~\ref{sec:bb-sha}).
\item\label{item:bb-sk-deep}
\emph{The inputs behind the inputs.} The Sha bound rests on Galois
representations with a truncated local--global compatibility
(Section~\ref{sec:bb-galois}). The admissible-family data rest on the
same Galois representations and on a covolume bound for lattices of
integral automorphic forms, which is where arithmetic geometry enters
(Section~\ref{sec:bb-covolume}).
\end{enumerate}
Given \ref{item:bb-sk-landau}--\ref{item:bb-sk-fam}, a zero $\rho$ of
$L(s,\psi)$ with $\operatorname{Re}\rho>1/2$ is excluded by choosing
$\eps$ so small that $1/2+9\eps<\operatorname{Re}\rho$. The passage
from primitive characters to all characters, and from the half-plane
to the critical line, is classical (Section~\ref{sec:bb-line}).

The dependencies are displayed in Figure~\ref{fig:dag}. The proof in
Section~\ref{sec:proof-main} reads the figure from the bottom right
to the top: Galois representations, the Sha bound, the supply of
hosts, the existence of tied families, zero-freeness for primitive
characters, and the passage to all characters.

\begin{figure}[htbp]
\centering
\input{alpoge/figs/v2/dag.tex}
\caption{Dependencies in the proof of Theorem~\ref{thm:main}.
Solid arrows point from inputs to results that use them. The
dashed arrow marks the tie between a family and its slope record.
A number after the sign \S\ is the number of the chapter that
develops the step.}
\label{fig:dag}
\end{figure}

Three of the boxes in Figure~\ref{fig:dag} carry almost all of the
mathematics that is specific to this argument. The first is Step~2,
the Chebyshev inequality, in which a chain of $m$ theta lifts is
compared on its archimedean and finite sides and $m$ is taken near
$\sqrt x$. The other two are the inputs at the bottom right: the
Sha bound on windows and the Galois representations behind it. The
assembly itself, as Remark~\ref{rem:content} records, is a chain of
implications.

\section{The vocabulary of the black boxes}
\label{sec:bb-vocabulary}

The black boxes are stated in terms that are defined in
Chapters~\ref{ch:landau}--\ref{ch:families}. This section recalls
those terms, in the order in which they are needed. The definitions
cited are the reference; where a definition is long, only the part
that the statements of this chapter use is recalled, and this is
said. Nothing in this section is new.

\subsection*{Smoothed Chebyshev functions}

Let $\Lambda$ be the von Mangoldt function and $\gamma$ Euler's
constant. For $x>0$ the smoothed Chebyshev function of $\zeta$ and
its normalised form are
\[
\Sz(x)=\sum_{1\le n\le x}\Lambda(n)\Bigl(\frac xn-1\Bigr),\qquad
\Gz(x)=\Sz(x)-x\log x+(1+\gamma)x .
\]
For a nonprincipal Dirichlet character $\chi$ and $y>0$ the complex
twisted function is
\[
g_\chi(y)=\sum_{1\le n\le y}\Lambda(n)\chi(n)\Bigl(\frac yn-1\Bigr)
+y\,\frac{L'(1,\chi)}{L(1,\chi)},
\]
and for the character modulo~$1$ one puts $g_1=\Gz$
(Definition~\ref{def:chebyshev} and
Proposition~\ref{prop:realcomplex}). The terms added to the prime
sums remove the contribution of $s=1$. A \emph{character family} $F$
is a finite list of pairs $(\chi_i,a_i)$, where each $\chi_i$ is a
Dirichlet character that is either the character modulo~$1$ or
nontrivial and each $a_i>0$ is a real weight; its weighted Chebyshev
function is $G_F(y)=\sum_ia_i\operatorname{Re}g_{\chi_i}(y)$
(Definition~\ref{def:charfamily}).

\subsection*{Levels}

A \emph{level} $\lev$ consists of a finite set $\Kept(\lev)$ of
primitive Dirichlet characters that contains the character modulo~$1$
and is closed under inversion, an optional exceptional character
$\psi_{\mathrm{exc}}\in\Kept(\lev)$, and a natural number $o(\lev)$.
The weight of $\psi\in\Kept(\lev)$ is $\wt(\psi)=1$ if
$\psi=\psi_{\mathrm{exc}}$ and $\wt(\psi)=(3o(\lev)+1)^{-1}$
otherwise, and the weighted Chebyshev function of the level is
\[
G_{\lev}(y)=\Gz(y)+\sum_{\psi\in\Kept(\lev),\ \psi\ne1}
\wt(\psi)\operatorname{Re}g_\psi(y).
\]
For a primitive character $\psi$ the \emph{cyclic level} $\lev_\psi$
has as kept set the primitive characters inducing the powers of
$\psi$, no exceptional character, and $o(\lev_\psi)$ equal to the
order of $\psi$ (Definition~\ref{def:level}). In particular
$\psi\in\Kept(\lev_\psi)$, and all weights of a level are positive.

\subsection*{Windows}

A set $S\subset\R$ \emph{has windows of exponent} $\lambda>1$ if for
every $y$ there is $Y\ge y$ with $[Y,Y^\lambda]\subset S$. A level
$\lev$ \emph{satisfies the window criterion at} $\lambda$ if for
every $\sigma_0\in(0,1)$ the following holds: if
$\{x:G_{\lev}(x)\le x^{\sigma_0}\}$ has windows of exponent
$\lambda$, then $L(s,\psi)\ne0$ for every $\psi\in\Kept(\lev)$ and
every $s$ with $\operatorname{Re}s>\sigma_0$, where $s\ne1$ when
$\psi$ is the character modulo~$1$ (Definition~\ref{def:windows}).

\subsection*{Scales and window pairs}

Throughout, $\eps\in(0,1/52)$. A \emph{scale} is a triple
$(x,\eps,J)$ with $x>1$ and $J\ge1$; one puts $b=\exp(x^\eps)$ and
$a=xb$. A \emph{system of scales} of exponent $\lambda$ is a $J$
together with a set $\Win\subset\R$ that has windows of exponent
$\lambda$. An integer $\nu$ is \emph{adapted} to $(\eps,x)$ if
$x^{2+8\eps}\le2^\nu\le2x^{2+8\eps}$. A pair of primes $(q,d)$ is a
\emph{window pair} for $(\eps,Y)$ if $q\equiv3\pmod8$ and
$Y^{\eps/5}/2<\log q\le Y^{\eps/5}$; $d\equiv7\pmod{16}$ and
$Y^{\eps/5}\log^2Y/2<\log d\le Y^{\eps/5}\log^2Y$; $q\mid d+1$; and
every odd prime $p\le Y^{\eps/5}$ divides $d+1$
(Definition~\ref{def:scales}). The upper bound $\eps<1/52$ is fixed
once and for all; only small values of $\eps$ are used in the end.

\subsection*{Hosts}

The geometry of the argument lives over an abelian CM field $E$
attached to $\psi$. A \emph{frame} $\phi$ fixes the shape of its
Galois group: an integer $\nu\ge0$, a finite abelian group $A$ with
a surjection onto $(\Z/2^{\nu+2}\Z)^\times/\{\pm1\}$, and a cyclic
subgroup $A_0\subset A$, the core, of index $o$; the Galois model is
$\GG=\{\pm1\}^2\times A$ (Definition~\ref{def:frame}).

A \emph{host}
$\mathcal H$ for $\phi$ consists of auxiliary primes $q\equiv3
\pmod8$ and $d\equiv7\pmod{16}$ with $(-d/q)=1$, an integer $Q_S$,
a number field $E$ that is Galois over $\Q$ with group $\GG$,
contains $K=\Q(\sqrt{-d})$ and embeds in $\Q(\zeta_{Q_Sqd})$, and an
Artin map $(\Z/Q_Sqd\Z)^\times\to\GG$, subject to the conditions of
Definition~\ref{def:host}. A character $\theta$ of $\GG$ then
determines a primitive Dirichlet character $\chi_\theta$. A
\emph{$\psi$-field} for $(\psi,\nu,q,d)$ is the subfield of a
cyclotomic field cut out by the conditions
$a\equiv\pm1\pmod{2^{\nu+2}}$, $\psi(a)=1$, $(a/q)=1$ and $(a/d)=1$
(Definition~\ref{def:psifield}); the host fields used for $\psi$ are
of this kind, with $\nu$ adapted to the scale, so their degree is
about $x^{2+8\eps}$.

A host is used together with a finite set $\Top\subset A$ of
\emph{top places}. The host \emph{satisfies the explicit-formula
condition} $\mathrm{EF}(\lev,\Top,x,\eps,C)$ if its index is
$o(\lev)$, the characters of $\Kept(\lev)$ are among the characters
$\chi_\theta$ in the way prescribed by Definition~\ref{def:shares-ef},
the set $\Top$ has prescribed shares in the cosets of $A_0$ up to
$x^{-1+\eps}$, and every \emph{bounded} character $\theta$ of $\GG$
satisfies
\[
|g_{\chi_\theta}(x)|\le x^{1/2+\eps}\qquad\text{and}\qquad
\Bigl|\frac{L'(1,\chi_\theta)}{L(1,\chi_\theta)}\Bigr|\le C\log x.
\]
The bounded characters are characters of $\GG$ trivial on the core;
they include every such character that is not kept
(Definition~\ref{def:bounded}). The pointwise bound is required at
the single scale $x$, not for all large arguments.

\subsection*{Host requirements}

Let $\lambda>1$. A \emph{host requirement} $\mathscr R$ is a property
of a scale $(x,\eps,J)$, a frame, a host $\mathcal H$ for that frame,
a finite set $\Top\subset A$ and a positive integer $t$. It is
\emph{supplied on windows of exponent} $\lambda$ if for every
$\eps\in(0,1/52)$ there are constants and a system of scales of
exponent $\lambda$ such that at every large scale of the system there
are ranks $\mathbf k=(k,m,t)$, a host $\mathcal H$ for $\mathbf k$
with top places $\Top_{\mathcal H}$ that is \emph{good} at the scale,
and $\mathscr R$ holds for them (Definition~\ref{def:host-supply}).
The goodness of a host at a scale is a list of size, splitting and
character conditions; it is part of the schematic content of that
definition and is not recalled here. For a primitive character
$\psi$, the \emph{requirement of} $\psi$ is the host requirement
$\mathscr R_\psi$: the host satisfies
$\mathrm{EF}(\lev_\psi,\Top,x,\eps,4)$, and its field $E$ has Sha
budget $(t,b^{1/2})$ in the sense recalled in
Section~\ref{sec:bb-sha}.

\subsection*{Packages, families and ties}

A \emph{ranks datum} for a frame is a triple $\mathbf k=(k,m,t)$ of
integers subject to the inequalities of Definition~\ref{def:ranks};
$m$ is the number of theta lifts. Over a host field there is a tower
of hermitian spaces: definite spaces $X_0,\dots,X_m$, whose unitary
groups carry algebraic automorphic forms, and spaces $W_1,\dots,W_m$
of signature $(p_j,q_j)=(k+j-1,m-j)$, whose unitary groups carry
Shimura varieties (Definition~\ref{def:spaces}). A \emph{package}
$\mathcal P$ is the automorphic input at one scale: a host, a ranks
datum, the tower, forms on the unitary groups of the tower with
their integral structures and theta lifts, and the hermitian
lattices and period maps through which the lifts are measured
(Definition~\ref{def:package}). A \emph{family at} $\eps$, written
$\Fam$, consists of a threshold, an unbounded set $\Win\subset\R$,
and for each $x$ a finite set of members, each a package with a size
datum at scale $x$, subject to uniformity conditions; it \emph{has
windows of exponent} $\lambda$ if $\Win$ does
(Definition~\ref{def:family}).

The inequalities that a package satisfies are collected in a
\emph{slope record}: a list of real numbers subject to eleven
inequalities (Definitions~\ref{def:record} and~\ref{def:size-data}).
A family is \emph{tied at} $x$, with constants $e_D,e_L,c_W,c_S,
e_K,C_L$, if there is a slope record at $x$ whose numbers are the
slopes, heights and local sizes of members of the family at $x$, one
member for each shift, and which satisfies the hypotheses listed in
Definition~\ref{def:tied}. The family satisfies the \emph{uniform
explicit-formula condition} for a level $\lev$ if, with one constant
$C$, the host of every member at every large $x$ satisfies
$\mathrm{EF}(\lev,\Top,x,\eps,C)$ (Definition~\ref{def:uniform-ef}).

\subsection*{Classes of families and the top existence statement}

Let $\lambda>1$ and let $\mathscr R$ be a host requirement. The
following terms are those of Definition~\ref{def:admissible}.
\begin{enumerate}[label=(\arabic*)]
\item A family \emph{satisfies the Sha hypothesis} if there is $x_0$
such that for every $x\ge x_0$ the host field of every member at $x$
has Sha budget $(t,b^{1/2})$, with $t$ the rank parameter and
$b=\exp(x^\eps)$.
\item The \emph{class of $\mathscr R$ at $\lambda$}, written
$\mathcal C_{\mathscr R}(\lambda)$, consists of the families with
windows of exponent $\lambda$ all of whose members beyond the
threshold meet $\mathscr R$.
\item A class $\mathcal C$ of families \emph{supplies tied families}
if there are real numbers $e_D,e_L$ and $\eps_1>0$ such that for
every $\eps\in(0,\eps_1)$ some family at $\eps$ in $\mathcal C$ is
tied at every large $x\in\Win$, with these $e_D,e_L$ and some
constants.
\item The \emph{top existence statement} is: for every $\lambda>1$
and every host requirement $\mathscr R$, if $\mathscr R$ is supplied
on windows of exponent $\lambda$ and every family of
$\mathcal C_{\mathscr R}(\lambda)$ satisfies the Sha hypothesis, then
$\mathcal C_{\mathscr R}(\lambda)$ supplies tied families.
\item \emph{Admissible-family data} for a class $\mathcal A$ of
families consist of realisation data, constants and a list of
propositions about the families of $\mathcal A$, in four groups:
automorphic forms, theta lifts and mass formulas; Galois
representations and local--global compatibility; Arakelov geometry
and covolume; counting and design.
\end{enumerate}
The definitions of packages, size data, families, the covolume
condition, host supply and admissible-family data are
\emph{schematic} in the sense of Section~\ref{sec:conventions}: they
list the parts of the data and the conditions that the argument
uses, and omit the remaining compatibility conditions. The
consequences for the statements of this chapter are collected in
Remark~\ref{rem:bb-schematic}.

\section{Black box: Galois representations}
\label{sec:bb-galois}

The first box supplies Galois representations attached to automorphic
representations over CM fields. Two notions are needed to state it;
both are defined in Chapter~\ref{ch:galois}, and both are weaker than
their counterparts in the literature.

A \emph{weak local Langlands map} for a number field $E$ and a finite
place $w$ is a map from the irreducible admissible representations of
the groups $\mathrm{GL}_n(E_w)$, $n\ge1$, to Weil--Deligne
representations of the Weil group of $E_w$ that satisfies a list of
laws: it preserves dimension, is local class field theory on
characters, is compatible with the Zelevinsky classification, and has
further properties on a class of \emph{good} representations
(Definition~\ref{def:weak-llc}). It is not required to be bijective,
and no condition on $L$- or $\epsilon$-factors is imposed.
\emph{Truncated local--global compatibility} is the statement that
for every CM field $E$, every prime $l$, every isomorphism
$\iota\colon\overline\Q_l\cong\C$, every $a\ge1$ and every regular
algebraic polarized cuspidal automorphic representation $\pi$ of
$\mathrm{GL}_a(\mathbb A_E)$ there is a \emph{compatible system} for
$\pi$: a continuous semisimple representation
$r\colon G_E\to\mathrm{GL}_a(\overline\Q_l)$ with prescribed
Hodge--Tate weights that matches $\pi_w$, through every weak local
Langlands map, at the finite places $w$ in the cases listed in
Definition~\ref{def:truncated-lgc}. Nothing is asserted about the
monodromy operator except where it vanishes, and nothing about
purity.

\begin{theorem}[Alp\"oge; restated from Theorem~\ref{thm:galois-reps}]
\label{bb:galois-reps}
\begin{enumerate}[label=(\alph*)]
\item For every number field $E$ and every finite
place $w$ of $E$, a weak local Langlands map for $(E,w)$ exists.
\item Truncated local--global compatibility
(Definition~\ref{def:truncated-lgc}) holds.
\end{enumerate}
\end{theorem}

\noindent\emph{What it says.} Regular algebraic polarized cuspidal
automorphic representations over CM fields have Galois
representations with the local properties listed in
Definition~\ref{def:truncated-lgc}, and local Langlands maps exist in
the weak sense above.

\noindent\emph{Why it is needed.} It is used twice. The Sha bound of
Section~\ref{sec:bb-sha} bounds the eigenspaces of the wrong parity
through automorphic points of ordinary deformation rings; these
points are the Galois representations attached to regular algebraic
polarized cuspidal representations of $\mathrm{GL}_3$ and
$\mathrm{GL}_4$ over imaginary quadratic fields, with the local
properties that part~\ref{item:gr-lgc} supplies
(Remark~\ref{thm:lgc-to-sha}). The admissible-family data of
Section~\ref{sec:bb-families} use the weak local Langlands maps of
part~\ref{item:gr-llc}, and their congruence link is listed with the
propositions on Galois representations and local--global
compatibility.

\noindent\emph{Where it is proved.} Chapter~\ref{ch:galois}, where
the construction of part~\ref{item:gr-lgc} is outlined at the level
of its named steps and compared with the published constructions;
the proof of part~\ref{item:gr-llc} is not outlined there. The
statement is not taken as a hypothesis from the literature
(Remark~\ref{rem:galois-inputs}).

\section{Black box: the Sha bound on windows}
\label{sec:bb-sha}

Let $E$ be a number field, $l$ a prime and $e\in\Z$. Let $S_l$ be the
set of places of $E$ above $l$ and $G_{E,S_l}$ the Galois group of
the maximal extension of $E$ unramified outside $S_l$ and the
infinite places. The \emph{Tate--Shafarevich group} of the twist $e$
is
\[
\Sha_l(E,e)=\Sha^1\bigl(G_{E,S_l},\Q_l/\Z_l(-e)\bigr),
\]
the group of classes in $H^1(G_{E,S_l},\Q_l/\Z_l(-e))$ whose
restrictions to the decomposition groups at $S_l$ vanish, and
$\shalen_l(e)\in\Z_{\ge0}\cup\{\infty\}$ is its length as a
$\Z_l$-module (Definition~\ref{def:sha}).

Let $t\ge1$ be an integer and $B\ge0$ real. The field $E$ has
\emph{Sha budget} $(t,B)$ if $\shalen_p(e)<\infty$ for every prime
$p$ and every $0<e<t$, and
\[
\sum_{l\in S}\log l\sum_{1\le e<t}\shalen_l(e)\le B
\]
for every finite set $S$ of primes. Let $C_t>0$. A primitive
Dirichlet character $\psi$ has \emph{Sha bounds on windows with
constant} $C_t$ if the following holds for every $\eps\in(0,1/52)$
and every function $\nu$ on $[1,\infty)$ such that $\nu(x)$ is
adapted to $(\eps,x)$ for every $x\ge1$. There are $J\ge1$ and $X_0$
such that for every $X\ge X_0$ and every window pair $(q,d)$ for
$(\eps,X^{1/J})$ there is a real $Y$ with $X^{1/J}\le Y$ and
$Y^2\le X$ such that, for every $x\in[Y,Y^2]$, every $\psi$-field for
$(\psi,\nu(x),q,d)$ has Sha budget
$\bigl(\lfloor x^{C_t}\rfloor,\exp(x^\eps)^{1/2}\bigr)$
(Definition~\ref{def:sha-windows}).

\begin{theorem}[Alp\"oge; restated from Theorem~\ref{thm:sha-windows}]
\label{bb:sha-windows}
Every primitive Dirichlet character has Sha bounds on windows with
constant $401$ (Definition~\ref{def:sha-windows}).
\end{theorem}

\noindent\emph{What it says.} On some window $[Y,Y^2]$ inside every
range $[X^{1/J},X]$, the CM fields of degree about $x^{2+8\eps}$
attached to $\psi$ satisfy
$\sum_l\log l\sum_{1\le e<x^{401}}\shalen_l(e)\le\exp(x^\eps/2)$. The
window may depend on $X$ and on the pair $(q,d)$, but the exponent
$J$ may not. By Poitou--Tate duality this is a bound, uniform in $l$,
on the subgroups $\Sha^2$ of the groups
$H^2(G_{E,S_l},\Z_l(e+1))$ (Remark~\ref{rem:poitou-tate}).

\noindent\emph{Why it is needed.} The congruence defects at the
twisted places of the families of Section~\ref{sec:bb-families} are
bounded by lengths of Selmer groups, and these by the Sha budget of
the host field. The budget enters the argument as part of the
requirement $\mathscr R_\psi$ on hosts, so the box is the hypothesis
of the supply of hosts (Section~\ref{sec:bb-hosts}); through the Sha
hypothesis on families it is then spent in the divisibility bound of
the upper link (Remark~\ref{rem:sha-upper-link}). The box is also the
reason for the windows. For a bounded number of exceptional Tate
twists the only bound available is doubly exponential, and a
pigeonhole argument gives the bound only on windows of scales
(Remark~\ref{rem:windows-origin}). The size of the bound is generous:
heuristics predict a total polynomial in $x$, far below
$\exp(x^\eps/2)$ (Remark~\ref{rem:budget}). The difficulty is
uniformity in the pair $(q,d)$ and in $l$, above all for the
eigenspaces of the wrong parity, where no $L$-value controls the
group.

\noindent\emph{Where it is proved.} Chapter~\ref{ch:sha}. The proof
splits the budget into a $2$-part, a minus part and a plus part, and
sorts the plus eigenspaces into six cases. Its automorphy input is
Theorem~\ref{bb:galois-reps}, part~\ref{item:gr-lgc}. The remaining
inputs are classical Galois cohomology and the theorems stated in
that chapter.

\section{Black box: the supply of hosts}
\label{sec:bb-hosts}

The host fields are produced in two stages, and the second stage is
the box that the proof of Theorem~\ref{thm:main} calls. The first
stage, in Chapter~\ref{ch:hosts}, shows that for every primitive
$\psi$, every large frame exponent and every suitable pair of
auxiliary primes $(q,d)$ there is a host whose characters trivial on
the core are the twists $\xi\chi_q^i\chi_d^j$ of the powers $\xi$ of
$\psi$ by the quadratic characters of $q$ and $d$
(Theorem~\ref{thm:fields-exist}). The bounded characters are twists
with $(i,j)\ne(0,0)$ (Proposition~\ref{prop:bounded-twists}), and the
explicit-formula condition asks for a bound on them at the scale
$x$. This is arranged by the choice of $(q,d)$. A window
pair $(q,d)$ for $(\eps,Y)$ with $\gcd(M,qd)=1$ is \emph{good} for
$(M,\eps,Y,X)$ if the twists of the characters modulo $M$ satisfy
$|g_\chi(y)|\le y^{1/2+\eps/2}$ and
$|L'(1,\chi)/L(1,\chi)|\le4\log y$ for every $y\in[Y,X]$, and the
twists of the real primitive characters of conductor dividing $8M$
have zero-free boxes of prescribed heights
(Definition~\ref{def:good-pair}).

\begin{theorem}[Alp\"oge; restated from Theorem~\ref{thm:good-pairs}]
\label{bb:good-pairs}
Let $M>0$ be an integer, $0<\eps<1/52$ and $\Delta\ge1$. There is
$Y_1$ such that for every $Y\ge Y_1$ some window pair for $(\eps,Y)$
is good for $(M,\eps,Y,Y^\Delta)$.
\end{theorem}

\noindent The proof of Theorem~\ref{bb:good-pairs}, outlined in
Chapter~\ref{ch:hosts}, is an averaging over the auxiliary primes:
since $\log q\approx Y^{\eps/5}$, a test point $y\in[Y,Y^\Delta]$ is
polylogarithmic in the conductors of the twists, and at such a scale
a twisted prime sum is small for most $q$ and $d$ by a mean-square
estimate. No hypothesis of Riemann type enters
(Remark~\ref{rem:host-input-strength}).

\begin{theorem}[Alp\"oge; restated from Theorem~\ref{thm:host-supply}]
\label{bb:host-supply}
Let $\psi$ be a primitive Dirichlet character. If $\psi$ has Sha
bounds on windows with constant $401$
(Definition~\ref{def:sha-windows}), then
$\mathscr R_\psi$ is supplied on windows of exponent $3/2$.
\end{theorem}

\noindent\emph{What it says.} Given the Sha bound for $\psi$, there
is, for every $\eps\in(0,1/52)$, a set of scales with windows
$[W,W^{3/2}]$ at every large scale of which there is a good host for
$\psi$ that satisfies the explicit-formula condition
$\mathrm{EF}(\lev_\psi,\Top,x,\eps,4)$ and whose field has Sha budget
$(t,b^{1/2})$.

\noindent\emph{Why it is needed.} It is one of the two hypotheses of
the zero-freeness theorem for primitive characters
(Theorem~\ref{bb:zero-free-primitive}): the families to which the
Chebyshev inequality is applied are built on these hosts, and the
explicit-formula condition and the Sha budget of the hosts become the
uniform explicit-formula condition and the Sha hypothesis of the
families.

\noindent\emph{Where it is proved.} Chapter~\ref{ch:families}, by an
outline that combines Theorem~\ref{thm:fields-exist},
Theorem~\ref{bb:good-pairs} and
Proposition~\ref{prop:bounded-twist-bounds} of
Chapter~\ref{ch:hosts} with the Sha bound. One pair $(q,d)$ serves
both the goodness and the Sha window. The Sha bound gives windows
$[Y,Y^2]$, and the supply of hosts is stated with windows
$[W,W^{3/2}]$. The step that turns the exponent $2$ into $3/2$ is
not reconstructed in this book. It is
Question~\ref{q:window-exponent}, which follows the proof of
Theorem~\ref{thm:host-supply} in Chapter~\ref{ch:families}.

\section{Black box: the covolume bound}
\label{sec:bb-covolume}

For a member of a family and a lift index $1\le j\le m$, let $M_j$ be
the lattice of integral forms on the indefinite unitary group
$U(W_j)$, of signature $(p_j,q_j)=(k+j-1,m-j)$, formed at a
\emph{point} of the member, that is, at a choice of auxiliary
integral data. The \emph{point conditions} are compatibility
conditions on such a choice. A point carries weights $w(j,\rho)$, a
period $\Omega$ and a positive integer $\deg$, its degree. It
satisfies the \emph{covolume inequality} with constant $C$ at $x$ if
for every lift index $j$
\[
\frac{\log\operatorname{covol}(M_j)}{\deg}\le\sum_{\rho\in\Phi}
\bigl(p_j-q_j\bigr)\,w(j,\rho)\log\Omega+C\,B_0\log\log x,
\]
where $\log\operatorname{covol}(M_j)$ is the normalised logarithmic
covolume of $M_j$ and $B_0$ is a size quantity of the member. The
\emph{covolume condition} is the statement that this inequality, for
all large $x$ and all large $C$, is a \emph{closed} property on a
fixed class $\mathcal C_{\mathrm{cov}}$ of families: for $\eps$ small
it holds for every family of the class whose constants exceed
thresholds chosen before the family (Definitions~\ref{def:closed}
and~\ref{def:covolume-condition}). The \emph{member hypothesis} of
that definition bounds the size of the coefficient field of a
package and fixes its local types at the two norm primes; the period
$\Omega$ is the CM period of the slope record, with
\[
\log\Omega=\gamma+\log(\pi/v)-\frac{L'(1,\chi_{-d})}{L(1,\chi_{-d})},
\]
where $v$ is the half root discriminant of the host field
(Remark~\ref{rem:stirling}).

\begin{theorem}[Alp\"oge; covolume bound; restated from
Theorem~\ref{thm:covolume}]
\label{bb:covolume}
There is $\eps_1>0$ such that for every
$\eps\in(0,\eps_1]$ there is a threshold $\kappa$ with the following
property. Let $\Fam$ be a family at $\eps$ in $\mathcal
C_{\mathrm{cov}}$ and $c$ a real number such that $\Fam$ is large for
$(\kappa,c)$. Then there is $C_0\ge0$ such that for all large $x$,
every member $(\mathcal P_{x,i},Z_{x,i})$ satisfying the member
hypothesis, every point of $\mathcal P_{x,i}$ satisfying the point
conditions and every $C\ge C_0$, the point satisfies the covolume
inequality with constant $C$ at $x$
(Definition~\ref{def:covolume-condition}\ref{item:cc-ineq}).
\end{theorem}

\noindent\emph{What it says.} The covolume condition
(Definition~\ref{def:covolume-condition}) holds: the property of a
family stated in the last sentence of the theorem is closed on
$\mathcal C_{\mathrm{cov}}$. For the lattice $M_j$ of integral cusp
forms of the weights $w(j,\rho)$ on the Shimura variety of $U(W_j)$,
at the level prescribed by the package, the normalised logarithmic
covolume is at most $\sum_{\rho\in\Phi}(p_j-q_j)\,w(j,\rho)
\log\Omega$, up to an error $O(B_0\log\log x)$, for all large $x$ and
uniformly over the members of a family whose field of coefficients is
not too large.

\noindent\emph{Why it is needed.} It is the one result through which
arithmetic geometry enters the argument
(Remark~\ref{rem:covol-supply}). Summed over $j$, its main terms give
the term $-ma_\tau\log\Omega$ of the divisor bound of the upper link
(Definition~\ref{def:upper-link}), which is the archimedean side of
the comparison. In the assembly it is an input to the
admissible-family data: the propositions of their Arakelov group are
obtained from the covolume condition
(Theorem~\ref{bb:packages}, part~\ref{item:packages-data}).

\noindent\emph{Where it is proved.} Chapter~\ref{ch:covolume}. The
proof is not reproduced in this book; the chapter lists its inputs,
which are cited theorems of arithmetic geometry (among them toroidal
compactifications, resolution of singularities, the arithmetic
Hilbert--Samuel theorem, a bigness theorem, vanishing theorems and
Grothendieck--Messing theory), and records the form in which each is
used (Remark~\ref{rem:ag-strength}).

\section{Black box: packages, Green's parametrisation, tied families}
\label{sec:bb-families}

Three results produce the families. The first is a theorem of the
literature on the characters of $\mathrm{GL}_n(\F_p)$.

\begin{theorem}[Green \cite{Green1955}; restated from
Theorem~\ref{thm:green}]
\label{bb:green}
Let $p$ be a prime and $n\ge1$. There is a map $P\mapsto\chi_P$ from
the monic irreducible polynomials $P\in\F_p[X]$ of degree $n$ with
$P(0)\ne0$ to the cuspidal irreducible characters of
$\mathrm{GL}_n(\F_p)$ such that
\begin{enumerate}[label=(\alph*)]
\item $\chi_P(1)=\prod_{i=1}^{n-1}(p^i-1)$ for every such $P$;
\item $P\mapsto\chi_P$ is injective;
\item every cuspidal irreducible character of $\mathrm{GL}_n(\F_p)$ is
$\chi_P$ for some such $P$.
\end{enumerate}
\end{theorem}

\noindent Green's theorem is not an analytic input, and the window
exponent $3/2$ does not come from it. It enters only through local
types at the norm primes $\varpi_0,\varpi_1$ of a host: it is used
to show that the rank defect of each
member of a family is small, and to label the local types in the ties
(Remark~\ref{rem:green}).

\begin{theorem}[Alp\"oge; restated from Theorem~\ref{thm:packages}]
\label{bb:packages}
Let $\eps\in(0,1/52)$.
\begin{enumerate}[label=(\alph*)]
\item For every scale $(x,\eps,J)$ with
$x$ sufficiently large there are packages
(Definition~\ref{def:package}) whose automorphic forms have the
weights, level and local types prescribed at that scale, on the tower
of unitary groups of Definition~\ref{def:spaces} over a CM field of
degree about $x^{2+8\eps}$.
\item The members of the families built
from these packages have points, and for each lift index $j$ the
lattice $M_j$ of the covolume condition
(Definition~\ref{def:covolume-condition}) is the lattice of integral
cusp forms on $U(W_j)$.
\item There is a class of families of
these packages that carries admissible-family data
(Definition~\ref{def:admissible}).
\end{enumerate}
\end{theorem}

\noindent\emph{What it says.} At every large scale there are
automorphic forms on the tower of unitary groups over a CM field of
degree about $x^{2+8\eps}$ with the prescribed weights, level and
local types. Their integral cusp forms give the lattices $M_j$ of the
covolume bound, and families of these packages satisfy the
propositions that make up admissible-family data.

\begin{theorem}[Alp\"oge; restated from Theorem~\ref{thm:tied-families}]
\label{bb:tied-families}
Suppose that some class $\mathcal A$ of families carries
admissible-family data (Definition~\ref{def:admissible}). Let
$\lambda>1$, and let $\mathscr R$ be a host requirement
(Definition~\ref{def:host-supply}) that is supplied on windows of
exponent $\lambda$ and such that every family of
$\mathcal C_{\mathscr R}(\lambda)$ satisfies the Sha hypothesis. Then
$\mathcal C_{\mathscr R}(\lambda)$ supplies tied families.
\end{theorem}

\noindent\emph{What it says.} Suppose some class of families carries
admissible-family data. Let $\mathscr R$ be a host requirement that
is supplied on windows of exponent $\lambda$, and suppose the host
fields of the families built on these hosts satisfy the Sha bound.
Then for every small $\eps$ there is a family of packages, built on
the supplied hosts, that is tied at every large scale of its windows.
For the families produced, $e_D=2$ and $e_L=0$.

\noindent\emph{Why they are needed.} Theorem~\ref{bb:packages},
part~\ref{item:packages-data}, provides the first hypothesis of
Theorem~\ref{bb:tied-families}, the class $\mathcal A$. The rest of
Theorem~\ref{bb:tied-families}, read for every $\lambda>1$ and every
host requirement $\mathscr R$, is the top existence statement of
Section~\ref{sec:bb-vocabulary}, and this statement is a hypothesis
of the zero-freeness theorem for primitive characters
(Theorem~\ref{bb:zero-free-primitive}). The top existence
statement does not mention the class $\mathcal A$, so one class
carrying admissible-family data suffices
(Remark~\ref{rem:quantifiers}). Supplying tied families cannot hold
vacuously: a closed property chooses its threshold before the family,
the family produced is large for that threshold, and an empty class
supplies no tied family.

\noindent\emph{Where they are proved.} Chapter~\ref{ch:families}.
For Theorem~\ref{bb:packages} the chapter indicates the standard
method for parts~\ref{item:packages-exist}
and~\ref{item:packages-points} and lists the inputs of
part~\ref{item:packages-data}; among them are the covolume bound
(Theorem~\ref{bb:covolume}) and the weak local Langlands maps of
Theorem~\ref{bb:galois-reps}. Theorem~\ref{bb:tied-families} is
proved in outline. The links of a family
(Definition~\ref{def:links}) are closed properties, so one threshold
serves for all of them. Green's parametrisation enters in the closed
form of the irreducibility link. The design statement of the data,
fed by the supplied hosts, gives a family that is large for the
threshold. The global link shows that this family is tied.

In terms of the objects of
Chapters~\ref{ch:theta}--\ref{ch:chebyshev}, the data encode nine
steps:
\begin{enumerate}[label=(\arabic*)]
\item build the packages and their lattices;
\item form the period map;
\item apply the slope inequality (Theorem~\ref{thm:frame});
\item bound the archimedean side;
\item bound the covolume (Theorem~\ref{bb:covolume});
\item bound the finite side;
\item balance the two sides (Theorem~\ref{bb:gap});
\item bound the defects at the twisted places by Sha lengths and
the Sha hypothesis;
\item supply near-orthogonality, good local types, masses and class
numbers.
\end{enumerate}

\section{Black box: the Chebyshev inequality on windows}
\label{sec:bb-chebyshev}

This box contains the main inequality of the argument. Its
hypotheses are stated in the
terms of Section~\ref{sec:bb-vocabulary}: a level, a family at
$\eps$, the uniform explicit-formula condition and the ties.

\begin{theorem}[Alp\"oge; Chebyshev inequality on windows; restated
from Theorem~\ref{thm:chebyshev-windows}]
\label{bb:chebyshev-windows}
Let $\lev$ be a level, $e_D$ and $e_L$ real numbers and $\eps>0$. Let
$\Fam$ be a family at $\eps$ that satisfies the uniform
explicit-formula condition for $\lev$ (Definition~\ref{def:uniform-ef}).
Suppose that there are real numbers $c_W,c_S,e_K,C_L$ and $x_0$ such
that $\Fam$ is tied (Definition~\ref{def:tied}) with constants
$e_D,e_L,c_W,c_S,e_K,C_L$ at every $x\in\Win$ with $x\ge x_0$. Then
there is $x_1$ such that
\[
G_\lev(x)\le x^{1/2+(e_D+e_L+7)\eps}
\]
for every $x\in\Win$ with $x\ge x_1$.
\end{theorem}

\noindent\emph{What it says.} If a family is tied at every large
point of its window set, with exponents $e_D$ and $e_L$ that do not
depend on the point, then the weighted Chebyshev function of the
level is at most $x^{1/2+O(\eps)}$ on that set. The bound is
one-sided and holds only for $x\in\Win$.

\noindent\emph{Why it is needed.} It is the only source of a
square-root bound in the argument. Applied to $\lev_\psi$ and to the
tied families of Theorem~\ref{bb:tied-families}, for which
$e_D+e_L+7=9$, it gives $G_{\lev_\psi}(x)\le x^{1/2+9\eps}$ on
windows of exponent $3/2$, which is the hypothesis of the windowed
Landau criterion. The theorem would say nothing for a family whose
window set had no large elements, or if $e_D$ and $e_L$ could grow as
$\eps\to0$. Neither happens: the window set of a family is unbounded,
and in the top existence statement $e_D$ and $e_L$ are fixed before
$\eps$ is sent to $0$ (Remark~\ref{rem:no-vacuity}).

\noindent\emph{Where it is proved.} Chapter~\ref{ch:chebyshev}, in
outline, from the four results restated below. The proof first
forgets the geometry: at each $x\in\Win$ the tie provides a slope
record, the arithmetic data of a host and auxiliary real numbers,
and after that only inequalities between real numbers are used.

\subsection*{Inside the box}

The inequality is assembled from an identity and three bounds. They
are restated here because the route through them is the core of the
argument; the proof of Theorem~\ref{thm:main} in
Section~\ref{sec:proof-main} calls only
Theorem~\ref{bb:chebyshev-windows}. Their statements use the
following terms. Let $H$ be a finite abelian group,
$V_4=(\Z/2\Z)^2$ and $\GG=V_4\times H$, with the index-two subgroup
$\Phi$ of the elements whose first coordinate is $0$, the \emph{CM
type}. \emph{Arithmetic data} $D$ for $\GG$ consist of natural
numbers $q_1$ and $d$ and, for each prime $p$, an inertia subgroup
$I_p\subset\GG$ and an element $\Frob_p\in\GG$; for a host they are
read off from the Artin map, with $H=A$ and $\Phi=\Gal(E/K)$
(Definitions~\ref{def:additive} and~\ref{def:arithdata}). For a
function $w$ on $\GG$ that is positive on $\Phi$, a point $\tau\in H$
and $a>0$, Definition~\ref{def:ask} attaches to $D$ a \emph{prime
sum}
\[
\begin{aligned}
\PS_{\tau,w}(a)={}&\sum_{p\le N\text{ inert}}2\log p
\sum_{\substack{1\le i\le N\\ i\Frob_p\notin\Phi}}
\Bigl(\frac{a}{p^i}-m_\tau(i\Frob_p)\Bigr)_+\\
&+(\text{a term for the prime }d),
\end{aligned}
\]
a \emph{main term}
\[
\MT_w(a)=a\Bigl(\log a-1-\gamma+\ell(\chi_K)-\frac1f
\sum_{\rho\in\Phi}\log w(\rho)\Bigr),
\]
and a \emph{defect} $\DF_{\tau,w}(a)$, which is a combination, with
coefficients $\overline{\theta(g)}\,m_\tau(g)/f$, of the values
$\GPsi_\theta(a/m_\tau(g))$ over all characters $\theta$ of $\GG$ and
all $g\notin\Phi$. Here $m_\tau$ is the partner weight, $f=|\Phi|$,
and $\GPsi_\theta$ is the arithmetic Chebyshev function attached to
$\theta$ by $D$ and a function $\ell$ on the characters of $\GG$. It
is built from the \emph{averaged values} $\theta_D(n)$ of $\theta$ at
the integers $n\ge1$ (Definition~\ref{def:arithdata}). We write
$\theta_0$ for the trivial character of $\GG$. The function
$\GPsi_{\theta_0}$ is $\Gz$. For a host one has
$\theta_D(n)=\chi_\theta(n)$, and with
$\ell(\theta)=L'(1,\chi_\theta)/L(1,\chi_\theta)$ the function
$\GPsi_\theta$ is $g_{\chi_\theta}$ for $\theta\ne\theta_0$.

\begin{proposition}[balance identity; restated from
Proposition~\ref{prop:balance}]
\label{bb:balance}
Let $D$ be arithmetic data for $\GG=V_4\times H$ that are split outside
$d$ on a finite set $S$ of primes, let $\ell$ be a function from the
characters of $\GG$ to $\C$, let $a>0$ and $\tau\in H$, let
$w\colon\GG\to\R$ be positive on $\Phi$, and let $N$ be an integer
with $a/w(\rho)\le N$ for every $\rho\in\Phi$. Then
\[
\PS_{\tau,w}(a)-\MT_w(a)=\DF_{\tau,w}(a).
\]
\end{proposition}

\noindent The prime sum has the shape of Legendre's formula for
$\log\lfloor a\rfloor!$, restricted to the primes inert in $K$ and
smoothed (Remark~\ref{rem:legendre}). The identity is finite Fourier
analysis on $\GG$ and is proved completely in
Chapter~\ref{ch:balance}. Through it, an inequality
$\PS\le\MT+(\text{small})$ is an upper bound for a combination of the
functions $\GPsi_\theta$. For a slope record with shifts $u_0\in H$,
the \emph{gap} at $u_0$ is the sum over the top places $\tau$ and
the lift indices $j<m$ of the differences $\PS-\operatorname{Re}\MT$
at $a=a_\tau$ and $w=w(u_0,j)$ (Definition~\ref{def:thick} and
Proposition~\ref{prop:main-terms}), and the \emph{error bracket}
$\Err$ is the explicit quantity of Definition~\ref{def:restsize}.

\begin{theorem}[Alp\"oge; gap bound; restated from
Theorem~\ref{thm:gap}]
\label{bb:gap}
Let $H$ be a finite abelian group of order $n$ and $\GG=V_4\times H$.
Let a slope record with shifts $U_0=H$, positive integers $m$ and $f$
and size constants be given (Definitions~\ref{def:record}
and~\ref{def:size-data}); write $x$, $\eps$, $a$, $b$ and $k$ for
the quantities of the record, and $c_{\mathrm{sp}}$,
$e_{\mathrm{div}}$ and $e_{\mathrm{loss}}$ for its size constants. Let
$c_W$ and $e_K$ be real numbers. Let $D$ be arithmetic data for $\GG$,
$\ell$ a function from the
characters of $\GG$ to $\C$, $N\ge1$ an integer, $\Top,\Top_2\subset
H$ finite sets, $a\colon H\to\R$ and
$w\colon H\times\{0,\dots,m-1\}\times\GG\to\R$.
Suppose that the record satisfies the size data
(Definition~\ref{def:size-data}), that the identification clause, the
window conditions with $c_W$ and $e_K$ and the bounds on $B_0$ and
$N_c$ (Definition~\ref{def:window-hyps}) hold, and that
\[
\log x\le x^\eps\qquad\text{and}\qquad
200(1+c_W)^2(1+c_{\mathrm{sp}})(1+e_K)\le x^\eps.
\]
Then for every $u_0\in H$
\[
\gap(u_0)\le x^{(e_{\mathrm{div}}+e_{\mathrm{loss}}+3)\eps}\,\Err,
\]
with $\Err$ the error bracket of Definition~\ref{def:restsize} for
$n=|H|$, $n_T=|\Top|$, $n_{T_2}=|\Top_2|$, the given $b$, $k$ and
$m$, and $a_E=a$.
\end{theorem}

\noindent The slope record is spent in the gap bound. The
record contains two bounds for one number, the frame slope $s$ of
the period lattice data. A lower bound comes from the divisor bound
and the quotient slopes; it carries the finite-side gain, which
bounds the prime sum, and the archimedean sum. An upper bound comes
from the slope inequality for the period maps into the test lattice.
Subtracting them removes the slope of the test lattice and the
heights, and leaves an inequality that one may read as ``finite-side
gain at most archimedean main term'' (Remark~\ref{rem:cancel}). The
rest of the proof is linear arithmetic under the window conditions.
The slope inequality has the form of the conclusion of the following
theorem, in which $\mu$,
$\mu_{\max}$, $\mu_{\min}$ are the slope and the extreme slopes of a
hermitian lattice, $\operatorname{ht}$ is the height of a map
(Definition~\ref{def:hlat}), and the frame operator and the
Hilbert--Schmidt defect are those of
Definition~\ref{def:frame-operator}.

\begin{theorem}[Alp\"oge; frame theorem; restated from
Theorem~\ref{thm:frame}]
\label{bb:frame}
Let $L$ be a totally complex number field and $\sigma\colon
L\hookrightarrow\C$ an embedding. Let $W$ be an $L$-vector space with a basis
indexed by a finite
set $J$, with a positive definite hermitian Gram matrix at each
infinite place, and let $R\subset W$ be a finitely generated
$\mathcal O_L$-submodule spanning $W$. Let $I$ be a nonempty finite
set, and for $i\in I$ let $V_i$ be an $L$-vector space with a basis
indexed by a finite set $\iota_i$ and positive definite hermitian
Gram matrices, $N_i\subset V_i$ a finitely generated $\mathcal
O_L$-submodule spanning $V_i$, $B_i\colon V_i\to W$ an $L$-linear map,
and $h_i$ a real number; assume $\sum_i|\iota_i|>0$. Let $J_1$ be a
finite set, and let $Y'$ and $Y_i$ ($i\in I$) be invertible complex
matrices indexed by $J$ and by $\iota_i$. For each $i$ let
$\widetilde B_i$ be a complex matrix with rows indexed by $J\sqcup J_1$
and columns by $\iota_i$, whose $J$-rows form the matrix
$Y'^{-1}[B_i]_\sigma Y_i$, where $[B_i]_\sigma$ is the matrix of $B_i$
in the given bases, embedded by $\sigma$. Let $c_i>0$ for $i\in I$, let $Q$ be
the frame operator of $(c_i,\widetilde B_i)$
(Definition~\ref{def:frame-operator}), and let $s,u,\Delta,\beta_0,
\beta_1,\beta_2$ be real numbers. Suppose:
\begin{enumerate}[label=(\alph*)]
\item $B_i(N_i)\subset R$ for all $i$, and $\operatorname{ht}(B_i)\le
h_i$ whenever $B_i\ne0$;
\item $s\sum_i\rk N_i\le\sum_i\rk N_i\,(\mu(N_i)-h_i)$; $u\ge0$ and
$\mu_{\max}(N_i)-h_i\le s+u$ for all $i$; $\Delta\ge0$ and
$s-\mu_{\min}(R)\le\Delta$;
\item $0\le\beta_0<1$ and $(1-\beta_0)|J\sqcup J_1|\le|J|$;
\item the Hilbert--Schmidt defect of $(c_i,\widetilde B_i)$ is at most
$\beta_2\sum_i\rk N_i$, with $0\le\beta_2\le1/4$; and
$\operatorname{tr}(Q^2)\le(1+\beta_1)(\operatorname{tr}Q)^2/
|J\sqcup J_1|$, with $\beta_1\ge0$.
\end{enumerate}
Then
\[
s-\log|I|-\frac{1}{1-\beta_0}\Bigl(4u\log\Bigl(1+\frac{\Delta}{u}
\Bigr)+\bigl(\beta_1+4\sqrt{\beta_2}\bigr)\Delta\Bigr)\le\mu(R),
\]
where the term $4u\log(1+\Delta/u)$ is read as $0$ when $u=0$.
\end{theorem}

\noindent The frame theorem compares a rank-weighted mean of the
slopes of many source lattices with the slope of one target lattice.
It bounds an average by the slope of the target, and not by its
maximal slope. So it is not Bost's slope inequality applied to each
period map. Bost's inequality is one step of its proof
(Remark~\ref{rem:not-bost}). No small sections are constructed. The
period map between two given lattices is shown to be integral and of
bounded height, and the slope inequality is applied to it
(Remark~\ref{rem:periodmap}). Chapter~\ref{ch:theta} states the
theorem and gives an account of the method.

The next two results turn the gap bound into a bound for $G_\lev$.
For arithmetic data $D$, a function $\ell$, a subgroup
$A_0\subset H$, a set $\Xi$ of characters of $\GG$, the kept set, and
a real number $\varkappa$, the \emph{extracted function} $G_A$ is the
combination of the functions $\operatorname{Re}\GPsi_\theta$,
$\theta\in\Xi$, with weights $1$ and $\varkappa$, and the \emph{rest}
$\Sigma_A$ collects the characters trivial on $0\times A_0$ that are
not in $\Xi$ (Definition~\ref{def:kept-chars}).

\begin{theorem}[Alp\"oge; core Chebyshev bound; restated from
Theorem~\ref{thm:core-bound}]
\label{bb:core-bound}
Let $C\ge0$, $\eps>0$, $\varkappa$ with $|\varkappa|\le1$, and $C_E,C_o,M$ be
real
numbers. There is $x_1$ such that the following holds for every
$x\ge x_1$. Let $H$ be a finite abelian group, $\GG=V_4\times H$, $m$ a
natural number, $D$ arithmetic data for $\GG$ that are split outside
$d$ on a finite set of primes, $\ell$ a function from the characters
of $\GG$ to $\C$, $N\ge1$ an integer, $A_0\subset H$ a subgroup,
$\Top,\Top_2\subset H$ finite sets, $a\colon H\to\R$,
$w\colon H\times\{0,\dots,m-1\}\times\GG\to\R$,
$b,b',\eta,X^\flat,\vartheta,k,a_E$ real numbers and $\Xi$ a set of
characters of $\GG$. Let $G_A$ and $\Sigma_A$ be the extracted
function and the rest (Definition~\ref{def:kept-chars}) for $D$,
$\ell$, $A_0$, $\Xi$ and $\varkappa$. Suppose that
\begin{enumerate}[label=(\roman*)]
\item $b>0$, $a_E\ge0$, $k>0$, $m\le x^{1/2+\eps/2}$ and $xm\le k$;
\item the size conditions and the thin-scale bound of
Definition~\ref{def:thick} hold, and $\Top_2\subset A_0$;
\item the cost conditions hold with $f=|\Phi|$, and the rest-size
conditions hold with $n=|H|$, $n_T=|\Top|$ and $n_{T_2}=|\Top_2|$
(Definition~\ref{def:restsize});
\item the character conditions of Definition~\ref{def:kept-chars} hold
with the given $\varkappa,M,C_o,C_E,\eps$ and $x$;
\item $\gap(u_0)\le x^{C\eps}\,\Err$ for every $u_0\in A_0$, with
$\Err$ the error bracket of Definition~\ref{def:restsize} for these
$n,n_T,n_{T_2},m,b,k,a_E$.
\end{enumerate}
Then $G_A(x)\le x^{1/2+(C+3)\eps}+\Sigma_A(x)$.
\end{theorem}

\begin{theorem}[Alp\"oge; restated from
Theorem~\ref{thm:averaged-hyps}]
\label{bb:averaged-hyps}
Let $\lev$ be a level, $\phi$ a frame, and take a host for $\phi$
(Definition~\ref{def:host}) with its arithmetic data $D$ for
$\GG=V_4\times A$. Let $\Top\subset A$ be finite and $x>0$, $\eps$ and
$C$ real, and suppose that the host satisfies the explicit-formula
condition $\mathrm{EF}(\lev,\Top,x,\eps,C)$
(Definition~\ref{def:shares-ef}). Let $\ell$ be a function from the
characters of $\GG$ to $\C$ such that $\ell(\theta)=L'(1,\chi)/L(1,\chi)$
whenever $\theta\ne\theta_0$ and $\chi$ is a primitive Dirichlet
character with $\chi(n)=\theta_D(n)$ for all $n\ge1$. Let $\Xi$ be the
set of characters $\theta$ of $\GG$ for which some primitive Dirichlet
character $\chi$ with $\chi(n)=\theta_D(n)$ for all $n\ge1$ lies in
$\Kept(\lev)$, let $A_0$ be the core of $\phi$, and let $G_A$ and
$\Sigma_A$ be the extracted function and the rest
(Definition~\ref{def:kept-chars}) for $D$, $\ell$, $A_0$, $\Xi$ and
$\varkappa=(3o(\lev)+1)^{-1}$.
Then:
\begin{enumerate}[label=(\alph*)]
\item $G_A(y)=G_\lev(y)$ for every $y>0$;
\item $\Sigma_A(x)\le4o(\lev)\,x^{1/2+\eps}$.
\end{enumerate}
\end{theorem}

\noindent In the core bound the gap is written through the balance
identity, the arguments $a_\tau/w$ are moved to $x$, and averages
over the shifts $u_0\in A_0$ and the top places $\tau\in\Top$
project onto the characters trivial on the core. The kept characters
come out with the positive weights $1$ and $(3o+1)^{-1}$. This is
why an upper bound for the gap gives an upper bound for $G_A$.
Theorem~\ref{bb:averaged-hyps} then identifies $G_A$ with $G_\lev$
for a host that satisfies the explicit-formula condition, and bounds
the rest by the pointwise bounds on the bounded characters. The
three results combine as follows. Since $e_{\mathrm{div}}\le e_D$ and
$e_{\mathrm{loss}}\le e_L$, Theorem~\ref{bb:gap} bounds the gap by
$x^{(e_D+e_L+3)\eps}\Err$. Theorem~\ref{bb:core-bound} with
$C=e_D+e_L+3$ gives $G_A(x)\le x^{1/2+(e_D+e_L+6)\eps}+\Sigma_A(x)$.
Theorem~\ref{bb:averaged-hyps} then gives
\[
G_\lev(x)\le x^{1/2+(e_D+e_L+6)\eps}+4o\,x^{1/2+\eps}
\le x^{1/2+(e_D+e_L+7)\eps}
\]
for large $x\in\Win$. The last step absorbs the term
$4o\,x^{1/2+\eps}$, and for this it uses $e_D+e_L\ge0$, which holds
because $e_D,e_L\ge0$ by the window conditions.

\begin{remark}[the square-root gain]\label{rem:shape}
The factorial comparison alone yields Chebyshev bounds of the
right order, without a square-root error. Here the gain comes
from the dependence of the slope-record errors on the number
of theta lifts. After normalisation, one part contributes
$x/m$ and another contributes $m$
(Proposition~\ref{prop:sqrtx}). Their balance is the source of
the exponent $1/2$, as explained in
Remark~\ref{rem:exponent-half}.
\end{remark}

\noindent The bookkeeping behind Remark~\ref{rem:shape} is
elementary. In it $m\approx\sqrt x$ is chosen.

\begin{proposition}[the bookkeeping behind $m\approx\sqrt x$;
restated from Proposition~\ref{prop:sqrtx}]
\label{bb:sqrtx}
Let $x\ge1$, $\eps>0$, $b>0$, $k>0$ and $a_E\ge0$ be real numbers and
$n,n_T,n_{T_2},m$ positive integers. Suppose that $a_E\le2xb$,
$n_T\ge x^{-\eps/8}n/2$, $x^{1/2-\eps}\le m\le x^{1/2+\eps/2}$,
$k\le n^2x^{1/2-\eps}$, $xm\le k$ and $n_{T_2}\le n$. Then the error
bracket $\Err$ of Definition~\ref{def:restsize} satisfies
\[
\frac{\Err}{n_Tmb}\le\frac{2x}m+2+4m\,x^{9\eps/8}+2x^{1/2+9\eps/8}
\le10\,x^{1/2+13\eps/8}.
\]
\end{proposition}

\noindent The constraints on $m$ and $k$ in
Proposition~\ref{bb:sqrtx} are hypotheses; for the families of
Theorem~\ref{bb:tied-families} they are part of the ties.

\section{Black box: the windowed Landau criterion}
\label{sec:bb-landau}

The analytic end of the argument is a converse to the explicit
formula. Landau's theorem on Laplace--Stieltjes integrals with a
monotone integrator (Theorem~\ref{thm:landau}) gives such a converse
when the bound $G_F(y)\le Cy^{\sigma_0}$ is assumed for all large
$y$. The version used here assumes it only on windows.

\begin{theorem}[Alp\"oge; restated from
Theorem~\ref{thm:windowed-landau}]
\label{bb:windowed-landau}
Let $F=((\chi_i,a_i))_i$ be a character family, and let
$0<\sigma_0<1$, $\delta>0$ and $C\in\R$. Suppose that for every $Y_0$
there is $Y\ge Y_0$ with $G_F(y)\le Cy^{\sigma_0}$ for all
$y\in[Y,Y^{1+\delta}]$. Then $L(s,\chi_i)\ne0$ for every $i$ and every
$s$ with $\operatorname{Re}s>\sigma_0$, where $s\ne1$ when $\chi_i$ is
the character modulo~$1$.
\end{theorem}

\begin{theorem}[Alp\"oge; restated from
Theorem~\ref{thm:level-criterion}]
\label{bb:level-criterion}
Let $\lev$ be a level and $\lambda>1$. Then $\lev$ satisfies the
window criterion at $\lambda$.
\end{theorem}

\noindent\emph{What they say.} An upper bound $x^{\sigma_0}$ for a
positive combination of smoothed Chebyshev functions, on windows
$[Y,Y^{1+\delta}]$ with $Y$ arbitrarily large, excludes zeros with
real part greater than $\sigma_0$ for every character in the
combination. Theorem~\ref{bb:level-criterion} is the case of the
combination $G_\lev$ of a level, with $\delta=\lambda-1$ and $C=1$.

\noindent\emph{Why they are needed.} They convert the Chebyshev
inequality into zero-freeness. Two features are used. Only an upper
bound is assumed. This suffices because, with positive weights, the
residues of the relevant Laplace transform at real zeros in $(0,1)$
all have one sign, so a hypothetical real zero cannot be cancelled
by another (Remark~\ref{rem:positivity}). The bound is also assumed
only on windows. The classical form would not be enough, because
the Chebyshev inequality is proved only on the windows that the Sha
bound provides (Remark~\ref{rem:landau-compare}). Any exponent
$\lambda>1$ is accepted, and the argument uses $\lambda=3/2$.

\noindent\emph{Where they are proved.} Chapter~\ref{ch:landau}.
Theorem~\ref{bb:windowed-landau} is proved in outline, by comparing
two estimates for the Taylor coefficients of a Laplace transform at a
point far to the right; its steps are elementary complex analysis.
Theorem~\ref{bb:level-criterion} is deduced from it there by a
complete proof. Whether the windowed criterion is new as stated is
left open (Question~\ref{q:landau-new}).

\section{Black box: from a half-plane to the critical line}
\label{sec:bb-line}

For a primitive Dirichlet character $\chi$ modulo $N$, let
$\mathfrak a\in\{0,1\}$ be given by $\chi(-1)=(-1)^{\mathfrak a}$,
and put
$\xi(s,\chi)=(N/\pi)^{(s+\mathfrak a)/2}\Gamma((s+\mathfrak a)/2)
L(s,\chi)$. We write (H) for the following hypothesis:
$L(s,\chi')\ne0$ for every primitive Dirichlet character $\chi'$ and
every $s$ with $\operatorname{Re}s>1/2$, where $s\ne1$ when $\chi'$
is the character modulo~$1$.

\begin{theorem}[Riemann, Hecke \cite{Davenport2000,Riemann1859};
restated from Theorem~\ref{thm:fe}]
\label{bb:fe}
Let $\chi$ be a primitive Dirichlet character modulo $N$. Then there is
$\epsilon(\chi)\in\C$ with $|\epsilon(\chi)|=1$ and
$\xi(s,\chi)=\epsilon(\chi)\xi(1-s,\overline\chi)$ for all $s\ne0,1$.
\end{theorem}

\begin{proposition}[restated from Proposition~\ref{prop:imprimitive}]
\label{bb:imprimitive}
Let $\chi$ be a Dirichlet character modulo $N$, induced by the
primitive character $\chi^*$ modulo $N^*$.
\begin{enumerate}[label=(\alph*)]
\item For $s\ne1$, and for every $s\in\C$ if
$\chi$ is nonprincipal,
$L(s,\chi)=L(s,\chi^*)\prod_{p\mid N}(1-\chi^*(p)p^{-s})$. In
particular, for such $s$ with $\operatorname{Re}s\ne0$, $L(s,\chi)=0$
if and only if $L(s,\chi^*)=0$.
\item The $L$-function of the character
modulo~$1$ is $\zeta$ on $s\ne1$.
\item Under (H), $L(s,\chi)\ne0$ for every $s$
with $\operatorname{Re}s>1/2$, where $s\ne1$ when $\chi$ is
principal.
\item Under (H), every zero of $L(s,\chi)$
with $0<\operatorname{Re}s<1$ has $\operatorname{Re}s=1/2$, and every
zero $s\ne1$ of $\zeta$ that is not in $\{-2,-4,-6,\dots\}$ has
$\operatorname{Re}s=1/2$.
\end{enumerate}
\end{proposition}

\noindent\emph{What they say.} Zero-freeness in
$\operatorname{Re}s>1/2$ for primitive characters implies the same
for all characters, and the functional equation then places the
zeros of the critical strip on the line $\operatorname{Re}s=1/2$.

\noindent\emph{Why they are needed.} The argument produces (H). The
statement of Theorem~\ref{thm:main} concerns all Dirichlet characters
and the critical line.

\noindent\emph{Where they are proved.} Chapter~\ref{ch:landau}.
Theorem~\ref{bb:fe} is cited, and
Proposition~\ref{bb:imprimitive} is proved completely. The step to
the Riemann hypothesis needs neither all Dirichlet characters nor
the property $|\epsilon(\chi)|=1$ of the root numbers. The statement
about $\zeta$ uses only (H) for the character modulo~$1$ and
Riemann's functional equation, and where the
functional equation of a primitive character is used the root number
multiplies zero (Remark~\ref{rem:zeta-suffices}).

\section{Proof of the main theorem}
\label{sec:proof-main}

All the boxes are now stated. They meet in one statement, which
turns the two existence results into zero-freeness for primitive
characters.

\begin{theorem}[Alp\"oge; restated from
Theorem~\ref{thm:zero-free-primitive}]
\label{bb:zero-free-primitive}
Let $\lambda>1$. Assume that the top existence statement
(Definition~\ref{def:admissible}) holds and that
$\mathscr R_\psi$ is supplied on windows of exponent $\lambda$ for every
primitive
Dirichlet character $\psi$. Then $L(s,\psi)\ne0$ for every primitive
$\psi$ and every $s$ with $\operatorname{Re}s>1/2$, except $s=1$ when
$\psi$ is trivial.
\end{theorem}

\noindent Theorem~\ref{bb:zero-free-primitive} is proved completely
in Chapter~\ref{ch:families}, from the Chebyshev inequality on
windows and the windowed Landau criterion. The argument is short,
and it is the place where the boxes of
Sections~\ref{sec:bb-chebyshev} and~\ref{sec:bb-landau} are joined
to those of Sections~\ref{sec:bb-hosts} and~\ref{sec:bb-families},
so we go through it here.

\begin{proof}[Proof of Theorem~\ref{bb:zero-free-primitive} from the
black boxes]
Fix a primitive $\psi$, and let
$\mathcal C=\mathcal C_{\mathscr R_\psi}(\lambda)$ be the class of
the requirement of $\psi$ at $\lambda$.

We check that the class meets the hypotheses of the top existence
statement. Every member of a family in $\mathcal C$ beyond the
threshold meets $\mathscr R_\psi$. The Sha budget $(t,b^{1/2})$ is
part of $\mathscr R_\psi$, so every family in $\mathcal C$ satisfies
the Sha hypothesis. The explicit-formula condition
$\mathrm{EF}(\lev_\psi,\Top,x,\eps,4)$ is also part of
$\mathscr R_\psi$, so every family in $\mathcal C$ satisfies the
uniform explicit-formula condition for $\lev_\psi$, with $C=4$. By
assumption $\mathscr R_\psi$ is supplied on windows of exponent
$\lambda$.

The top existence statement therefore applies
to $\lambda$ and $\mathscr R_\psi$, and $\mathcal C$ supplies tied
families: there are $e_D,e_L$ and $\eps_1>0$ such that for every
$\eps\in(0,\eps_1)$ some family $\Fam$ at $\eps$ in $\mathcal C$ is
tied at every large $x\in\Win$, with these $e_D,e_L$ and some
constants. The numbers $e_D$ and $e_L$ do not depend on $\eps$.

Now let $s$ satisfy $\operatorname{Re}s>1/2$, with $s\ne1$ when
$\psi$ is trivial, put
$\mu=\min(\operatorname{Re}s-1/2,1/2)>0$, and choose
$\eps\in(0,\eps_1)$ with $(e_D+e_L+7)\eps\le\mu/2$. By the previous
paragraph there is a family $\Fam$ at this $\eps$ in $\mathcal C$
and there are real numbers $c_W,c_S,e_K,C_L$ and $x_0$ such that
$\Fam$ is tied with constants $e_D,e_L,c_W,c_S,e_K,C_L$ at every
$x\in\Win$ with $x\ge x_0$. The family $\Fam$ satisfies the uniform
explicit-formula condition for $\lev_\psi$, since it lies in
$\mathcal C$. By Theorem~\ref{bb:chebyshev-windows}, applied to
$\lev_\psi$ and $\Fam$, there is $x_1$ such that
\[
G_{\lev_\psi}(x)\le x^{1/2+(e_D+e_L+7)\eps}\le x^{\sigma_0},
\qquad \sigma_0=\tfrac12+\tfrac\mu2,
\]
for every $x\in\Win$ with $x\ge x_1$.

It remains to apply the Landau criterion. Since $\Fam$ lies in
$\mathcal C$, the
set $\Win$ has windows of exponent $\lambda$. Hence the set
$\{x:G_{\lev_\psi}(x)\le x^{\sigma_0}\}$, which contains every large
element of $\Win$, has windows of exponent $\lambda$, and
$\sigma_0\in(0,1)$. By Theorem~\ref{bb:level-criterion}, $\lev_\psi$
satisfies the window criterion at $\lambda$, so $L(s',\chi)\ne0$ for
every $\chi\in\Kept(\lev_\psi)$ and every $s'$ with
$\operatorname{Re}s'>\sigma_0$, where $s'\ne1$ when $\chi$ is the
trivial character. The character $\psi$ is kept by $\lev_\psi$, and
$\operatorname{Re}s>\sigma_0$ because
$\mu\le\operatorname{Re}s-1/2$. Hence $L(s,\psi)\ne0$.
\end{proof}

\noindent The construction depends on the order of the quantifiers
in this proof. The exponents $e_D,e_L$ and the window exponent
$\lambda$ are fixed first. The zero $s$ is given next, and it
determines $\eps$. The family, its window set and all remaining
constants depend on $\eps$. The bound $x^{1/2+(e_D+e_L+7)\eps}$ is used only
on the window set of that one family.

\begin{proof}[Proof of Theorem~\ref{thm:main}]
We chain the black boxes.
\begin{enumerate}[label=(\arabic*)]
\item\label{item:main-s1}
By Theorem~\ref{thm:galois-reps}, weak local Langlands maps exist
and the truncated local--global compatibility of
Definition~\ref{def:truncated-lgc} holds for them. This is
Theorem~\ref{bb:galois-reps}.
\item\label{item:main-s2}
By Theorem~\ref{thm:sha-windows}, whose automorphy input is
step~\ref{item:main-s1} (Remark~\ref{thm:lgc-to-sha}), every
primitive Dirichlet character $\psi$ has Sha bounds on windows with
constant $401$. This is Theorem~\ref{bb:sha-windows}.
\item\label{item:main-s3}
By Theorem~\ref{thm:host-supply}, applied to each primitive $\psi$
with step~\ref{item:main-s2} as its hypothesis, $\mathscr R_\psi$ is
supplied on windows of exponent $3/2$ for every primitive $\psi$.
This is Theorem~\ref{bb:host-supply}; the good pairs of
Theorem~\ref{bb:good-pairs} are used in its proof.
\item\label{item:main-s4}
By Theorem~\ref{thm:packages}\ref{item:packages-data}, whose proof
uses the covolume bound (Theorem~\ref{thm:covolume}) and the maps of
step~\ref{item:main-s1}, some class of families carries
admissible-family data. By Theorem~\ref{thm:tied-families}, applied
to that class, the top existence statement holds; Green's
parametrization (Theorem~\ref{thm:green}) enters its proof. These
are Theorems~\ref{bb:covolume},
\ref{bb:green}, \ref{bb:packages} and~\ref{bb:tied-families}.
\item\label{item:main-s5}
By Theorem~\ref{thm:zero-free-primitive} with $\lambda=3/2$, whose
two hypotheses are steps~\ref{item:main-s4}
and~\ref{item:main-s3}, $L(s,\psi)\ne0$ for every primitive $\psi$
and every $s$ with $\operatorname{Re}s>1/2$, except $s=1$ when
$\psi$ is trivial. This is Theorem~\ref{bb:zero-free-primitive},
whose proof above uses Theorems~\ref{bb:chebyshev-windows}
and~\ref{bb:level-criterion}. The conclusion is the hypothesis (H)
of Section~\ref{sec:bb-line}.
\item\label{item:main-s6}
By Proposition~\ref{prop:imprimitive}, the nonvanishing extends to
all Dirichlet characters, and the functional equation
(Theorem~\ref{thm:fe}) puts the zeros in $0<\operatorname{Re}s<1$ on
the line $\operatorname{Re}s=1/2$, for $\zeta$ and for every
$L(s,\chi)$. In detail, let $\chi$ be a Dirichlet character.
Part~\ref{item:imprim-all} of Proposition~\ref{bb:imprimitive} gives
$L(s,\chi)\ne0$ for $\operatorname{Re}s>1/2$, with $s=1$ excluded
when $\chi$ is principal. Part~\ref{item:imprim-line} gives that
every zero of $L(s,\chi)$ with $0<\operatorname{Re}s<1$ has real
part $1/2$, and that every zero of $\zeta$ outside
$\{-2,-4,-6,\dots\}$ has real part $1/2$. The point $s=1$ is a pole
of $\zeta$ and not a zero.
\end{enumerate}
This proves Theorem~\ref{thm:main} relative to the black boxes.
Section~\ref{sec:bb-status} records their status.
\end{proof}

\begin{remark}[where the content is]\label{rem:content}
The deduction of Theorem~\ref{thm:main} from its ingredients
carries no substantial argument of its own: the proof above is a
chain of implications. The substantive inputs
are the windowed Landau criterion (Theorem~\ref{thm:windowed-landau}),
the frame theorem (Theorem~\ref{thm:frame}), the covolume bound
(Theorem~\ref{thm:covolume}), the Sha bound on windows
(Theorem~\ref{thm:sha-windows}, with Theorem~\ref{thm:galois-reps}),
the supply of hosts (Theorems~\ref{thm:good-pairs}
and~\ref{thm:host-supply}), and the existence of packages and tied
families (Theorems~\ref{thm:packages} and~\ref{thm:tied-families}).
The analytic core that turns these into $G_\lev(x)\le
x^{1/2+O(\eps)}$ on windows is the gap bound,
Theorem~\ref{thm:gap}, with the core bound,
Theorem~\ref{thm:core-bound}. Once the gap bound is in place, the
Chebyshev inequality on windows (Theorem~\ref{thm:chebyshev-windows})
is linear arithmetic.
\end{remark}

\begin{remark}[the constants of the chain]\label{rem:bb-constants}
Five numbers pass from box to box, and each has one origin.
\begin{enumerate}[label=(\arabic*)]
\item The bound $\eps<1/52$ is fixed once, in the definition of a
scale. Only small values of $\eps$ are used in the end.
\item The constant $401$ of the Sha bound makes the budget cover
$\lfloor x^{401}\rfloor$ Tate twists; the number of twists needed is
the rank parameter $t$, which satisfies $t\le x^{400}$
(Remark~\ref{rem:windows-origin}).
\item The window exponent is $2$ in the Sha bound, and $3/2$ in the
supply of hosts (step~\ref{item:main-s3}) and in the application of
the Landau criterion (step~\ref{item:main-s5}). How the $2$ becomes
$3/2$ is Question~\ref{q:window-exponent}. The exponent $3/2$ does
not come from Green's theorem, which enters at
step~\ref{item:main-s4}, between the two. It is the window exponent
of the supply of hosts, and Green's theorem is not an analytic input
(Remark~\ref{rem:green}). The Landau criterion
accepts any exponent greater than $1$.
\item The constant $4$ in $\mathrm{EF}(\lev_\psi,\Top,x,\eps,4)$ is
the constant of the pointwise clause of a good pair
(Definition~\ref{def:good-pair}).
\item The exponent $1/2+(e_D+e_L+7)\eps$ has $e_D=2$ and $e_L=0$ for
the families of Theorem~\ref{bb:tied-families}, so it is
$1/2+9\eps$. Of the $7$, the gap bound contributes $3$, the core
bound $3$, and the absorption of the rest term $4o\,x^{1/2+\eps}$
one more.
\end{enumerate}
\end{remark}

\section{Chebyshev's comparison and the function-field model}
\label{sec:bb-analogy}

Chebyshev compared the archimedean size of a factorial, given by
Stirling's formula, with its prime factorisation, given by
Legendre's formula. The comparison bounds a prime sum because
the two calculations concern the same integer. Here the objects
are periods of successive theta lifts of algebraic automorphic
forms. Their archimedean sizes contain a Stirling term and a CM
period; their finite sizes record divisibility at primes inert
in an imaginary quadratic field. A map between lattices of forms
and its slope estimates take the place of the factorial identity.
Table~\ref{tab:analogy} records the parallel and the point at which
the present argument gains its sharper error. Each row of its right
column lies inside one of the black boxes. The first five rows lie
in the Chebyshev inequality on windows, through
Proposition~\ref{bb:balance} and Theorems~\ref{bb:gap}
and~\ref{bb:frame}. The row on multiplicity lies in
Proposition~\ref{bb:sqrtx}, and the last two rows lie in
Theorems~\ref{bb:chebyshev-windows} and~\ref{bb:windowed-landau}.

\begin{table}[htbp]
\centering
\input{alpoge/figs/v2/analogy.tex}
\caption{Chebyshev's factorial comparison and the corresponding
parts of the present argument. The slope comparison and the
inert-prime gain are developed in
Chapters~\ref{ch:theta}--\ref{ch:covolume}.}
\label{tab:analogy}
\end{table}

Chapter~\ref{ch:function-fields} proves the Riemann hypothesis for
a function field over a finite field by the method of Stepanov and
Bombieri, and Table~\ref{tab:bb-model} places each part of that
proof beside the black box that plays its role. The table is a
guide to roles. It asserts no implication between its two columns.

\begin{table}[htbp]
\centering\small
\renewcommand{\arraystretch}{1.25}
\begin{tabular}{@{}>{\raggedright\arraybackslash}p{4.3cm}
  >{\raggedright\arraybackslash}p{4.5cm}
  >{\raggedright\arraybackslash}p{2.9cm}@{}}
\toprule
In Chapter~\ref{ch:function-fields} & In the assembled argument &
Black box\\
\midrule
The centred smoothed sum $G(T)$ of the point counts
(Definition~\ref{def:smoothed-sum}) &
the weighted Chebyshev function $G_\lev$ &
Section~\ref{sec:bb-vocabulary}\\
A bound on a set containing arbitrarily long intervals forces
$|\alpha_i|=\sqrt q$ (Proposition~\ref{prop:landau},
Theorem~\ref{thm:route2prime}) &
a bound on windows excludes zeros with
$\operatorname{Re}s>\sigma_0$ &
Theorem~\ref{bb:windowed-landau}\\
The Stepanov--Bombieri inequality, from an auxiliary section
vanishing at the Frobenius places
(Theorem~\ref{thm:aux-section}) &
the gap bound, from a slope inequality for period maps; no
auxiliary section is constructed &
Theorems~\ref{bb:gap}, \ref{bb:frame}\\
The error $(g+1)q/M+2gM$, balanced at $M=\sqrt q$
(Proposition~\ref{prop:balancing}, Theorem~\ref{thm:bombieri}) &
the errors $x/m$ and $m$, balanced at $m\approx\sqrt x$ &
Proposition~\ref{bb:sqrtx}\\
Averaging over a Galois group to pass between upper and lower bounds
(Theorem~\ref{thm:averaging}) &
Fourier analysis on $\Gal(E/\Q)$, with positive weights for the kept
characters &
Proposition~\ref{bb:balance}, Theorem~\ref{bb:core-bound}\\
The functional equation places the reciprocal roots on
$|\alpha|=\sqrt q$ (Theorem~\ref{thm:rh}) &
the functional equation places the zeros on
$\operatorname{Re}s=1/2$ &
Theorem~\ref{bb:fe}, Proposition~\ref{bb:imprimitive}\\
The $M$-th power map (Frobenius), the source of multiplicity &
the chain of $m$ theta lifts; $\Z$ has no such endomorphism (item
(B1) of Section~\ref{sec:dictionary}) &
Chapter~\ref{ch:theta}, Proposition~\ref{bb:sqrtx}\\
No counterpart &
hosts on windows, the Sha bound, Galois representations, the
covolume bound, packages and tied families &
Sections~\ref{sec:bb-galois}--\ref{sec:bb-families}\\
\bottomrule
\end{tabular}
\caption{The function-field model of
Chapter~\ref{ch:function-fields} beside the black boxes of this
chapter.}
\label{tab:bb-model}
\end{table}

Three points of the comparison need care.
Section~\ref{sec:dictionary} of Chapter~\ref{ch:function-fields}
lists further places where the analogy breaks.

The first point is the direction of the inequality. The Chebyshev
inequality on windows is an upper bound for $G_\lev$. In the model,
an upper bound for the centred sum $G(T)$ corresponds to a
\emph{lower} bound for the point counts
(Proposition~\ref{prop:which-bound}), so the signs are opposite to
those of the Stepanov--Bombieri upper bound for point counts. The
inequality that the assembled argument proves is therefore the
counterpart of the hypothesis of Theorem~\ref{thm:route2prime}. It
does not correspond to the inequality that the Stepanov--Bombieri
method supplies. In Chapter~\ref{ch:function-fields} the lower bound
is obtained from the upper bound by Galois averaging
(Theorem~\ref{thm:averaging}). In the assembled argument the upper
bound for $G_\lev$ is proved directly, and the Landau criterion does
the rest.

The second point is the auxiliary function. The Stepanov--Bombieri
proof
constructs a section with prescribed vanishing and counts its zeros
against its poles. The assembled argument constructs no space of
small sections for its slope inequality. The inequality is applied
to the period map between two given lattices
(Remark~\ref{rem:periodmap}).

The third point is the sign-only argument. The second argument of
Chapter~\ref{ch:function-fields}, through the identity
$qN=M^k+R$, is conditional on a hypothesis that is not proved there
(Theorem~\ref{thm:sign}, hypothesis~\ref{H2}). It gets its bound
from a sign that prevents cancellation. No identity of this form
occurs among the black boxes of this chapter. In its place the
assembly has the positivity of the weights of a level. Because the
weights are positive, the contributions of real zeros cannot
cancel, and a one-sided bound suffices
(Remark~\ref{rem:positivity}).

\section{The status of the inputs}
\label{sec:bb-status}

The proof of Theorem~\ref{thm:main} in Section~\ref{sec:proof-main}
is complete relative to its black boxes. This section records what
this book provides for each box, and the caveats attached to them in
the chapters where they are treated.
\begin{figure}[tbp]
\centering
\input{book/figs/ch03-status-map.tex}
\end{figure}

\begin{remark}[attributions and kinds of proof]
\label{rem:bb-attribution}
The conventions are those of Section~\ref{sec:conventions}. A
statement marked (Alp\"oge) is a statement of Alp\"oge's argument.
Alp\"oge claims it, and it is stated to be formalized in his Lean
archive \cite{AlpogeArchive2026}; the archive was not built or
checked for this book, and nothing here establishes the statement
(Remark~\ref{rem:status}). This book states it and explains it. A
statement carrying another name is cited from the literature, and an
unattributed statement is proved in the text. Each proof block is of
one of four kinds.
\begin{enumerate}[label=(\arabic*)]
\item A \emph{proof} is complete.
\item A \emph{proof (outline)} gives Alp\"oge's route with its steps
named.
\item A block \emph{on the proof} gives the inputs of the proof and
the shape of its route, without an argument.
\item An \emph{idea of proof} is an explanatory account of a
standard method. It is not Alp\"oge's proof, and nothing later
depends on it.
\end{enumerate}
Not every box is closed by a complete proof in the later chapters.
Figure~\ref{fig:bb-status-map} places every box on the dependency structure of the proof of Theorem~\ref{thm:main}. Only the boxes marked ``proof'' in Table~\ref{tab:bb-status} are
proved completely in this book. For the others the book gives what
the table says, and for a proof it can only refer to the archive,
which it has not checked, or to the cited source.
\end{remark}

\begin{table}[htbp]
\centering\small
\renewcommand{\arraystretch}{1.15}
\begin{tabular}{@{}lll
  >{\raggedright\arraybackslash}p{4.2cm}@{}}
\toprule
Black box & Source & Chapter & What the book gives\\
\midrule
Theorem~\ref{bb:galois-reps} & Alp\"oge & \ref{ch:galois} &
proof (outline) of~\ref{item:gr-lgc}; \ref{item:gr-llc} not outlined\\
Theorem~\ref{bb:sha-windows} & Alp\"oge & \ref{ch:sha} &
on the proof\\
Theorem~\ref{bb:good-pairs} & Alp\"oge & \ref{ch:hosts} &
proof (outline)\\
Theorem~\ref{bb:host-supply} & Alp\"oge & \ref{ch:families} &
proof (outline)\\
Theorem~\ref{bb:covolume} & Alp\"oge & \ref{ch:covolume} &
on the proof; not reproduced\\
Theorem~\ref{bb:green} & Green & \ref{ch:families} & cited\\
Theorem~\ref{bb:packages} & Alp\"oge & \ref{ch:families} &
idea of proof, on the proof\\
Theorem~\ref{bb:tied-families} & Alp\"oge & \ref{ch:families} &
proof (outline)\\
Theorem~\ref{bb:chebyshev-windows} & Alp\"oge & \ref{ch:chebyshev} &
proof (outline)\\
Proposition~\ref{bb:balance} & --- & \ref{ch:balance} & proof\\
Theorem~\ref{bb:gap} & Alp\"oge & \ref{ch:chebyshev} &
proof (outline)\\
Theorem~\ref{bb:frame} & Alp\"oge & \ref{ch:theta} & idea of proof\\
Theorem~\ref{bb:core-bound} & Alp\"oge & \ref{ch:balance} &
proof (outline)\\
Theorem~\ref{bb:averaged-hyps} & Alp\"oge & \ref{ch:balance} &
idea of proof\\
Proposition~\ref{bb:sqrtx} & --- & \ref{ch:chebyshev} & proof\\
Theorem~\ref{bb:windowed-landau} & Alp\"oge & \ref{ch:landau} &
proof (outline)\\
Theorem~\ref{bb:level-criterion} & Alp\"oge & \ref{ch:landau} &
proof\\
Theorem~\ref{bb:fe} & Riemann, Hecke & \ref{ch:landau} & cited\\
Proposition~\ref{bb:imprimitive} & --- & \ref{ch:landau} & proof\\
Theorem~\ref{bb:zero-free-primitive} & Alp\"oge &
\ref{ch:families} & proof\\
\bottomrule
\end{tabular}
\caption{The black boxes of this chapter, the chapter in which each
is treated, and the kind of proof given there
(Remark~\ref{rem:bb-attribution}). The proofs marked ``proof'' are
complete relative to the results they quote.}
\label{tab:bb-status}
\end{table}

\begin{remark}[schematic statements]\label{rem:bb-schematic}
The definitions of packages, size data, the upper and lower links,
families, the covolume condition, host supply, the links and
admissible-family data
(Definitions~\ref{def:package}, \ref{def:size-datum},
\ref{def:upper-link}, \ref{def:lower-link}, \ref{def:family},
\ref{def:covolume-condition}, \ref{def:host-supply},
\ref{def:links} and~\ref{def:admissible}) are schematic: each lists
the parts of the data that the argument uses and the conditions on
them that the text uses, and omits the remaining compatibility
conditions. So not every black box can be checked against its
definitions line by line. Theorems~\ref{bb:host-supply},
\ref{bb:covolume},
\ref{bb:packages}, \ref{bb:tied-families},
\ref{bb:chebyshev-windows} and~\ref{bb:zero-free-primitive} are
statements about data given by schematic definitions, and are
schematic to the same extent; where the later chapters give an
informal statement in standard terms, it is repeated here under
``what it says''. The other black boxes of this chapter are stated
in full, apart from two clauses of Definition~\ref{def:weak-llc}
that are summarised there.
\end{remark}

\begin{remark}[no hypothesis is stated]\label{rem:bb-unconditional}
Theorem~\ref{thm:main} is stated without hypothesis. In particular,
none of the following appears as an assumption of a black box: a
bound of Riemann type for the auxiliary characters, local--global
compatibility for the Galois representations, or the local Langlands
correspondence. This is a statement about the form of the argument. The
proofs of the boxes not marked ``proof'' in
Table~\ref{tab:bb-status} are not checked in this book, so nothing
here establishes that the boxes hold (Remark~\ref{rem:status}).
\begin{enumerate}[label=(\arabic*)]
\item The analytic input on the twists is the explicit-formula
condition at the single scale $x$, which is polylogarithmic in the
conductor. It follows from the averaging over auxiliary primes, and
neither it nor the zero-free boxes of the real twists has the
strength of the Riemann hypothesis near $\operatorname{Re}s=1/2$
(Remark~\ref{rem:host-input-strength}). The classical inputs to the
averaging are listed in Remark~\ref{rem:classical-inputs}; none
assumes a hypothesis of Riemann type.
\item The Galois-theoretic inputs are the two parts of
Theorem~\ref{bb:galois-reps}, which is a statement of the argument
and not an assumption (Remark~\ref{rem:galois-inputs}). Truncated
compatibility asks for less than the published theorems, and the
local Langlands correspondence is replaced by the existence of weak
local Langlands maps (Remark~\ref{rem:llc}).
\item Every result of Chapter~\ref{ch:sha} is stated without
hypothesis; its only input beyond classical Galois cohomology and
the theorems stated there is Theorem~\ref{bb:galois-reps},
part~\ref{item:gr-lgc} (Remark~\ref{thm:lgc-to-sha}).
\end{enumerate}
\end{remark}

\begin{remark}[the scan of the archive for assumptions]\label{rem:ai-scan}
The $4081$ naive hits of a scan of the archive for assumptions are not
assumptions. The scans find no axiom declaration, no proof placeholder
and no unsafe or kernel-bypassing construct. The hits are macros,
imports and meta-programming commands. The scans also count $37$
declarations with conclusion \emph{False}, $124$ uses of
kernel-checked decision and $11$ declarations of a trivially true
type, of which only a sample was read. The scans are static: the
archive was not built (Section~\ref{sec:ai-formalization}).
\end{remark}

\begin{remark}[imported results and the forms in which they are used]
\label{rem:bb-imported}
The black boxes cite the literature at the following places.
\begin{enumerate}[label=(\arabic*)]
\item The functional equation (Theorem~\ref{bb:fe}) and Green's
parametrisation (Theorem~\ref{bb:green}) are used as stated.
Landau's theorem (Theorem~\ref{thm:landau}) is the model for the
windowed criterion and is not itself used.
\item The covolume bound rests on cited theorems of arithmetic
geometry. Each is used in a specific formulation, not necessarily
in the generality of the theorem it is named after, ranging from
near-general to special or weakened (Remark~\ref{rem:ag-strength}).
Each formulation is implied by the
cited theorem, with two qualifications recorded there: for a nef
line bundle the asymptotic Hilbert--Samuel formula is due to
Moriwaki, and for the bigness theorem the implication needs an
axiomatic intersection pairing to be the arithmetic intersection
pairing; this is not checked, and the implication is not claimed.
\item Bost's slope inequality (Theorem~\ref{thm:bost}) is one step
of the proof of the frame theorem; the slope clause of a slope
record is not Bost's inequality (Remark~\ref{rem:not-bost}).
\item Poitou--Tate duality, the finiteness of local Galois
cohomology and local Tate duality are used in the Sha bound. The
existence of packages and of admissible-family data draws on the
Siegel--Weil formula, Howe duality, results of M\'inguez on theta
integrals, the Casselman--Wallach theory and Minkowski's bound on
finite subgroups of $\mathrm{GL}_n(\Z)$
(Definition~\ref{def:admissible} and the discussion of
Theorem~\ref{thm:packages}).
\item Some named results are not used. Clozel's purity lemma, the
main theorem of complex multiplication and the theorem of Deligne
and Ribet do not enter the Sha bound
(Remark~\ref{rem:unused-inputs}), and no height of a CM point and no
Colmez-type sum of logarithmic derivatives enters the main term
(Remark~\ref{rem:no-colmez}).
\end{enumerate}
\end{remark}

\begin{remark}[transcribed theorems in the archive]\label{rem:ai-transcribed}
The statements of the archive that transcribe published theorems are
neither all genuine nor all hollow. Of $30$ sampled, $27$ are genuine
statements about Mathlib objects. Of the $195$ theorems that produce
them, $182$ cite theorems proved inside the archive.
\end{remark}

\begin{remark}[imports in the archive]\label{rem:ai-imports}
Long import chains do not show that the named theorems of the archive
are logically used. All $175{,}187$ modules are import-reachable from
the main file, so import reachability says nothing about term
dependency, which is unknown.
\end{remark}

\begin{remark}[points left open in the later chapters]
\label{rem:bb-open}
The chapters that treat the black boxes leave the following points
open, and the reader of the assembly should know of them.
\begin{enumerate}[label=(\arabic*)]
\item Whether the windowed Landau criterion is new as stated is a
question (Question~\ref{q:landau-new}).
\item In the supply of good pairs, the constant $4$ of the pointwise
clause must be uniform in the exponent $\Delta$. The route described
in Chapter~\ref{ch:hosts} gives a constant depending on $\Delta$,
and the sharper bound on $L'(1,\chi)/L(1,\chi)$ that uniformity
requires is not reconstructed there.
\item The constant $401$ leaves a margin of one in the exponent
over the $x^{400}$ twists that are needed; no step that uses the
margin is identified (Remark~\ref{rem:windows-origin}).
\item The bounds for the eigenspaces of the wrong parity inside the
Sha bound are of the type of Vandiver's and Greenberg's conjectures,
and Chapter~\ref{ch:sha} records that no published result of this
strength is known (Remark~\ref{rem:vandiver}). A vanishing theorem of Kato
type used there (Theorem~\ref{thm:bloch-kato}) is stated without an
outline of proof, and whether the twist it needs lies in the range
that Kato treats is not settled there.
\item The proof of the covolume bound is not reproduced, the proof
of part~\ref{item:gr-llc} of Theorem~\ref{bb:galois-reps} is not
outlined, and the compact-group statement used in the construction
of the Galois representations
(Proposition~\ref{prop:compact-group}) is stated without an outline
of proof.
\item The step from the window exponent $2$ of the Sha bound to the
exponent $3/2$ of the supply of hosts is
Question~\ref{q:window-exponent}; see also
Remark~\ref{rem:window-exponent}.
\end{enumerate}
\end{remark}
\clearpage
\clearpage
\chapter{Smoothed Chebyshev functions and a windowed Landau criterion}
\label{ch:landau}\label{sec:landau}

This chapter supplies the analytic end of the argument. Its main
result is the windowed Landau criterion,
Theorem~\ref{thm:windowed-landau}: a one-sided bound
$G_F(y)\le Cy^{\sigma_0}$ for a positive combination $G_F$ of
smoothed twisted Chebyshev functions, assumed only on windows
$[Y,Y^{1+\delta}]$ with $Y$ arbitrarily large, forces the
$L$-functions in the combination to have no zero in
$\operatorname{Re}s>\sigma_0$. Theorem~\ref{thm:level-criterion}
specialises the criterion to the combinations $G_\lev$ attached to
levels, which are the functions that the later chapters bound, and
Proposition~\ref{prop:imprimitive} passes from primitive characters
and a half-plane to all Dirichlet characters and the critical line.
The proof of Theorem~\ref{thm:main} in
Chapter~\ref{ch:assembly} (Section~\ref{sec:proof-main}) uses both
as black boxes. Theorem~\ref{thm:level-criterion} enters through
Theorem~\ref{thm:zero-free-primitive} of Chapter~\ref{ch:families},
the step that turns growth bounds on windows into zero-freeness.
Proposition~\ref{prop:imprimitive}, with the functional equation
(Theorem~\ref{thm:fe}), is the last step.
Section~\ref{sec:conventions} fixes the conventions for statements
and proofs used in Chapter~\ref{ch:assembly} and in
Chapters~\ref{ch:landau}--\ref{ch:families}.

The argument proves zero-freeness through one-sided growth. Its
analytic end is a converse to the explicit formula: if a suitable
smoothed prime sum is bounded above by $x^{\sigma_0}$, then the
corresponding $L$-functions have no zero in $\operatorname{Re}s>
\sigma_0$. Two features of the converse are used throughout. It needs
only an \emph{upper} bound, provided
the prime sum is a combination of twisted sums with positive weights.
It needs the bound only on \emph{windows} $[Y,Y^\lambda]$ for
arbitrarily large $Y$, not for all large $x$.

\section[Conventions]%
{Conventions for Chapters~\ref{ch:assembly}--\ref{ch:families}}
\label{sec:conventions}

The conventions of this section hold in Chapter~\ref{ch:assembly}
and in Chapters~\ref{ch:landau}--\ref{ch:families}.
A statement marked (Alp\"oge) is a statement of
Alp\"oge's argument. Alp\"oge claims it, and it is stated to be
formalized in his Lean archive \cite{AlpogeArchive2026}; the
archive was not built or checked for this book, and nothing here
establishes the statement (Remark~\ref{rem:status}). We state it
and explain it. A statement
carrying another name is cited from the literature, and an
unattributed statement is proved in the text.

Each proof block carries one of four labels.
\begin{enumerate}[label=(\roman*)]
\item \emph{Proof}: a complete proof.
\item \emph{Proof (outline)}: Alp\"oge's route, with its steps named;
steps that are standard, or that are cited theorems, are said to be.
\item \emph{On the proof}: the inputs of the proof and the shape of
its route, without an argument.
\item \emph{Idea of proof}: our explanatory account of a standard
method for statements of this kind. It is not Alp\"oge's proof, and
nothing later depends on it.
\end{enumerate}

Some definitions in
Chapters~\ref{ch:theta}, \ref{ch:covolume} and~\ref{ch:families}
are marked \emph{schematic}. A schematic definition lists the parts of
the data that the argument uses, and the conditions on them that the
text uses; the remaining compatibility conditions among the data are
omitted. A result about data given by a schematic definition is stated
schematically, and an informal statement in standard terms accompanies
it.

\section{Smoothed Chebyshev functions}

\begin{definition}[smoothed Chebyshev functions]\label{def:chebyshev}
Let $\Lambda$ be the von Mangoldt function, $\gamma$ Euler's constant
and $\Ch(t)=\sum_{n\le t}\Lambda(n)$ Chebyshev's function. Let $x,y>0$
and let $\chi$ be a Dirichlet character modulo $N$.
\begin{enumerate}[label=(\arabic*)]
\item The \emph{smoothed Chebyshev function} of $\zeta$ is
\[
\Sz(x)=\sum_{1\le n\le x}\Lambda(n)\Bigl(\frac{x}{n}-1\Bigr),
\]
and its \emph{normalised form} is $\Gz(x)=\Sz(x)-x\log x+(1+\gamma)x$.
\item The \emph{real twisted function} of $\chi$ is
\[
G_\chi(x)=\sum_{1\le n\le x}\Lambda(n)\operatorname{Re}\chi(n)
\Bigl(\frac{x}{n}-1\Bigr)+x\operatorname{Re}
\frac{L'(1,\chi)}{L(1,\chi)}.
\]
\item The \emph{complex form} for $\zeta$ is
\[
G(y)=\sum_{n\le y}\Lambda(n)\max\Bigl(\frac yn-1,0\Bigr)
-y(\log y-1-\gamma).
\]
\item The \emph{complex twisted function} of $\chi$ is $g_\chi(y)=G(y)$
if $\chi$ is principal (we then write $\chi=1$), and, for $\chi\ne1$,
\[
g_\chi(y)=\sum_{1\le n\le y}\Lambda(n)\chi(n)\max\Bigl(\frac{y}{n}-1,0
\Bigr)+y\,\frac{L'(1,\chi)}{L(1,\chi)}.
\]
\end{enumerate}
\end{definition}

The letter $\psi$ is kept for characters throughout, which is why
Chebyshev's function is written $\Ch$. For a nonprincipal $\chi$ the
value $L(1,\chi)$ is nonzero, so the term $yL'(1,\chi)/L(1,\chi)$ is
defined. The extra terms
in $\Gz$, $G_\chi$ and $g_\chi$ remove the contribution of $s=1$, as
Remark~\ref{rem:explicit} explains.

\begin{proposition}\label{prop:realcomplex}
Let $y>0$. Then:
\begin{enumerate}[label=(\alph*)]
\item\label{item:rc-zeta} $G(y)=\Gz(y)$;
\item\label{item:rc-twist} if $\chi\ne1$ is a Dirichlet character,
then $\operatorname{Re}g_\chi(y)=G_\chi(y)$;
\item\label{item:rc-integral} if $y\ge1$, then
$\Sz(y)=y\int_1^y\Ch(t)t^{-2}\,dt$.
\end{enumerate}
\end{proposition}

\begin{proof}
\noindent\emph{Proof of \ref{item:rc-zeta}.}
For $1\le n\le y$ we have $y/n-1\ge0$, so $\max(y/n-1,0)=y/n-1$. This
gives $G(y)=\Sz(y)-y\log y+(1+\gamma)y=\Gz(y)$.

\noindent\emph{Proof of \ref{item:rc-twist}.} Since $\Lambda(n)$ and
$y/n-1$ are real, the real part of each summand of $g_\chi(y)$ is
$\Lambda(n)\operatorname{Re}\chi(n)(y/n-1)$, and
$\operatorname{Re}(yL'(1,\chi)/L(1,\chi))=y\operatorname{Re}
(L'(1,\chi)/L(1,\chi))$. Hence $\operatorname{Re}g_\chi(y)=G_\chi(y)$.

\noindent\emph{Proof of \ref{item:rc-integral}.} The function
$\Ch(t)=\sum_{n\le t}\Lambda(n)$ is a finite sum on $[1,y]$,
so
\[
y\int_1^y\Ch(t)t^{-2}\,dt=y\sum_{n\le y}\Lambda(n)\int_n^y
t^{-2}\,dt
=\sum_{n\le y}\Lambda(n)\Bigl(\frac{y}{n}-1\Bigr)=\Sz(y).\qedhere
\]
\end{proof}

\begin{remark}[what the functions measure]\label{rem:explicit}
Perron's formula with the kernel $x^s/(s(s-1))$ gives the explicit
formula for $\Ch(x)$ \cite{Davenport2000,MontgomeryVaughan2007} in
integrated form:
\[
\Sz(x)=\frac{1}{2\pi i}\int_{(2)}\Bigl(-\frac{\zeta'}{\zeta}(s)\Bigr)
\frac{x^s}{s(s-1)}\,ds\qquad(x>1).
\]
The double pole at $s=1$ contributes exactly $x\log x-(1+\gamma)x$,
the pole at $s=0$ contributes $\log2\pi$, and the trivial zeros
contribute $O(1/x)$. Hence
\[
\Gz(x)=\log2\pi-\sum_\rho\frac{x^\rho}{\rho(\rho-1)}+O(1/x),
\]
the sum running over the nontrivial zeros of $\zeta$ and converging
absolutely. For $\chi\ne1$ the term $xL'(1,\chi)/L(1,\chi)$ removes
the simple pole of the integrand at $s=1$, and
$g_\chi(x)=-\sum_\rho x^\rho/(\rho(\rho-1))+O(\log x)$, the sum now
over the nontrivial zeros of $L(s,\chi)$. Consequently the Riemann
hypothesis for $L(s,\chi)$ gives $g_\chi(x)=O(x^{1/2})$, and the
Riemann hypothesis gives $\Gz(x)=O(x^{1/2})$. The rest of this chapter
is about the converse.
\end{remark}

Figure~\ref{fig:gz} shows $\Gz(x)/\sqrt x$ for $x\le10^7$ against the
size that the Riemann hypothesis predicts.

\begin{figure}[htbp]
\centering
\input{alpoge/figs/v2/gz.tex}
\caption{The function $\Gz(x)=\sum_{n\le x}\Lambda(n)(x/n-1)-x\log x
+(1+\gamma)x$, centred at its constant $\log2\pi$ and divided by
$\sqrt x$, for $10\le x\le10^7$. The line is computed from $\Lambda$,
the triangles from the explicit formula with the first $500$ zeros of
$\zeta$. The band $\pm c_0$, with $c_0=\sum_\rho1/|\rho(\rho-1)|$, is
the envelope that the explicit formula gives when every zero lies on
the critical line. Numerics at finite $x$ do not bear on the
asymptotic statement.}
\label{fig:gz}
\end{figure}

\section{The windowed Landau criterion}

The criterion is stated for finite positive combinations of the
functions $g_\chi$.

\begin{definition}[character families]\label{def:charfamily}
A \emph{character family} $F$ is a finite list of pairs
$(\chi_i,a_i)$, where each $\chi_i$ is a Dirichlet character modulo
some $N_i$ that is either the character modulo~$1$ or nontrivial, and
each $a_i>0$ is a real \emph{weight}. Its \emph{weighted Chebyshev
function} is $G_F(y)=\sum_ia_i\operatorname{Re}g_{\chi_i}(y)$.
\end{definition}

The classical model for a one-sided converse is Landau's theorem on
Laplace--Stieltjes integrals with a monotone integrator.

\begin{theorem}[Landau \cite{Landau1905}]\label{thm:landau}
Let $\alpha$ be a nondecreasing function on $[0,\infty)$ and suppose
that $F(s)=\int_0^\infty e^{-su}\,d\alpha(u)$ has a finite abscissa of
convergence $\sigma_c$. Then $F$ has no analytic continuation to any
neighbourhood of $s=\sigma_c$.
\end{theorem}

\noindent See also \cite{Widder1941} and, for Dirichlet series,
\cite{MontgomeryVaughan2007}. Applied to
$\int_1^\infty(Cy^{\sigma_0}-G_F(y))y^{-s-1}\,dy$, it gives the
statement below with the bound assumed for \emph{all} large $y$. The
version used in the argument assumes it only on windows.
\begin{figure}[tbp]
\centering
\input{book/figs/ch04-gamma-kernel.tex}
\end{figure}

\begin{theorem}[Alp\"oge]\label{thm:windowed-landau}
Let $F=((\chi_i,a_i))_i$ be a character family, and let
$0<\sigma_0<1$, $\delta>0$ and $C\in\R$. Suppose that for every $Y_0$
there is $Y\ge Y_0$ with $G_F(y)\le Cy^{\sigma_0}$ for all
$y\in[Y,Y^{1+\delta}]$. Then $L(s,\chi_i)\ne0$ for every $i$ and every
$s$ with $\operatorname{Re}s>\sigma_0$, where $s\ne1$ when $\chi_i$ is
the character modulo~$1$.
\end{theorem}

Step~(2) of the proof below uses how much of the mass of a Gamma
kernel lies in one window. The next two statements give the correct
rate and show that it does not improve as the expansion point moves
to the right.

For $k\ge0$ and $\lambda>0$ let $\mu_{k,\lambda}$ be the probability
measure $\frac{\lambda^{k+1}}{k!}u^ke^{-\lambda u}\,du$ on $(0,\infty)$.
It is the law of $Z/\lambda$, where $Z$ has the Gamma distribution
with density $z^ke^{-z}/k!$; its mean is $(k+1)/\lambda$.

\begin{proposition}[Gamma tails]\label{prop:gamma-tails}
Let $k\ge0$, $\lambda>0$ and $0<t<1$, and put $h(t)=t-\log(1+t)$.
Then $h(t)>0$ and
\[
\mu_{k,\lambda}\Bigl(\Bigl\{u:\Bigl|\frac{\lambda u}{k+1}-1\Bigr|
\ge t\Bigr\}\Bigr)\le2e^{-(k+1)h(t)}.
\]
\end{proposition}

\begin{proof}
The function $h$ vanishes at $0$ and has derivative $t/(1+t)>0$, so
$h(t)>0$. The variable $Z=\lambda u$ has moment generating function
$\int_0^\infty e^{\theta z}z^ke^{-z}\,dz/k!=(1-\theta)^{-(k+1)}$ for
$\theta<1$. Put $n=k+1$. For $0<\theta<1$, Markov's inequality applied
to $e^{\theta Z}$ gives
$\mu_{k,\lambda}(Z\ge n(1+t))\le(1-\theta)^{-n}e^{-\theta n(1+t)}$;
with $\theta=t/(1+t)$ the right side is $e^{-nh(t)}$. Applied to
$e^{-\theta Z}$ with $\theta=t/(1-t)$, it gives
$\mu_{k,\lambda}(Z\le n(1-t))\le(1+\theta)^{-n}e^{\theta n(1-t)}
=e^{-n(-t-\log(1-t))}$. The power series
$-t-\log(1-t)=\sum_{j\ge2}t^j/j$ and $t-\log(1+t)=\sum_{j\ge2}
(-1)^jt^j/j$ show $-t-\log(1-t)\ge h(t)$. Adding the two tails gives
the claim.
\end{proof}

The rate $e^{-(k+1)h(t)}$ does not depend on $\lambda$, and nothing
better is true: the next proposition shows that the mass outside a
window is not $O(\lambda^{-k})$, which is what an earlier version of
step~(2) asserted.

\begin{proposition}\label{prop:gamma-no-sigma-gain}
For every $\lambda>0$, $U>0$ and $k\ge1$,
\[
\mu_{k,\lambda}\bigl((0,\infty)\setminus[U,2U]\bigr)\ge8^{-(k+1)}.
\]
In particular, for fixed $\lambda>8$ there is no constant $C'$ such
that, for arbitrarily large $U$, some $k\ge U$ has
$\mu_{k,\lambda}((0,\infty)\setminus[U,2U])\le C'\lambda^{-k}$.
\end{proposition}

\begin{proof}
Put $a=\lambda U$; in the variable $z=\lambda u$ the claim is that
$Z\notin[a,2a]$ with probability at least $8^{-(k+1)}$. If
$a\ge k/2$, then
\[
\mu(Z<a)\ge\frac1{k!}\int_0^{k/2}z^ke^{-z}\,dz
\ge\frac{e^{-k/2}(k/2)^{k+1}}{(k+1)!}\ge e^{-k/2}4^{-(k+1)}
\ge8^{-(k+1)},
\]
using $(k+1)!\le(k+1)^{k+1}\le(2k)^{k+1}$ and $e^{-1/2}\ge1/2$. If
$a<k/2$, then $2a<k$ and
\[
\mu(Z>2a)\ge\frac1{k!}\int_k^{k+1}z^ke^{-z}\,dz
\ge\frac{k^k}{k!}e^{-(k+1)}\ge e^{-(k+1)}\ge8^{-(k+1)} .
\]
For the second sentence, $8^{-(k+1)}\le C'\lambda^{-k}$ forces
$C'\ge(\lambda/8)^k/8$, which is unbounded as $k\ge U\to\infty$ when
$\lambda>8$.
\end{proof}

The proof compares two estimates for the Taylor coefficients of a
Laplace transform at a point far to the right: one coming from the
bound on windows, the other from a hypothetical zero. Its steps are
elementary complex analysis.

\begin{proof}[Proof (outline)]
Put $g(u)=G_F(e^u)$ and, for $\operatorname{Re}s>2$,
\[
\widehat g_F(s)=\int_0^\infty g(u)e^{-su}\,du
=\int_1^\infty G_F(y)y^{-s-1}\,dy .
\]
\begin{enumerate}[label=(\arabic*)]
\item \emph{Continuation and poles.} For a single nontrivial $\chi$, a
direct computation gives
$\int_1^\infty g_\chi(y)y^{-s-1}\,dy=-\frac{L'}{L}(s,\chi)
\frac{1}{s(s-1)}+\frac{L'}{L}(1,\chi)\frac1{s-1}$, in which the
pole at $s=1$ cancels; for $\chi=1$ the correction
$y(\log y-1-\gamma)$ in $G$ cancels the double pole of
$-\frac{\zeta'}{\zeta}(s)/(s(s-1))$. Since
$\operatorname{Re}g_\chi=(g_\chi+g_{\overline\chi})/2$ and no
$L(s,\chi_i)$ vanishes on $\operatorname{Re}s\ge1$, the transform
$\widehat g_F$ continues meromorphically to $\operatorname{Re}s>0$,
with simple poles exactly at the zeros of the $L(s,\chi_i)$ and
$L(s,\overline\chi_i)$. At a zero $\rho$ the residue is
$-c_\rho/(\rho(\rho-1))$, where $c_\rho>0$ is the sum of the weights
times the multiplicities. For real $\rho\in(0,1)$ this residue is
positive.
\item \emph{The Gamma kernel.} Up to the sign $(-1)^k$, the $k$-th
Taylor coefficient of $\widehat g_F$ at a point $s_0$ with
$\operatorname{Re}s_0=\sigma>2$ is
$c_k(s_0)=\frac1{k!}\int_0^\infty g(u)u^ke^{-s_0u}\,du$. The
kernel $u^ke^{-\sigma u}$ peaks at $u=k/\sigma$. Fix the window
$W=[U,(1+\delta)U]$ with $U\ge1$; in the variable $y=e^u$ it is the
window $[Y,Y^{1+\delta}]$ with $Y=e^U$
(Figure~\ref{fig:gamma-kernel}). Put $\lambda=\sigma-2$, assume
$\lambda\ge4/\delta$, and choose $k$ with
$k+1=\lceil\lambda(1+\delta/2)U\rceil$, so that $k+1\ge\lambda U$.
With $t=\delta/(4+3\delta)$, every $u\notin W$ has
$|\lambda u/(k+1)-1|\ge t$, so by Proposition~\ref{prop:gamma-tails}
\[
\mu_{k,\lambda}\bigl((0,\infty)\setminus W\bigr)\le
2e^{-(k+1)h(t)} .
\]
The rate depends on $\delta$ but not on $\sigma$; by
Proposition~\ref{prop:gamma-no-sigma-gain} no bound of the form
$C'\sigma^{-k}$ holds.
\item\label{item:landau-s3}
\emph{The window bound.} We may take $C\ge0$. By
Proposition~\ref{prop:realcomplex}\ref{item:rc-integral} and
$\Ch(t)\le2t$, there is $A$ with $|g(u)|\le Ae^{2u}$ for $u\ge0$.
On $W$ we have $g\le R(u):=Ce^{\sigma_0u}$, hence
$|g|+g\le2R$ there, and off $W$ we have $|g|+g\le2Ae^{2u}$. Since
$g$ is real,
\[
|c_k(\sigma+it_0)|+\operatorname{Re}c_k(\sigma)
\le\frac1{k!}\int_0^\infty(|g|+g)u^ke^{-\sigma u}\,du
\le\frac{2C}{(\sigma-\sigma_0)^{k+1}}
+\frac{4Ae^{-(k+1)h(t)}}{(\sigma-2)^{k+1}} .
\]
If $\sigma$ is so large in terms of $\delta$ that
$e^{-h(t)}(\sigma-\sigma_0)<\sigma-2$, the right side is at most
$C''(\sigma-\sigma_0)^{-(k+1)}$ with $C''=2C+4A$.
\item \emph{The contradiction.} A zero with $\operatorname{Re}\rho>
\sigma_0$ gives, through the pole of $\widehat g_F$ nearest to
$\sigma$ (real, or an isolated nearest complex pole after a choice
of $t_0$), a lower bound $c\,r^{-(k+1)}$ on the Taylor coefficients
for all large $k$, with $c>0$ and $r$ the distance from the expansion
point to the pole used. That pole has real part $>\sigma_0$, so
$r<\sigma-\sigma_0$, and the lower bound contradicts
step~\ref{item:landau-s3} once $k$ is large, that is, once $U$ is
large ($k+1\ge(\sigma-2)U$). One window per coefficient suffices,
because the
lower bound holds for all large $k$ and not only along a subsequence.
\qedhere
\end{enumerate}
\end{proof}

\begin{remark}[a correction]\label{rem:gamma-correction}
An earlier version of step~(2) asserted that all but a proportion
$C'/\sigma^k$ of the kernel's mass lies in the window. That is false
(Proposition~\ref{prop:gamma-no-sigma-gain}). Hilaf Hasson found the
error, with an AI-assisted reading of the earlier version, and gave
the counterexample of Proposition~\ref{prop:gamma-no-sigma-gain}.
Steps~(2) and~(3) above use the correct rate of
Proposition~\ref{prop:gamma-tails} instead. The rate
$e^{-(k+1)h(t)}$ suffices because $k+1$ is of size $\sigma U$, so it
beats the trivial growth $e^{2u}$ of $g$ once $\sigma$ is large in
terms of~$\delta$. Steps~(1) and~(4) are unchanged, and the outline
remains an outline: step~(4), including the choice of $t_0$ when the
nearest pole is not real, is not written out here.
\end{remark}

\begin{remark}[why the weights must be positive]\label{rem:positivity}
Theorem~\ref{thm:windowed-landau} assumes only an upper bound for
$G_F$. This suffices because, with positive weights, the residues of
$\widehat g_F$ at real zeros in $(0,1)$ all have one sign, so a
hypothetical real zero cannot be cancelled by another. The argument
must therefore produce an upper bound for a \emph{positive}
combination of $\Gz$ and the functions $\operatorname{Re}g_\psi$. The
weights that occur are $1$ and $(3o+1)^{-1}$
(Definition~\ref{def:level}); they arise from the shares of
Proposition~\ref{prop:shares}.
\end{remark}

Figure~\ref{fig:landau} pictures the mechanism.

\begin{figure}[htbp]
\centering
\input{alpoge/figs/v2/landau.tex}
\caption{The mechanism of Theorem~\ref{thm:windowed-landau}. The Laplace
transform $\widehat g_F$ of the proof has poles at the zeros of the
$L(s,\chi_i)$; the crosses are the first zeros of $\zeta$. A
hypothetical real zero $\beta$ with $\frac12<\sigma_0<\beta$ (hollow
point) would be a pole with residue $-c_\beta/(\beta(\beta-1))>0$,
where $c_\beta>0$ is the sum of the weights times the multiplicities
of $\beta$ (step~(1) of the proof), and positive weights
give all such residues one sign.
The Taylor expansion at a point $\sigma$ far to the right then
has every coefficient bounded below, which contradicts the bound on
windows. The disc is drawn as an ellipse because the two axes have
different scales.}
\label{fig:landau}
\end{figure}

\begin{remark}[windows are needed]\label{rem:landau-compare}
The classical, non-windowed form of
Theorem~\ref{thm:windowed-landau}, which follows from
Theorem~\ref{thm:landau} (see the paragraph after it), does not
suffice for the argument. The growth bounds that the argument proves
hold only on windows: the Sha
bound of Chapter~\ref{ch:sha} controls the congruence defects only
for $x$ in a set with windows $[Y,Y^2]$
(Theorem~\ref{thm:windows-pigeonhole}), and the Chebyshev inequality
of Chapter~\ref{ch:chebyshev} is proved only there. So the windowed
form is the one used. The closest result in the literature known to
us, \cite[Thm.~3.1]{BrouckeRevesz2026landau}, assumes a pointwise
bound on $|F(s)|$ near $\sigma=1$, not a one-sided bound on windows.
\end{remark}

Crossley applies Alp\"oge's pole argument to the same family at three dyadic scales \cite[\texttt{compactDifference\_zeroFree}]{Crossley2026Window}. If the bound on $[Y,Y^{1+\delta}]$ in Theorem~\ref{thm:windowed-landau} holds for $-G_F(y)+3G_F(y/2)-2G_F(y/4)$, the same zero-free conclusion follows. The multiplier $(1-2^{-s})(2^{1-s}-1)$ has no zeros in $0<\operatorname{Re}s<1$ and is positive on the real interval $(0,1)$. It therefore preserves the poles in the strip and the signs of the real residues used in the one-sided argument.

\begin{question}\label{q:landau-new}
Is the windowed criterion of Theorem~\ref{thm:windowed-landau} new as
stated?
\end{question}

\section{Levels}

The positive combinations that the argument controls are attached to
finite sets of primitive characters closed under inversion.

\begin{definition}[levels]\label{def:level}
\begin{enumerate}[label=(\arabic*)]
\item A \emph{level} $\lev$ consists of a finite set $\Kept(\lev)$ of
primitive Dirichlet characters that contains the character
modulo~$1$ and is closed under inversion, an optional
\emph{exceptional character} $\psi_{\mathrm{exc}}\in\Kept(\lev)$, and
a natural number $o(\lev)$. The elements of $\Kept(\lev)$ are its
\emph{kept characters}.
\item The \emph{weight} of $\psi\in\Kept(\lev)$ is $\wt(\psi)=1$ if
$\psi=\psi_{\mathrm{exc}}$ and $\wt(\psi)=(3o(\lev)+1)^{-1}$
otherwise.
\item The \emph{weighted Chebyshev function} of $\lev$ is
\[
G_{\lev}(y)=\Gz(y)+\sum_{\psi\in\Kept(\lev),\ \psi\ne1}
\wt(\psi)\operatorname{Re}g_\psi(y).
\]
\item For a primitive character $\psi$, the \emph{cyclic level}
$\lev_\psi$ has as kept set the primitive characters inducing the
powers of $\psi$, no exceptional character, and $o(\lev_\psi)$ equal
to the order of $\psi$.
\end{enumerate}
\end{definition}

All weights of a level are positive, and the trivial character enters
with weight~$1$ through $\Gz$. Thus $G_{\lev_\psi}=\Gz+(3o+1)^{-1}
\sum_{\chi}\operatorname{Re}g_\chi$, the sum over the primitive
characters $\chi\ne1$ that induce powers of $\psi$, where $o$ is the
order of $\psi$.

\begin{definition}[windows]\label{def:windows}
Let $S\subset\R$, $\lambda>1$, and let $\lev$ be a level.
\begin{enumerate}[label=(\arabic*)]
\item $S$ \emph{has windows of exponent} $\lambda$ if for every $y$
there is $Y\ge y$ with $[Y,Y^\lambda]\subset S$.
\item $\lev$ \emph{satisfies the window criterion at} $\lambda$ if for
every $\sigma_0\in(0,1)$ the following holds: if
$\{x:G_{\lev}(x)\le x^{\sigma_0}\}$ has windows of exponent
$\lambda$, then $L(s,\psi)\ne0$ for every $\psi\in\Kept(\lev)$ and
every $s$ with $\operatorname{Re}s>\sigma_0$, where $s\ne1$ when
$\psi$ is the character modulo~$1$.
\end{enumerate}
\end{definition}

\begin{theorem}[Alp\"oge]\label{thm:level-criterion}
Let $\lev$ be a level and $\lambda>1$. Then $\lev$ satisfies the
window criterion at $\lambda$.
\end{theorem}

\begin{proof}
Let $F$ be the character family whose pairs are $(1,1)$, for the
character modulo~$1$, and $(\psi,\wt(\psi))$ for
$\psi\in\Kept(\lev)$, $\psi\ne1$. A primitive character is either the
character modulo~$1$ or nontrivial, and every weight is positive, so
$F$ is a character family. Since $g_1=G=\Gz$ on $y>0$
(Proposition~\ref{prop:realcomplex}\ref{item:rc-zeta}), $G_F=G_{\lev}$. Let
$\sigma_0\in(0,1)$ and suppose that $\{x:G_{\lev}(x)\le x^{\sigma_0}\}$
has windows of exponent $\lambda$. Then for every $Y_0$ there is
$Y\ge Y_0$ with $G_F(y)\le y^{\sigma_0}$ on $[Y,Y^\lambda]$. This is
the hypothesis of Theorem~\ref{thm:windowed-landau} with
$\delta=\lambda-1$ and $C=1$, whose conclusion is the required
nonvanishing.
\end{proof}

The windows that occur are pictured in Figure~\ref{fig:windows}: the
Sha bound gives exponent $2$, the supply of hosts is stated with
$3/2$ (Question~\ref{q:window-exponent}), and
Theorem~\ref{thm:level-criterion} accepts any exponent $\lambda>1$.

\begin{figure}[htbp]
\centering
\input{alpoge/figs/v2/windows.tex}
\caption{Windows on a $\log\log x$ axis, on which $[Y,Y^\lambda]$ has
fixed width $\log\lambda$. The Sha bound gives windows $[Y,Y^2]$: in
every range $[X^{1/J},X]$ a pigeonhole argument over the at most $R$
exceptional Tate twists finds a window on which all of them are small
(Theorem~\ref{thm:windows-pigeonhole},
Remark~\ref{rem:windows-origin}). The supply of hosts is stated
with windows $[W,W^{3/2}]$; a window $[Y,Y^2]$ contains
$[Y,Y^{3/2}]$, and the step between the two statements is not
reconstructed (Question~\ref{q:window-exponent}).
Theorem~\ref{thm:windowed-landau} needs only windows
$[Y,Y^{1+\delta}]$ for arbitrarily large $Y$. Here $R=3$; positions
and widths are illustrative.}
\label{fig:windows}
\end{figure}

\section{From a half-plane to the critical line}

The passage from zero-freeness in $\operatorname{Re}s>1/2$ to the
critical line is classical. For a primitive Dirichlet character $\chi$
modulo $N$, let $\mathfrak a\in\{0,1\}$ be given by
$\chi(-1)=(-1)^{\mathfrak a}$, and put
$\xi(s,\chi)=(N/\pi)^{(s+\mathfrak a)/2}\Gamma((s+\mathfrak a)/2)
L(s,\chi)$.
We write (H) for the following hypothesis: $L(s,\chi')\ne0$ for every
primitive Dirichlet character $\chi'$ and every $s$ with
$\operatorname{Re}s>1/2$, where $s\ne1$ when $\chi'$ is the character
modulo~$1$.

\begin{theorem}[Riemann, Hecke \cite{Davenport2000,Riemann1859}]
\label{thm:fe}
Let $\chi$ be a primitive Dirichlet character modulo $N$. Then there is
$\epsilon(\chi)\in\C$ with $|\epsilon(\chi)|=1$ and
$\xi(s,\chi)=\epsilon(\chi)\xi(1-s,\overline\chi)$ for all $s\ne0,1$.
\end{theorem}

\noindent For $N=1$ this is Riemann's functional equation.

\begin{proposition}\label{prop:imprimitive}
Let $\chi$ be a Dirichlet character modulo $N$, induced by the
primitive character $\chi^*$ modulo $N^*$.
\begin{enumerate}[label=(\alph*)]
\item\label{item:imprim-euler} For $s\ne1$, and for every $s\in\C$ if
$\chi$ is nonprincipal,
$L(s,\chi)=L(s,\chi^*)\prod_{p\mid N}(1-\chi^*(p)p^{-s})$. In
particular, for such $s$ with $\operatorname{Re}s\ne0$, $L(s,\chi)=0$
if and only if $L(s,\chi^*)=0$.
\item\label{item:imprim-zeta} The $L$-function of the character
modulo~$1$ is $\zeta$ on $s\ne1$.
\item\label{item:imprim-all} Under (H), $L(s,\chi)\ne0$ for every $s$
with $\operatorname{Re}s>1/2$, where $s\ne1$ when $\chi$ is
principal.
\item\label{item:imprim-line} Under (H), every zero of $L(s,\chi)$
with $0<\operatorname{Re}s<1$ has $\operatorname{Re}s=1/2$, and every
zero $s\ne1$ of $\zeta$ that is not in $\{-2,-4,-6,\dots\}$ has
$\operatorname{Re}s=1/2$.
\end{enumerate}
\end{proposition}

\begin{proof}[Proof of \ref{item:imprim-euler}]
For $\operatorname{Re}s>1$ the Euler products give
$L(s,\chi)=\prod_p(1-\chi(p)p^{-s})^{-1}$ with $\chi(p)=\chi^*(p)$ for
$p\nmid N$ and $\chi(p)=0$ for $p\mid N$, which is the identity. Both
sides are meromorphic on $\C$ and holomorphic off $s=1$, so the
identity theorem extends it to $s\ne1$. If $\chi$ is nonprincipal,
then $\chi^*$ is nontrivial, $L(s,\chi)$ and $L(s,\chi^*)$ are entire,
and the identity theorem on $\C$ gives the identity for every $s$. A
factor $1-\chi^*(p)p^{-s}$ vanishes only if $|p^{-s}|=1$, that is,
$\operatorname{Re}s=0$.
\end{proof}

\begin{proof}[Proof of \ref{item:imprim-zeta}]
The character modulo~$1$ takes the value $1$ at every $n\ge1$, so the
two Dirichlet series coincide on $\operatorname{Re}s>1$. Conclude by
the identity theorem.
\end{proof}

\begin{proof}[Proof of \ref{item:imprim-all}]
Let $\operatorname{Re}s>1/2$, with $s\ne1$ if $\chi$ is principal. If
$\chi$ is nonprincipal, then $\chi^*$ is nontrivial, so
$L(s,\chi^*)\ne0$ by (H), also at $s=1$. If $\chi$ is principal, then
$\chi^*$ is the character modulo~$1$ and $s\ne1$, so again
$L(s,\chi^*)\ne0$. In both cases part~\ref{item:imprim-euler} applies
at $s$, and since $\operatorname{Re}s\ne0$ it gives $L(s,\chi)\ne0$.
\end{proof}

\begin{proof}[Proof of \ref{item:imprim-line}]
First let $\chi$ be primitive, $0<\operatorname{Re}s<1$ and
$L(s,\chi)=0$. The factor $(N/\pi)^{(s+\mathfrak a)/2}
\Gamma((s+\mathfrak a)/2)$ is finite and nonzero, because
$\operatorname{Re}((s+\mathfrak a)/2)>0$; the same holds at $1-s$.
Apply Theorem~\ref{thm:fe} to $\overline\chi$ at the point $1-s$:
$\xi(1-s,\overline\chi)=\epsilon(\overline\chi)\xi(s,\chi)=0$, so
$L(1-s,\overline\chi)=0$, and $\overline\chi$ is primitive. If
$\operatorname{Re}s>1/2$ this contradicts (H) for $\chi$; if
$\operatorname{Re}s<1/2$ it contradicts (H) for $\overline\chi$ at
$1-s$ (note $s\ne1$ and $1-s\ne1$). So $\operatorname{Re}s=1/2$. For
an arbitrary $\chi$, a zero of $L(s,\chi)$ in the strip is a zero of
$L(s,\chi^*)$ by part~\ref{item:imprim-euler}, so it lies on the line.

Now let $\zeta(s)=0$ with $s\ne1$ and $s\notin\{-2,-4,\dots\}$. By
part~\ref{item:imprim-zeta} and (H) for the character modulo~$1$,
$\operatorname{Re}s\le1/2$. If $0<\operatorname{Re}s<1$, then
$\operatorname{Re}s=1/2$ by the first paragraph. If
$\operatorname{Re}s\le0$, then $\operatorname{Re}(1-s)\ge1$, so
$\zeta(1-s)\ne0$ by (H). For $s\ne0$ the functional equation
$\pi^{-s/2}\Gamma(s/2)\zeta(s)=\pi^{-(1-s)/2}\Gamma((1-s)/2)
\zeta(1-s)$ then has a nonzero finite right-hand side, so
$\zeta(s)=0$ forces a pole of $\Gamma(s/2)$, that is,
$s\in\{-2,-4,\dots\}$, which was excluded. The case $s=0$ does not
occur, as $\zeta(0)=-1/2$.
\end{proof}

\begin{remark}[the step to the Riemann hypothesis]
\label{rem:zeta-suffices}
The step from the half-plane to the Riemann hypothesis for $\zeta$
does not use all Dirichlet characters, their functional equations or
the property $|\epsilon(\chi)|=1$ of the root numbers. The statement
about $\zeta$ in
Proposition~\ref{prop:imprimitive}\ref{item:imprim-line} uses only the
hypothesis for the character modulo~$1$ and Riemann's functional
equation. Where the functional equation of a primitive character is
used, in the first paragraph of the proof of
part~\ref{item:imprim-line}, it is applied to $\overline\chi$ at a
zero of $\xi(s,\chi)$, so the root number multiplies zero and none of
its properties is used.
\end{remark}
\clearpage
\clearpage
\chapter{Host fields and averaging over auxiliary primes}
\label{ch:hosts}\label{sec:hosts}

This chapter constructs the host fields attached to a primitive
character $\psi$ and bounds the characters of their Galois groups
other than the kept ones. It has two outputs.
Theorem~\ref{thm:fields-exist} gives, for every primitive $\psi$ and
suitable auxiliary primes $q,d$, a host with twist form for $\psi$.
Theorem~\ref{thm:good-pairs} gives, at every large scale, a window
pair $(q,d)$ for which every twist obeys a pointwise bound at the
test scale and the real twists have a zero-free box. Through
Proposition~\ref{prop:bounded-twist-bounds} this is the pointwise
clause of the explicit-formula condition
(Definition~\ref{def:shares-ef}), the hypothesis through which this
chapter reaches Chapter~\ref{ch:balance}. These three results are
combined in the supply of hosts, Theorem~\ref{thm:host-supply} of
Chapter~\ref{ch:families}. In Chapter~\ref{ch:assembly}, the supply
of hosts is a step of the proof of Theorem~\ref{thm:main}
(Section~\ref{sec:proof-main}), and Theorem~\ref{thm:good-pairs} is
one of the inputs that Remark~\ref{rem:content} lists as
substantive. No hypothesis of Riemann type is used
(Remark~\ref{rem:host-input-strength}).

The geometry of Chapters~\ref{ch:theta}--\ref{ch:chebyshev} lives over
an abelian CM field $E$ attached to the primitive character $\psi$
whose $L$-function is to be controlled. The Galois group of $E$ has
many characters besides the powers of $\psi$, and the balance identity
of Chapter~\ref{ch:balance} brings all of them in. Every character
of $\Gal(E/\Q)$ that is trivial on a certain subgroup is a twist
$\xi\chi_q^i\chi_d^j$ of a power $\xi$ of $\psi$ by the quadratic
characters of two auxiliary primes $q$ and $d$. The twists with
$(i,j)\ne(0,0)$ need to be bounded only at a test scale $x$ that is
polylogarithmic in their conductor, and at such a scale a twisted
prime sum is small for most choices of $q$ and $d$, by a mean-square
estimate over the auxiliary prime.

Two letters are fixed for this chapter. The $2$-adic exponent of a
frame is $\nu$; the bare letter $e$ is kept for Tate twists in
Chapter~\ref{ch:sha}. The exponent of the interval $[Y,Y^\Delta]$ on
which the auxiliary primes are tested is $\Delta$.

\section{Frames and hosts}

\begin{definition}[frames]\label{def:frame}
For an integer $\nu\ge0$ let $H_\nu$ be the quotient of
$(\Z/2^{\nu+2}\Z)^\times$ by $\{\pm1\}$.
\begin{enumerate}[label=(\arabic*)]
\item A \emph{frame} $\phi$ consists of an integer $\nu\ge0$, a finite
abelian group $A$, written additively, a surjection $h_0\colon A\to
H_\nu$, and an element $\alpha_0\in A$ with $\langle\alpha_0\rangle
\cap\ker h_0=0$.
\item Its \emph{core} is $A_0=\langle\alpha_0\rangle$, with order
$n_0=|A_0|$, and its \emph{index} is $o=[A:A_0]$.
\item Its \emph{Galois model} is $\GG=\{\pm1\}^2\times A$, with the
elements $c=((-1,-1),0)$ and $\iota=((-1,1),0)$.
\end{enumerate}
\end{definition}

The group $\GG$ will be the Galois group of the host field, with $c$
complex conjugation. The first sign records the action on
$\sqrt{-d}$ and the second the action on a second quadratic subfield;
the group $A$ carries the
$2$-power cyclotomic part through $h_0$ and the part cut out by
$\psi$.

\begin{definition}[hosts]\label{def:host}
Let $\phi$ be a frame.
\begin{enumerate}[label=(\arabic*)]
\item A \emph{host} $\mathcal H$ for $\phi$ consists of:
\begin{enumerate}[label=(\roman*)]
\item primes $q\equiv3\pmod 8$ and $d\equiv7\pmod{16}$ with
$(-d/q)=1$, the \emph{auxiliary primes};
\item an integer $Q_S>0$ with $2^{\nu+2}\mid Q_S$ and
$\gcd(Q_S,qd)=1$, every odd prime divisor of which splits in
$K=\Q(\sqrt{-d})$;
\item a number field $E$, Galois over $\Q$, with an isomorphism
$\Gal(E/\Q)\cong\GG$, containing $\sqrt{-d}$, on which $\GG$ acts
through the first sign of $\{\pm1\}^2$, and with an embedding
$E\hookrightarrow\Q(\zeta_{Q_Sqd})$;
\item an Artin map $(\Z/Q_Sqd\Z)^\times\to\GG$, compatible with the
cyclotomic action, sending $-1$ to $c$, compatible with $h_0$ on the
$2$-power part, with inertia at $d$ of order at most $2$ and inertia
at $q$ of exponent $2$.
\end{enumerate}
\item The \emph{modulus} of $\mathcal H$ is $N_{\mathcal H}=Q_Sqd$.
\item For a character $\theta$ of $\GG$, $\chi_\theta$ is the
primitive Dirichlet character that agrees with $\theta$ composed with
the Artin map on the integers prime to $N_{\mathcal H}$.
\item $\omega_q$ is the character of $\GG$ given by the second sign.
\item The \emph{half root discriminant} of $\mathcal H$ is
$v=\frac12|\operatorname{disc}E|^{1/[E:\Q]}$.
\end{enumerate}
\end{definition}

From these conditions the Artin map is surjective, $E$ is a CM field
and $c$ is complex conjugation. The subgroup $\Phi=\Gal(E/K)$ has
index~$2$; it is the CM type induced from $K$
(Definition~\ref{def:additive}). Figure~\ref{fig:hosttower} summarises
the fields and the group.

\begin{figure}[htbp]
\centering
\input{alpoge/figs/v2/hosttower.tex}
\caption{The host field of Definitions~\ref{def:frame}
and~\ref{def:host}. The field $E$ is abelian over $\Q$, contains
$\sqrt{-d}$ and embeds in $\Q(\zeta_{Q_Sqd})$, where $q\equiv3\pmod8$
and $d\equiv7\pmod{16}$ are the auxiliary primes. The group
$\Phi=\Gal(E/K)$ has index~$2$ and is the CM type induced from $K$;
$E$ is a CM field, and $c$ is complex conjugation.}
\label{fig:hosttower}
\end{figure}

The host field is cut out of a cyclotomic field by explicit congruence
conditions.

\begin{definition}[$\psi$-fields]\label{def:psifield}
Let $\psi$ be a primitive character of conductor $N_\psi$, let
$\nu\ge0$, and let $q,d$ be primes. A number field $E$ is a
\emph{$\psi$-field} for $(\psi,\nu,q,d)$ if there are an integer $N$
with $2^{\nu+2},N_\psi,q,d\mid N$ and an embedding
$E\hookrightarrow\Q(\zeta_N)$ whose image is the fixed field of the
group $U$ of classes $a\in(\Z/N\Z)^\times$ with
$a\equiv\pm1\pmod{2^{\nu+2}}$, $\psi(a)=1$, $(a/q)=1$ and $(a/d)=1$.
\end{definition}

\begin{proposition}\label{prop:psifield-chars}
Let $\psi$ be a primitive character, $\nu\ge0$ an integer and $q,d$
primes. Let $E$ be a $\psi$-field for $(\psi,\nu,q,d)$, with integer
$N$ and embedding $E\hookrightarrow\Q(\zeta_N)$. Then the characters of
$\Gal(E/\Q)$, viewed through the embedding as Dirichlet characters
modulo $N$, are exactly
the products $\chi_{(2)}\,\xi\,\chi_q^i\chi_d^j$ with $\chi_{(2)}$ an
even character modulo $2^{\nu+2}$, $\xi\in\langle\psi\rangle$,
$\chi_q=(\cdot/q)$, $\chi_d=(\cdot/d)$ and $i,j\in\{0,1\}$.
\end{proposition}

\begin{proof}
Let $U\subset(\Z/N\Z)^\times$ be the group of classes $a$ with
$a\equiv\pm1\pmod{2^{\nu+2}}$, $\psi(a)=1$, $(a/q)=1$ and $(a/d)=1$,
so that $E$ is the fixed field of $U$ (Definition~\ref{def:psifield}).
By the Galois
correspondence for $\Q(\zeta_N)$, $\Gal(E/\Q)\cong(\Z/N\Z)^\times/U$, so its
characters are the Dirichlet characters
modulo $N$ trivial on $U$ \cite{Washington1997}. Let $X$ be the group
generated by the even characters modulo $2^{\nu+2}$, $\psi$, $\chi_q$
and $\chi_d$. A class $a$ lies in the kernel of every even character
modulo $2^{\nu+2}$ if and only if $a\equiv\pm1\pmod{2^{\nu+2}}$,
because these characters are the characters of
$(\Z/2^{\nu+2}\Z)^\times/\{\pm1\}$. Hence $U=\bigcap_{\chi\in X}\ker
\chi$. For a finite abelian group $G$ and a subgroup $X$ of its dual,
the characters trivial on $\bigcap_{\chi\in X}\ker\chi$ form $X$, by
duality of finite abelian groups. Every element of $X$ has the stated
form, because $X$ is abelian, generated by these characters, and
$\chi_q^2=\chi_d^2=1$.
\end{proof}

\section{Kept and bounded characters}

The balance identity of Chapter~\ref{ch:balance} produces a
combination of the functions $g_{\chi_\theta}$ over characters
$\theta$ of $\GG$ trivial on $A_0$. Those that belong to the level are
kept; the others must be bounded.

\begin{definition}[kept and bounded characters]\label{def:bounded}
Let $\mathcal H$ be a host for a frame $\phi$, and let $S$ be a finite
set of primitive characters.
\begin{enumerate}[label=(\arabic*)]
\item $X_+$ is the set of characters $\theta$ of $\GG$ trivial on
$A_0$ with $\theta(\iota)=1$, and $X_-$ the set of those with
$\theta(\iota)=-1$.
\item A character $\theta\in X_+\cup X_-$ is \emph{kept} if
$\chi_\theta\in S$.
\item A character $\theta$ is \emph{bounded} if $\theta\in X_+$ and
$\theta$ is not kept, or $\theta\in X_-$.
\item $\mathcal H$ has \emph{kept characters in $X_+$} if every
element of $S$ is $\chi_\theta$ for some $\theta\in X_+$.
\item Let $\psi_{\mathrm{exc}}$ be an optional element of $S$. The
host \emph{keeps the exceptional character} $\psi_{\mathrm{exc}}$ if
$\psi_{\mathrm{exc}}=\chi_{\omega_q}$ when $\psi_{\mathrm{exc}}$ is
given, and $\omega_q$ is not kept when it is not given.
\end{enumerate}
\end{definition}

For a level $\lev$ we take $S=\Kept(\lev)$ and $\psi_{\mathrm{exc}}$
the exceptional character of $\lev$
(Definition~\ref{def:level}).

\begin{definition}[twist form]\label{def:twistform}
Let $\psi$ be a primitive character and $\mathcal H$ a host for a frame
$\phi$, with auxiliary primes $q,d$, and put $\chi_q=(\cdot/q)$ and
$\chi_d=(\cdot/d)$. The host has \emph{twist form for $\psi$} if every
character
$\theta$ of $\GG$ trivial on $A_0$ admits $\xi\in\langle\psi\rangle$
and $i,j\in\{0,1\}$ with $\theta(\iota)=(-1)^j$ and
$\theta(a)=\xi(a)\chi_q(a)^i\chi_d(a)^j$ for every integer $a$ prime
to $N_{\mathcal H}$, $\theta$ being composed with the Artin map.
\end{definition}

\begin{proposition}\label{prop:bounded-twists}
Let $\psi$ be a primitive character and $\mathcal H$ a host with twist
form for $\psi$, with auxiliary primes $q,d$, and put
$\chi_q=(\cdot/q)$ and $\chi_d=(\cdot/d)$. Let $\theta$ be bounded for
$S=\Kept(\lev_\psi)$.
Let $\xi\in\langle\psi\rangle$ and $i,j\in\{0,1\}$ satisfy
$\theta(\iota)=(-1)^j$ and $\theta(a)=\xi(a)\chi_q(a)^i\chi_d(a)^j$
for every integer $a$ prime to $N_{\mathcal H}$. Then
$(i,j)\ne(0,0)$.
\end{proposition}

\begin{proof}
Suppose $(i,j)=(0,0)$. Then $\theta(\iota)=1$, so $\theta\notin X_-$,
and therefore $\theta$ is in $X_+$ and not kept. But $\theta$ agrees
with $\xi$ on the integers prime to $N_{\mathcal H}$, so
$\chi_\theta$ is the primitive character inducing
$\xi\in\langle\psi\rangle$, which lies in $\Kept(\lev_\psi)$. Thus
$\theta$ is kept, a contradiction.
\end{proof}

Table~\ref{tab:chargrid} lists the four kinds of characters.

\begin{table}[htbp]
\centering
\input{alpoge/figs/v2/chargrid.tex}
\caption{Characters $\theta$ of $\GG$ trivial on $A_0$, for a host with
twist form for $\psi$. Each is $\xi\chi_q^i\chi_d^j$ with
$\xi\in\langle\psi\rangle$ and $\theta(\iota)=(-1)^j$
(Definition~\ref{def:twistform}). By
Proposition~\ref{prop:bounded-twists} a bounded character has
$(i,j)\ne(0,0)$. Conversely, $(0,1)$ and $(1,1)$ are bounded because
$\theta(\iota)=-1$, and $(1,0)$ is bounded when $q\nmid N_\psi$, since
it is then not kept (Definition~\ref{def:bounded}). The bounded
characters are controlled by the pointwise bound at the test scale,
clause~\ref{item:gp-pointwise} of a good pair (Definition~\ref{def:good-pair});
for real
twists, clause~\ref{item:gp-box} adds a zero-free box.}
\label{tab:chargrid}
\end{table}

Hosts with all these properties exist for every primitive $\psi$.

\begin{theorem}[Alp\"oge]\label{thm:fields-exist}
Let $\psi$ be a primitive character of conductor $N_\psi$. There are
an integer $c_\psi\ge1$ and $\nu_0$ with the following
property. Let $\nu\ge\nu_0$,
let $z\ge4N_\psi$, and let $q,d$ be primes with $q\equiv3\pmod8$,
$q\ne3$, $d\equiv7\pmod{16}$, $(-d/q)=1$ and $\gcd(N_\psi,qd)=1$, such
that every odd prime $p\le z$ splits in $\Q(\sqrt{-d})$. Then there
are a frame $\phi$ with exponent $\nu$ and $|A|=c_\psi2^\nu$, and a host
$\mathcal H$ for $\phi$ with auxiliary primes $q,d$, such that, for
$S=\Kept(\lev_\psi)$: the index of $\phi$ is $o(\lev_\psi)$, the host
has kept characters in $X_+$, it keeps the exceptional character of
$\lev_\psi$, and it has twist form for $\psi$.
\end{theorem}

\begin{proof}[Idea of proof]
We describe the standard route. One takes for $E$ a $\psi$-field
(Definition~\ref{def:psifield}) inside $\Q(\zeta_N)$ for a suitable
$N$ divisible by $2^{\nu+2}$, $N_\psi$, $q$ and $d$. By
Proposition~\ref{prop:psifield-chars} the character group of
$\Gal(E/\Q)$ is generated by the even characters modulo $2^{\nu+2}$,
$\psi$, $\chi_q$ and $\chi_d$, and the frame is read off from this
description: the two signs come from $\chi_q$ and $\chi_d$ (the
first sign is the action on $\sqrt{-d}$), and $A$ carries the
$2$-power and $\psi$-parts, so that the characters trivial on $A_0$
are exactly the twists of Definition~\ref{def:twistform}.
The splitting in $K$ of the odd primes $p\le z$ is a hypothesis of
the theorem. The conditions
on $Q_S$ are congruence conditions, handled by the Chinese remainder
theorem and quadratic reciprocity. All of this is standard class field theory
over
$\Q$. In the application the splitting hypothesis is arranged by the
choice of a window pair (Definition~\ref{def:scales} and the paragraph
after it).
\end{proof}

\section{Scales and the explicit-formula condition}

The bounded characters are needed only at one scale $x$, for top
places with prescribed shares.

\begin{definition}[scales and window pairs]\label{def:scales}
Let $\eps\in(0,1/52)$.
\begin{enumerate}[label=(\arabic*)]
\item A \emph{scale} is a triple $(x,\eps,J)$ with $x>1$ and $J\ge1$.
Put $b=\exp(x^\eps)$ and $a=xb$.
\item A \emph{system of scales} of exponent $\lambda$ is a $J$ together
with a set $\Win\subset\R$ that has windows of exponent $\lambda$
(Definition~\ref{def:windows}).
\item An integer $\nu$ is \emph{adapted} to $(\eps,x)$ if
$x^{2+8\eps}\le2^\nu\le2x^{2+8\eps}$.
\item A pair of primes $(q,d)$ is a \emph{window pair} for $(\eps,Y)$
if $q\equiv3\pmod8$ and $Y^{\eps/5}/2<\log q\le Y^{\eps/5}$;
$d\equiv7\pmod{16}$ and $Y^{\eps/5}\log^2Y/2<\log d\le
Y^{\eps/5}\log^2Y$; $q\mid d+1$; and every odd prime $p\le Y^{\eps/5}$
divides $d+1$.
\end{enumerate}
\end{definition}

The upper bound $\eps<1/52$ is fixed once and for all and recurs in
the statements below. Only small values of $\eps$ are used in the end:
the proof of Theorem~\ref{thm:zero-free-primitive} lets $\eps\to0$.

For a window pair, $d\equiv-1\pmod p$ for every odd prime
$p\le Y^{\eps/5}$ and for $p=q$, so all these primes split in
$\Q(\sqrt{-d})$, and $(-d/q)=1$. Thus a window pair meets the
conditions of Theorem~\ref{thm:fields-exist} with $z=Y^{\eps/5}$ once
$Y$ is large. Since $\log q\approx Y^{\eps/5}$, a test point
$y\in[Y,Y^\Delta]$ is polylogarithmic in $q$, and hence in the
conductors of the twists (Figure~\ref{fig:hosts-scales}).

A host also carries two further primes, used in the congruence
modules of Chapters~\ref{ch:theta}--\ref{ch:chebyshev}. They are
fixed here because their sizes depend on the scale.

\begin{definition}[norm primes]\label{def:norm-primes}
Let $\mathcal H$ be a host with field $E$ and auxiliary primes $q,d$,
let $K=\Q(\sqrt{-d})\subset E$, let $k\ge1$ be an integer and let $b>1$ be real.
\begin{enumerate}[label=(\arabic*)]
\item A pair of \emph{norm primes} of $\mathcal H$ at $(k,b)$ is a
pair of distinct primes $\varpi_0,\varpi_1$ such that, for
$i\in\{0,1\}$, $(kb)^8\le\varpi_i\le2(kb)^8$ and
$\varpi_i=\pi_i\overline{\pi_i}$ for some $\pi_i\in\mathcal O_K$ with
$\pi_i\equiv1\pmod{N_{\mathcal H}\mathcal O_E}$.
\item For such a pair, $\lgv=\log(\varpi_0\varpi_1)$.
\end{enumerate}
\end{definition}

In applications $k$ is the first entry of a ranks datum
(Definition~\ref{def:ranks}) and $b=\exp(x^\eps)$ is the parameter of
the scale. From $\pi_i\equiv1$ modulo $N_{\mathcal H}$ one gets
$\varpi_i\equiv1\pmod{N_{\mathcal H}}$, so $\varpi_i$ splits
completely in $\Q(\zeta_{N_{\mathcal H}})$ and hence in $E$. The norm
primes are level primes at which the local type of the automorphic
forms is prescribed; they enter the later chapters only through
$\lgv$, in the upper link (Definition~\ref{def:upper-link}) and in the
ties (Definition~\ref{def:tied}). They are distinct from the auxiliary
primes $q,d$ and from the core generator $\alpha_0$ of the frame.

\begin{definition}[shares and the explicit-formula condition]
\label{def:shares-ef}
Let $\lev$ be a level, $\mathcal H$ a host for a frame $\phi$ with
core $A_0$ and index $o$, $\Top\subset A$ a finite set, and
$x,\eps,C,\delta$ real.
\begin{enumerate}[label=(\arabic*)]
\item The \emph{share} of $\alpha\in A$ is $\lambda_\phi(\alpha)=
4/(3o+1)$ if $\alpha\in A_0$ and $3/(3o+1)$ otherwise.
\item $\Top$ \emph{has shares up to} $\delta$ if for every
$\alpha\in A$,
\[
\bigl|\#\{\tau\in\Top:\tau-\alpha\in A_0\}-\lambda_\phi(\alpha)|\Top|
\bigr|\le\delta|\Top|.
\]
\item The host \emph{satisfies the explicit-formula condition}
$\mathrm{EF}(\lev,\Top,x,\eps,C)$ if $o=o(\lev)$; the host has kept
characters in $X_+$ and keeps the exceptional character of $\lev$,
for $S=\Kept(\lev)$; $\Top$ has shares up to $x^{-1+\eps}$; and every
bounded $\theta$ satisfies $|g_{\chi_\theta}(x)|\le x^{1/2+\eps}$ and
$|L'(1,\chi_\theta)/L(1,\chi_\theta)|\le C\log x$.
\end{enumerate}
\end{definition}

The shares produce the weights of a level: averaging the
characters of $A/A_0$ over a set $\Top$ with these shares gives $1$
for the trivial character and $(3o+1)^{-1}$ for the others
(Proposition~\ref{prop:shares}). In Chapter~\ref{ch:balance} the set
$\Top$ is the set of top places. The rest of this chapter shows how
the pointwise clause is obtained.

\section{Averaging over the auxiliary prime}

The pointwise bounds for the twists are obtained by choosing $q$ and
$d$ at random. The objects of the averaging are the following.

\begin{definition}[head sums, prime sets and boxes]\label{def:headsum}
Let $Z,y>0$, let $q\ge1$ be an integer and let
$\kappa\colon\Z_{\ge0}\to\C$.
\begin{enumerate}[label=(\arabic*)]
\item The \emph{head sum} is
\[
S(Z,y,q,\kappa)=\sum_{1\le n\le Z}\Lambda(n)\min\Bigl(1,\frac yn
\Bigr)\kappa(n)\Bigl(\frac nq\Bigr),
\]
with $(\frac nq)$ the Jacobi symbol.
\item For integers $m_0,a_0$ and $U>0$, $P(m_0,a_0,U)$ is the set of
primes $q$ with $U^{1/2}<q\le U$ and $q\equiv a_0\pmod{m_0}$.
\item For $\eps,Y,U>0$ the pair $(m_0,a_0)$ is \emph{admitted} if
$m_0>0$, $8\mid m_0$, $\gcd(a_0,m_0)=1$,
$\frac14(\log Y)^2\log m_0\le\log U$ and
$Y^{\eps/6}\le\log U\le Y^{\eps/5}(\log Y)^2$.
\item For $C,Z\in\R$, a finite set $P$ of integers and $I\subset\R$,
the \emph{mean-square bound} $\mathrm{MeanSq}(C,Z,P,I)$ holds if
\[
\sum_{q\in P}|S(Z,y,q,\kappa)|^2\le C\,y\log y\,|P|
\]
for every $y\in I$ and every $\kappa$ with $|\kappa(n)|\le1$ for all
$n$.
\item A Dirichlet character $\chi$ modulo $N$ has a \emph{zero-free
box of height} $Q$ if $L(s,\chi)\ne0$ whenever
$\operatorname{Re}s>1-1/400$ and $|\operatorname{Im}s|\le Q$.
\item For $Q,C,\tau,Y,X\in\R$, a Dirichlet character $\chi$ modulo $N$
\emph{has a box} $\mathrm{Box}(Q,C,\tau,Y,X)$ if $\chi\ne1$, $N\le
Q$, $\chi$ has a zero-free box of height $Q$,
$|L'(1,\chi)/L(1,\chi)|\le C\log Y$, and $|g_\chi(y)|\le
y^{1/2+\tau}$ for every $y\in[Y,X]$.
\end{enumerate}
\end{definition}

The head sum is a twisted prime sum truncated at $Z$, with the weight
$\min(1,y/n)$ in place of $\max(y/n-1,0)$. For $\chi(n)=\kappa(n)
(\frac nq)$ nonprincipal and $y\le Z$, it is related to $g_\chi$ by
an exact identity. Write $\sum_{n\le y}\Lambda(n)\chi(n)(y/n-1)=
y\sum_{n\le y}\Lambda(n)\chi(n)/n-\sum_{n\le y}\Lambda(n)\chi(n)$,
and use $-L'(1,\chi)/L(1,\chi)=\sum_n\Lambda(n)\chi(n)/n$, a
convergent series by the prime number theorem in progressions. Then
\begin{equation}\label{eq:head-tail}
g_\chi(y)=-S(Z,y,q,\kappa)-y\sum_{n>Z}\frac{\Lambda(n)\chi(n)}{n}.
\end{equation}
So $g_\chi(y)$ is small when the head sum is small and the tail is
small, and the tail is controlled by the zeros of $L(s,\chi)$ near
$\operatorname{Re}s=1$.

\begin{definition}[good pairs]\label{def:good-pair}
Let $M\ge1$ be an integer, $\eps>0$ and $Y\le X$.
\begin{enumerate}[label=(\roman*)]
\item The \emph{box heights} are
$T_j=8M\exp(Y^{\eps/5})\exp(Y^{\eps/5}\log^2Y)^j$ for $j\in\{0,1\}$.
\item A window pair $(q,d)$ for $(\eps,Y)$ with $\gcd(M,qd)=1$ is
\emph{good} for $(M,\eps,Y,X)$ if it satisfies
clauses~\ref{item:gp-pointwise} and~\ref{item:gp-box} below, where
$\chi_q=(\cdot/q)$ and $\chi_d=(\cdot/d)$.
\end{enumerate}
\begin{enumerate}[label=(\arabic*)]
\item\label{item:gp-pointwise} For every character $\xi$ modulo $M$, every
$(i,j)\ne(0,0)$ and the primitive character $\chi$ inducing
$\xi\chi_q^i\chi_d^j$ (with respect to the modulus $2Mqd$), every
$y\in[Y,X]$ satisfies $|g_\chi(y)|\le y^{1/2+\eps/2}$ and
$|L'(1,\chi)/L(1,\chi)|\le4\log y$.
\item\label{item:gp-box}
For every real primitive character $\varrho$ with conductor
dividing $8M$, every $(i,j)\ne(0,0)$ and the primitive character
$\chi$ inducing $\varrho\chi_q^i\chi_d^j$, $\chi$ has a zero-free box
of height $T_j$.
\end{enumerate}
\end{definition}

The letter $T$ is the usual one for the height of a box of zeros; the
letter $H$ is reserved for the groups $H_\nu$ of
Definition~\ref{def:frame} and $H$ of Definition~\ref{def:additive}.
Since $Y^{\eps/5}/2<\log q\le
Y^{\eps/5}$, the height $T_0$ lies in $[8Mq,8Mq^2)$, and similarly
$T_1\in[8Mqd,8Mq^2d^2)$. So the zero-free boxes in clause~\ref{item:gp-box}
reach a
fixed power of the conductor of $\chi$, whose $q$- and $d$-parts are
at most $q$ and $d$.
\begin{figure}[tbp]
\centering
\input{book/figs/ch05-scale-ladder.tex}
\end{figure}

The averaging rests on four estimates. The first is a mean-square
bound over primes in a progression.

\begin{proposition}[Alp\"oge]\label{prop:meansq}
There is $C>0$ such that for all $\eps\in(0,1]$ and $\Delta\ge1$
there is $Y_0$ with the following property. Let $Y\ge Y_0$ and $U>0$,
and let $(m_0,a_0)$ be admitted for $(\eps,Y,U)$, with every prime
divisor $p$ of $m_0$ satisfying $p\le Y^{\eps/5}$ or
$p>Y^{400\Delta}$. Then $\mathrm{MeanSq}(C,Y^{400\Delta},
P(m_0,a_0,U),[Y,Y^\Delta])$ holds.
\end{proposition}

\begin{proof}[Idea of proof]
This is our account of the standard method. Put $Z=Y^{400\Delta}$
and $w(n)=\Lambda(n)\min(1,y/n)$, and expand the square:
\[
\sum_{q\in P}|S(Z,y,q,\kappa)|^2=\sum_{n,n'\le Z}w(n)w(n')\kappa(n)
\overline{\kappa(n')}\sum_{q\in P}\Bigl(\frac{nn'}q\Bigr).
\]
The diagonal terms, with $nn'$ a square, contribute
$\ll|P|\sum_nw(n)^2\ll|P|\,y\log y$. For $nn'$ not a square,
quadratic reciprocity makes $q\mapsto(\frac{nn'}q)$ a character of
modulus $\ll Z^2$. The sum over $q\in P(m_0,a_0,U)$ is then a sum
over primes in a progression to a modulus $\ll Z^2m_0$. Its logarithm
is $o(\log U)$, since $\log Z^2\ll\Delta\log Y$, $\log U\ge Y^{\eps/6}$
and $\log m_0\le4\log U/(\log Y)^2$. It need not be polylogarithmic
in $U$, because of the factor $m_0$. In this range the prime number
theorem in progressions, with the Deuring--Heilbronn phenomenon for a
possible exceptional zero (Remark~\ref{rem:classical-inputs}), gives
cancellation, unless the character is determined by the
progression. The hypothesis on $m_0$ says that $m_0$
has no prime factor in $(Y^{\eps/5},Z]$. The terms $n=p^r$ whose
symbol is constant on the progression ($p\le Y^{\eps/5}$, or $r$
even) contribute $O(y^{1/2}+Y^{\eps/5}\log Y)$ to the head sum,
which is within the bound. This is a large-sieve type argument.
\end{proof}

\begin{remark}[a heuristic]\label{rem:indep-signs}
For $n\le(\log q)^{O(1)}$ the values $(n/q)$ behave like independent
signs as $q$ runs over primes, by the prime number theorem in
progressions to the modulus $\prod_{p\le n}p$, which is small on the
logarithmic scale of $q$. This is why the head sums are small on
average. In this step no circularity with the Riemann hypothesis
appears: the moduli of the progressions are far smaller than $q$.
\end{remark}

The second converts the mean square into a bound for all but a small
proportion of the moduli, uniformly in $y$.

\begin{proposition}[Alp\"oge]\label{prop:bad-moduli}
Let $C\ge0$, $\tau>0$ and $r_0>0$. There is $Y_0$ such that for all
$Y\ge Y_0$, all $X,Z\in\R$, all finite sets $P$ of natural numbers,
all subsets $B\subseteq P$ and all $\kappa$ with $|\kappa(n)|\le1$ for every
$n$ the
following holds. If $\mathrm{MeanSq}(C,Z,P,[Y,X])$ holds and every
$q\in B$ has some $y\in[Y,X]$ with $|S(Z,y,q,\kappa)|>y^{1/2+\tau}/2$,
then $|B|\le r_0|P|$.
\end{proposition}

\begin{proof}[Idea of proof]
At a single $y$, Chebyshev's inequality applied to the mean square
bounds the number of $q$ with $|S(Z,y,q,\kappa)|>y^{1/2+\tau}/2$ by
$4C\log y\,y^{-2\tau}|P|$. A bad modulus needs a large head sum only
at \emph{some} $y\in[Y,X]$, so one must also pass from single points
to intervals. The function $y\mapsto S(Z,y,q,\kappa)$ is piecewise
linear, with derivative $\sum_{y<n\le Z}\Lambda(n)\kappa(n)
(\frac nq)/n$ off the integers. On each dyadic block $[W,2W]$ a
Sobolev-type inequality bounds $\sup|S|^2$ by its mean over the block
plus an integral of $|S|\,|\partial_yS|$, both of which are
controlled by mean squares. Summing over the blocks, the bounds
$\ll\log W\,W^{-2\tau}$ form a convergent series whose sum tends to
$0$ as $Y\to\infty$. This account is not a complete proof.
\end{proof}

The third counts the moduli for which a box fails.

\begin{proposition}[Alp\"oge]\label{prop:zero-box}
There is $C$ with the following property. Let $Q\ge3$ and $r_1$ be
real, and let $B$ be a finite set of primes such that each $q\in B$ has
$q>r_1$ and admits a positive integer $r\le r_1$ with $rq\le Q$ and a
primitive character modulo $rq$ without a zero-free box of height $Q$.
Then $|B|\le CQ^{1/4}$.
\end{proposition}

\begin{proof}[Idea of proof]
The input is a zero-density estimate for boxes near
$\operatorname{Re}s=1$ \cite{MontgomeryVaughan2007}: the number of
primitive characters of conductor at most $Q$ with a zero in
$\operatorname{Re}s>1-1/400$, $|\operatorname{Im}s|\le Q$, is a small
power of $Q$. Since $q>r_1\ge r$, distinct $q\in B$ give characters
of distinct conductors $rq$, so $|B|$ is at most that count.
\end{proof}

The fourth turns a bound for the head sum into a box.

\begin{proposition}[Alp\"oge]\label{prop:box-from-headsum}
There is $C>0$ such that for all $\eps\in(0,1/4]$, $\tau>0$ and
$\Delta\ge1$ there is $Y_0$ with the following property. Let $Y\ge
Y_0$, $Q\in\R$, $N\ge1$, let $\chi$ be a Dirichlet character modulo
$N$, $q\ge0$ an integer and $\kappa\colon\Z_{\ge0}\to\C$. Suppose that
$Y^{\eps/6}\le\log Q\le Y^{\eps/4}$; that $\chi\ne1$, $N\le Q$ and
$\chi$ has a zero-free box of height $Q$; that $\chi(n)=\kappa(n)
(\frac nq)$ for every $n\ge0$; and that
$|S(Y^{400\Delta},y,q,\kappa)|\le y^{1/2+\tau}/2$ for every
$y\in[Y,Y^\Delta]$. Then $\chi$ has a box
$\mathrm{Box}(Q,C,\tau,Y,Y^\Delta)$.
\end{proposition}

\begin{proof}[Idea of proof]
By the identity~\eqref{eq:head-tail} with $Z=Y^{400\Delta}$, it
suffices to show that the tail $y\sum_{n>Z}\Lambda(n)\chi(n)/n$ is at
most $y^{1/2+\tau}/2$ for $y\le Y^\Delta$, and to bound
$L'(1,\chi)/L(1,\chi)=-\sum_{n\le Z}\Lambda(n)\chi(n)/n
-\sum_{n>Z}\Lambda(n)\chi(n)/n$. Both follow from a truncated
explicit formula for the tail, truncated at height $Q$. Since
$L(s,\chi)$ has no zero with $\operatorname{Re}s>1-1/400$ and
$|\operatorname{Im}s|\le Q$, the contour can be moved to
$\operatorname{Re}s=1-1/400$ below height $Q$, at a cost of
$Z^{-1/400}=Y^{-\Delta}$ times factors polynomial in $\log Q\le
Y^{\eps/4}$; the truncation error is controlled because $Q$ is large
compared with $\log Z$. The head of $L'/L(1,\chi)$ is at most $\log Z\ll
\Delta\log Y$. Our route therefore gives $C$ depending on $\Delta$;
the statement has $C$ uniform in $\Delta$, and we do not reconstruct
the sharper bound on $L'(1,\chi)/L(1,\chi)$ that this requires. This
is of
Littlewood type
\cite{Littlewood1928}, with the zero-free box in place of the
generalised Riemann hypothesis. Littlewood's bound is for
$L(1,\chi)$, and the bound here is for $L'(1,\chi)/L(1,\chi)$.
\end{proof}

With these four estimates good pairs exist at every large scale (Figure~\ref{fig:averaging-flow}).

\begin{theorem}[Alp\"oge]\label{thm:good-pairs}
Let $M>0$ be an integer, $0<\eps<1/52$ and $\Delta\ge1$. There is
$Y_1$ such that for every $Y\ge Y_1$ some window pair for $(\eps,Y)$
is good for $(M,\eps,Y,Y^\Delta)$.
\end{theorem}

\begin{proof}[Proof (outline)]
Most pairs are good. For $y\in[Y,Y^\Delta]$ and
$\log q\approx Y^{\eps/5}$ the test point $y$ is polylogarithmic in
the conductor. By~\eqref{eq:head-tail} with $Z=Y^{400\Delta}$,
$g_\chi(y)$ is minus the head sum $S(Z,y,q,\kappa)$ minus a tail. For the
twists with $(i,j)=(1,0)$ one takes $\kappa=\xi$ and averages over
$q$. The head sums have mean square $\ll y\log y$ over the primes $q$
in a progression (Proposition~\ref{prop:meansq}), so all but a
proportion $r_0$ of them satisfy the bound for every $y$ in the range
(Proposition~\ref{prop:bad-moduli}). The tail is negligible when
$L(s,\chi)$ has a zero-free box, which fails for $\ll Q^{1/4}$
primes (Proposition~\ref{prop:zero-box}), a negligible proportion.
Proposition~\ref{prop:box-from-headsum} then gives the bounds of
clause~\ref{item:gp-pointwise} for most $q$, and the box count gives
clause~\ref{item:gp-box} for the
real twists. One chooses $q$ in the good majority, and then $d$ in the
same way with $q$ fixed, taking $\kappa=\xi\chi_q^i$ for the twists
with $j=1$, in the class $d\equiv-1$ modulo
$q\prod_{3\le p\le Y^{\eps/5}}p$, so that $(q,d)$ is a window pair
(Definition~\ref{def:scales}). The constant $4$ in
clause~\ref{item:gp-pointwise} uses that the constant $C$ of
Proposition~\ref{prop:box-from-headsum} does not depend on $\Delta$.
\end{proof}

Good pairs give the pointwise clause of
Definition~\ref{def:shares-ef} at the scale $x$, with $C=4$.
\begin{figure}[tbp]
\centering
\input{book/figs/ch05-averaging-flow.tex}
\end{figure}

\begin{proposition}\label{prop:bounded-twist-bounds}
Let $\psi$ be a primitive character of conductor $N_\psi$ and
$\mathcal H$ a host with twist form for $\psi$, with $N_\psi\mid Q_S$
and every prime divisor of $Q_S$ dividing $2N_\psi$. Let $\eps>0$,
$x\ge1$ and $Y\le x$, and suppose that the auxiliary primes $(q,d)$ of
$\mathcal H$ are good for $(N_\psi,\eps,Y,x)$. Then every $\theta$
that is bounded for $\Kept(\lev_\psi)$ satisfies
$|g_{\chi_\theta}(x)|\le x^{1/2+\eps}$ and
$|L'(1,\chi_\theta)/L(1,\chi_\theta)|\le4\log x$.
\end{proposition}

\begin{proof}
Since $\mathcal H$ has twist form for $\psi$
(Definition~\ref{def:twistform}), there are $\xi\in\langle\psi\rangle$
and $i,j\in\{0,1\}$ with $\theta(\iota)=(-1)^j$ and
$\theta(a)=\xi(a)\chi_q(a)^i\chi_d(a)^j$ for every integer $a$ prime
to $N_{\mathcal H}$. By
Proposition~\ref{prop:bounded-twists}, $(i,j)\ne(0,0)$. An integer $a$
prime to $2N_\psi qd$ is prime to $N_{\mathcal H}=Q_Sqd$, since every
prime divisor of $Q_S$ divides $2N_\psi$. So $\chi_\theta(a)=
\theta(a)=\xi(a)\chi_q(a)^i\chi_d(a)^j$ for such $a$, and
$\chi_\theta$ is the primitive character inducing
$\xi\chi_q^i\chi_d^j$ with respect to $2N_\psi qd$.
Clause~\ref{item:gp-pointwise} of
Definition~\ref{def:good-pair} with $M=N_\psi$, at $y=x\in[Y,x]$,
gives $|g_{\chi_\theta}(x)|\le x^{1/2+\eps/2}\le x^{1/2+\eps}$, as
$x\ge1$, and the bound for $L'(1,\chi_\theta)/L(1,\chi_\theta)$.
\end{proof}

\section{The strength of the analytic input}

\begin{remark}[the analytic input is not of Riemann type]
\label{rem:host-input-strength}
The argument does not need, for every character of the host field
that is not kept, a pointwise bound for $g_\chi$ of
Riemann-hypothesis quality at all large $y$. The analytic
input on the twists is the explicit-formula condition of
Definition~\ref{def:shares-ef}, whose pointwise clause is required at
the single scale $x$, which is polylogarithmic in the conductor. It is
not assumed: it follows from Theorem~\ref{thm:good-pairs} and
Proposition~\ref{prop:bounded-twist-bounds}, through the supply of
hosts in Theorem~\ref{thm:host-supply}. In terms of zeros, it gives
pointwise bounds at polylogarithmic scales and zero-freeness of the
real twists in the box $\operatorname{Re}s>1-1/400$,
$|\operatorname{Im}s|\le T_j$, whose height is a power of the
conductor. Neither statement has the strength of the Riemann
hypothesis near $\operatorname{Re}s=1/2$.
\end{remark}

\begin{remark}[classical inputs]\label{rem:classical-inputs}
The classical inputs to the averaging are of the following kinds.
\begin{itemize}
\item The zero-free region with at most one exceptional zero
\cite{Page1935}.
\item Siegel's theorem \cite{Siegel1935}, with its Siegel--Walfisz
form in progressions \cite{Walfisz1936}.
\item The Deuring--Heilbronn phenomenon in the Linnik range
\cite{Linnik1944}.
\item A log-free zero-density estimate \cite{Gallagher1970}, behind
Proposition~\ref{prop:zero-box}.
\item Mean squares of real character sums
\cite{Montgomery1971,HeathBrown1995}, behind
Proposition~\ref{prop:meansq}.
\item Mertens-type bounds.
\end{itemize}
None of
Propositions~\ref{prop:meansq}--\ref{prop:box-from-headsum} assumes a
hypothesis of Riemann type.
\end{remark}
\clearpage
\clearpage
\chapter{The balance identity and Fourier analysis on the Galois group}
\label{ch:balance}\label{sec:balance}

This chapter proves the balance identity,
Proposition~\ref{prop:balance}, and outlines the proof of the core
Chebyshev bound, Theorem~\ref{thm:core-bound}. It also treats the
averaging results
(Proposition~\ref{prop:shares}, Proposition~\ref{prop:averaged} and
Theorem~\ref{thm:averaged-hyps}) that identify the function extracted
from the balance identity with the weighted Chebyshev function
$G_\lev$ of a level (Definition~\ref{def:level}). In
Chapter~\ref{ch:assembly} the balance identity organises the
finite-side calculation. The core bound is half of the analytic
core of the proof of Theorem~\ref{thm:main}: with the gap bound,
Theorem~\ref{thm:gap} of Chapter~\ref{ch:chebyshev}, it gives the
Chebyshev inequality on windows, Theorem~\ref{thm:chebyshev-windows}
(Remark~\ref{rem:content}).

Chebyshev obtained bounds of the right order for $\Ch(x)$ from
Legendre's formula for the prime factorisation of factorials. The
argument of these notes uses a variant of Legendre's count in which
only the primes inert in an imaginary quadratic field
$K=\Q(\sqrt{-d})$ appear, weighted by a function on a CM type. This
chapter shows that such a count equals an explicit main term plus a
combination of twisted Chebyshev functions, and that averaging over
suitable translates isolates exactly the weighted Chebyshev function
$G_\lev$ of a level, with its positive weights.

The identity (Proposition~\ref{prop:balance}) is finite Fourier
analysis on the Galois group $\GG$ of the host field: orthogonality of
characters sorts the prime powers, and the terms linear in the
variable assemble into the main term. Averaging over shifts and over
top places with prescribed shares (Proposition~\ref{prop:shares})
then extracts the characters trivial on a cyclic subgroup, with
weights $1$ and $(3o+1)^{-1}$. The chapter ends with the core bound
(Theorem~\ref{thm:core-bound}): an upper bound for the prime sum by
its main term, up to the error bracket, gives an upper bound
$x^{1/2+O(\eps)}$ for $G_\lev(x)$. Chapters~\ref{ch:theta}
to~\ref{ch:chebyshev} supply such an upper bound for the prime sum.

\section{The additive model and arithmetic data}

\begin{definition}[the additive model]\label{def:additive}
Let $V_4=(\Z/2\Z)^2$, let $H$ be a finite abelian group of order $n$,
and let $\GG=V_4\times H$, written additively. Write an element as
$g=((g_1,g_2),h)$ with $g_1,g_2\in\Z/2\Z$ and $h\in H$.
\begin{enumerate}[label=(\arabic*)]
\item The \emph{CM type} is the subgroup $\Phi=\{g\in\GG:g_1=0\}$, and
$f=|\Phi|=2n$.
\item For $\tau\in H$, the \emph{partner map} is
$\pt_\tau\colon\GG\to\GG$, $\pt_\tau(g)=g+((1,1),\tau)$.
\item The characters $\chi_K$ and $\omega$ of $\GG$ are
$\chi_K(g)=(-1)^{g_1}$ and $\omega(g)=(-1)^{g_2}$, and
$s_0=((1,0),0)$.
\end{enumerate}
\end{definition}

\noindent For the rest of this chapter $H$ is a finite abelian group
of order $n$ and $\GG=V_4\times H$, with the CM type $\Phi$, the
number $f$, the partner maps $\pt_\tau$, the characters $\chi_K$ and
$\omega$ and the element $s_0$ of Definition~\ref{def:additive}.
We write $\theta$ for a character of $\GG$ and $\theta_0$
for the trivial character. For a host (Definition~\ref{def:host}) we
take $H=A$. The multiplicative model $\{\pm1\}^2\times A$ of
Definition~\ref{def:frame} is matched with the additive model
coordinate by coordinate: writing $\bar s\in\Z/2\Z$ for the
class with $s=(-1)^{\bar s}$, the element $((s_1,s_2),h)$ goes to
$((\bar s_1,\bar s_2),h)$. Thus $g_1$ records the action on
$\sqrt{-d}$, so
$\Phi=\Gal(E/K)$ is the CM type induced from $K$ and $\chi_K$ is the
quadratic character of $K$. Under the dictionary,
$\iota\mapsto s_0$, $c\mapsto((1,1),0)$, and the character $\omega_q$
of Definition~\ref{def:host} becomes $\omega$; the conditions
$\theta(\iota)=\pm1$ of Definition~\ref{def:bounded} become
$\theta(s_0)=\pm1$.

\begin{definition}[arithmetic data]\label{def:arithdata}
Let $H$ be a finite abelian group and $\GG=V_4\times H$, with its CM
type $\Phi$ and element $s_0$ (Definition~\ref{def:additive}).
\begin{enumerate}[label=(\arabic*)]
\item \emph{Arithmetic data} $D$ for $\GG$ consist of natural numbers
$q_1$ and $d$ and, for each prime $p$, a subgroup $I_p\subset\GG$ (the
\emph{inertia} at $p$) and an element $\Frob_p\in\GG$.
\item For arithmetic data $D$, a character $\theta$ of $\GG$ and an
integer $n\ge1$, the \emph{averaged value} of $\theta$ at $n$ is
\[
\theta_D(n)=\prod_{p\mid n}\frac{1}{|I_p|}\sum_{z\in I_p}
\theta\bigl(v_p(n)\Frob_p+z\bigr).
\]
\item Let $\ell$ be a function from the characters of $\GG$ to $\C$.
The \emph{arithmetic Chebyshev function} of $\theta$ is
$\GPsi_\theta(y)=\Gz(y)$ if $\theta=\theta_0$, and otherwise
\[
\GPsi_\theta(y)=\sum_{1\le n\le y}\Lambda(n)\,\theta_D(n)
\Bigl(\frac yn-1\Bigr)+y\,\ell(\theta)\qquad(y>0).
\]
\item A prime $p$ is \emph{inert} for $D$ if $p\notin\{2,q_1,d\}$ and
$\Frob_p\notin\Phi$.
\item Let $S$ be a finite set of primes. The data $D$ are \emph{split
outside $d$ on $S$} if $d$ is prime, $2,q_1\in S$, $d\notin S$,
$I_p=0$ for every prime $p\notin S\cup\{d\}$, $I_p\subset\Phi$ and
$\Frob_p\in\Phi$ for every $p\in S$, and $I_d=\{0,s_0\}$.
\item The data $D$ satisfy the \emph{host facts} if $q_1$ and $d$ are
distinct odd primes and $D$ is split outside $d$ on $\{2,q_1\}$.
\item Let $\mathcal H$ be a host (Definition~\ref{def:host}) with
modulus $N=N_{\mathcal H}$, auxiliary primes $q,d$ and Artin map
$\mathrm{Art}\colon(\Z/N\Z)^\times\to\GG$, with $H=A$. The
\emph{arithmetic data of} $\mathcal H$ have $q_1=q$ and the given
$d$. For a prime $p\nmid N$, $I_p=0$ and $\Frob_p=\mathrm{Art}(p)$.
For a prime $p\mid N$, write $(\Z/N\Z)^\times=(\Z/p^{v}\Z)^\times
\times(\Z/N'\Z)^\times$ with $v=v_p(N)$ and $p\nmid N'$, let $I_p$ be
the image under $\mathrm{Art}$ of the first factor (the inertia at
$p$), and let $\Frob_p$ be the image of the class that is $1$ in the
first factor and $p$ in the second.
\end{enumerate}
\end{definition}

\noindent For the arithmetic data of a host, $\theta_D(n)$ is the
value at $n$ of the primitive Dirichlet character $\chi_\theta$
attached to $\theta$: averaging over the inertia group at $p$ passes
from $\theta$ composed with the Artin map to the primitive
character. For a host we take
\[
\ell(\theta)=\frac{L'(1,\chi_\theta)}{L(1,\chi_\theta)}\qquad
(\theta\ne\theta_0),
\]
which is defined because $\chi_\theta\ne1$ for $\theta\ne\theta_0$, so
that $\GPsi_\theta=g_{\chi_\theta}$ (Definition~\ref{def:chebyshev}).
The ties of Definition~\ref{def:tied} impose this choice. For the
identity below, $\ell$ is arbitrary.

\begin{definition}[prime sum, main term and defect]\label{def:ask}
Let $D$ be arithmetic data for $\GG=V_4\times H$
(Definition~\ref{def:arithdata}), let $\ell$ be a function from the
characters of $\GG$ to $\C$, let $N\ge1$ be an integer, $a>0$ and
$\tau\in H$, and let $w\colon\GG\to\R$ be positive on $\Phi$.
\begin{enumerate}[label=(\arabic*)]
\item\label{item:ask-partner}
The \emph{partner weight} is $m_\tau(g)=w(\pt_\tau(g))$ for
$g\in\GG$.
\item The \emph{prime sum} is
\[
\begin{aligned}
\PS_{\tau,w}(a)={}&\sum_{p\le N\text{ inert}}2\log p
\sum_{\substack{1\le i\le N\\ i\Frob_p\notin\Phi}}
\Bigl(\frac{a}{p^i}-m_\tau(i\Frob_p)\Bigr)_+\\
&+\log d\sum_{1\le i\le N}\ \sum_{\substack{g\notin\Phi\\
g-i\Frob_d\in I_d}}\Bigl(\frac{a}{d^i}-m_\tau(g)\Bigr)_+.
\end{aligned}
\]
\item The \emph{main term} is
\[
\MT_w(a)=a\Bigl(\log a-1-\gamma+\ell(\chi_K)-\frac1f
\sum_{\rho\in\Phi}\log w(\rho)\Bigr).
\]
\item The \emph{defect} is
\[
\DF_{\tau,w}(a)=\frac1f\sum_{g\notin\Phi}\sum_\theta
\overline{\theta(g)}\,m_\tau(g)\,\GPsi_\theta\bigl(a/m_\tau(g)\bigr),
\]
the inner sum over all characters $\theta$ of $\GG$.
\end{enumerate}
\end{definition}

\noindent The prime sum is a sum over the primes $p\le N$ inert for
$D$ and over the exponents $i\le N$ with $i\Frob_p\notin\Phi$, plus a
term for the prime $d$; in the case of a host these are the primes
inert in $K$ and, for such a prime, the odd exponents. The weights
$m_\tau(g)$ at the nodes $g\notin\Phi$ are the values of $w$ on $\Phi$,
moved by the partner map (Figure~\ref{fig:partner}).

\begin{proposition}[balance identity]\label{prop:balance}
Let $D$ be arithmetic data for $\GG=V_4\times H$ that are split outside
$d$ on a finite set $S$ of primes, let $\ell$ be a function from the
characters of $\GG$ to $\C$, let $a>0$ and $\tau\in H$, let
$w\colon\GG\to\R$ be positive on $\Phi$, and let $N$ be an integer
with $a/w(\rho)\le N$ for every $\rho\in\Phi$. Then
\[
\PS_{\tau,w}(a)-\MT_w(a)=\DF_{\tau,w}(a).
\]
\end{proposition}

\noindent Written out, the identity says that the Legendre-type count
$\PS_{\tau,w}(a)$ over inert primes equals
$a\bigl(\log a-1-\gamma+\ell(\chi_K)-f^{-1}\sum_{\rho\in\Phi}\log
w(\rho)\bigr)$ plus a combination, with coefficients
$\overline{\theta(g)}\,m_\tau(g)/f$, of the values
$\GPsi_\theta(a/m_\tau(g))$ over all characters $\theta$ and all
$g\notin\Phi$.

\begin{proof}
We expand $\DF_{\tau,w}(a)$ by orthogonality. Its prime-power terms
reassemble into $\PS_{\tau,w}(a)$, and its terms linear in $a$
reassemble into $\MT_w(a)$.

Translation by $((1,1),\tau)$ changes $g_1$,
so $\pt_\tau$ maps $\GG\smallsetminus\Phi$ bijectively onto $\Phi$.
Hence for $g\notin\Phi$ the weight $m_g:=m_\tau(g)$ is positive and
$y_g:=a/m_g\le N$. Also $|\GG\smallsetminus\Phi|=4n-2n=f$.

Next we split $\GPsi$. Since $\theta_{0,D}(n)=1$ for every $n$, we
may write, for every $\theta$ and $y>0$,
\[
\GPsi_\theta(y)=\sum_{n\le y}\Lambda(n)\,\theta_D(n)\Bigl(\frac yn-1
\Bigr)-\kappa_\theta(y),
\]
with $\kappa_{\theta_0}(y)=y\log y-(1+\gamma)y$ and
$\kappa_\theta(y)=-y\,\ell(\theta)$ for $\theta\ne\theta_0$.

Consider the prime powers. If $\Lambda(n)\ne0$, then $n=p^i$ with $i\ge1$,
and $\theta_D(p^i)=|I_p|^{-1}\sum_{z\in I_p}\theta(i\Frob_p+z)$. By
orthogonality, $\sum_\theta\overline{\theta(g)}\theta(h)=|\GG|$ if
$g=h$ and $0$ otherwise, so
\[
\sum_\theta\overline{\theta(g)}\,\theta_D(p^i)=
\begin{cases}|\GG|/|I_p|&\text{if }g\in i\Frob_p+I_p,\\
0&\text{otherwise.}\end{cases}
\]
Also $m_g(y_g/p^i-1)=a/p^i-m_g$, which is nonnegative exactly
when $p^i\le y_g$; and $p^i\le y_g\le N$ forces $p\le N$ and $i\le N$.
So the prime-power part of $\DF_{\tau,w}(a)$ equals
\[
\frac1f\sum_{g\notin\Phi}\ \sum_{\substack{p\le N\\ 1\le i\le N}}
\log p\cdot\frac{|\GG|}{|I_p|}\cdot\mathbf 1[g\in i\Frob_p+I_p]\cdot
\Bigl(\frac a{p^i}-m_g\Bigr)_+.
\]
We sort the primes. If $p\in S$, then $i\Frob_p+I_p\subset\Phi$
contains no $g\notin\Phi$, so $p$ contributes nothing. Also $p$ is not
inert, since $\Frob_p\in\Phi$. If $p\notin S\cup\{d\}$, then $I_p=0$
and the coefficient is $|\GG|/f=2$. The only $g$ is $i\Frob_p$, and
$i\Frob_p\notin\Phi$ for some $i$ only if $\Frob_p\notin\Phi$, that is,
only if $p$ is inert (recall $2,q_1\in S$ and $p\ne d$). These primes
give the first sum in $\PS_{\tau,w}(a)$. If $p=d$, then $|I_d|=2$ and
the coefficient is $(4n/2)/f=1$, which gives the second sum. The three
cases are disjoint, because $2,q_1\in S$ and $d\notin S$. Hence the
prime-power part of $\DF_{\tau,w}(a)$ is $\PS_{\tau,w}(a)$.

It remains to treat the linear terms. The rest of $\DF_{\tau,w}(a)$ is
$-f^{-1}\sum_{g\notin\Phi}\sum_\theta\overline{\theta(g)}\,m_g\,
\kappa_\theta(y_g)$. For $\theta_0$ it equals
\[
-\frac1f\sum_{g\notin\Phi}\Bigl(a\log\frac a{m_g}-(1+\gamma)a\Bigr)
=-a(\log a-1-\gamma)+\frac af\sum_{\rho\in\Phi}\log w(\rho),
\]
by the bijection $\pt_\tau$. For $\theta\ne\theta_0$ it equals
$f^{-1}a\,\ell(\theta)\sum_{g\notin\Phi}\overline{\theta(g)}$. Since
$\sum_{g\in\GG}\overline{\theta(g)}=0$, the last sum is
$-\sum_{g\in\Phi}\overline{\theta(g)}$, which is $-f$ if $\theta$ is
trivial on $\Phi$ and $0$ otherwise. The only nontrivial character
trivial on the index-two subgroup $\Phi$ is $\chi_K$, so these terms
contribute $-a\,\ell(\chi_K)$. Adding up,
$\DF_{\tau,w}(a)=\PS_{\tau,w}(a)-\MT_w(a)$.
\end{proof}

\begin{figure}[htbp]
\centering
\input{alpoge/figs/v2/partner.tex}
\caption{The partner map $\pt_\tau(g)=g+((1,1),\tau)$ on
$\GG=V_4\times H$, here with $H=\Z/3$ and $\tau=1$. It changes the
coordinate $g_1$, so it maps $\GG\smallsetminus\Phi$ bijectively onto
$\Phi$. Hence the partner weights $m_\tau(g)$ at the nodes outside
$\Phi$ run over the weights $w(\rho)$, $\rho\in\Phi$. This is why
$f^{-1}\sum_{\rho\in\Phi}\log w(\rho)$ appears in the main term
$\MT_w$.}
\label{fig:partner}
\end{figure}

\begin{remark}[Legendre's formula]\label{rem:legendre}
Legendre's formula gives $\log\lfloor a\rfloor!=\sum_p\log p
\sum_{i\ge1}\lfloor a/p^i\rfloor$ \cite{Legendre1808}. The prime sum
$\PS_{\tau,w}(a)$ has the smoothed shape
$\sum_p\log p\sum_i(a/p^i-w)_+$ of this count, restricted to the primes
inert in $K$ (for CM by $K$, the supersingular primes) and to the
exponents $i$ with $i\Frob_p\notin\Phi$, which for an inert prime are
the odd ones. Through Proposition~\ref{prop:balance}, an inequality
$\PS\le\MT+(\text{small})$ is therefore a Chebyshev-type upper bound
\cite{Chebyshev1852} for a combination of the functions $\GPsi_\theta$.
\end{remark}

\begin{example}[inert primes for $\Q(\sqrt{-7})$]\label{ex:legendre}
Let $K=\Q(\sqrt{-7})$ and $\chi=\chi_{-7}$. The simplest instance of
the shape of Proposition~\ref{prop:balance}, outside the formalism of
arithmetic data, takes the weight $w\equiv1$:
\[
A(a)=\sum_{n\le a}\Lambda(n)\bigl(1-\chi(n)\bigr)\Bigl(\frac an-1
\Bigr),\qquad
M(a)=a\Bigl(\log a-1-\gamma+\frac{L'(1,\chi)}{L(1,\chi)}\Bigr).
\]
For a prime $p$ inert in $K$, that is, $p\bmod 7\in\{3,5,6\}$, we have
$1-\chi(p^i)=2$ for odd $i$ and $0$ for even $i$. For a split prime,
$1-\chi(p^i)=0$, and $1-\chi(7^i)=1$. So $A(a)$ is the Legendre-type
count over the inert primes and odd powers, at weight $2\log p$, plus
the ramified prime $7$ at weight $\log7$. Directly from
Definition~\ref{def:chebyshev},
\[
A(a)-M(a)=\Gz(a)-g_\chi(a),
\]
a main term plus a combination of two Chebyshev functions, which is
the shape of the balance identity. Here
$L'(1,\chi)/L(1,\chi)=0.0156357\ldots$, and
\[
\begin{array}{rrrr}
a&A(a)&M(a)&(A(a)-M(a))/\sqrt a\\\hline
100&305.081&304.359&0.0722\\
200&748.989&747.347&0.1161\\
1000&5346.323&5346.175&0.0047
\end{array}
\]
Figure~\ref{fig:inert} plots the normalised difference for $10\le
a\le10^5$; it stays small. Figure~\ref{fig:legendre} compares, for
$a=200$, the per-prime terms of Legendre's count with those of the
inert count.
\end{example}

\noindent The computation illustrates the shape of the main term only.
It does not test the inequality of Theorem~\ref{thm:core-bound}, whose
content lies in the size of the error.

\begin{figure}[htbp]
\centering
\input{alpoge/figs/v2/legendre.tex}
\caption{For $a=200$ and $K=\Q(\sqrt{-7})$, per prime $p\le60$: the
smoothed Legendre count $\log p\sum_{p^i\le a}a/p^i$ of
$\log\lfloor a\rfloor!$, and the count of the type of the prime sum
$\PS$ (Definition~\ref{def:ask}), which keeps only the primes inert in
$K$ (those with $p\bmod7\in\{3,5,6\}$) and the odd powers, at weight
$2\log p$. The count uses no weights and no ramified term at $d$.}
\label{fig:legendre}
\end{figure}

\begin{figure}[htbp]
\centering
\input{alpoge/figs/v2/inert.tex}
\caption{For $K=\Q(\sqrt{-7})$ and weight $w\equiv1$:
$A(a)=\sum_{n\le a}\Lambda(n)(1-\chi_{-7}(n))(a/n-1)$, the
Legendre-type sum over the primes inert in $K$ (odd powers) and the
ramified prime, against
$M(a)=a(\log a-1-\gamma+L'(1,\chi_{-7})/L(1,\chi_{-7}))$, where
$L'(1,\chi_{-7})/L(1,\chi_{-7})=0.015635\ldots$. The normalised
difference stays small. The computation illustrates the shape of the
main term $\MT$. It does not test the inequality of
Theorem~\ref{thm:core-bound}.}
\label{fig:inert}
\end{figure}

\section{Thick and thin weights}

The weights $w$ in the prime sum will come from the geometry of
Chapter~\ref{ch:theta}, one weight function for each shift $u_0$ and
each lift index $j<m$. They are of two kinds. At the \emph{thick}
nodes they are close to a large number $b$, so the corresponding
arithmetic Chebyshev functions are evaluated near $a_\tau/b\approx x$.
At the \emph{thin} nodes they are at most $2b'$, far below $b$, so the
functions are evaluated at arguments of size at least $X^\flat/2$,
where the prime number theorem makes them negligible.

\begin{definition}[thick and thin weights]\label{def:thick}
Let $D$ be arithmetic data for $\GG=V_4\times H$, $\ell$ a function
from the characters of $\GG$ to $\C$, $N\ge1$ and $m\ge1$ integers,
$A_0\subset H$ a subgroup, and $\Top,\Top_2\subset H$ finite sets. Let
$a\colon H\to\R$, written $\tau\mapsto a_\tau$, and
$w\colon H\times\{0,\dots,m-1\}\times\GG\to\R$; write $w(u_0,j)$ for
the function $\rho\mapsto w(u_0,j,\rho)$ on $\GG$. Let
$x,b,b',\eta,X^\flat,\vartheta$ be real numbers.
\begin{enumerate}[label=(\arabic*)]
\item A node $\rho\in\Phi$ is \emph{thick at} $u_0\in H$ if
$\rho=((0,1),t+u_0)$ for some $t\in\Top_2$, and \emph{thin at} $u_0$
otherwise.
\item The \emph{size conditions} hold if $\eta\ge0$, $X^\flat>0$,
$b'>0$, $b'X^\flat\le xb$ and $\vartheta\ge0$; if $|a_\tau/b-x|\le1$
for every $\tau\in\Top$; and if for every $u_0\in A_0$, $j<m$ and
$\rho\in\Phi$:
\begin{enumerate}[label=(\alph*)]
\item $w(u_0,j,\rho)>0$;
\item $b\le w(u_0,j,\rho)\le b+\eta+1$ if $\rho$ is thick at $u_0$;
\item $w(u_0,j,\rho)\le2b'$ if $\rho$ is thin at $u_0$;
\item $a_\tau/w(u_0,j,\rho)\le N$ for every $\tau\in\Top$.
\end{enumerate}
\item The \emph{thin-scale bound} holds if
$|\GPsi_\theta(y)|\le\vartheta y$ for every character $\theta$ of
$\GG$ and every $y\ge X^\flat/3$.
\item For $u_0\in H$ such that $w(u_0,j)$ is positive on $\Phi$ for
every $j<m$, the \emph{gap} at $u_0$ is
\[
\gap(u_0)=\sum_{\tau\in\Top}\sum_{j<m}\Bigl(
\PS_{\tau,w(u_0,j)}(a_\tau)-\operatorname{Re}\MT_{w(u_0,j)}(a_\tau)
\Bigr),
\]
with $\PS$ and $\MT$ the prime sum and the main term of
Definition~\ref{def:ask} for $D$, $\ell$ and $N$.
\end{enumerate}
\end{definition}

\noindent At the sizes used later, $b=\exp(x^\eps)$, $a_\tau\approx
xb$, $X^\flat=\exp(x^{\eps/3})$ and $b'=a/X^\flat$
(Definition~\ref{def:scales}). A thick weight therefore tests
$\GPsi_\theta$ at $a_\tau/w\approx x$, and a thin weight tests it at
$a_\tau/w\ge X^\flat/2$ approximately, where the thin-scale bound, an
estimate of prime number theorem type, applies (Figure~\ref{fig:thick-thin}). The bound is required
from $X^\flat/3$ on because the exact lower bound is
$a_\tau/w\ge(x-1)b/(2b')$, and the factor $(x-1)/x$ that this loses
against $X^\flat/2$ is absorbed by passing from $1/2$ to $1/3$.
\begin{figure}[tbp]
\centering
\input{book/figs/ch06-thick-thin.tex}
\end{figure}

\begin{definition}[kept characters and extracted functions]
\label{def:kept-chars}
Let $D$ be arithmetic data for $\GG=V_4\times H$, let $\ell$ be a
function from the characters of $\GG$ to $\C$, let $A_0\subset H$ be a
subgroup, $\Top\subset H$ a finite set and $x>0$. Let $\Psi(A_0)$ be
the set of characters of $\GG$ trivial on $0\times A_0$, let $\Xi$ be
a set of characters of $\GG$ (the \emph{kept set}), and let
$\varkappa,M,C_o,C_E,\eps$ be real numbers.
\begin{enumerate}[label=(\arabic*)]
\item The \emph{weight} of a character $\theta$ is $\wt_A(\theta)=1$ if
$\theta$ is trivial on $0\times H$ and $\wt_A(\theta)=\varkappa$ otherwise.
\item The \emph{extracted function} and the \emph{rest} are
\[
G_A(y)=\sum_{\theta\in\Xi}\wt_A(\theta)\operatorname{Re}
\GPsi_\theta(y),\qquad
\Sigma_A(y)=\sum_{\theta\in\Psi(A_0)\smallsetminus\Xi}
|\GPsi_\theta(y)|.
\]
\item The \emph{character conditions} hold if: the trivial character
lies in $\Xi$; every $\theta\in\Xi$ is trivial on $0\times A_0$ and
satisfies $\theta(s_0)=1$; $|\Psi(A_0)|\le M$; for every $\theta\in\Xi$
nontrivial on $0\times H$,
\[
\Bigl|\sum_{\tau\in\Top}\theta(0,\tau)-\varkappa\,|\Top|\Bigr|\le
Mx^{-1+\eps}|\Top|;
\]
$|\ell(\theta)|\le C_o$ for every nontrivial $\theta\in\Xi$; and
$|\ell(\theta)|\le C_E\log x$ for every
$\theta\in\Psi(A_0)\smallsetminus\Xi$.
\end{enumerate}
\end{definition}

\begin{definition}[rest sizes, costs and the error bracket]
\label{def:restsize}
Let $\eps,x,b,k,a_E,\eta,\vartheta,f$ be real numbers and
$n,n_T,n_{T_2},m$ natural numbers with $n\ge1$.
\begin{enumerate}[label=(\arabic*)]
\item The \emph{rest-size conditions} hold if $a_E\le2xb$,
$x^{-\eps/8}n/2\le n_T$, $x^{-\eps}\le n_{T_2}/n$, $x^{1/2-\eps}\le m$
and $k\le n^2x^{1/2-\eps}$.
\item The \emph{cost conditions} hold if $2f(x+1)\vartheta\le1/x$ and
$6f(x+1)(\eta+1)\le b/x$.
\item The \emph{error bracket} is
\[
\Err=n_Ta_E\Bigl(1+\frac{m^2}{k}\Bigr)+m\Bigl(2nmbx^\eps+
\frac{n_{T_2}kbx^\eps}{n^2x^{-\eps}}\Bigr).
\]
\end{enumerate}
\end{definition}

\noindent Proposition~\ref{prop:sqrtx} shows that, under the rest-size
conditions, $xm\le k$ and $m\le x^{1/2+\eps/2}$, the bracket divided by
$n_Tmb$ is of size $x^{1/2+O(\eps)}$. The square root in the final
bound comes from that computation.

\section{Averaging over shifts and top places}

The kept characters must be weighted by the weights of a level. The
weights come from the distribution of the top places among the cosets
of $A_0$ for a frame $\phi$ (Definition~\ref{def:frame}).

\begin{proposition}[shares give the weights]\label{prop:shares}
Let $\phi$ be a frame with group $A$, core $A_0$ and index
$o=[A:A_0]$, let $\delta$ be real, let $\Top\subset A$ be a finite set
with shares up to $\delta$ (Definition~\ref{def:shares-ef}), and let
$\theta$ be a character of $A$ trivial on $A_0$. Put
$\varkappa_\theta=1$ if $\theta$ is trivial and
$\varkappa_\theta=(3o+1)^{-1}$ otherwise. Then
\[
\Bigl|\sum_{\tau\in\Top}\theta(\tau)-\varkappa_\theta|\Top|\Bigr|\le
o\,\delta\,|\Top|.
\]
\end{proposition}

\begin{proof}
Choose representatives $r_1=0,r_2,\dots,r_o$ of the cosets of $A_0$ in
$A$. Since $\theta$ is constant on cosets,
\[
\sum_{\tau\in\Top}\theta(\tau)=\sum_{k=1}^o\theta(r_k)\,
\#\{\tau\in\Top:\tau-r_k\in A_0\}.
\]
By the shares hypothesis and $|\theta|=1$, this differs from
$|\Top|\sum_k\theta(r_k)\lambda_\phi(r_k)$ by at most
$o\,\delta\,|\Top|$. Now
$\sum_k\theta(r_k)\lambda_\phi(r_k)=\bigl(4+3\sum_{k\ge2}
\theta(r_k)\bigr)/(3o+1)$. If $\theta$ is trivial, this is
$(4+3(o-1))/(3o+1)=1$. Otherwise $\theta$ is a nontrivial character of
$A/A_0$, so $\sum_{k=1}^o\theta(r_k)=0$ and
$\sum_{k\ge2}\theta(r_k)=-1$, which gives $(4-3)/(3o+1)$.
\end{proof}

\noindent With $\delta=x^{-1+\eps}$, Proposition~\ref{prop:shares} is
the share clause of the character conditions
(Definition~\ref{def:kept-chars}) with $\varkappa=(3o+1)^{-1}$ and any
$M\ge o$. The weights of a level (Definition~\ref{def:level}) come
from this proposition, and so does their positivity, which the
windowed Landau criterion needs (Remark~\ref{rem:positivity}).
Figure~\ref{fig:shares} shows the shares
and the averages for $o=3$.

\begin{figure}[htbp]
\centering
\input{alpoge/figs/v2/shares.tex}
\caption{(a) The shares $\lambda_\phi$ of
Definition~\ref{def:shares-ef}: $4/(3o+1)$ on $A_0$ and $3/(3o+1)$ on
each other coset, here with $o=3$. (b) The averages over $\Top$ of the
characters of $A/A_0$, for top places with these shares: $1$ for the
trivial character and $(3o+1)^{-1}$ for each nontrivial one, up to
$o\,\delta$ (Proposition~\ref{prop:shares}). These are the positive
weights of a level (Definition~\ref{def:level}).}
\label{fig:shares}
\end{figure}

\begin{proposition}[the extracted function is $G_\lev$]
\label{prop:averaged}
Let $H$ be a finite abelian group of order $n$, $A_0\subset H$ a
subgroup of order $n_0$, $D$ arithmetic data for $\GG=V_4\times H$,
$\ell$ a function from the characters of $\GG$ to $\C$, and $\Xi$ a set
of characters of $\GG$. Let $\lev$ be a level with no exceptional
character and with $o(\lev)=o=[H:A_0]$, let $\varkappa=(3o+1)^{-1}$, and form
$\Psi(A_0)$, $G_A$ and $\Sigma_A$ (Definition~\ref{def:kept-chars}).
Suppose that $\Xi$ contains the trivial character and that each
nontrivial $\theta\in\Xi$ is nontrivial on $0\times H$ and is given a
primitive character $\chi_\theta$, such that $\theta\mapsto\chi_\theta$
is a bijection from the nontrivial elements of $\Xi$ onto
$\Kept(\lev)\smallsetminus\{1\}$ and $\GPsi_\theta(y)=
g_{\chi_\theta}(y)$ for $y\ge1$.
\begin{enumerate}[label=(\alph*)]
\item\label{item:averaged-eq} $G_A(y)=G_\lev(y)$ for every $y\ge1$.
\item\label{item:averaged-rest} Let $x>0$ and $\eps\in\R$. If
$|\GPsi_\theta(x)|\le x^{1/2+\eps}$ for every
$\theta\in\Psi(A_0)\smallsetminus\Xi$, then
$\Sigma_A(x)\le4o\,x^{1/2+\eps}$.
\end{enumerate}
\end{proposition}

\begin{proof}[Proof of \ref{item:averaged-eq}]
The trivial character lies in $\Xi$, has weight $1$ and contributes
$\operatorname{Re}\GPsi_{\theta_0}=\Gz$. Every other
$\theta\in\Xi$ is nontrivial on $0\times H$, so it has weight $\varkappa$ and
contributes $\varkappa\operatorname{Re}g_{\chi_\theta}$. Since $\lev$ has no
exceptional character, all its weights are $(3o+1)^{-1}=\varkappa$
(Definition~\ref{def:level}), and the bijection matches these terms
with the terms of $G_\lev$.
\end{proof}

\begin{proof}[Proof of \ref{item:averaged-rest}]
The set $\Psi(A_0)$ is the dual of $\GG/(0\times A_0)$, which has order
$4n/n_0=4o$. So $\Sigma_A(x)$ has at most $4o$ terms, each at most
$x^{1/2+\eps}$.
\end{proof}

\noindent The hypothesis that $\Xi$ contains the trivial character is
the first of the character conditions; without it $G_A$ would lack the
term $\Gz$.

\begin{theorem}[Alp\"oge]\label{thm:averaged-hyps}
Let $\lev$ be a level, $\phi$ a frame, and take a host for $\phi$
(Definition~\ref{def:host}) with its arithmetic data $D$ for
$\GG=V_4\times A$. Let $\Top\subset A$ be finite and $x>0$, $\eps$ and
$C$ real, and suppose that the host satisfies the explicit-formula
condition $\mathrm{EF}(\lev,\Top,x,\eps,C)$
(Definition~\ref{def:shares-ef}). Let $\ell$ be a function from the
characters of $\GG$ to $\C$ such that $\ell(\theta)=L'(1,\chi)/L(1,\chi)$
whenever $\theta\ne\theta_0$ and $\chi$ is a primitive Dirichlet
character with $\chi(n)=\theta_D(n)$ for all $n\ge1$. Let $\Xi$ be the
set of characters $\theta$ of $\GG$ for which some primitive Dirichlet
character $\chi$ with $\chi(n)=\theta_D(n)$ for all $n\ge1$ lies in
$\Kept(\lev)$, let $A_0$ be the core of $\phi$, and let $G_A$ and
$\Sigma_A$ be the extracted function and the rest
(Definition~\ref{def:kept-chars}) for $D$, $\ell$, $A_0$, $\Xi$ and
$\varkappa=(3o(\lev)+1)^{-1}$.
Then:
\begin{enumerate}[label=(\alph*)]
\item\label{item:ah-eq} $G_A(y)=G_\lev(y)$ for every $y>0$;
\item\label{item:ah-rest} $\Sigma_A(x)\le4o(\lev)\,x^{1/2+\eps}$.
\end{enumerate}
\end{theorem}

\noindent Theorem~\ref{thm:averaged-hyps} says that for a host the
hypotheses of Proposition~\ref{prop:averaged} are met, so that the
extracted function of the host is the weighted Chebyshev function of
the level and the rest is of square-root size.
\begin{figure}[tbp]
\centering
\input{book/figs/ch06-kept-grid.tex}
\end{figure}

\begin{proof}[Idea of proof]
We indicate why the statement holds, from the definitions. Both
conclusions are term-by-term comparisons:
steps~\ref{item:ah-s1}--\ref{item:ah-s3} give
\ref{item:ah-eq}, and step~\ref{item:ah-s4} gives \ref{item:ah-rest}.
\begin{enumerate}[label=(\arabic*)]
\item\label{item:ah-s1}
The condition on $\ell$, together with the identification of
$\theta_D(n)$ with the value at $n$ of the primitive character attached
to $\theta$, gives $\GPsi_\theta=g_{\chi_\theta}$ for $\theta\ne\theta_0$.
\item The Artin map of a host is surjective, so $\theta$ is determined by
$\chi_\theta$. The explicit-formula condition asks that every kept
character be $\chi_\theta$ for some $\theta\in X_+$; hence
$\theta\mapsto\chi_\theta$ is a bijection from $\Xi$ onto $\Kept(\lev)$,
and $\Xi\subset X_+\subset\Psi(A_0)$.
\item\label{item:ah-s3}
The characters of $\GG$ trivial on $0\times A$ are $\theta_0$,
$\omega$, $\chi_K$ and $\chi_K\omega$, and only $\theta_0$ and $\omega$
take the value $1$ at $s_0$ (Figure~\ref{fig:kept-grid}). So the kept characters of weight $1$ are
$\theta_0$ and, when the level has an exceptional character, $\omega$,
with $\chi_\omega=\chi_{\omega_q}=\psi_{\mathrm{exc}}$ because the host
keeps the exceptional character. This matches the weights of
Definition~\ref{def:level}, and $G_A=G_\lev$.
\item\label{item:ah-s4} A character trivial on $A_0$ and not kept is bounded
(Definition~\ref{def:bounded}), so the explicit-formula condition gives
$|g_{\chi_\theta}(x)|\le x^{1/2+\eps}$ for it. Since $o(\phi)=o(\lev)$,
the count of Proposition~\ref{prop:averaged}\ref{item:averaged-rest}
gives $\Sigma_A(x)\le4o(\lev)\,x^{1/2+\eps}$.
\end{enumerate}
\end{proof}

The core bound converts an upper bound for the gap into an upper bound
for the extracted function. By Proposition~\ref{prop:balance}, for
weights positive on $\Phi$ and $N$ large enough,
\[
\gap(u_0)=\sum_{\tau\in\Top}\sum_{j<m}\operatorname{Re}
\DF_{\tau,w(u_0,j)}(a_\tau),
\]
since $\PS$ is real. So an upper bound for the gap is an upper bound
for a weighted combination of the functions $\GPsi_\theta$, evaluated
at the points $a_\tau/w(u_0,j,\rho)$.

\begin{theorem}[Alp\"oge; core Chebyshev bound]\label{thm:core-bound}
Let $C\ge0$, $\eps>0$, $\varkappa$ with $|\varkappa|\le1$, and $C_E,C_o,M$ be
real
numbers. There is $x_1$ such that the following holds for every
$x\ge x_1$. Let $H$ be a finite abelian group, $\GG=V_4\times H$, $m$ a
natural number, $D$ arithmetic data for $\GG$ that are split outside
$d$ on a finite set of primes, $\ell$ a function from the characters
of $\GG$ to $\C$, $N\ge1$ an integer, $A_0\subset H$ a subgroup,
$\Top,\Top_2\subset H$ finite sets, $a\colon H\to\R$,
$w\colon H\times\{0,\dots,m-1\}\times\GG\to\R$,
$b,b',\eta,X^\flat,\vartheta,k,a_E$ real numbers and $\Xi$ a set of
characters of $\GG$. Let $G_A$ and $\Sigma_A$ be the extracted
function and the rest (Definition~\ref{def:kept-chars}) for $D$,
$\ell$, $A_0$, $\Xi$ and $\varkappa$. Suppose that
\begin{enumerate}[label=(\roman*)]
\item $b>0$, $a_E\ge0$, $k>0$, $m\le x^{1/2+\eps/2}$ and $xm\le k$;
\item the size conditions and the thin-scale bound of
Definition~\ref{def:thick} hold, and $\Top_2\subset A_0$;
\item the cost conditions hold with $f=|\Phi|$, and the rest-size
conditions hold with $n=|H|$, $n_T=|\Top|$ and $n_{T_2}=|\Top_2|$
(Definition~\ref{def:restsize});
\item the character conditions of Definition~\ref{def:kept-chars} hold
with the given $\varkappa,M,C_o,C_E,\eps$ and $x$;
\item $\gap(u_0)\le x^{C\eps}\,\Err$ for every $u_0\in A_0$, with
$\Err$ the error bracket of Definition~\ref{def:restsize} for these
$n,n_T,n_{T_2},m,b,k,a_E$.
\end{enumerate}
Then $G_A(x)\le x^{1/2+(C+3)\eps}+\Sigma_A(x)$.
\end{theorem}

\begin{proof}[Proof (outline)]
The argument has five elementary steps (Figure~\ref{fig:core-steps}). Its only analytic input is
the classical prime number theorem behind the thin-scale bound.
\begin{enumerate}[label=(\arabic*)]
\item Write $\gap(u_0)$ through the balance identity, as displayed
above. The thin nodes contribute values of $\GPsi_\theta$ at
arguments at least $X^\flat/3$, which the thin-scale bound and the
cost conditions make negligible.
\item At the thick nodes the weight is within $\eta+1$ of $b$, and
$|a_\tau/b-x|\le1$. Moving each argument $a_\tau/w$ to exactly $x$
costs $O(\log x+|\ell(\theta)|)$ per term, by the Lipschitz behaviour
of $\GPsi_\theta$; the character conditions bound $|\ell(\theta)|$.
\item Averaging over the shifts $u_0\in A_0$ projects onto the $4o$
characters of $\GG$ trivial on $0\times A_0$. The kept characters
enter with the right sign: a node $g\notin\Phi$ whose partner
$\pt_\tau(g)=((0,1),t+u_0)$ is thick satisfies
$g=s_0+((0,0),t+u_0-\tau)$, and $\theta(s_0)=1$ for $\theta\in\Xi$
(Definition~\ref{def:kept-chars}). So the coefficient
$\overline{\theta(g)}$ of $\GPsi_\theta$ in the gap is
$\overline{\theta((0,0),t+u_0-\tau)}$, and after the averages over
$u_0$ and $\tau$ the kept characters carry the positive weights that
Remark~\ref{rem:positivity} requires. This is why an \emph{upper}
bound for the gap gives an \emph{upper} bound for $G_A$.
\item Summing over $\tau\in\Top$ produces the weights $1$ and $\varkappa$, by
the share clause of the character conditions
(compare Proposition~\ref{prop:shares}); the characters outside $\Xi$
are collected in $\Sigma_A(x)$.
\item Divide by $|\Top|\,m\,b\,\delta_2/2$, where
$\delta_2=|\Top_2|/n$ is the proportion of thick shifts. By
Proposition~\ref{prop:sqrtx} and the rest-size conditions, the
bracket $x^{C\eps}\Err$ then becomes at most $x^{1/2+(C+3)\eps}$ for
large $x$.
\end{enumerate}
\end{proof}

\noindent Theorem~\ref{thm:core-bound}, combined with
Theorem~\ref{thm:averaged-hyps}, turns a gap bound into the bound
$G_\lev(x)\le x^{1/2+(C+3)\eps}+4o\,x^{1/2+\eps}$. The gap bound is
Theorem~\ref{thm:gap}, and the combination is
Theorem~\ref{thm:chebyshev-windows}.
\begin{figure}[tbp]
\centering
\input{book/figs/ch06-core-steps.tex}
\end{figure}

\FloatBarrier
\clearpage
\clearpage
\chapter{Hermitian lattices, the theta tower and slope records}
\label{ch:theta}\label{sec:slopes}

This chapter sets up the geometric side of the argument: the tower of
hermitian spaces and the theta chain over a host field, hermitian
lattices and their slopes, packages, and the slope record, a list of
real numbers and inequalities through which the geometry enters the
analytic chapters. Its main theorem is the frame theorem,
Theorem~\ref{thm:frame}, a slope comparison between many source
lattices and one target lattice. The slope clause of the size data
(Definition~\ref{def:size-data}\ref{item:sd-slope}) has the shape of
its conclusion. The upper and lower links
(Definitions~\ref{def:upper-link} and~\ref{def:lower-link}) record
the archimedean and finite sides. In Chapter~\ref{ch:assembly} the
frame theorem is one of the inputs that Remark~\ref{rem:content}
lists as substantive. The slope record and the links are used in
Chapters~\ref{ch:covolume}, \ref{ch:chebyshev}
and~\ref{ch:families}, where packages are constructed and tied to
slope records.

Chebyshev's proof of $\Ch(x)\asymp x$ compares two evaluations of
$\log n!$: the archimedean one, by Stirling's formula, and the $p$-adic
ones, by Legendre's formula. Alpöge's argument replaces $n!$ by the
periods of theta lifts of algebraic automorphic forms of huge weight
$a\approx xb$. On the archimedean side the periods carry Stirling's term
$a\log a-a$ and the $a$-th power of a CM period $\Omega$, whose logarithm
contains $L'(1,\chi_{-d})/L(1,\chi_{-d})$. On the finite side they carry
a Legendre-type gain at every prime inert in $K$. A slope inequality for
the period map between two hermitian lattices links the two sides, and a
chain of $m\approx\sqrt x$ lifts makes the total error of relative size
about $x^{-1/2}$.

The geometric objects enter the analytic argument of
Chapters~\ref{ch:balance} and~\ref{ch:chebyshev} only through the
slope record. The lattices reappear in Chapter~\ref{ch:covolume},
where the arithmetic geometry is concentrated.

\section{The tower of unitary groups and the theta chain}

\begin{definition}[ranks]\label{def:ranks}
Let a frame be given (Definition~\ref{def:frame}), whose group has
order $\operatorname{ord}$.
\begin{enumerate}[label=(\arabic*)]
\item\label{item:rk-datum} A \emph{ranks datum} for the frame is a triple
$\mathbf k=(k,m,t)$ of integers with $k$ even, $m\ge2$ even and $t$
odd, such that
\[
k\ge m+3,\quad t\ge k+m+4,\quad t>3m,\quad
m\le2\operatorname{ord}\le k,\quad m^2\le k+m(t+1).
\]
\item The \emph{derived quantities} of $\mathbf k$ are $n=
\operatorname{ord}$, $f=2n$, $r=k+m-1$, $n_V=r+t$, $r_S=m/2-1$ and
$N(j)=k+j(t+1)$ for $j\ge0$. The \emph{signature at level} $j$ is
$(p_j,q_j)=(k+j-1,m-j)$.
\end{enumerate}
\end{definition}

\noindent The last two conditions of~\ref{item:rk-datum} read $m\le f\le k$ and
$m^2\le N(m)$. Since $k\ge m+3$, the difference $p_j-q_j=k-1-m+2j$ is
positive for every $j\ge0$.

\begin{definition}[the spaces of the tower]\label{def:spaces}
Let $E$ be a host field (Definition~\ref{def:host}) with maximal
totally real subfield $F$, and let $\mathbf k=(k,m,t)$ be a ranks
datum with derived quantities $r,n_V,r_S,N(j)$
(Definition~\ref{def:ranks}). The \emph{spaces} of $\mathbf k$ over
$E/F$ consist of the following hermitian spaces over $E/F$, each with a
unimodular base lattice.
\begin{enumerate}[label=(\arabic*)]
\item For $0\le j\le m$, a definite space $X_j$ of rank $N(j)$.
\item A space $X'$ of rank $r+mt$, an isometric embedding
$X'\hookrightarrow X_m$, and an anisotropic vector $z\in X_m$
orthogonal to its image.
\item\label{item:sp-y}
For $0\le j\le m$, a space $Y_j$ of rank $r+jt$, with isometric
embeddings $X_j\hookrightarrow Y_j$ for $j<m$, an isometry
$X'\cong Y_m$, and a decomposition
$Y_0=X_0\oplus\langle-T\rangle\oplus\mathbb H^{r_S}$, where $\mathbb H$
is a hyperbolic plane and $T$ is a prime.
\item For $1\le j\le m$, a space $W_j$ of rank $r$ and signature
$(p_j,q_j)=(k+j-1,m-j)$ at every archimedean place, locally isometric
to $Y_0$ at every finite place.
\item For $1\le j\le m$, a space $V_j$ of rank $n_V$ carrying a lattice
with an orthonormal basis up to sign.
\end{enumerate}
\end{definition}

\noindent The hermitian forms, the lattices and the isometries are part
of the data. We do not write them out. Nothing below depends on more
than the ranks and signatures listed. The spaces $X_j$ are definite, so
their unitary groups $U(X_j)$ carry algebraic automorphic forms in the
sense of Gross~\cite{Gross1999}. The spaces $W_j$ are indefinite for
$j<m$, and their unitary groups carry the Shimura varieties of
Chapter~\ref{ch:covolume}.

\begin{remark}[the top of the tower]\label{rem:topspace}
The embeddings $X_j\hookrightarrow Y_j$ of
Definition~\ref{def:spaces}\ref{item:sp-y} stop below the top of the
tower: there is no embedding $X_m\hookrightarrow Y_m$. The ranks
satisfy
\[
\rk Y_j-\rk X_j=(r+jt)-N(j)=m-1-j.
\]
An isometric embedding of a nondegenerate space is injective, so such
an embedding exists only for $j\le m-1$; at $j=m-1$ the two ranks are
equal. At the top the direction reverses: $Y_m\cong X'$ embeds in
$X_m$, whose rank $N(m)=\rk Y_m+1$ is one larger. By the rank count,
$X_m$ is the orthogonal sum of the image of $X'$ and the line spanned
by the anisotropic vector $z$.
\end{remark}

\noindent Figure~\ref{fig:tower} shows the spaces for small ranks.

\begin{figure}[htbp]
\centering
\input{alpoge/figs/v2/tower.tex}
\caption{The spaces of Definition~\ref{def:spaces} for the small ranks
$(k,m,t)=(6,2,13)$. They satisfy the conditions of
Definition~\ref{def:ranks}: $k$ and $m$ are even, $t$ is odd,
$m+3\le k$, $k+m+4\le t$, $3m<t$, $m^2\le k+m(t+1)$, and
$m\le2\,\mathrm{ord}\le k$ for $\mathrm{ord}\in\{1,2,3\}$. Here
$r=k+m-1=7$, $N(j)=k+j(t+1)$, $\rk Y_j=r+jt$ and $n_V=r+t$. The space
$W_j$ has signature $(k+j-1,m-j)$ and is definite at $j=m$; $Y_0$ has
signature $(6,1)$. The embeddings $X_j\hookrightarrow Y_j$ exist only
for $j<m$ (Remark~\ref{rem:topspace}); at $j=m$ the space $X'\cong
Y_m$ embeds in $X_m$, whose rank is one larger. The arrows $\Theta_j$
are the theta kernels of Definition~\ref{def:column}.}
\label{fig:tower}
\end{figure}

\begin{definition}[forms, theta kernels and the column; schematic]
\label{def:column}
Let $\mathbf k=(k,m,t)$ be a ranks datum and fix spaces of
$\mathbf k$ over $E/F$ (Definition~\ref{def:spaces}). For $0\le i\le
m$ let $\operatorname{cls}_i$ be the class set of the genus of a
unimodular lattice in $X_i$, weighted by the masses
$1/\#\operatorname{Stab}$.
\begin{enumerate}[label=(\arabic*)]
\item A \emph{form at level} $i$ is a function
$\operatorname{cls}_i\times Y_i\to\C$.
\item The \emph{theta kernel} $\Theta_i$, for $0\le i<m$, is the
linear map from forms at level $i$ to forms at level $i+1$ given by
\[
(\Theta_i\phi)(\Lambda')=\sum_{\Lambda\in\operatorname{cls}_i}
\frac{\upsilon_i(\Lambda)\,\upsilon'_i(\Lambda')}
{\#\operatorname{Stab}(\Lambda)}\,K_i(\Lambda,\Lambda')\bigl(
\phi(\Lambda)\bigr)\qquad(\Lambda'\in\operatorname{cls}_{i+1}),
\]
where $\upsilon_i,\upsilon'_i$ are functions of absolute value $1$ on
the class sets (the splitting characters), and $K_i(\Lambda,\Lambda')$
is a linear map from functions on $Y_i$ to functions on $Y_{i+1}$,
the theta kernel between the lattices $\Lambda$ and $\Lambda'$,
composed on both sides with multiplication by fixed nonvanishing
functions (the archimedean normalisations).
\item The \emph{column} is the $m$-fold composite
$\operatorname{col}=\Theta_{m-1}\circ\dots\circ\Theta_0$, the
\emph{theta chain} from $U(X_0)$ to $U(X_m)$.
\item Let $\mathrm L_z$ be the map that evaluates a form at level $m$
at the anisotropic vector $z\in X_m$ of the spaces and then applies a
fixed linear lowering map to a space of row vectors, and let
$\langle\ ,\ \rangle$ be a fixed bilinear pairing of row vectors with
test forms. The \emph{period} of a form $\phi$ at level $m$ against a
test form $F$ is $P(\phi,F)=\langle\mathrm L_z\phi,F\rangle$. The
\emph{period map} of a form $\phi$ at level $0$ is
$F\mapsto P(\operatorname{col}\phi,F)$.
\end{enumerate}
\end{definition}

\noindent Thus $\Theta_i$ lifts forms on $U(X_i)$ to forms on
$U(X_{i+1})$, and the period map pairs the lift of a form on $U(X_0)$
to the top of the tower with test forms. In terms of Arthur parameters
each step of the chain may be pictured as
$\pi\mapsto\pi\boxplus\mu[t+1]$. This picture is only a guide. The
kernels $K_i$, the lowering map and the pairing are not written out
(Section~\ref{sec:conventions}); the argument uses the column only
through the period maps $B_u$ of Definition~\ref{def:period-lattice}.

\section{Hermitian lattices and slopes}

\begin{definition}[hermitian lattices, slopes and heights]
\label{def:hlat}
Let $L$ be a number field.
\begin{enumerate}[label=(\arabic*)]
\item A \emph{hermitian vector bundle} $\overline{\mathcal E}$ over
$\operatorname{Spec}\mathcal O_L$ is a finitely generated projective
$\mathcal O_L$-module $\mathcal E$ with a hermitian norm on
$\mathcal E\otimes_\sigma\C$ for each embedding $\sigma\colon
L\hookrightarrow\C$, the norms at $\sigma$ and $\bar\sigma$ being
conjugate.
\item A \emph{hermitian lattice} over $\mathcal O_L$ is a
finite-dimensional $L$-vector space with a basis, a full-rank
$\mathcal O_L$-submodule, and a positive hermitian Gram matrix at each
infinite place. It is a hermitian vector bundle over
$\operatorname{Spec}\mathcal O_L$.
\item The \emph{normalised Arakelov degree} is
$\widehat{\deg}_n=\widehat{\deg}/[L:\Q]$. The \emph{slope} of
$\overline{\mathcal E}\ne0$ is $\hat\mu(\overline{\mathcal E})=
\widehat{\deg}_n(\overline{\mathcal E})/\rk\mathcal E$, and its
\emph{maximal} and \emph{minimal slopes}
$\hat\mu_{\max}(\overline{\mathcal E})$ and
$\hat\mu_{\min}(\overline{\mathcal E})$ are those of its
Harder--Narasimhan filtration~\cite{Stuhler1976,Grayson1984}. For a
hermitian lattice $N$ we write $\mu(N)$, $\mu_{\max}(N)$,
$\mu_{\min}(N)$, and $N^\vee$ for the dual lattice.
\item\label{item:hl-height}
For hermitian vector bundles $\overline{\mathcal E}$,
$\overline{\mathcal F}$ and an $L$-linear map $\varphi\colon\mathcal
E_L\to\mathcal F_L$, the \emph{height} of $\varphi$ is
\[
h(\overline{\mathcal E},\overline{\mathcal F},\varphi)=\frac1{[L:\Q]}
\Bigl(\sum_{\mathfrak p}\log\|\varphi\|_{\mathfrak p}
+\sum_{\sigma\colon L\hookrightarrow\C}\log\|\varphi\|_\sigma\Bigr),
\]
the first sum over the finite places of $L$. For hermitian lattices
$N,R$ and $\varphi\ne0$ we write $\operatorname{ht}(\varphi)$.
\end{enumerate}
\end{definition}

\begin{definition}[period lattice data]\label{def:period-lattice}
Let $E$ be a host field and $L\supset E$ a CM field, and let
$\Words$ be a finite nonempty set, whose elements we call
\emph{words}.
\begin{enumerate}[label=(\arabic*)]
\item \emph{Period lattice data} over $L$ indexed by $\Words$ consist
of a hermitian lattice $R$ over $\mathcal O_L$ (the \emph{test forms});
for each $u\in\Words$ a hermitian lattice $N_u$ (the \emph{ramified
forms}), an $L$-linear map $B_u\colon N_u\otimes L\to R^\vee\otimes L$
(the \emph{period map}) and a real number $h_u$; and a real number
$\Delta_R$ (the \emph{spread}). They are required to satisfy
$B_u(N_u)\subset R^\vee$, and $\operatorname{ht}(B_u)\le h_u$ whenever
$B_u\ne0$ (Definition~\ref{def:hlat}\ref{item:hl-height}).
\item The \emph{frame slope} of the data is
\[
s=\frac{\sum_{u\in\Words}\rk(N_u)\bigl(\mu(N_u)-h_u\bigr)}
{\sum_{u\in\Words}\rk(N_u)}.
\]
\end{enumerate}
\end{definition}

\noindent The maps $B_u$ are not required to be injective.

\begin{definition}[packages; schematic]\label{def:package}
A \emph{package} $\mathcal P$ consists of the following data.
\begin{enumerate}[label=(\arabic*)]
\item A ranks datum $\mathbf k=(k,m,t)$
(Definition~\ref{def:ranks}) and a host
(Definition~\ref{def:host}) with $\sqrt{-q}\in E$, together with:
a \emph{shift} $u_0\in A$; size bounds; small primes split in $K$; an ideal
$\mathfrak b$ with
$\mathfrak d_{E/K}=\mathfrak b\mathfrak b^c$; the norm primes
$\varpi_0,\varpi_1$ of the host (Definition~\ref{def:norm-primes});
an ideal of $\mathcal O_K$ supported
on split primes; the top places $\Top$ with top sizes $a_\tau$
($\tau\in\Top$); the set $\Top_2$; weights $w(j,\rho)>0$ for the lift
indices $1\le j\le m$ and $\rho\in\Phi$; and a finite set $\Words$ of
words, each with a level prime $T(u)$.
\item Spaces of $\mathbf k$ over $E/F$ (Definition~\ref{def:spaces}).
\item Coefficient data: a CM field $L\supset E$, archimedean places,
local twist types, finite-dimensional Hilbert spaces with lattices and
idempotents, spaces of forms on $U(W_j)$ with integral structures,
theta maps from the definite groups $U(V_j)$, eigenvalues, nonnegative
congruence invariants, and nonnegative rational numbers
$\shalen(l,e)$ (the lengths used in Chapter~\ref{ch:sha}).
\item For each word $u$, a column (Definition~\ref{def:column}) with
its period map.
\item Period lattice data over $L$ indexed by $\Words$
(Definition~\ref{def:period-lattice}).
\item Form lattices, with isometric embeddings of the lattices $N_u$
into them.
\item Local $p$-adic and Galois-side data, a theta jet, and
compatibility conditions among all of these data.
\end{enumerate}
\end{definition}

\noindent The definition is schematic (Section~\ref{sec:conventions}):
the compatibility conditions among these data are omitted. Informally,
a package is the automorphic input at one scale: a host field, the
tower of hermitian spaces over it, forms on the unitary groups of the
tower with their integral structures and theta lifts, and the lattices
and period maps through which the lifts are measured. A package
contains no inequality of the kind collected in
the slope record (Definition~\ref{def:size-data}). The links of
Definitions~\ref{def:upper-link} and~\ref{def:lower-link} attach the
inequalities to a package.

\begin{definition}[size data; schematic]\label{def:size-datum}
Let $\mathcal P$ be a package, $x>1$ and $\eps>0$. A \emph{size datum}
for $\mathcal P$ at scale $x$ and $\eps$ consists of a scale
$(x,\eps,J)$ (Definition~\ref{def:scales}) and of local sizes of
$\mathcal P$ at the places of a finite set of places of its
coefficient field, the \emph{support}.
\end{definition}

\noindent The local sizes are real numbers attached to the places
$\mathfrak p$ of the support and vanishing off it: for each lift index
$j$ a \emph{gain} $\operatorname{gain}_j(\mathfrak p)$, a \emph{field
term} $F_j(\mathfrak p)$ and, for each word $u$, a \emph{congruence
term} $\Cg_{u,j}(\mathfrak p)$; and a \emph{loss}
$\mathrm{ta}_{\mathfrak p}$ and an \emph{orbit term}
$\mathrm{io}_{\mathfrak p}$. A local quantity $y$ at $\mathfrak p$
contributes the global size
\[
\operatorname{sz}_{\mathfrak p}(y)=2e_{\mathfrak p}\,
\frac{\log N\mathfrak p}{[L:\Q]}\,y,
\]
where $L$ is the coefficient field of $\mathcal P$ and $e_{\mathfrak
p}$ the ramification index of $\mathfrak p$ over $\Q$. The
\emph{field term} and the \emph{finite-side gain} of the size datum are
\[
\Fo=\sum_j\sum_{\mathfrak p}\operatorname{sz}_{\mathfrak p}
\bigl(F_j(\mathfrak p)\bigr),\qquad
\Ga=\sum_{\mathfrak p}\operatorname{sz}_{\mathfrak p}\Bigl(
\sum_j\operatorname{gain}_j(\mathfrak p)-\mathrm{ta}_{\mathfrak p}
\Bigr).
\]
The constraints that a size datum imposes on its local sizes are
omitted (Section~\ref{sec:conventions}).

The comparison of slopes that the argument rests on goes back to Bost.

\begin{theorem}[Bost \cite{Bost2001}]\label{thm:bost}
Let $L$ be a number field, let $\overline{\mathcal E}$ and
$\overline{\mathcal F}$ be hermitian vector bundles over
$\operatorname{Spec}\mathcal O_L$, and let $\varphi\colon\mathcal
E_L\to\mathcal F_L$ be an injective $L$-linear map. Then
\[
\widehat{\deg}_n(\overline{\mathcal E})\le\rk\mathcal E\,
\bigl(\hat\mu_{\max}(\overline{\mathcal F})+h(\overline{\mathcal E},
\overline{\mathcal F},\varphi)\bigr).
\]
\end{theorem}

\noindent Equivalently, if $\mathcal E\ne0$, then
$\hat\mu(\overline{\mathcal E})\le\hat\mu_{\max}(\overline{\mathcal
F})+h(\overline{\mathcal E},\overline{\mathcal F},\varphi)$. See
\cite[\S4.1.4, Prop.~4.5, (4.10)--(4.11)]{Bost2001}. The
inequality originates in \cite[\S4.3.1, Prop.~4.3(1)]{Bost1996}, where
$\varphi$ is a morphism of $\mathcal O_L$-modules and only archimedean
norms appear; see also~\cite{Chen2010}.

\begin{remark}[the period map as the auxiliary map]
\label{rem:periodmap}
In transcendence proofs, the map $\varphi$ of Theorem~\ref{thm:bost}
evaluates small sections that one first constructs. Here its role is
played by the period map $B_u\colon N_u\to R^\vee$ of
Definition~\ref{def:period-lattice}. No small sections are constructed.
Instead one shows that the period map between two given lattices is
integral and of bounded height, and applies a slope inequality to it.
\end{remark}

\section{The slope record}

The geometric objects enter the analytic argument only through the
following list of real numbers and inequalities. Definitions
\ref{def:record}--\ref{def:record-main} are shared with
Chapters~\ref{ch:balance} and~\ref{ch:chebyshev}.

\begin{definition}[slope records]\label{def:record}
Let $m$ and $f$ be positive integers.
\begin{enumerate}[label=(\arabic*)]
\item A \emph{slope record} consists of finite sets $U_0$ (the
\emph{shifts}), $\Words(u_0)\ne\emptyset$ for $u_0\in U_0$ (the
\emph{words} at $u_0$) and $\Top$ (the \emph{top places}),
together with the following real numbers:
\begin{enumerate}[label=(\alph*)]
\item $x,\eps,a,b,b',k,\Upsilon,\Upsilon_a,B_0,N_c,\gamma^*,\ell_K,v,
\lgv$ and $\bar\beta_0,\bar\beta_1,\bar\beta_2$;
\item the \emph{top sizes} $a_\tau$ for $\tau\in\Top$, and the
\emph{weights} $w(u_0,j,\rho)$ for $u_0\in U_0$, $0\le j<m$ and
$0\le\rho<f$;
\item for $u_0\in U_0$ and $u\in\Words(u_0)$: $\rk(u)$, $T(u)$,
$\mu_N(u)$, $\mu_N^{\max}(u)$, $h(u)$, $\mu_M(u)$, $\mu_M^{\max}(u)$,
and $\Cg_j(u)$ for $0\le j<m$;
\item for $u_0\in U_0$: $\mu_R(u_0)$, $\mu_R^{\min}(u_0)$, $\Fo(u_0)$,
$\Ga(u_0)$, $\PS(u_0)$, $\IO(u_0)$, $\Delta_R(u_0)$,
$\beta_0(u_0)$, $\beta_1(u_0)$ and $\beta_2(u_0)$.
\end{enumerate}
\item The \emph{top count} and the \emph{top total} of the record are
$n_T=|\Top|$ and $\Sigma a=\sum_{\tau\in\Top}a_\tau$.
\item The \emph{congruence size} of a word $u$ is
$\Cg(u)=\sum_{j<m}\Cg_j(u)$.
\item The \emph{frame slope} at $u_0\in U_0$ is
\[
s(u_0)=\frac{\sum_{u\in\Words(u_0)}\rk(u)\bigl(\mu_N(u)-h(u)\bigr)}
{\sum_{u\in\Words(u_0)}\rk(u)},
\]
defined when the denominator is nonzero.
\item The \emph{excess} at $u_0\in U_0$ is
$u_F(u_0)=\max_{u\in\Words(u_0)}\bigl(\mu_N^{\max}(u)-h(u)\bigr)
-s(u_0)$.
\item\label{item:rec-period}
The \emph{period} of the record is the real number $\Omega>0$
with $\log\Omega=\gamma^*+\log(\pi/v)-\ell_K$, defined when $v>0$.
\end{enumerate}
\end{definition}

\noindent The record is a list of real numbers, but each entry stands
for a quantity of the geometry, and the clauses of
Definition~\ref{def:size-data} below compare these quantities. Once
the record is tied to a family (Definition~\ref{def:tied}), the
entries have the meanings listed in Table~\ref{tab:record-glossary}.
Two conventions fix the indexing. The record index $j$, with
$0\le j<m$, corresponds to the lift index $j+1$ of a package, whose
weights $w(j+1,\rho)$ are indexed by $1\le j+1\le m$. The index
$\rho<f$ of the record corresponds to a node of $\Phi$ through the
bijection $e_P$ of the identification clause
(Definition~\ref{def:window-hyps}). Likewise the top places of the
record correspond to the top places $\Top$ of a host through the
bijection $e_T$ of that clause, and we write $\Top$ for both.

\begin{table}[htbp]
\centering
\small
\begin{tabular}{@{}p{0.2\linewidth}p{0.5\linewidth}p{0.2\linewidth}@{}}
\toprule
symbol & meaning once tied & clauses\\
\midrule
$a_\tau$, $n_T$, $\Sigma a$ & top sizes, the number of top places,
their total & (ii)--(v)\\
$w(u_0,j,\rho)$ & weight of the member at the shift $u_0$, lift
index $j+1$, node $e_P(\rho)$ & (v)\\
$\rk(u)$, $T(u)$ & rank of the ramified lattice $N_u$; level prime of
the word $u$ & (ii), (iv), (v)\\
$\mu_N(u)$, $\mu_N^{\max}(u)$ & slope and maximal slope of $N_u$ &
(vii), (ix)\\
$h(u)$ & the height bound $h_u$ for the period map $B_u$ & (ii),
(iv), (ix)\\
$\mu_M(u)$, $\mu_M^{\max}(u)$ & slope and maximal slope of the form
lattice of $u$ & (vi), (vii)\\
$\Cg_j(u)$, $\Cg(u)$ & congruence size of $u$ at lift index $j+1$,
and its sum over $j$ & (iii), (vii)\\
$\mu_R$, $\mu_R^{\min}$ & slope and minimal slope of $R^\vee$, the
dual of the lattice of test forms & (ii), (ix), (x)\\
$\Fo$ & field term of the size datum & (ii), (vii)\\
$\Ga$ & finite-side gain of the size datum, the Legendre-type gain at
the inert places & (vii), (viii)\\
$\PS$ & the prime sum $\sum_\tau\sum_j\PS_{\tau,w}(a_\tau)$ of the
arithmetic data & (viii)\\
$\IO$ & orbit term: the sum of the global sizes of the orbit terms
$\mathrm{io}_{\mathfrak p}$ & (vii), (viii)\\
$\gamma^*$, $\ell_K$, $v$ & Euler's constant, $L'(1,\chi_{-d})/
L(1,\chi_{-d})$, the half root discriminant; they give $\Omega$ & (ii)\\
$\lgv$ & $\log(\varpi_0\varpi_1)$ for the norm primes & (iii), (v)\\
$\Upsilon$, $\Upsilon_a$, $B_0$, $N_c$ & size quantities of the size
datum & (ii)--(vii)\\
$\Delta_R$, $\beta_1$, $\beta_2$ & spread and frame defects; set by
no tie & (ix)--(xi)\\
$\beta_0$ & rank defect of the member & (ix), (xi)\\
\bottomrule
\end{tabular}
\caption{The entries of a slope record (Definition~\ref{def:record})
and their meanings once the record is tied
(Definition~\ref{def:tied}), with the clauses of
Definition~\ref{def:size-data} in which they occur.}
\label{tab:record-glossary}
\end{table}
\FloatBarrier

\begin{definition}[size data of a slope record]\label{def:size-data}
Let a slope record be given (Definition~\ref{def:record}), with
positive integers $m$ and $f$.
\begin{enumerate}[label=(\arabic*)]
\item \emph{Size constants} are real numbers
\[
c_{\mathrm{div}},\ e_{\mathrm{div}},\ c_{\mathrm{cg}},\ c_{\mathrm{ht}},
\ c_{\mathrm{ct}},\ c_{\mathrm{fs}},\ c_{\mathrm{qs}},\
c_{\mathrm{loss}},\ e_{\mathrm{loss}},\ c_{\mathrm{sp}}\quad
\text{and}\quad X^\sharp.
\]
\item The \emph{losses} of the record for given size constants are
\[
\begin{aligned}
\Lo_P&=c_{\mathrm{loss}}\,n_Tmb\,x^{e_{\mathrm{loss}}\eps},\\
\Lo_Q(u)&=c_{\mathrm{qs}}\IO(u_0)+c_{\mathrm{qs}}\,n_Tm\Bigl(\frac ak+b'
\Bigr)\log T(u)+c_{\mathrm{qs}}X^\sharp+1
\end{aligned}
\]
for $u_0\in U_0$ and $u\in\Words(u_0)$.
\item The record \emph{satisfies the size data} for given size
constants if for every $u_0\in U_0$ and all $u,u'\in\Words(u_0)$ the
following clauses hold, where $s$, $u_F$, $\mu_R$, $\mu_R^{\min}$,
$\Fo$, $\Ga$, $\PS$, $\IO$, $\Delta_R$ and $\beta_0,\beta_1,\beta_2$
are taken at $u_0$.
\begin{enumerate}[label=(\roman*)]
\item\label{item:sd-pos}
(Positivity.) $x,b,k,v,f>0$; $a,b'\ge0$; $a_\tau\ge0$ for every
$\tau$; $w(u_0,j,\rho)>0$ for all $j,\rho$; $\rk(u)>0$; $T(u)>0$;
$\mu_N(u)\le\mu_N^{\max}(u)$; $c_{\mathrm{qs}}\ge0$; $\bar\beta_0<1$;
$\bar\beta_1\ge0$; and $0\le\bar\beta_2\le1/4$.
\item\label{item:sd-div} (Divisor bound.)
\[
\begin{aligned}
\mu_R+h(u)-\Fo\le{}&\sum_{\tau\in\Top}\Bigl(m\bigl(a_\tau\log
a_\tau-a_\tau-a_\tau\log\Omega\bigr)-a_\tau\log T(u)\Bigr)\\
&+c_{\mathrm{div}}\,n_Ta\Upsilon_a
+c_{\mathrm{div}}\,m(B_0+N_c)\,x^{e_{\mathrm{div}}\eps}.
\end{aligned}
\]
\item\label{item:sd-cg} (Congruence sizes.) $|\Cg(u)-m\,\Sigma a\,\lgv|\le
c_{\mathrm{cg}}\,mB_0$.
\item\label{item:sd-ht} (Heights.) $\bigl|\bigl(h(u)+\Sigma a\log T(u)\bigr)
-\bigl(h(u')+\Sigma a\log T(u')\bigr)\bigr|\le
c_{\mathrm{ht}}\,(n_Ta\Upsilon_a+mB_0)$.
\item\label{item:sd-ct} (Counts.)
\[
\sum_{j<m}\frac1f\sum_{\rho<f}\log\frac{w(u_0,j,\rho)\,v}{\pi}
+m\lgv-c_{\mathrm{ct}}\Bigl(\Upsilon_a+\frac{m^2\Upsilon_a}{k}\Bigr)
\le\log T(u)\le c_{\mathrm{ct}}\,m\Upsilon,
\]
and $|\log T(u)-\log T(u')|\le c_{\mathrm{ct}}\log k$.
\item\label{item:sd-fs}
(Form slopes.) $-c_{\mathrm{fs}}\,n_Ta\Upsilon_a(1+m/k)\le
\mu_M(u)$ and $\mu_M^{\max}(u)\le c_{\mathrm{fs}}\,n_Ta\Upsilon_a$.
\item\label{item:sd-qs}
(Quotient slopes.) $\mu_N^{\max}(u)\le\mu_M^{\max}(u)+\Ga+\Fo
-\Cg(u)+c_{\mathrm{qs}}\,mB_0$ and
$\mu_M(u)+\Ga+\Fo-\Cg(u)-\Lo_Q(u)\le\mu_N(u)$.
\item\label{item:sd-loss} (Losses.) $\PS-\Ga\le\Lo_P$ and $\IO\le\Lo_P$.
\item\label{item:sd-slope}
(Slope inequality.) If $0\le\beta_0<1$, $\beta_1\ge0$,
$0\le\beta_2\le1/4$, $\Delta_R>0$ and $s-\mu_R^{\min}\le\Delta_R$, then
\[
s-\log|\Words(u_0)|-\frac{1}{1-\beta_0}\Bigl(4u_F\log\Bigl(1+
\frac{\Delta_R}{u_F}\Bigr)+\bigl(\beta_1+4\sqrt{\beta_2}\bigr)\Delta_R
\Bigr)\le\mu_R,
\]
where the term $4u_F\log(1+\Delta_R/u_F)$ is read as $0$ when
$u_F=0$.
\item\label{item:sd-spread} (Spread.) $s-\mu_R^{\min}\le\Delta_R$ and
$0<\Delta_R\le x^{c_{\mathrm{sp}}}b$.
\item\label{item:sd-def} (Defects.) $0\le\beta_0\le\bar\beta_0$,
$0\le\beta_1\le\bar\beta_1$ and $0\le\beta_2\le\bar\beta_2$.
\end{enumerate}
\end{enumerate}
\end{definition}

\noindent Once the record is tied to a family
(Definition~\ref{def:tied}), $\mu_R$ and $\mu_R^{\min}$ are the slope
and the minimal slope of $R^\vee$, $\mu_N(u)$ and $\mu_N^{\max}(u)$
are the slope and maximal slope of $N_u$, and $h(u)=h_u$. The numbers
$\IO$ and $\Cg_j(u)$ become sums of local sizes over the
support of a size datum, and $\PS$ becomes the sum of the prime sums
$\PS_{\tau,w}$ of Definition~\ref{def:ask}. The spread $\Delta_R$ and
the defects $\beta_1,\beta_2$ are set by no tie. They are constrained
only by the clauses above. Figure~\ref{fig:clause-map} shows which entries of the record each clause connects. The hypotheses of clause~\ref{item:sd-slope} follow
from
clauses~\ref{item:sd-pos}, \ref{item:sd-spread} and~\ref{item:sd-def}.
\begin{figure}[tbp]
\centering
\input{book/figs/ch07-clause-map.tex}
\end{figure}

\begin{definition}[main term of a slope record]\label{def:record-main}
Let a slope record be given (Definition~\ref{def:record}) whose top
sizes $a_\tau$ and weights $w(u_0,j,\rho)$ are positive. Its
\emph{main term} at $u_0\in U_0$ is
\[
\Main(u_0)=\sum_{j<m}\sum_{\tau\in\Top}a_\tau\Bigl(\log a_\tau
-1-\gamma^*+\ell_K-\frac1f\sum_{\rho<f}\log w(u_0,j,\rho)\Bigr).
\]
\end{definition}

\section{The frame theorem}

The slope clause~\ref{item:sd-slope} has the shape of the conclusion of the
following
theorem. It compares a rank-weighted mean of the slopes of many source
lattices with the slope of one target lattice.

\begin{definition}[frame operator]\label{def:frame-operator}
Let $J$ and $\iota_i$, for $i$ in a finite set $I$, be finite sets.
For each $i\in I$ let $\widetilde B_i$ be a complex matrix with rows
indexed by $J$ and columns by $\iota_i$, and let $c_i$ be a real
number.
\begin{enumerate}[label=(\arabic*)]
\item The \emph{frame operator} of the weighted family
$(c_i,\widetilde B_i)_{i\in I}$ is $Q=\sum_{i\in I}c_i\,\widetilde
B_i\widetilde B_i^*$, an operator on $\C^J$.
\item The \emph{Hilbert--Schmidt defect} of the family is
$\sum_{i\in I}\|c_i\widetilde B_i^*\widetilde B_i-1\|_{\mathrm{HS}}^2$,
where $1$ is the identity of $\C^{\iota_i}$ and
$\|A\|_{\mathrm{HS}}^2=\operatorname{tr}(A^*A)$.
\end{enumerate}
\end{definition}

\begin{theorem}[Alp\"oge; frame theorem]\label{thm:frame}
Let $L$ be a totally complex number field and $\sigma\colon
L\hookrightarrow\C$ an embedding. Let $W$ be an $L$-vector space with a basis
indexed by a finite
set $J$, with a positive definite hermitian Gram matrix at each
infinite place, and let $R\subset W$ be a finitely generated
$\mathcal O_L$-submodule spanning $W$. Let $I$ be a nonempty finite
set, and for $i\in I$ let $V_i$ be an $L$-vector space with a basis
indexed by a finite set $\iota_i$ and positive definite hermitian
Gram matrices, $N_i\subset V_i$ a finitely generated $\mathcal
O_L$-submodule spanning $V_i$, $B_i\colon V_i\to W$ an $L$-linear map,
and $h_i$ a real number; assume $\sum_i|\iota_i|>0$. Let $J_1$ be a
finite set, and let $Y'$ and $Y_i$ ($i\in I$) be invertible complex
matrices indexed by $J$ and by $\iota_i$. For each $i$ let
$\widetilde B_i$ be a complex matrix with rows indexed by $J\sqcup J_1$
and columns by $\iota_i$, whose $J$-rows form the matrix
$Y'^{-1}[B_i]_\sigma Y_i$, where $[B_i]_\sigma$ is the matrix of $B_i$
in the given bases, embedded by $\sigma$. Let $c_i>0$ for $i\in I$, let $Q$ be
the frame operator of $(c_i,\widetilde B_i)$
(Definition~\ref{def:frame-operator}), and let $s,u,\Delta,\beta_0,
\beta_1,\beta_2$ be real numbers. Suppose:
\begin{enumerate}[label=(\alph*)]
\item $B_i(N_i)\subset R$ for all $i$, and $\operatorname{ht}(B_i)\le
h_i$ whenever $B_i\ne0$;
\item $s\sum_i\rk N_i\le\sum_i\rk N_i\,(\mu(N_i)-h_i)$; $u\ge0$ and
$\mu_{\max}(N_i)-h_i\le s+u$ for all $i$; $\Delta\ge0$ and
$s-\mu_{\min}(R)\le\Delta$;
\item $0\le\beta_0<1$ and $(1-\beta_0)|J\sqcup J_1|\le|J|$;
\item the Hilbert--Schmidt defect of $(c_i,\widetilde B_i)$ is at most
$\beta_2\sum_i\rk N_i$, with $0\le\beta_2\le1/4$; and
$\operatorname{tr}(Q^2)\le(1+\beta_1)(\operatorname{tr}Q)^2/
|J\sqcup J_1|$, with $\beta_1\ge0$.
\end{enumerate}
Then
\[
s-\log|I|-\frac{1}{1-\beta_0}\Bigl(4u\log\Bigl(1+\frac{\Delta}{u}
\Bigr)+\bigl(\beta_1+4\sqrt{\beta_2}\bigr)\Delta\Bigr)\le\mu(R),
\]
where the term $4u\log(1+\Delta/u)$ is read as $0$ when $u=0$.
\end{theorem}

\noindent So $(W,R)$ and the $(V_i,N_i)$ are hermitian lattices in
the sense of Definition~\ref{def:hlat}, and the slopes $\mu$,
$\mu_{\max}$, $\mu_{\min}$ and the heights are computed from the given
Gram matrices. The matrices $Y'$ and $Y_i$ are arbitrary: the
hypotheses (c) and (d) concern only the matrices $\widetilde B_i$, and
the theorem holds for every choice of $Y'$ and $Y_i$ for which they
are satisfied. A natural choice builds them from the Gram matrices at
$\sigma$, so that $\widetilde B_i$ is the matrix of $B_i$ in bases
that are orthonormal at $\sigma$. The rows indexed by $J_1$ are padding rows,
and $\beta_0$ bounds their
proportion by (c). In the
theorem $s$ plays the role of the frame slope of
Definition~\ref{def:period-lattice}, $u$ is the excess of the most
unstable source over the mean, $\Delta$ is the spread of the mean
over the target, and $\beta_0,\beta_1,\beta_2$ measure how far the
normalised maps at $\sigma$ are from a tight frame (Figure~\ref{fig:near-tight}): $\beta_0$ the
proportion of extra coordinates, $\beta_1$ the spread of the
eigenvalues of $Q$, and $\beta_2$ the distance of each weighted
$\widetilde B_i$ from an isometry.
\begin{figure}[tbp]
\centering
\input{book/figs/ch07-near-tight.tex}
\end{figure}

\begin{proof}[Idea of proof]
This is our account of the method. Alp\"oge's proof is not
reproduced (Section~\ref{sec:conventions}).
Theorem~\ref{thm:bost} is applied, as one step, to the
Harder--Narasimhan pieces of the sources $N_i$, and bounds the slope of
each piece by the slope of its image in $R$ plus the height of $B_i$.
A linear-algebra statement about near-tight frames then controls how
much of $W$ the images miss: if the weighted maps $c_i^{1/2}\widetilde
B_i$ have Hilbert--Schmidt defect at most $\beta_2\sum\rk N_i$ and the
trace condition with $\beta_1$ holds, then the images of subspaces of
small codimension in the $V_i$ span all of $\C^{J\sqcup J_1}$ except a
proportion bounded by $\beta_1+4\sqrt{\beta_2}$ plus four times the
codimension. In this route, a layer-cake summation over the slopes between $s$
and $s+u$ produces the term $4u\log(1+\Delta/u)$, the missing
proportion times the spread $\Delta$ produces the term
$(\beta_1+4\sqrt{\beta_2})\Delta$, and the factor $(1-\beta_0)^{-1}$
comes from passing from $J\sqcup J_1$ to $J$. The term $\log|I|$ is
the cost of averaging over the index set.
\end{proof}

\noindent Figure~\ref{fig:frame} shows the configuration.

\begin{figure}[htbp]
\centering
\input{alpoge/figs/v2/frame.tex}
\caption{The frame theorem (Theorem~\ref{thm:frame}): source lattices $N_i$
map integrally into one target lattice $R$. If the maps are jointly
close to an isometry at one infinite place, the mean source slope $s$
is at most the target slope plus explicit corrections. Bost's
inequality (Theorem~\ref{thm:bost}) compares one source with the
maximal slope of the target. It is one step of the proof of the frame
theorem (Remark~\ref{rem:not-bost}).}
\label{fig:frame}
\end{figure}

\begin{remark}[the slope clause is not Bost's inequality]
\label{rem:not-bost}
The slope input of the argument is neither Theorem~\ref{thm:bost}
applied to each period map $B_u$ (when $B_u$ is injective), which
bounds $\mu(N_u)$ by $\mu_{\max}(R^\vee)+h_u$, nor an average of these
bounds. Clause~\ref{item:sd-slope} of
Definition~\ref{def:size-data} bounds the \emph{average} $s$ of
$\mu(N_u)-h_u$ by the \emph{slope} $\mu(R^\vee)$, not by the maximal
slope, up to a correction controlled by the spread
$s-\mu_{\min}(R^\vee)\le\Delta_R$ and by the defects
$\beta_0,\beta_1,\beta_2$. It has the form of the conclusion of
Theorem~\ref{thm:frame}, with target $R^\vee$, sources $N_u$ and maps
$B_u$. Bost's inequality is one step of the proof of that theorem.
Clause~\ref{item:sd-slope} is a lower bound for $\mu(R^\vee)$.
Clause~\ref{item:sd-div} is an
upper bound for $\mu(R^\vee)+h(u)-\Fo$, so the two clauses together
determine $\mu(R^\vee)$ up to the error terms. This is the cancellation
of Remark~\ref{rem:cancel}. In the construction of tied families the
spread is taken to be $\Delta_R=x^{C_H}b$, and clause~\ref{item:sd-spread}
follows from
the bound $s+\mu_{\max}(R)\le x^{C_H}b$ of
Definition~\ref{def:upper-link}\ref{item:ul-link}\ref{item:ul-int}, since
$\mu_{\min}(R^\vee)=-\mu_{\max}(R)$.
\end{remark}

\section{The archimedean side}

\begin{definition}[the upper link; schematic]\label{def:upper-link}
Let $\mathcal P$ be a package (Definition~\ref{def:package}) with
ranks datum $(k,m,t)$, period lattice data $R,N_u,B_u,h_u$ and frame
slope $s$, norm primes $\varpi_0,\varpi_1$
(Definition~\ref{def:norm-primes}), top places $\Top$ with top sizes
$a_\tau$ and level primes $T(u)$, and put $\lgv=\log(\varpi_0
\varpi_1)$. Let a size datum for $\mathcal P$ be given
(Definition~\ref{def:size-datum}), with the scale and size quantities
$x,\eps,a,b,\Upsilon_a,B_0,N_c$, the real threshold $X^\sharp$ and the
field term $\Fo$, and with the local sizes $\Cg_j(u)$ of the
congruence modules of the words $u$ at the lift indices $j$. Let
$\Omega$ be the period with
$\log\Omega=\gamma+\log(\pi/v)-L'(1,\chi_{-d})/L(1,\chi_{-d})$, where
$v$ is the half root discriminant of the host
(Definition~\ref{def:host}), and let $C,C_{\mathrm{abs}},C_H$ be real
numbers.
\begin{enumerate}[label=(\arabic*)]
\item The \emph{archimedean main term} of a word $u$ is
\[
A(u)=\sum_{\tau\in\Top}\Bigl(m\bigl(a_\tau\log a_\tau-a_\tau-a_\tau
\log\Omega\bigr)-a_\tau\log T(u)\Bigr).
\]
\item\label{item:ul-link}
The package \emph{satisfies the upper link} with constants
$C,C_{\mathrm{abs}},C_H$ if the following hold.
\begin{enumerate}[label=(\alph*)]
\item\label{item:ul-int}
(Integrality.) $B_u(N_u)\subset R^\vee$ for every word $u$, $R$
has full rank, and $s+\mu_{\max}(R)\le x^{C_H}b$; and the local sizes
vanish above $X^\sharp$: at every place $\mathfrak p$ of residue
characteristic $p>X^\sharp$ with $p\notin\{T(u),\varpi_0,\varpi_1,d\}$
and $p$ prime to the orders of the stabilisers in the class sets, the
gains and the congruence terms of $u$ at $\mathfrak p$ are $0$.
\item\label{item:ul-bessel}
(Bessel identity.) The periods are rational, and there are a
unitary map $V$ of the space of test forms and, for each word $u$, a
unitary map $U_u$ of the space of forms of $u$ and a scalar
$\alpha_u\ne0$, such that $B_u(\phi)(F)=\alpha_u\,
b_u(U_u\phi,VF)$, where $b_u$ is the pairing of Bessel type carried by
the package.
\item (Heights.) $\operatorname{ht}(B_u)\le h_u$ whenever $B_u\ne0$,
for the lattice structures of $N_u$ and $R^\vee$.
\item\label{item:ul-div} (Divisor bound.) For every word $u$,
\[
\mu(R^\vee)+h_u-\Fo\le A(u)+C|\Top|a\Upsilon_a
+Cm(B_0+N_c)x^{C_{\mathrm{abs}}\eps}.
\]
\item\label{item:ul-cg}
(Congruences and heights.) For every word $u$ and lift index $j$,
$\bigl|\Cg_j(u)-\bigl(\sum_\tau a_\tau\bigr)\lgv\bigr|\le CB_0$; and
for all words $u,u'$,
\[
\Bigl|\Bigl(h_u+\sum_\tau a_\tau\log T(u)\Bigr)-\Bigl(h_{u'}+\sum_\tau
a_\tau\log T(u')\Bigr)\Bigr|\le C\bigl(|\Top|a\Upsilon_a+mB_0\bigr).
\]
\end{enumerate}
\end{enumerate}
\end{definition}

\noindent In words, the Bessel identity
\ref{item:ul-link}\ref{item:ul-bessel} says that the periods agree, up
to a nonzero scalar and unitary changes of variables, with a pairing
of Bessel type. The pairings $b_u$
are part of the schematic data of a package
(Section~\ref{sec:conventions}). The
divisor bound~\ref{item:ul-link}\ref{item:ul-div} has the shape of
clause~\ref{item:sd-div} of
Definition~\ref{def:size-data}, and~\ref{item:ul-link}\ref{item:ul-cg} has the
shape of
clauses~\ref{item:sd-cg} and~\ref{item:sd-ht}. Three main terms occur and are
kept apart.
The main term of the balance identity for one weight is $\MT_w(a)$
(Definition~\ref{def:ask}). The main term of a record is $\Main(u_0)$
(Definition~\ref{def:record-main}); under the identification clause it
is the sum of the $\operatorname{Re}\MT$ over $\tau$ and $j$
(Proposition~\ref{prop:main-terms}). The archimedean side of the upper
link is $A(u)$.

\begin{remark}[Stirling and the CM period]\label{rem:stirling}
Heuristically, the archimedean norm of a vector of weight $a$ is
governed by Stirling's formula \cite[(5.11.1)]{DLMF2026}, which gives
$a\log a-a$, normalised by the $a$-th power of a CM period $\Omega$
with
\[
\log\Omega=\gamma+\log(\pi/v)-\frac{L'(1,\chi_{-d})}{L(1,\chi_{-d})},
\]
where $v$ is the half root discriminant of the host field
(Definition~\ref{def:host}). This is a period of Chowla--Selberg or
Kronecker-limit type for CM by
$K$; compare \cite{ChowlaSelberg1967}, which gives the product formula
for CM periods. The $m$ lifts contribute $m$ copies of this term. The
quotient $L'(1,\chi_{-d})/L(1,\chi_{-d})$ enters the argument at this
point.
\end{remark}

\begin{remark}[no CM heights]\label{rem:no-colmez}
The argument does not use heights of CM points on a Shimura variety
or a Colmez-type sum of $L'/L(1,\chi)$ over many characters
\cite{Colmez1993}. Only
$L'(1,\chi_{-d})/L(1,\chi_{-d})$ enters the main term. The other
characters enter only through prime sums at inert primes weighted by
$w$. No CM point is evaluated, and the slope inequality is applied to a
period map. Arakelov geometry is used to bound the covolume of a
lattice on the indefinite side (Theorem~\ref{thm:covolume}).
\end{remark}

\section{The finite side}

\begin{definition}[the lower link; schematic]\label{def:lower-link}
Let $\mathcal P$ be a package with ranks datum $(k,m,t)$, so with
$m$ lift indices, top places $\Top$ and top sizes $a_\tau$, and let
$f=|\Phi|$ and $N(m)=k+m(t+1)$. Let a size datum for $\mathcal P$ be
given (Definition~\ref{def:size-datum}), with scale $x$, total size
$a$ and real threshold $X^\sharp$. For each place $\mathfrak p$ of its
support and each lift index $j$ let $\operatorname{gain}_j(\mathfrak
p)$ and $\mathrm{ta}_{\mathfrak p}$ be the local sizes of the size
datum. Let $C$ be a constant of the link, put
\[
\lambda_{\mathrm{fl}}(p)=Cf\,N(m)^2\log_pa+Cf(\log_pa)^2
\]
for a prime $p$, and let $\lambda_{\mathrm{lo}}(p)$ be a lower-order
term of the same kind: $C$ times size quantities of the datum times
$\log_p(2a)$, plus, for inert $p<X^\sharp$, correction terms, each a
multiple of $\log x$ by size quantities of the scale, whose exact form
is omitted (Section~\ref{sec:conventions}).
\begin{enumerate}[label=(\arabic*)]
\item For a lift index $j$, $\tau\in\Top$ and a prime $p$, the
\emph{Legendre sum} is
\[
L_{j,\tau}(p)=\sum_{i\ \mathrm{odd}}\Bigl(\frac{a_\tau}{p^i}
-w^{\pt}_{j,\tau}\Bigr)_+,
\]
where $w^{\pt}_{j,\tau}$ is the partner weight of the package at
level $j$ (compare Definition~\ref{def:ask}\ref{item:ask-partner}).
\item The package satisfies the \emph{Legendre-type gain condition} if
at every place $\mathfrak p$ above a prime $p\le X^\sharp$ inert in $K$, and for
every lift index $j$,
\begin{gather*}
0\le\operatorname{gain}_j(\mathfrak p)\le\frac{4|\Top|a}{p},\\
\sum_j\sum_\tau\bigl(L_{j,\tau}(p)-\lambda_{\mathrm{lo}}(p)\bigr)
-m\,\lambda_{\mathrm{fl}}(p)\le\sum_j\operatorname{gain}_j(\mathfrak
p)-\mathrm{ta}_{\mathfrak p},
\end{gather*}
the inner sum over the top places attached to $\mathfrak p$.
\item\label{item:ll-link}
The package satisfies the \emph{lower link} if it satisfies the
Legendre-type gain condition and further local conditions: bounds for
local lengths at places whose residue characteristic is not $T(u)$,
and bounds for sums over orbits of places.
\end{enumerate}
\end{definition}

\noindent The
further local conditions in~\ref{item:ll-link} are part of the data and are not
written out (Section~\ref{sec:conventions}). The gain condition
is the finite side of the argument:
each inert prime $p\le X^\sharp$ contributes the Legendre sum of
$a_\tau/p^i$ over odd $i$, which mirrors the contribution of $p$ to
the prime sum $\PS_{\tau,w}$ of Definition~\ref{def:ask}.

\begin{remark}[the balance of $n!$]\label{rem:factorial}
Summing the gains over places gives approximately
$\sum_{p\text{ inert}}(a/p)\log p+\cdots$, a prime sum over a
Frobenius class. Set against Stirling's $a\log a-a$
(Remark~\ref{rem:stirling}), this is the balance between the
archimedean and the finite sizes of $\log n!$ used by
Chebyshev~\cite{Chebyshev1852}. The method is Chebyshev's factorial
argument, twisted by the characters of $\Gal(E/\Q)$ through CM periods
and sharpened by a chain of $m\approx\sqrt x$ theta lifts. The chain
is the multiplicity device. The role that vanishing of an auxiliary
polynomial to high order plays in Stepanov's
method~\cite{Stepanov1969} is played here by $p$-adic divisibility of
lifted periods at all inert $p$ at once, forced by weights $a\ge p$.
\end{remark}
\clearpage
\clearpage
\chapter{Families and the covolume bound}
\label{ch:covolume}\label{sec:covolume}

This chapter collects the packages of Chapter~\ref{ch:theta} into
families (Definition~\ref{def:family}), introduces the closed
properties through which constants are fixed before a family is chosen
(Definition~\ref{def:closed}), and states the covolume bound,
Theorem~\ref{thm:covolume}, the single result through which arithmetic
geometry enters the argument. The covolume bound is not proved here:
its inputs are listed after the statement, and
Remark~\ref{rem:ag-strength} records the precise form in which each of
them is used, including one input for which the implication from the
cited theorem is not claimed. Chapter~\ref{ch:assembly} uses
Theorem~\ref{thm:covolume} as a black box. In the proof of the main
theorem it enters through the existence of admissible-family data,
Theorem~\ref{thm:packages}\ref{item:packages-data} of
Chapter~\ref{ch:families}, and, summed over the lifts, it supplies
the main term of the divisor bound of the upper link
(Definition~\ref{def:upper-link}).

The packages of Chapter~\ref{sec:slopes} are assembled into families,
one package for each member at each scale. The covolume bound is
a bound for the covolumes of lattices of integral forms on the
indefinite unitary groups $U(W_j)$. The chapter ends with
the forms in which the theorems of arithmetic geometry are used.

\section{Families and closed properties}

\begin{definition}[families; schematic]\label{def:family}
Let $\eps>0$.
\begin{enumerate}[label=(\arabic*)]
\item\label{item:fam-def}
A \emph{family at} $\eps$, written $\Fam$, consists of:
\begin{enumerate}[label=(\alph*)]
\item a threshold $x_1>1$ and an unbounded set $\Win\subset\R$;
\item for each real $x$ a finite index set $I_x$, nonempty for
$x\in\Win$, and for each $i\in I_x$ a \emph{member}: a package
$\mathcal P_{x,i}$ together with a size datum $Z_{x,i}$ for
$\mathcal P_{x,i}$ at scale $x$ and $\eps$
(Definition~\ref{def:size-datum}) when $x>x_1$, such that the
third entry $J$ of the scale $(x,\eps,J)$ of $Z_{x,i}$ is independent
of $x$ and $i$;
\item\label{item:fam-unif}
uniformity conditions: for each $x$, all members at $x$ share
their ranks (Definition~\ref{def:ranks}), the primes $q$ and $d$ and
the integer $Q_S$ of their hosts (Definition~\ref{def:host}), the set
$\Top_2$, the Artin data, the auxiliary primes $\varpi_0,\varpi_1$ and
their sets of words, through fixed identifications; and for every
member and every shift there is a member with that shift and the same
top sizes;
\item\label{item:fam-param}
parameter choices and tie data, and period lattice data
(Definition~\ref{def:period-lattice}) whose heights $h_u$ are given by
an explicit formula.
\end{enumerate}
\item A family $\Fam$ \emph{has windows of exponent} $\lambda>1$ if
$\Win$ has windows of exponent $\lambda$
(Definition~\ref{def:windows}).
\end{enumerate}
\end{definition}
\noindent The local sizes of the size data, the remaining uniformity
conditions in~\ref{item:fam-def}\ref{item:fam-unif}, the parameter choices and
the formula for the
heights in~\ref{item:fam-def}\ref{item:fam-param} are part of the data of a
family. They are not written out
here (Section~\ref{sec:conventions}). Later chapters use them
through the ties of
Definition~\ref{def:tied} and the hypotheses of
Definitions~\ref{def:covolume-condition} and~\ref{def:admissible}.

In plain terms, a property is \emph{closed} if, for $\eps$ small, it
holds for every family whose constants exceed thresholds chosen before
the family. The next definition makes this precise.

\begin{definition}[thresholds and closed properties]\label{def:closed}
Let $\eps>0$. Each family $\Fam$ at $\eps$
(Definition~\ref{def:family}) carries finite lists of real constants
$c_{\mathrm{pre}}(\Fam)$, $c_{\mathrm{out}}(\Fam)$,
$c_{\mathrm{last}}(\Fam)$ and, for each real number $c$, a list
$c_{\mathrm{mid}}(\Fam,c)$ whose first entry is $c$.
\begin{enumerate}[label=(\arabic*)]
\item A \emph{threshold} $\kappa$ consists of a real number
$\kappa_{\mathrm{pre}}$ and three functions $\kappa_t$,
$\kappa_{\mathrm{mid}}$, $\kappa_{\mathrm{last}}$ from finite lists of
real numbers to real numbers.
\item A family $\Fam$ is \emph{large for} $(\kappa,c)$ if every entry
of $c_{\mathrm{pre}}(\Fam)$ is at least $\kappa_{\mathrm{pre}}$, every
entry of $c_{\mathrm{mid}}(\Fam,c)$ is at least
$\kappa_{\mathrm{mid}}(c_{\mathrm{out}}(\Fam))$, and every entry of
$c_{\mathrm{last}}(\Fam)$ is at least $\kappa_{\mathrm{last}}$ of the
list $c_{\mathrm{out}}(\Fam)$ followed by $c_{\mathrm{mid}}(\Fam,c)$.
\item\label{item:cl-closed}
Let $\mathcal C$ be a class of families and $P$ a property of a
family and a real number. Then $P$ is \emph{closed} on $\mathcal C$ if
there is $\eps_1>0$ such that for every $\eps\in(0,\eps_1]$ there is a
threshold $\kappa$ for which $P(\Fam,c)$ holds for every family
$\Fam$ at $\eps$ in $\mathcal C$ and every real $c$ such that $\Fam$
is large for $(\kappa,c)$.
\end{enumerate}
\end{definition}

\noindent The entries of the constant lists are constants of the
family and are not listed here. In~\ref{item:cl-closed} the threshold is chosen
\emph{before} the family. Several closed properties can therefore be combined by
combining their thresholds, and a single family that is large for the
combined threshold is then supplied by the existence theorem of
Chapter~\ref{sec:assembly}.

\section{The covolume bound}

\begin{definition}[the covolume condition; schematic]
\label{def:covolume-condition}
Let $\mathcal C_{\mathrm{cov}}$ be the class of families whose
packages carry the integral structures used below, let $\eps>0$, and
let $\Fam$ be a family at $\eps$ in $\mathcal C_{\mathrm{cov}}$ with members
$(\mathcal P_{x,i},Z_{x,i})$.
\begin{enumerate}[label=(\arabic*)]
\item\label{item:cc-member}
A member $(\mathcal P,Z)$ satisfies the \emph{member hypothesis}
if the following two conditions hold.
\begin{enumerate}[label=(\alph*)]
\item (Field size.) Let $F_O\subset\C$ be the field generated over
$\Q$ by the image of the coefficient ring of $\mathcal P$ under its
fixed complex embedding, and put $P_O=(19fr^2+198f^2)\,x^\eps$,
where $f$ and $r$ are the derived quantities of the ranks of
$\mathcal P$ (Definition~\ref{def:ranks}) and $(x,\eps,J)$ is the
scale of $Z$. Then $\log[F_O:\Q]\le P_O$ and
$\log|\operatorname{disc}F_O|/[F_O:\Q]\le P_O$.
\item (Local types at the norm primes.) For $k=0,1$, write
$G_k=\mathrm{GL}_r(\F_{\varpi_k})$. There is a labelling of the
irreducible characters of $G_k$ by polynomials over
$\F_{\varpi_k}$, of the kind used in Green's parametrization
(Theorem~\ref{thm:green}) and built on the character of
$\F_{\varpi_k}^\times$ that $\mathcal P$ fixes at $\varpi_k$, such
that for every polynomial $P$ the idempotent of $\C[G_k]$ that
$\mathcal P$ attaches to $P$ is the sum of the central idempotents
of the irreducible characters with label $P$.
\end{enumerate}
\item A \emph{point} of a member is a choice of the auxiliary integral
data (lattices, integral structures and constants) at which the
lattices $M_j$ below are formed, and the \emph{point conditions} are
compatibility conditions on such a choice; both are schematic. A
point $\mathfrak z$ carries a positive integer $\deg$, the degree by
which covolumes are normalised.
\item\label{item:cc-ineq}
Let $\mathfrak z$ be a point of a member $\mathcal P$, with weights
$w(j,\rho)$ and period $\Omega$, and let $C\ge0$ and $x>1$. For a lift
index $1\le j\le m$, let $M_j$ be the lattice of integral forms on the
indefinite unitary group $U(W_j)$ attached to $\mathfrak z$, and let
$(p_j,q_j)=(k+j-1,m-j)$ be the signature of $W_j$
(Definition~\ref{def:spaces}). The point satisfies the \emph{covolume
inequality} with constant $C$ at $x$ if for every lift index $j$
\[
\frac{\log\operatorname{covol}(M_j)}{\deg}\le\sum_{\rho}\bigl(p_j-q_j
\bigr)\,w(j,\rho)\log\Omega+C\,B_0\log\log x,
\]
where $\log\operatorname{covol}(M_j)$ is the normalised logarithmic
covolume of $M_j$, $\deg$ is the degree of the point, the sum runs over
$\rho\in\Phi$, and $B_0$ is the size quantity of the member
(Definition~\ref{def:size-datum}).
\item The \emph{covolume condition} is the statement that the
following property of a family $\Fam$ and a real number $c$ is closed
on $\mathcal C_{\mathrm{cov}}$ (Definition~\ref{def:closed}): there is
$C_0\ge0$ such that for all large $x$, every index $i$ for which
$(\mathcal P_{x,i},Z_{x,i})$ satisfies the member hypothesis, every
point $\mathfrak z$ of $\mathcal P_{x,i}$ that satisfies the point
conditions, and every $C\ge C_0$, the point $\mathfrak z$ satisfies the
covolume inequality with constant $C$ at $x$.
\end{enumerate}
\end{definition}

\noindent The class $\mathcal C_{\mathrm{cov}}$ is fixed once and for
all. It is distinct from the classes of admissible families of
Chapter~\ref{sec:assembly} (see the paragraph after
Definition~\ref{def:admissible}). The class $\mathcal
C_{\mathrm{cov}}$ and the points of a member with their point
conditions are schematic (Section~\ref{sec:conventions}). The period
$\Omega$ is that of the slope
record (Definition~\ref{def:record}\ref{item:rec-period}), with
$\gamma^*=\gamma$ and
$\ell_K=L'(1,\chi_{-d})/L(1,\chi_{-d})$ (Remark~\ref{rem:stirling}).
The lattices $M_j$ lie on the indefinite side, one for each lift index.
They are not the test lattice $R$ of
Definition~\ref{def:period-lattice}.

\begin{remark}[the member hypothesis is not a tie]
\label{rem:member-hyp}
Chapter~\ref{sec:chebwin} ties a family at a scale to a slope record
(Definition~\ref{def:tied}). The member hypothesis of
Definition~\ref{def:covolume-condition}\ref{item:cc-member} is not a
condition on the slope record of a member. The member hypothesis is a bound on
the size of the field $F_O$ together with a condition on the local
types of the package at the two norm primes. Both conditions involve
only the package and its size datum. They involve no point of the
member and no slope record. The
covolume bound is independent of the analytic data of
Chapters~\ref{sec:balance} and~\ref{sec:chebwin}.
\end{remark}

\begin{theorem}[Alp\"oge; covolume bound]\label{thm:covolume}
There is $\eps_1>0$ such that for every
$\eps\in(0,\eps_1]$ there is a threshold $\kappa$ with the following
property. Let $\Fam$ be a family at $\eps$ in $\mathcal
C_{\mathrm{cov}}$ and $c$ a real number such that $\Fam$ is large for
$(\kappa,c)$. Then there is $C_0\ge0$ such that for all large $x$,
every member $(\mathcal P_{x,i},Z_{x,i})$ satisfying the member
hypothesis, every point of $\mathcal P_{x,i}$ satisfying the point
conditions and every $C\ge C_0$, the point satisfies the covolume
inequality with constant $C$ at $x$
(Definition~\ref{def:covolume-condition}\ref{item:cc-ineq}).
\end{theorem}

\noindent Written out, the conclusion is that for every lift index $j$
\[
\frac{\log\operatorname{covol}(M_j)}{\deg}\le\sum_{\rho\in\Phi}
\bigl(p_j-q_j\bigr)\,w(j,\rho)\log\Omega+C\,B_0\log\log x,
\]
where $M_j$, $(p_j,q_j)$, $\deg$, $w(j,\rho)$, $\Omega$ and $B_0$ are
the lattice, signature, degree, weights, period and size quantity
attached to the point and the member by
Definition~\ref{def:covolume-condition}\ref{item:cc-ineq}. So
Theorem~\ref{thm:covolume} says that the covolume condition holds
(Figure~\ref{fig:covolume-quantifiers} shows the order of its quantifiers).

\noindent \emph{Informal statement.} For the lattice $M_j$ of
integral cusp forms of the weights $w(j,\rho)$ on the Shimura variety
of $U(W_j)$, at the level prescribed by the package
(Theorem~\ref{thm:packages}\ref{item:packages-points}), the normalised
logarithmic covolume is at most $\sum_{\rho\in\Phi}(p_j-q_j)\,
w(j,\rho)\log\Omega$, up to an error $O(B_0\log\log x)$, for all
large $x$ and uniformly over the members of a family whose field of
coefficients is not too large.
\begin{figure}[tbp]
\centering
\input{book/figs/ch08-quantifiers.tex}
\end{figure}

\begin{proof}[On the proof]
We do not reproduce the proof. The bound rests on the following
inputs, all of which are cited theorems of arithmetic geometry in the
forms described in Remark~\ref{rem:ag-strength} and
Figure~\ref{fig:ag-inputs}.
\begin{enumerate}[label=(\arabic*)]
\item Toroidal compactifications of the unitary Shimura varieties of
the groups $U(W_j)$ and of Siegel modular varieties
\cite{FaltingsChai1990,Lan2013}, and resolution of singularities in
characteristic zero \cite{Hironaka1964}.
\item The arithmetic Hilbert--Samuel theorem
\cite{GilletSoule1992,AbbesBouche1995} for a smooth metric, applied
once, as an upper bound for a counting volume on the projective image
of a map given by cusp forms with a Fubini--Study metric. Also
Moriwaki's bigness theorem \cite{Moriwaki2009}.
\item A comparison of the Fubini--Study metric with the Petersson
metric through the Bergman kernel and the comparison of sup and $L^2$
norms of cusp forms. Geometric degrees from a Chern--Weil
statement for good metrics of Mumford type \cite{Mumford1977}.
\item Kawamata--Viehweg vanishing \cite{Kawamata1982,Viehweg1982} and
Nadel vanishing \cite{Nadel1990}.
\item Integral models of the Shimura varieties, through
Grothendieck--Messing deformation theory \cite{Messing1972}.
\end{enumerate}
\end{proof}

\noindent Figure~\ref{fig:covolume-main} shows the coefficient of the
main term along the tower, and how the bound feeds the divisor bound
of a slope record.

\begin{figure}[htbp]
\centering
\input{alpoge/figs/v2/covolume-main.tex}
\caption{(a) The signature $(p_j,q_j)=(k+j-1,m-j)$ of the space $W_j$ for
$k=8$ and $m=4$; $t$ plays no role. Since $k\ge m+3$, the coefficient
$p_j-q_j=k-1-m+2j$ of the main term of Theorem~\ref{thm:covolume} is
positive for every $j$ and grows with $j$. (b) The theorems of
arithmetic geometry feed only the covolume bound
(Remark~\ref{rem:covol-supply}). The dashed arrow: the main term
$\log\Omega$ of the covolume bound corresponds to the term
$-ma_\tau\log\Omega$ of the divisor bound, clause~\ref{item:sd-div} of
Definition~\ref{def:size-data}.}
\label{fig:covolume-main}
\end{figure}

\noindent Summed over $j$, the main terms of
Theorem~\ref{thm:covolume} give the term $-ma_\tau\log\Omega$ of the
divisor bound of the upper link
(Definition~\ref{def:upper-link}\ref{item:ul-link}\ref{item:ul-div}):
the lattice of rows is built from the $M_j$.
\begin{figure}[tbp]
\centering
\input{book/figs/ch08-ag-inputs.tex}
\end{figure}

\section{How arithmetic geometry enters}

\begin{remark}[one output of the arithmetic geometry]
\label{rem:covol-supply}
The theorems of arithmetic geometry used by the argument (resolution
of singularities, arithmetic Hilbert--Samuel, Moriwaki's bigness
theorem, toroidal compactifications, Kawamata--Viehweg and Nadel
vanishing, Grothendieck--Messing theory) do not enter the argument at
separate steps. All of them enter through the one covolume bound,
Theorem~\ref{thm:covolume}.
Nothing in the argument bounds heights of CM points
(Remark~\ref{rem:no-colmez}).
\end{remark}

\begin{remark}[smooth metrics suffice]\label{rem:smooth-metrics}
The Hodge bundle with its Petersson metric has logarithmic
singularities along the boundary of a toroidal compactification
\cite{BurgosKramerKuhn2005}. The argument does not need an arithmetic
Hilbert--Samuel theorem for such log-singular metrics. The
Hilbert--Samuel statements used require smooth metrics, and four
devices avoid the singularity. First, only
cusp forms with $L^2$ Petersson norms enter. Second, cusp forms are
identified with sections of an adjoint-type bundle on a smooth
projective resolution. Third, for the height bound the Petersson
metric is replaced by a Fubini--Study metric on a projective image,
with Bergman-kernel and sup-versus-$L^2$ comparisons. Fourth,
geometric degrees come from a Chern--Weil statement for good metrics
of Mumford type \cite{Mumford1977}. The comparisons of the second and
third devices cost at most the unit error $C\,B_0\log\log x$ of each
lift in Theorem~\ref{thm:covolume}.
\end{remark}

\begin{remark}[the heights need no further geometry]\label{rem:one-point}
The heights $h_u$ of the period lattice data
(Definition~\ref{def:period-lattice}) need no arithmetic geometry
beyond the covolume bound. The height used depends only on a number
computed from the rows of the package. All the geometric content of
the argument is in the covolume bound,
Theorem~\ref{thm:covolume}.
\end{remark}

\begin{remark}[the forms of the arithmetic-geometry inputs]
\label{rem:ag-strength}
Each input of Theorem~\ref{thm:covolume} is used in a specific
formulation, which need not have the generality of the theorem it is
named after. The formulations range from near-general to special or
weakened.
\begin{enumerate}[label=(\arabic*)]
\item Near-general.
\begin{enumerate}[label=(\alph*)]
\item Resolution of singularities in characteristic zero
\cite{Hironaka1964}: a quasi-compact, separated, reduced scheme
locally of finite type over a field of characteristic zero admits a
sequence of blow-ups to a smooth scheme that is an isomorphism over
the smooth locus and whose exceptional locus has simple normal
crossings.
\item Arithmetic Hilbert--Samuel \cite{GilletSoule1992,AbbesBouche1995}:
for a nef hermitian line bundle $\overline L$ with smooth metric on an
arithmetic variety of dimension $D+1$, whose self-intersection number
$a$ enters as a hypothesis, and any hermitian line bundle $\overline
N$, the logarithm of the number of small sections of $\overline
L^{\otimes n}\otimes\overline N$ is $a\,n^{D+1}/(D+1)!+o(n^{D+1})$.
\end{enumerate}
\item Special or weakened.
\begin{enumerate}[label=(\alph*)]
\item Moriwaki's bigness theorem \cite{Moriwaki2009}: a lower bound for
the arithmetic volume by the self-intersection number, stated on
multi-projective coordinate rings, for an intersection pairing that is
specified by axioms.
\item Toroidal compactifications \cite{FaltingsChai1990,Lan2013}: for
unitary Shimura varieties of signature $(p,q)$ attached to a CM field,
a smooth model with simple normal crossings boundary and a logarithmic
Kodaira--Spencer isomorphism; and, in a separate integral statement,
for Siegel modular varieties of level $n\ge3$ odd over
$\Z[\zeta_n,1/n]$, a smooth projective compactification with a
semi-abelian extension of the universal abelian scheme and an ample
modification of a power of the Hodge bundle.
\item Kawamata--Viehweg vanishing \cite{Kawamata1982,Viehweg1982}
for a nef and big line bundle on a
smooth projective variety over an algebraically closed field of
characteristic zero, without $\Q$-divisors.
\item Nadel vanishing \cite{Nadel1990} in the form: there is \emph{some} ideal
sheaf
whose cosupport is the non-klt locus and for which the twisted higher
cohomology vanishes. The ideal sheaf is not identified with the
multiplier ideal.
\item Grothendieck--Messing theory \cite{Messing1972} over a
divided-power thickening
with nilpotent divided powers, along which $p$ is nilpotent.
\end{enumerate}
\end{enumerate}
Each formulation is implied by the cited theorem, with two
qualifications. For~(1)(b) with a nef, rather than ample, line bundle
the asymptotic formula is due to Moriwaki \cite{Moriwaki2009}, as a
consequence of the continuity of arithmetic volumes; Gillet--Soul\'e
and Abbes--Bouche treat ample bundles with positive metrics.
For~(2)(a) the implication needs the axiomatic intersection pairing to
be the arithmetic intersection pairing of \cite{Moriwaki2009}. This
is not checked here, so for~(2)(a) the implication is not claimed.
\end{remark}
\clearpage
\clearpage
\chapter{The Chebyshev inequality on windows}
\label{ch:chebyshev}\label{sec:chebwin}

This chapter proves the Chebyshev inequality on windows,
Theorem~\ref{thm:chebyshev-windows}: for a family that satisfies the
uniform explicit-formula condition and is tied at every large scale of
its window set, $G_\lev(x)\le x^{1/2+(e_D+e_L+7)\eps}$ on the windows.
The proof combines the slope records of Chapter~\ref{ch:theta} with
the balance identity and the core bound of Chapter~\ref{ch:balance},
through the gap bound, Theorem~\ref{thm:gap}. The proof of the gap
bound is given in outline: its steps are written out, but its last
step, an estimation of error terms by linear arithmetic in which the
largeness of $x$ absorbs the constants, is not. In
Chapter~\ref{ch:assembly} the Chebyshev inequality on windows is a
black box. It is applied in the last step of the proof of the main
theorem, inside Theorem~\ref{thm:zero-free-primitive} of
Chapter~\ref{ch:families}, whose output goes into the windowed Landau
criterion for levels (Theorem~\ref{thm:level-criterion},
Chapter~\ref{ch:landau}).

A slope record
(Definitions~\ref{def:record} and~\ref{def:size-data}) is a list of
real numbers subject to eleven inequalities. The prime sum and the
gap of Chapter~\ref{sec:balance} are attached to arithmetic data. Two
hypotheses connect them: an identification clause, which says that the
record's prime sum, top sizes and weights are those of the arithmetic
data (Definition~\ref{def:window-hyps}), and the ties of a family
(Definition~\ref{def:tied}), which say that the record's numbers are
the slopes, heights and local sizes of actual packages. Under these
hypotheses the eleven clauses give an upper bound for the gap
(Theorem~\ref{thm:gap}), the core bound of Chapter~\ref{sec:balance}
turns it into an upper bound for $G_\lev$, and one obtains
\[
G_\lev(x)\le x^{1/2+(e_D+e_L+7)\eps}
\]
for all large $x$ in the window set of the family
(Theorem~\ref{thm:chebyshev-windows}).

In the computation the slope of the test lattice cancels
(Remark~\ref{rem:cancel}). The exponent $1/2$ comes from the
bookkeeping of the error bracket, in which the number $m$ of theta
lifts is balanced against $x/m$
(Proposition~\ref{prop:sqrtx}).

\section{Window hypotheses and ties}

\begin{definition}[window and identification hypotheses]
\label{def:window-hyps}
Let $H$ be a finite abelian group of order $n$ and $\GG=V_4\times H$.
Let a slope record with shifts $U_0=H$ and positive integers $m$ and
$f$ be given (Definition~\ref{def:record}), together with size
constants (Definition~\ref{def:size-data}) and real numbers $c_W$,
$e_K$ and $C_L$. In~\ref{item:wh-window} the letters
$x,\eps,a,b,b',k,\Upsilon,\Upsilon_a,B_0,N_c,\bar\beta_0,\bar\beta_1,
\bar\beta_2$ and $n_T$ denote the quantities of the record, and the
$c$'s and $e$'s with subscripts denote its size constants.
Let $D$ be arithmetic data for $\GG$, $\ell$ a
function from the characters of $\GG$ to $\C$, $N\ge1$ an integer,
$\Top,\Top_2\subset H$ finite sets, $a\colon H\to\R$ and
$w\colon H\times\{0,\dots,m-1\}\times\GG\to\R$.
\begin{enumerate}[label=(\arabic*)]
\item\label{item:wh-window} The \emph{window conditions} hold if
\begin{enumerate}[label=(\alph*)]
\item $1\le x$, $0<\eps\le1$, $1\le b$, $xb\le a$ and $1\le n_T$;
\item $m\le x^{1/2+\eps/2}$ and $xm\le k$;
\item $0\le e_K$, $1\le k\le x^{e_K}$ and $k\le a$;
\item $0\le\Upsilon\le2x^\eps$, $0\le\Upsilon_a\le2x^\eps$, $0\le b'$,
$kb'\le a$ and $0\le X^\sharp\le a$;
\item $0\le a_\tau\le2a$ for every top place $\tau$ of the record, and
$B_0,N_c\ge0$;
\item $\log|\Words(u_0)|\le(k-3)\log2$ for every $u_0\in U_0$;
\item $\bar\beta_0\le1/2$, $\bar\beta_1\ge0$ and $\bar\beta_2\ge0$;
\item $c_{\mathrm{div}},c_{\mathrm{cg}},c_{\mathrm{ht}},c_{\mathrm{ct}},
c_{\mathrm{fs}},c_{\mathrm{qs}},c_{\mathrm{loss}}\in[0,c_W]$ and
$e_{\mathrm{div}},e_{\mathrm{loss}},c_{\mathrm{sp}}\ge0$;
\item\label{item:wh-defect} (\emph{defect smallness})
$\bigl(\bar\beta_1+4\sqrt{\bar\beta_2}\bigr)\,x^{c_{\mathrm{sp}}}b\le
c_W\,x^\eps\,n_T\,a$.
\end{enumerate}
\item The \emph{identification clause} holds if there are bijections
$e_T$ from the top places of the record onto $\Top$ and
$e_P\colon\{0,\dots,f-1\}\to\Phi$
such that the record's top size at $\tau$ is $a_{e_T(\tau)}$ and its
weight at $(u_0,j,\rho)$ is $w(u_0,j,e_P(\rho))$; $\gamma^*=\gamma$;
$\ell_K=\operatorname{Re}\ell(\chi_K)$; and for every $u_0\in H$
\[
\PS(u_0)=\sum_{\tau\in\Top}\sum_{j<m}\PS_{\tau,w(u_0,j)}(a_\tau).
\]
\item The \emph{bounds on $B_0$ and $N_c$} hold if
\[
B_0\le2nmbx^\eps\qquad\text{and}\qquad
N_c\le\frac{|\Top_2|\,kbx^\eps}{n^2x^{-\eps}}.
\]
\item The \emph{values at $x$} are bounded if
$|\GPsi_\theta(x)|\le x^{1/2+\eps}$ and $|\ell(\theta)|\le C_L\log x$
for $\theta\in\{\omega,\chi_K,\chi_K\omega\}$.
\end{enumerate}
\end{definition}

\noindent The defect
smallness~\ref{item:wh-window}\ref{item:wh-defect} is
where the size of the defects in clause~\ref{item:sd-slope} of
Definition~\ref{def:size-data} is fixed; in the families of
Chapter~\ref{sec:assembly}, $\bar\beta_1=1/x$.

\begin{proposition}[the main terms]\label{prop:main-terms}
Let $H$ be a finite abelian group and $\GG=V_4\times H$. Let a slope
record with shifts $U_0=H$ and positive integers $m$ and $f$ be given
(Definition~\ref{def:record}). Let $D$ be arithmetic data for $\GG$,
$\ell$ a function from the characters of $\GG$ to $\C$, $N\ge1$ an
integer, $\Top,\Top_2\subset H$ finite sets, $a\colon H\to\R$ with
$a_\tau>0$ for $\tau\in\Top$, and
$w\colon H\times\{0,\dots,m-1\}\times\GG\to\R$ with
$w(u_0,j,\rho)>0$ for all $u_0\in H$, $j<m$ and $\rho\in\Phi$.
Suppose that the identification clause
(Definition~\ref{def:window-hyps}) holds, and let $u_0\in H$. Then the
main term $\Main$ (Definition~\ref{def:record-main}), the term $\MT$
(Definition~\ref{def:ask}) and the gap $\gap$
(Definition~\ref{def:thick}) satisfy:
\begin{enumerate}[label=(\alph*)]
\item\label{item:mt-main} $\displaystyle\Main(u_0)=\sum_{\tau\in\Top}
\sum_{j<m}\operatorname{Re}\MT_{w(u_0,j)}(a_\tau)$;
\item\label{item:mt-gap} $\gap(u_0)=\PS(u_0)-\Main(u_0)$.
\end{enumerate}
\end{proposition}

\begin{proof}
\noindent\emph{Proof of \ref{item:mt-main}.} Since $a_\tau$ and the weights are
real and positive,
\[
\operatorname{Re}\MT_w(a)=a\Bigl(\log a-1-\gamma+\operatorname{Re}
\ell(\chi_K)-\frac1f\sum_{\rho\in\Phi}\log w(\rho)\Bigr),
\]
where $f=|\Phi|$; the bijection $e_P$ shows that this $f$ is the
record's $f$. By the identification clause, $\gamma^*=\gamma$,
$\ell_K=\operatorname{Re}\ell(\chi_K)$, and $e_T$ and $e_P$ carry the
record's top sizes and weights to $a_\tau$ and $w(u_0,j,\rho)$.
Summing over $\tau$ and $j$ and comparing with
Definition~\ref{def:record-main} gives the identity.

\noindent\emph{Proof of \ref{item:mt-gap}.} This follows from the
definition of $\gap(u_0)$, part~\ref{item:mt-main} and the clause of
the identification for $\PS(u_0)$.
\end{proof}

\begin{definition}[uniform explicit-formula condition]
\label{def:uniform-ef}
Let $\lev$ be a level, $\eps>0$ and $\Fam$ a family at $\eps$
(Definition~\ref{def:family}). The family satisfies the \emph{uniform
explicit-formula condition} for $\lev$ if there are real numbers $C$
and $x_0$ such that for every $x\ge x_0$ and every $i\in I_x$ the host
of the package $\mathcal P_{x,i}$, with its top places $\Top$,
satisfies $\mathrm{EF}(\lev,\Top,x,\eps,C)$
(Definition~\ref{def:shares-ef}).
\end{definition}

\begin{definition}[ties]\label{def:tied}
Let $\eps>0$, let $\Fam$ be a family at $\eps$
(Definition~\ref{def:family}), and let $e_D,e_L,c_W,c_S,e_K,C_L$ and
$x$ be real numbers. The family is \emph{tied at $x$} with these
constants if there are
\begin{itemize}
\item an index $i_0\in I_x$, with $A$ the group of the frame of the
package $\mathcal P_{x,i_0}$ and $\GG=V_4\times A$;
\item positive integers $m$ and $f$, and a map $\mathrm{mem}\colon A\to
I_x$, the \emph{choice of members};
\item a slope record with shifts $U_0=A$ whose words $\Words(u_0)$ are
the words of the package $\mathcal P_{x,\mathrm{mem}(u_0)}$, and size
constants for it;
\item a function $\ell$ from the characters of $\GG$ to $\C$, an
integer $N\ge1$, finite sets $\Top,\Top_2\subset A$, maps
$a\colon A\to\R$ and
$w\colon A\times\{0,\dots,m-1\}\times\GG\to\R$, and real numbers
$\eta,X^\flat,\vartheta$,
\end{itemize}
such that the following hold for every $u_0\in A$.
\begin{enumerate}[label=(\arabic*)]
\item (Members.) The member $\mathcal P_{x,\mathrm{mem}(u_0)}$ has the
frame, the integers $m$ and $f$, and the primes $q$ and $d$ of
$\mathcal P_{x,i_0}$; its host has shift $u_0$; and the sizes that its
size datum assigns to a top place do not depend on $u_0$.
\item (Global ties.) The record's $a$, $b$, $b'$ and $\Upsilon$ are
those of the scale of the size datum $Z_{x,\mathrm{mem}(u_0)}$; $k$ is
the rank $k$ of the package (Definition~\ref{def:ranks}); $\Upsilon_a$,
$B_0$, $N_c$ and $X^\sharp$ are those of the size datum, and so are
$\eta$ and $X^\flat$; $\gamma^*=\gamma$; $\ell_K$ is the logarithmic
derivative at $1$ recorded by the size datum; $v$ is the half root
discriminant of the host (Definition~\ref{def:host}); and
$\lgv=\log(\varpi_0\varpi_1)$, with $\varpi_0,\varpi_1$
the norm primes of the host (Definition~\ref{def:norm-primes}).
\item (Local ties.) For every word $u\in\Words(u_0)$: $\rk(u)$ is the
rank of the lattice $N_u$ of the period lattice data
(Definition~\ref{def:period-lattice}); $T(u)$ is the level prime of
$u$; $\mu_N(u)$ and $\mu_N^{\max}(u)$ are the slope and the maximal
slope of $N_u$ (Definition~\ref{def:hlat}); $h(u)=h_u$; $\mu_M(u)$ and
$\mu_M^{\max}(u)$ are the slope and the maximal slope of the form
lattice of $u$; and
\[
\Cg_j(u)=\sum_{\mathfrak p}\operatorname{sz}_{\mathfrak p}
\bigl(\text{congruence invariant of $u$ at lift index $j+1$}\bigr),
\]
the sum over the places $\mathfrak p$ of the support of the size
datum, with $\operatorname{sz}_{\mathfrak p}$ its local sizes.
For the shift itself, $\mu_R(u_0)$ and $\mu_R^{\min}(u_0)$ are the slope and the
minimal slope of $R^\vee$; $\Fo(u_0)$ and $\Ga(u_0)$ are those of the
size datum; $\IO(u_0)$ is a sum over the support of
local quantities of the size datum; and $\beta_0(u_0)$ is the rank
defect of the member.
\item (Arithmetic ties.) $D$ are the arithmetic data built from the
Artin map of the host of $\mathcal P_{x,i_0}$; $\Top$, $\Top_2$ and
$a$ are those of the members; $w(u_0,j,\cdot)$ are the weights of
$\mathcal P_{x,\mathrm{mem}(u_0)}$ at lift index $j+1$; and
$\ell(\theta)=L'(1,\chi)/L(1,\chi)$ whenever $\theta\ne\theta_0$ and
$\chi$ is a primitive Dirichlet character with $\chi(n)=\theta_D(n)$
for every $n\ge1$.
\item (Constants.) The record's $x$ and $\eps$ are $x$ and $\eps$,
$e_{\mathrm{div}}\le e_D$, $e_{\mathrm{loss}}\le e_L$ and
$c_{\mathrm{sp}}\le c_S$.
\item\label{item:tie-hyp} (Hypotheses.) The record satisfies the size data
(Definition~\ref{def:size-data}); the window conditions with $c_W$ and
$e_K$, the identification clause and the bounds on $B_0$ and $N_c$
(Definition~\ref{def:window-hyps}); the size conditions of
Definition~\ref{def:thick} for every $u_0\in A$, and the thin-scale
bound; the bound on the values at $x$ with $C_L$; the cost conditions
with $f=|\Phi|$; and the rest-size conditions with $n=|A|$,
$n_T=|\Top|$, $n_{T_2}=|\Top_2|$ and $a_E=a$
(Definition~\ref{def:restsize}).
\end{enumerate}
\end{definition}

\noindent The record's prime sum $\PS(u_0)$ is fixed by the
identification clause in~\ref{item:tie-hyp}. The spread $\Delta_R$ and the
defects
$\beta_1,\beta_2$ are fixed by no tie; they are constrained only by
the hypotheses collected in~\ref{item:tie-hyp}. The definition names quantities
of
packages and of size data (shifts, level primes, form lattices,
congruence invariants, local sizes, the rank defect) that belong to the data of
a
package (Definition~\ref{def:package}) and of a family
(Definition~\ref{def:family}); the definitions there are schematic
(Section~\ref{sec:conventions}). Through
the ties, the window conditions and the rest-size conditions, a tied
family has $x^{1/2-\eps}\le m\le x^{1/2+\eps/2}$ and $xm\le k$ at
every scale at which it is tied.

\section{The gap bound}
\begin{figure}[tbp]
\centering
\input{book/figs/ch09-gap-flow.tex}
\end{figure}

\begin{theorem}[Alp\"oge; gap bound]\label{thm:gap}
Let $H$ be a finite abelian group of order $n$ and $\GG=V_4\times H$.
Let a slope record with shifts $U_0=H$, positive integers $m$ and $f$
and size constants be given (Definitions~\ref{def:record}
and~\ref{def:size-data}); write $x$, $\eps$, $a$, $b$ and $k$ for
the quantities of the record, and $c_{\mathrm{sp}}$,
$e_{\mathrm{div}}$ and $e_{\mathrm{loss}}$ for its size constants. Let
$c_W$ and $e_K$ be real numbers. Let $D$ be arithmetic data for $\GG$,
$\ell$ a function from the
characters of $\GG$ to $\C$, $N\ge1$ an integer, $\Top,\Top_2\subset
H$ finite sets, $a\colon H\to\R$ and
$w\colon H\times\{0,\dots,m-1\}\times\GG\to\R$.
Suppose that the record satisfies the size data
(Definition~\ref{def:size-data}), that the identification clause, the
window conditions with $c_W$ and $e_K$ and the bounds on $B_0$ and
$N_c$ (Definition~\ref{def:window-hyps}) hold, and that
\[
\log x\le x^\eps\qquad\text{and}\qquad
200(1+c_W)^2(1+c_{\mathrm{sp}})(1+e_K)\le x^\eps.
\]
Then for every $u_0\in H$
\[
\gap(u_0)\le x^{(e_{\mathrm{div}}+e_{\mathrm{loss}}+3)\eps}\,\Err,
\]
with $\Err$ the error bracket of Definition~\ref{def:restsize} for
$n=|H|$, $n_T=|\Top|$, $n_{T_2}=|\Top_2|$, the given $b$, $k$ and
$m$, and $a_E=a$.
\end{theorem}

\begin{proof}[Proof (outline)]
By Proposition~\ref{prop:main-terms}\ref{item:mt-gap},
$\gap(u_0)=\PS(u_0)-\Main(u_0)$,
so it suffices to bound $\PS(u_0)$ by $\Main(u_0)$ plus errors (Figure~\ref{fig:gap-flow}). Fix
$u_0$; all record quantities below are taken at $u_0$, and for a
function $X$ on the words we write
\[
\langle X\rangle=\frac{\sum_{u\in\Words(u_0)}\rk(u)X(u)}
{\sum_{u\in\Words(u_0)}\rk(u)},
\]
so that $s=\langle\mu_N-h\rangle$. Put
\[
\begin{aligned}
A(u)&=\sum_{\tau\in\Top}\Bigl(m\bigl(a_\tau\log a_\tau
-a_\tau-a_\tau\log\Omega\bigr)-a_\tau\log T(u)\Bigr),\\
E_{\mathrm{div}}&=c_{\mathrm{div}}n_Ta\Upsilon_a+c_{\mathrm{div}}
m(B_0+N_c)x^{e_{\mathrm{div}}\eps},
\end{aligned}
\]
and let $\sigma=\log|\Words(u_0)|+(1-\beta_0)^{-1}\bigl(4u_F\log(1+
\Delta_R/u_F)+(\beta_1+4\sqrt{\beta_2})\Delta_R\bigr)$ be the slack of
clause~\ref{item:sd-slope}. The clauses below are those of
Definition~\ref{def:size-data}.
\begin{enumerate}[label=(\arabic*)]
\item\label{item:gap-s1} Clause~\ref{item:sd-loss} gives $\PS\le\Ga+\Lo_P$.
\item The second inequality of clause~\ref{item:sd-qs} and
clause~\ref{item:sd-div} give, for
every word $u$,
$\mu_N(u)-h(u)\ge\Ga+\mu_R+\mu_M(u)-\Cg(u)-\Lo_Q(u)-A(u)
-E_{\mathrm{div}}$; the term $\Fo$ cancels. Averaging,
\begin{equation}\label{eq:sandwich-lower}
s\ge\Ga+\mu_R-\bigl\langle A+\Cg+\Lo_Q-\mu_M\bigr\rangle
-E_{\mathrm{div}}.
\end{equation}
The hypotheses of clause~\ref{item:sd-slope} hold by clauses~\ref{item:sd-pos},
\ref{item:sd-spread} and~\ref{item:sd-def}, so
\begin{equation}\label{eq:sandwich-upper}
s\le\mu_R+\sigma.
\end{equation}
Subtracting, $\Ga\le\sigma+E_{\mathrm{div}}+\langle A+\Cg+\Lo_Q
-\mu_M\rangle$, an inequality in which neither $\mu_R$ nor the
heights $h(u)$ appear (Remark~\ref{rem:cancel}).
\item\label{item:gap-s3}
With $\log\Omega=\gamma^*+\log(\pi/v)-\ell_K$, the lower bound
for $\log T(u)$ in clause~\ref{item:sd-ct}, multiplied by $\Sigma a$, and the
bound
$\Cg(u)\le m\,\Sigma a\,\lgv+c_{\mathrm{cg}}mB_0$ of clause~\ref{item:sd-cg}
give
\[
A(u)+\Cg(u)\le\Main(u_0)+c_{\mathrm{cg}}mB_0+c_{\mathrm{ct}}\,\Sigma a
\Bigl(\Upsilon_a+\frac{m^2\Upsilon_a}{k}\Bigr):
\]
the terms $m\,\Sigma a\log(\pi/v)$ and $m\,\Sigma a\log(v/\pi)$
cancel, the terms $\pm m\,\Sigma a\,\lgv$ cancel, and what remains of
the archimedean sum and of the weights in clause~\ref{item:sd-ct} is
$\Main(u_0)$. Clause~\ref{item:sd-fs} bounds $-\mu_M(u)$ by
$c_{\mathrm{fs}}n_Ta\Upsilon_a(1+m/k)$.
\item Combining \ref{item:gap-s1}--\ref{item:gap-s3},
\[
\begin{aligned}
\PS\le{}&\Main(u_0)+\Lo_P+\sigma+E_{\mathrm{div}}
+c_{\mathrm{cg}}mB_0+\langle\Lo_Q\rangle\\
&+c_{\mathrm{ct}}\Sigma a\Bigl(\Upsilon_a+\frac{m^2\Upsilon_a}k
\Bigr)+c_{\mathrm{fs}}n_Ta\Upsilon_a\Bigl(1+\frac mk\Bigr).
\end{aligned}
\]
The error terms on the right are then estimated under the window
conditions and the bounds on $B_0$ and $N_c$ (with
$\IO\le\Lo_P$ from clause~\ref{item:sd-loss} and $\log T(u)\le c_{\mathrm{ct}}
m\Upsilon$ from clause~\ref{item:sd-ct} inside $\Lo_Q$). This estimation, linear
arithmetic in which the largeness of $x$ absorbs the constants, gives
the bracket $x^{(e_{\mathrm{div}}+e_{\mathrm{loss}}+3)\eps}\Err$.
\end{enumerate}
\end{proof}

\begin{remark}[the test lattice cancels]\label{rem:cancel}
The slope $\mu(R^\vee)$ of the test lattice, and with it its
covolume, does not appear in the final inequality, although the
periods live in the lattice $R^\vee$. The
two lattice inequalities are the bounds~\eqref{eq:sandwich-lower}
and~\eqref{eq:sandwich-upper} for the frame slope $s$:
\[
\Ga+\mu_R-\bigl\langle A+\Cg+\Lo_Q-\mu_M\bigr\rangle-E_{\mathrm{div}}
\ \le\ s\ \le\ \mu_R+\sigma.
\]
The left bound comes from the divisor bound (clause~\ref{item:sd-div}) and the
quotient slopes (clause~\ref{item:sd-qs}), the right one from the slope
inequality (clause~\ref{item:sd-slope}). Both are bounds on the one number
$\mu_R=\mu(R^\vee)$ in terms of $s$, and subtracting them removes it
(Figure~\ref{fig:sandwich}).
The heights $h(u)$ cancel as well, since they enter $s$ and
clause~\ref{item:sd-div} with opposite signs. What survives is an inequality
between the arithmetic quantity $\Ga$, which bounds the prime sum, and
the archimedean main term; one may read it as ``finite-side gain at
most archimedean main term''.
\end{remark}

\begin{figure}[htbp]
\centering
\input{alpoge/figs/v2/sandwich.tex}
\caption{Clause~\ref{item:sd-div} of Definition~\ref{def:size-data} bounds the
slope
$\mu(R^\vee)$ of the test lattice from above; here
$A(u)=\sum_\tau\bigl(m(a_\tau\log a_\tau-a_\tau-a_\tau\log\Omega)
-a_\tau\log T(u)\bigr)$ is the archimedean sum. Clause~\ref{item:sd-slope}
bounds it
from below in terms of the frame slope $s$. Subtracting, $\mu(R^\vee)$
and the heights cancel (Remark~\ref{rem:cancel}). Bost's inequality
would compare with $\mu_{\max}(R^\vee)$ instead
(Remark~\ref{rem:not-bost}). Positions on the line are
illustrative.}
\label{fig:sandwich}
\end{figure}

\section{The square-root bookkeeping}

\begin{proposition}[the bookkeeping behind $m\approx\sqrt x$]
\label{prop:sqrtx}
Let $x\ge1$, $\eps>0$, $b>0$, $k>0$ and $a_E\ge0$ be real numbers and
$n,n_T,n_{T_2},m$ positive integers. Suppose that $a_E\le2xb$,
$n_T\ge x^{-\eps/8}n/2$, $x^{1/2-\eps}\le m\le x^{1/2+\eps/2}$,
$k\le n^2x^{1/2-\eps}$, $xm\le k$ and $n_{T_2}\le n$. Then the error
bracket $\Err$ of Definition~\ref{def:restsize} satisfies
\[
\frac{\Err}{n_Tmb}\le\frac{2x}m+2+4m\,x^{9\eps/8}+2x^{1/2+9\eps/8}
\le10\,x^{1/2+13\eps/8}.
\]
\end{proposition}

\begin{proof}
Divide each of the three terms of $\Err$ by $n_Tmb$. Using
$a_E\le2xb$ and $xm\le k$,
\[
\frac{a_E(1+m^2/k)}{mb}\le\frac{2x}m+\frac{2xm}k\le\frac{2x}m+2.
\]
Using $1/n_T\le2x^{\eps/8}/n$,
\[
\frac{2nmx^\eps}{n_T}\le4m\,x^{9\eps/8}.
\]
Using also $n_{T_2}\le n$ and $k\le n^2x^{1/2-\eps}$,
\[
\frac{n_{T_2}\,k\,x^{2\eps}}{n^2n_T}\le2x^{1/2+9\eps/8}.
\]
This is the first inequality. For the second, $m\ge x^{1/2-\eps}$
gives $2x/m\le2x^{1/2+\eps}$, and $m\le x^{1/2+\eps/2}$ gives
$4mx^{9\eps/8}\le4x^{1/2+13\eps/8}$; since $x\ge1$, each of the four
terms is at most a multiple of $x^{1/2+13\eps/8}$, with multiples
$2,2,4,2$.
\end{proof}

\noindent The exponent $1/2$ is the balance point of the terms $x/m$
and $m$: the largest of $x/m$, $m$ and $x^{1/2}$ is smallest at
$m=x^{1/2}$, and it is $x^{1/2}$ there (Figure~\ref{fig:exponents}). The
constraints on $m$ and $k$
are hypotheses: for the families of Chapter~\ref{sec:assembly} they
are part of the ties (Definition~\ref{def:tied}). The proof of
Theorem~\ref{thm:core-bound} divides by $n_Tmb\,\delta_2/2$ with
$\delta_2=n_{T_2}/n$ instead of $n_Tmb$. Since $x^{-\eps}\le
n_{T_2}/n$ by the rest-size conditions, the extra factor satisfies
$2/\delta_2\le2x^\eps$, which changes only the multiple of $\eps$ in
the exponent.

\begin{figure}[htbp]
\centering
\input{alpoge/figs/v2/exponents.tex}
\caption{Exponents of $x$ in the three terms of
$\Err/(n_Tmb)\ll(x/m+m+x^{1/2})x^{O(\eps)}$
(Proposition~\ref{prop:sqrtx}), with $m=x^\mu$ and the factors
$x^{O(\eps)}$ ignored. The largest term is smallest at $m=x^{1/2}$.
The shaded band is the range $x^{1/2-\eps}\le m\le x^{1/2+\eps/2}$,
drawn at $\eps=0.05$; this range is part of the hypotheses of
Theorem~\ref{thm:chebyshev-windows}, through the ties of
Definition~\ref{def:tied}, and it is met by the families of
Theorem~\ref{thm:tied-families}.}
\label{fig:exponents}
\end{figure}

\section{The Chebyshev inequality on windows}
\begin{figure}[tbp]
\centering
\input{book/figs/ch09-exponent-ledger.tex}
\end{figure}

\begin{theorem}[Alp\"oge; Chebyshev inequality on windows]
\label{thm:chebyshev-windows}
Let $\lev$ be a level, $e_D$ and $e_L$ real numbers and $\eps>0$. Let
$\Fam$ be a family at $\eps$ that satisfies the uniform
explicit-formula condition for $\lev$ (Definition~\ref{def:uniform-ef}).
Suppose that there are real numbers $c_W,c_S,e_K,C_L$ and $x_0$ such
that $\Fam$ is tied (Definition~\ref{def:tied}) with constants
$e_D,e_L,c_W,c_S,e_K,C_L$ at every $x\in\Win$ with $x\ge x_0$. Then
there is $x_1$ such that
\[
G_\lev(x)\le x^{1/2+(e_D+e_L+7)\eps}
\]
for every $x\in\Win$ with $x\ge x_1$.
\end{theorem}

\begin{proof}[Proof (outline)]
The first step discards the geometric data (Figure~\ref{fig:exponent-ledger}). At each $x\in\Win$
with $x\ge x_0$, the
tie provides a slope record, arithmetic data of a host and the data
$\ell,N,\Top,\Top_2,a,w,\eta,X^\flat,\vartheta$; dropping the ties
themselves and the bound on the values at $x$ leaves hypotheses on
real numbers only. Since $e_{\mathrm{div}}\le e_D$ and
$e_{\mathrm{loss}}\le e_L$, Theorem~\ref{thm:gap} bounds the gap by
$x^{(e_D+e_L+3)\eps}\Err$ at every shift, once $x$ is large in terms
of $c_W,c_S,e_K$. Theorem~\ref{thm:core-bound}, with $C=e_D+e_L+3$,
then gives $G_A(x)\le x^{1/2+(e_D+e_L+6)\eps}+\Sigma_A(x)$, where the
constants of the character conditions come from the uniform
explicit-formula condition and from the level, so that one $x_1$
serves for all $x$. Theorem~\ref{thm:averaged-hyps} identifies $G_A$
with $G_\lev$ and bounds $\Sigma_A(x)$ by $4o\,x^{1/2+\eps}$. So
\[
G_\lev(x)\le x^{1/2+(e_D+e_L+6)\eps}+4o\,x^{1/2+\eps},
\]
which is at most $x^{1/2+(e_D+e_L+7)\eps}$ for large $x$, since
$e_D,e_L\ge0$ by the window conditions.
\end{proof}

\begin{remark}[no vacuity]\label{rem:no-vacuity}
Theorem~\ref{thm:chebyshev-windows} would say nothing for a family whose
window set has no large elements, or if the exponents $e_D$ and $e_L$
could grow as $\eps\to0$. Neither happens. The window set of a family
is unbounded (Definition~\ref{def:family}). The same families and the
same ties occur in the existence statement
Theorem~\ref{thm:tied-families}, where the window set is required to
have windows of a fixed exponent, and the exponents $e_D$ and $e_L$
are fixed before $\eps$ is sent to $0$ (Figure~\ref{fig:eps}).
\end{remark}

\begin{remark}[where the exponent $1/2$ comes from]
\label{rem:exponent-half}
Chebyshev's identity alone gives $\Ch(x)\asymp x$, so a one-sided bound
with square-root error must come from somewhere else. In the present
argument it comes from the way the errors of the slope record depend
on the number $m$ of theta lifts. The main term $\Main(u_0)$ is a sum
of $m$ copies of an archimedean term, one for each lift, and so is of
size about $m\,n_Ta\log a$. The error bracket $\Err$ has two kinds of
term. The first, $n_Ta_E(1+m^2/k)$, does not grow with $m$ once
$k\ge xm$; it collects errors such as $c_{\mathrm{div}}n_Ta\Upsilon_a$
in the divisor bound (clause~\ref{item:sd-div}), which do not depend on $m$. The
other kind, of the form $m(B_0+N_c)$, grows like $m^2$ through
$B_0\le2nmbx^\eps$, while its part $mN_c$ gives, through the bound on
$N_c$ and $k\le n^2x^{1/2-\eps}$, the term of size $x^{1/2}$ of
Proposition~\ref{prop:sqrtx}. After division by $n_Tmb$ the first kind
and the $B_0$ part become the terms $x/m$ and $m$ of
Proposition~\ref{prop:sqrtx}, and their balance at $m\approx x^{1/2}$
is the exponent $1/2$. If the errors of the first kind grew linearly
in $m$, the first term would be of size $x$, the size given by
Chebyshev's identity.

The inequality that the argument proves is, through
Proposition~\ref{prop:balance} and Remark~\ref{rem:legendre}, a
one-sided bound for a Legendre-type sum over the primes inert in $K$,
with the exact main term and an error $x^{1/2+O(\eps)}$ after
normalisation. The main term enters with coefficient~$1$ through an
identity (Proposition~\ref{prop:balance}) and through the archimedean
side of clause~\ref{item:sd-div}; no sieve is used.
\end{remark}

\FloatBarrier
\clearpage
\clearpage
\chapter{Tate--Shafarevich groups of Tate twists}
\label{ch:sha}\label{sec:sha}

This chapter states a uniform bound for the Tate--Shafarevich groups
of the Tate twists $\Q_l/\Z_l(-e)$ over the host fields, the Sha bound
on windows (Theorem~\ref{thm:sha-windows}), and the results from which
it is put together: a case split of the eigenspaces, proved in full
(Proposition~\ref{prop:cases-exhaustive}), a bound for the minus part
(Theorem~\ref{thm:minus}), bounds for the plus part
(Theorems~\ref{thm:plus-gk}, \ref{thm:plus}
and~\ref{thm:remaining-counts}), two auxiliary theorems
(Theorems~\ref{thm:bloch-kato} and~\ref{thm:h2-bound}) and a
pigeonhole theorem that produces the windows
(Theorem~\ref{thm:windows-pigeonhole}). Most of these theorems are
stated with an outline or a description of the route of proof only;
each proof environment says which. In Chapter~\ref{ch:assembly} the
Sha bound on windows is a black box: it is the second step of the
proof of the main theorem, its automorphy input is
Theorem~\ref{thm:galois-reps} of Chapter~\ref{ch:galois}
(Remark~\ref{thm:lgc-to-sha}), and its output is the hypothesis of the
supply of hosts, Theorem~\ref{thm:host-supply} of
Chapter~\ref{ch:families}.

The congruence defects at the twisted places of the families of
Chapter~\ref{sec:assembly} are bounded by lengths of Selmer groups, and
these in turn by a bound on the Tate--Shafarevich groups of the Tate
twists $\Q_l/\Z_l(-e)$ over the host fields. The bound is needed
uniformly in the auxiliary pair $(q,d)$, and it is generous in size:
its total is allowed to be $\exp(x^\eps/2)$, while the classical
heuristics predict a size polynomial in $x$. The difficulty is
uniformity, and above all the eigenspaces of the wrong parity, where no
$L$-value controls the group. For a bounded number of exceptional Tate
twists the only bound available is doubly exponential, and a pigeonhole
argument then gives the bound only on windows of scales. This is where
the windows of Chapter~\ref{sec:landau} come from.

The logical structure of the chapter is shown in
Figure~\ref{fig:sha-deps}. No result in it carries an unproved
hypothesis in its statement. The proofs are Alp\"oge's claims
(Remark~\ref{rem:status}). The automorphy that the plus-part bounds
use is Theorem~\ref{thm:galois-reps}\ref{item:gr-lgc}, proved in
Chapter~\ref{sec:galois}, and the output feeds the supply of hosts in
Chapter~\ref{sec:assembly}.

\begin{figure}[htbp]
\centering
\input{alpoge/figs/v2/sha-deps.tex}
\caption{Dependencies in Chapter~\ref{sec:sha}. Shaded boxes are
results stated in this chapter; the white box collects classical
inputs. An arrow points from a result to a result whose proof uses
it.}
\label{fig:sha-deps}
\end{figure}

Throughout this chapter $l$ denotes a prime number, the coefficient
prime; it is kept apart from the function $\ell$ of
Chapter~\ref{sec:balance}.

\section{The groups and the bound}

\begin{definition}[Tate--Shafarevich groups of Tate twists]
\label{def:sha}
Let $E$ be a number field with an algebraic closure $\overline E$, let
$l$ be a prime, and let $e\in\Z$. Let $S_l$ be the set of places of
$E$ above $l$, and let $G_{E,S_l}$ be the Galois group of the maximal
extension of $E$ in $\overline E$ unramified outside $S_l$ and the
infinite places. Choose a decomposition group at each place of $S_l$;
the group defined below, and its length, do not depend on these
choices or on the choice of $\overline E$, up to isomorphism.
\begin{enumerate}[label=(\arabic*)]
\item The \emph{Tate--Shafarevich group} of the twist $e$ is
\[
\Sha_l(E,e)=\Sha^1\bigl(G_{E,S_l},\Q_l/\Z_l(-e)\bigr),
\]
the group of classes in $H^1(G_{E,S_l},\Q_l/\Z_l(-e))$ whose
restrictions to the chosen decomposition groups at $S_l$ vanish.
\item Its \emph{length} $\shalen_l(e)\in\Z_{\ge0}\cup\{\infty\}$ is the
length of $\Sha_l(E,e)$ as a $\Z_l$-module.
\end{enumerate}
\end{definition}

\begin{definition}[Sha budgets and Sha bounds on windows]
\label{def:sha-windows}
Let $t\ge1$ be an integer and $B\ge0$ a real number, and let $\psi$
be a primitive Dirichlet character.
\begin{enumerate}[label=(\arabic*)]
\item A number field $E$ has \emph{Sha budget} $(t,B)$ if
$\shalen_p(e)<\infty$ for every prime $p$ and every
$0<e<t$, and
\[
\sum_{l\in S}\log l\sum_{1\le e<t}\shalen_l(e)\le B
\]
for every finite set $S$ of primes.
\item Let $C_t>0$. The character $\psi$ has \emph{Sha bounds on windows
with constant $C_t$} if the following holds for every
$\eps\in(0,1/52)$ and every function $\nu$ on $[1,\infty)$ such that
$\nu(x)$ is adapted to $(\eps,x)$ for every $x\ge1$
(Definition~\ref{def:scales}). There are $J\ge1$ and $X_0$ such that
for every $X\ge X_0$ and every window pair $(q,d)$ for
$(\eps,X^{1/J})$ (Definition~\ref{def:scales}) there is a real $Y$
with $X^{1/J}\le Y$ and $Y^2\le X$ such that, for every
$x\in[Y,Y^2]$, every $\psi$-field for $(\psi,\nu(x),q,d)$
(Definition~\ref{def:psifield}) has Sha budget
$\bigl(\lfloor x^{C_t}\rfloor,\exp(x^\eps)^{1/2}\bigr)$.
\item The \emph{Sha bound on windows} is the statement that every
primitive Dirichlet character has Sha bounds on windows with constant
$401$.
\end{enumerate}
\end{definition}

\noindent In words: on some window $[Y,Y^2]$ inside every range
$[X^{1/J},X]$, the CM fields of degree about $x^{2+8\eps}$ attached to
$\psi$ satisfy
$\sum_l\log l\sum_{1\le e<x^{401}}\shalen_l(e)\le\exp(x^\eps/2)$. The
window $[Y,Y^2]$ may depend on $X$ and on the pair $(q,d)$, but the
exponent $J$ may not.

\begin{remark}[duality and $K$-theory]\label{rem:poitou-tate}
By Poitou--Tate duality \cite{Tate1963,Poitou1967} the group
$\Sha_l(E,e)$ is dual to
$\Sha^2(G_{E,S_l},\Z_l(e+1))$, which is a subgroup of
$H^2(G_{E,S_l},\Z_l(e+1))$. For odd $l$ the Quillen--Lichtenbaum
comparison identifies this $H^2$ with the $l$-part of
$K_{2e}(\mathcal O_E)$; the comparison rests on the theorem of
Voevodsky and Rost (the Bloch--Kato conjecture, proved for odd $l$ in
\cite{Voevodsky2011}). So the Sha bound on windows is, in classical
terms, a bound on the subgroup $\Sha^2$ of the $l$-part of
$K_{2e}(\mathcal O_E)$, the kernel of localisation at the places
above $l$, uniform in $l$; it is not a bound on the whole $l$-part.
Figure~\ref{fig:duality} sets out these
identifications.
\end{remark}

\begin{figure}[htbp]
\centering
\input{alpoge/figs/v2/duality.tex}
\caption{The classical dictionary for the groups bounded by the Sha
bound on windows (Definition~\ref{def:sha-windows}): Poitou--Tate
duality and, dashed, the Quillen--Lichtenbaum identification for odd
$l$, which rests on the theorem of Voevodsky and Rost
(Remark~\ref{rem:poitou-tate}). The size estimate is heuristic.}
\label{fig:duality}
\end{figure}

\begin{remark}[the size of the budget]\label{rem:budget}
Heuristically $\log|K_{2e}(\mathcal O_E)|\approx e[E:\Q]\log(e\,
\mathrm{rd}_E)$, where $\mathrm{rd}_E$ is the root discriminant. With
$[E:\Q]\approx x^{2+8\eps}$ and $e<x^{401}$ the total is polynomial in
$x$, far below the allowed $\exp(x^\eps/2)$. The bound therefore has
a wide margin in size. It must hold uniformly in $(q,d)$, and for
that the eigenspaces of the wrong parity must be controlled. This step can be
compared with Vandiver's
conjecture (Remark~\ref{rem:vandiver}).
\end{remark}

\section{Eigenspaces and the case split}

The host fields are abelian over $\Q$, so their Tate--Shafarevich
groups split into eigenspaces for the characters of the Galois group,
and complex conjugation acts on each eigenspace by a sign. The bound is
proved eigenspace by eigenspace, and the eigenspaces fall into a minus
part and six kinds of plus part.

As elsewhere in these notes, $\psi$ is the primitive character from
which a host is built. In this section and the next, $\eta$ denotes
a Dirichlet character of the host modulus whose eigenspace is
considered, an \emph{eigen-character}.

\begin{definition}[abelian data and their fields]
\label{def:abelian-datum}
Let $N_M\ge1$ be an integer.
\begin{enumerate}[label=(\arabic*)]
\item An \emph{abelian datum} $M$ of modulus $N_M$ is a subgroup
$H_M\le(\Z/N_M\Z)^\times$. Its \emph{degree} is the index
$[(\Z/N_M\Z)^\times:H_M]$.
\item A \emph{prime pair} for $M$ is a pair of primes $(q,d)$ with
$q\equiv3\pmod8$, $d\equiv7\pmod{16}$ and $\gcd(qd,N_M)=1$.
\item Let $\nu\ge0$ be an integer and $(q,d)$ a pair of primes. A
number field $E$ is a \emph{field for} $(M,\nu,q,d)$ if there are an
integer $N$ divisible by $2^{\nu+2}$, $N_M$, $q$ and $d$ and an embedding
$E\hookrightarrow\Q(\zeta_N)$ whose image is the fixed field of the
classes $a\in(\Z/N\Z)^\times$ with $a\equiv\pm1\pmod{2^{\nu+2}}$,
$a\bmod N_M\in H_M$, $(a/q)=1$ and $(a/d)=1$.
\item The \emph{conductor} of an abelian number field $E$ is the least
$N\ge1$ such that $E$ embeds into $\Q(\zeta_N)$.
\end{enumerate}
\end{definition}

\noindent For a primitive character $\psi$ of conductor $N_\psi$ and the
datum $M=(N_\psi,\ker\psi)$, the fields for $(M,\nu,q,d)$ are the
$\psi$-fields of Definition~\ref{def:psifield}.

\begin{definition}[eigenspaces and cases]\label{def:cases}
Let $E$ be a CM field, $l$ a prime and $e\in\Z$.
\begin{enumerate}[label=(\arabic*)]
\item The \emph{plus part} and the \emph{minus part} of
$\Sha_l(E,e)$ are its eigenspaces for complex conjugation with
eigenvalues $1$ and $-1$.
\item A character $\eta$ of $\Gal(E/\Q)$ is \emph{critical} for $e$ if
$\eta(-1)=(-1)^{e+1}$, and \emph{plus} for $e$ if $\eta(-1)=(-1)^e$.
\item A \emph{host datum} $\mathfrak h$ consists of a primitive Dirichlet
character $\psi$ of conductor $N_\psi$ and order $o_\psi$, an integer
$\nu\ge2$ and a prime pair $(q,d)$ for $(N_\psi,\ker\psi)$. Its
\emph{modulus} is $N_{\mathfrak h}=\operatorname{lcm}(2^{\nu+2},N_\psi)\,qd$,
and
$H_{\mathfrak h}\le(\Z/N_{\mathfrak h}\Z)^\times$ is the subgroup of classes $u$
with
$u\equiv\pm1\pmod{2^{\nu+2}}$, $\psi(u)=1$, and $u$ a square modulo
$q$ and modulo $d$.
\item Let $p_K$ be the least prime with $p_K\equiv3\pmod4$, $p_K\ge7$
and $p_K\nmid N_\psi$; let $h^*$ be the class number of
$\Q(\sqrt{-p_K})$; and let $p_1$ be the least prime that is
$\equiv1$ modulo the odd part of $o_\psi$ and prime to $2N_\psi p_K$.
The \emph{bad set} $B_\psi$ is the set of odd primes dividing
$o_\psi N_\psi p_Kh^*p_1$, together with $3$.
\item The host datum $\mathfrak h$ is \emph{large for} an integer $t\ge1$ if
$2t+1<q$, $2t+1<d$, and each of $N_\psi$, $o_\psi$, $p_K$, $h^*$, $p_1$
is less than $q$ and less than $d$.
\item Let $N\ge1$ and let $r$ be a prime. Write
$(\Z/N\Z)^\times=(\Z/r^{v_r(N)}\Z)^\times\times(\Z/N'\Z)^\times$ with
$r\nmid N'$, and let $\operatorname{pr}_r$ be the endomorphism of
$(\Z/N\Z)^\times$ that is the identity on the first factor and trivial
on the second. For a character $\theta$ modulo $N$ put
$\theta|_r=\theta\circ\operatorname{pr}_r$. For a character $\eta$
modulo $N$, let $\eta^{(l')}$ be the prime-to-$l$ part of $\eta$ (the component
of order prime to $l$ in
the decomposition of $\eta$ into its $l$-part and its prime-to-$l$
part), and put
\[
\chi'_l(\eta)=\eta^{(l')}\cdot(\eta^{(l')}|_l)^{-1}.
\]
\item Let $\mathfrak h$ be a host datum, $t\ge1$, $\eta$ a character modulo
$N_{\mathfrak h}$ trivial on $H_{\mathfrak h}$, $l$ an odd prime and $0<e<t$.
The triple
$(l,e,\eta)$ is in
\begin{enumerate}[label=(\alph*)]
\item \emph{Case~I} if $l\notin B_\psi$ and $(l-1)\nmid2e$;
\item \emph{Case~Z} if $l\notin B_\psi$, $(l-1)\mid2e$ and $\eta=1$;
\item \emph{Case~II} if $e\ge2$, $\chi'_l(\eta)$ has even order, and
$l\in B_\psi$ or $(l-1)\mid2e$;
\item \emph{Case~III} if $l\notin B_\psi$, $(l-1)\mid2e$, and
$\eta\ne1$ has odd order;
\item \emph{Case~IV} if $l\in B_\psi$, $e\ge2$ and $\chi'_l(\eta)$ has
odd order;
\item \emph{Case~V} if $l\in B_\psi$ and $e=1$.
\end{enumerate}
\end{enumerate}
\end{definition}

\noindent The projection $\operatorname{pr}_r$ keeps the component in
the factor $(\Z/r^{v_r(N)}\Z)^\times$ of the Chinese remainder
decomposition. It is not the projection onto the Sylow $r$-subgroup:
the factor $(\Z/r^{v_r(N)}\Z)^\times$ has order prime to $r$ in
general. So $\chi'_l(\eta)$ is the prime-to-$l$ part of $\eta$ with its
component at the modulus $l^{v_l(N)}$ removed.

\begin{proposition}\label{prop:cases-exhaustive}
Let $\mathfrak h$ be a host datum with primitive character $\psi$, let $t\ge1$
be an integer for which $\mathfrak h$ is large, let $\eta$ be a character modulo
$N_{\mathfrak h}$ trivial on $H_{\mathfrak h}$, let $l$ be an odd prime and let
$0<e<t$.
\begin{enumerate}[label=(\alph*)]
\item\label{item:cases-good} If $l\notin B_\psi$, then $l$ does not
divide the order of $\eta$ and $\eta^{(l')}=\eta$; if in addition
$(l-1)\mid2e$, then $e\ge2$, $l\nmid N_{\mathfrak h}$ and $\chi'_l(\eta)=\eta$.
The triple $(l,e,\eta)$ is in Case~I, Z, II or~III.
\item\label{item:cases-bad} If $l\in B_\psi$, the triple $(l,e,\eta)$
is in Case~II, IV or~V.
\end{enumerate}
\end{proposition}

\noindent The prime $l=2$ lies outside the case split, and the minus
part is treated separately (Theorem~\ref{thm:minus}). The
$2$-parts of the lengths are bounded by the first clause of
Theorem~\ref{thm:windows-pigeonhole}. Figure~\ref{fig:case-plane} shows, for one $\psi$, where the cases lie in the plane of $(l,e)$, and Figure~\ref{fig:cases} shows the
case split with the result that bounds each case.
\begin{figure}[tbp]
\centering
\input{book/figs/ch10-case-plane.tex}
\end{figure}

\begin{proof}[Proof of \ref{item:cases-good}]
Let $l$ be an odd prime not in $B_\psi$. Since $3\in B_\psi$, we have
$l\ge5$, and $l$ divides none of $o_\psi$, $N_\psi$.

We first show that the order of $\eta$ divides $2^{\nu+1}o_\psi$. The
group $H_{\mathfrak h}$ is the intersection of the kernels of four homomorphisms
on
$(\Z/N_{\mathfrak h}\Z)^\times$, namely reduction modulo $2^{\nu+2}$ followed by
the
quotient by $\{\pm1\}$, the character $\psi$, and the quadratic
residue symbols modulo $q$ and $d$. Hence $(\Z/N_{\mathfrak
h}\Z)^\times/H_{\mathfrak h}$
embeds into a product of groups of exponents dividing $2^\nu$,
$o_\psi$, $2$ and $2$, and its exponent divides
$\operatorname{lcm}(2^\nu,o_\psi,2)$, hence $2^{\nu+1}o_\psi$. A character
trivial on $H_{\mathfrak h}$ factors through this
quotient. As $l$ is odd and $l\nmid o_\psi$, $l$ does not divide the
order of $\eta$, so the $l$-part of $\eta$ is trivial and
$\eta^{(l')}=\eta$.

Suppose now that $(l-1)\mid2e$. Since $l\ge5$ and $e>0$, we get
$4\le l-1\le2e$, so $e\ge2$; and $l-1\le2e<2t$, so $l\le2t<q$ and
$l<d$ because $\mathfrak h$ is large for $t$. The odd prime divisors of
$N_{\mathfrak h}$
divide $N_\psi qd$, and $l$ divides none of $N_\psi$, $q$, $d$; hence
$l\nmid N_{\mathfrak h}$. Then the factor $(\Z/l^{v_l(N_{\mathfrak
h})}\Z)^\times$ is trivial,
$\operatorname{pr}_l$ is trivial, $\eta^{(l')}|_l=1$ and
$\chi'_l(\eta)=\eta^{(l')}=\eta$.

Now sort the triple. If $(l-1)\nmid2e$, it is in Case~I. Otherwise
$(l-1)\mid2e$. If $\eta=1$, the triple is in Case~Z. If $\eta\ne1$
has odd order, it is in Case~III. If $\eta$ has even order, then
$\chi'_l(\eta)=\eta$ has even order and $e\ge2$, so it is in
Case~II.
\end{proof}

\begin{proof}[Proof of \ref{item:cases-bad}]
Let $l\in B_\psi$. If $e=1$ the triple is in Case~V. Otherwise
$e\ge2$, and the triple is in Case~II if $\chi'_l(\eta)$ has even order
and in Case~IV if it has odd order.
\end{proof}

\begin{remark}[the parity test is on $\chi'_l(\eta)$]
\label{rem:cases-parity}
Cases~II and~IV test the parity of the order of $\chi'_l(\eta)$.
They do not test the order of $\eta$ or of its $l$-part. For odd $l$ the
$l$-part of $\eta$ has odd order, so
a test on it could never hold in Case~II, and the list would miss the
triples in the proof of Proposition~\ref{prop:cases-exhaustive}
\ref{item:cases-good} with $(l-1)\mid2e$ and $\eta$ of even order.
Only Case~III tests the order of $\eta$ itself.
\end{remark}

\begin{figure}[htbp]
\centering
\input{alpoge/figs/v2/cases.tex}
\caption{The case split of Definition~\ref{def:cases} for the plus part
($l$ an odd prime, $0<e<t$, the host large, $\eta$ trivial on $H_{\mathfrak
h}$),
with the result that bounds each case. The cases are exhaustive
(Proposition~\ref{prop:cases-exhaustive}). The prime $l=2$ lies
outside the split, and the minus part is bounded separately
(Theorem~\ref{thm:minus}).}
\label{fig:cases}
\end{figure}

\begin{example}[the projection is not the Sylow projection]
\label{ex:crt}
Let $N=63$ and $l=3$, so that
$(\Z/63\Z)^\times\cong(\Z/9\Z)^\times\times(\Z/7\Z)^\times\cong
C_6\times C_6$, and write the characters as $\alpha^a\beta^b$, with
$\alpha$ a generator of the characters of the first factor and $\beta$
of the second. Take $\eta=\alpha$, of order $6$. Its prime-to-$3$ part
is $\eta^{(l')}=\alpha^3$, a character of the factor $(\Z/9\Z)^\times$,
so $\eta^{(l')}|_3=\alpha^3$ and $\chi'_3(\eta)=1$ has odd order, while
$\eta$ has even order. Projecting to the Sylow $3$-subgroup instead
would return $\eta^{(l')}=\alpha^3$, of even order. Table~\ref{tab:crt}
lists further characters; the host conditions are ignored here.
\end{example}

\begin{table}[htbp]
\centering
\input{alpoge/figs/v2/tab-crt.tex}
\caption{For $N=63$ and $l=3$:
$(\Z/63)^\times\cong(\Z/9)^\times\times(\Z/7)^\times\cong C_6\times
C_6$, with characters $\alpha^a\beta^b$. Here $\eta^{(l')}$ is the
prime-to-$3$ part of $\eta$, $\eta^{(l')}|_l$ its composite with the
projection to the factor $(\Z/9)^\times$, and
$\chi'_l(\eta)=\eta^{(l')}(\eta^{(l')}|_l)^{-1}$. In the first row $\eta$
has even order and $\chi'_l(\eta)$ is trivial, so the parity test of
Cases~II and~IV differs from a test on $\eta$. The last two columns
show the result of projecting to the Sylow $3$-subgroup instead, which
returns $\eta^{(l')}$. The example ignores the host conditions.}
\label{tab:crt}
\end{table}

\section{The minus part}

\begin{theorem}[Alp\"oge]\label{thm:minus}
There is a constant $C$ with the following property. Let $M$ be an
abelian datum of degree $o$, let $\nu\ge2$ and $t\ge2$ be integers, let
$(q,d)$ be a prime pair for $M$, and let $E$ be a CM field that is a
field for $(M,\nu,q,d)$, of degree $d_E=[E:\Q]$ and conductor $N$.
For an odd prime $l$ and $0<e<t$ let $\shalen^-_l(e)$ be the length of
the minus part of $\Sha_l(E,e)$. Then $\shalen^-_l(e)<\infty$ for every
odd prime $l$ and every $0<e<t$, and for every finite set $S$ of odd
primes
\[
\sum_{l\in S}\log l\sum_{1\le e<t}\shalen_l^-(e)\le
3d_Et^2\log(tN)+C\sum_{\substack{l\mid o\\ l\text{ odd}}}
l\,d_E^2t^2\log(tlN)^2.
\]
\end{theorem}

\begin{proof}[Idea of proof]
We describe the standard route to a bound of this kind.
The minus part of $\Sha_l(E,e)$ is the sum of the
eigenspaces for the critical characters $\eta$ of $\Gal(E/\Q)$. For an
odd $l$ with $l\nmid[E:\Q]$ and $\eta$ critical for $e$, the length
of the $\eta$-eigenspace is bounded by the $l$-adic valuation of a
generalised Bernoulli number attached to $e+1$ and $\eta$. This bound
comes from Poitou--Tate duality and an orthogonality statement for
unramified classes: a system of elements of cyclotomic fields with
norm relations (Gauss sums), the valuations given by Stickelberger's
theorem, and derivative classes of Kolyvagin type produce enough
classes to bound the eigenspace from one side. Summing over $l$ uses
that for a nonzero integer $B$,
\[
\sum_l v_l(B)\log l=\log|B|,
\]
so one bound on the size of a product of Bernoulli numbers controls
every $l$ at once. The second term of the bound accounts for the odd
primes dividing the degree $o$, where $l\mid[E:\Q]$ is possible.
\end{proof}

\begin{remark}[the minus part is a one-sided bound]
\label{rem:minus-mech}
The bound on the minus part does not use the main conjecture of
Iwasawa theory for abelian fields \cite{MazurWiles1984}. It needs only
an inequality in one direction, which comes from duality, the
orthogonality of unramified classes and the Stickelberger valuations of
Gauss sums. No equality and no characteristic ideal enter. The
situation is classical: on the critical eigenspaces an $L$-value is
available and bounds the group.
\end{remark}

\section{The plus part}

\begin{definition}[coefficient rings, orbits and the hypotheses GD]
\label{def:gd}
Let $l$ be a prime, $N\ge1$ an integer and $\overline\Q$ an algebraic
closure of $\Q$.
\begin{enumerate}[label=(\arabic*)]
\item A \emph{coefficient ring} for $l$ is a $\Z_l$-subalgebra
$\mathcal O$ of an algebraic closure of $\Q_l$ that is a discrete
valuation ring and a finitely generated $\Z_l$-module; $e_{\mathcal O}$ is its
ramification index over $\Z_l$ and $F_{\mathcal O}$ its field of fractions.
\item For a subgroup $H'\le(\Z/N\Z)^\times$, the \emph{orbit}
$\mathcal X_{\mathcal O}(H')$ is the finite set of Dirichlet characters
modulo $N$ with values in $\mathcal O$ whose kernel is exactly $H'$.
For a host datum $\mathfrak h$ and $e\ge0$, $\mathcal X^+_{\mathcal O}(\mathfrak
h,e)$ is the set of characters
$\eta$ modulo $N_{\mathfrak h}$ with values in $\mathcal O$ that are trivial on
$H_{\mathfrak h}$ and plus for $e$.
\item Let $\eta$ be a character modulo $N$ with values in
$\mathcal O$, viewed as a character of $\Gal(\overline\Q/\Q)$, let
$\omega_l$ be the $l$-adic cyclotomic character and let $e\ge0$. The
\emph{eigen-length} $h^1_{\mathcal O}(e,\eta)\in
\Z_{\ge0}\cup\{\infty\}$ is the $\mathcal O$-length of the group of
classes in
$H^1\bigl(\Gal(\overline\Q/\Q),(F_{\mathcal O}/\mathcal
O)(\eta^{-1}\omega_l^{-e})
\bigr)$ that are unramified at the primes not dividing $lN$ and
trivial on the decomposition groups at the primes dividing $lN$. Its
\emph{normalised length} is
$\bar h^1_{\mathcal O}(e,\eta)=h^1_{\mathcal O}(e,\eta)/e_{\mathcal O}$
when $h^1_{\mathcal O}(e,\eta)$ is finite. Neither length depends on
the choice of $\overline\Q$.
\item Let $\mathbb K\subset\overline\Q$ be a number field, $e$ an integer
and $\eta$ a character modulo $N$. The \emph{hypotheses}
$\mathrm{GD}(\mathbb K,l,e,\eta)$ hold if: $l$ is odd; $e\ge1$;
$N$ is even; $l$ does not divide the order of $\eta$;
$(\eta|_l)^2=1$; $\eta$ is plus for $e$; $\mathbb K$ is an imaginary
quadratic field; $l\nmid|\operatorname{disc}\mathbb K|\,h_{\mathbb K}$;
$\gcd(|\operatorname{disc}\mathbb K|,N)=1$; $(l-1)\nmid2e$; and the
intersection of the kernel of the mod-$l$ cyclotomic character of
$\Gal(\overline\Q/\Q)$ with the kernel of $\chi'_l(\eta)$ is not contained
in $\Gal(\overline\Q/\mathbb K)$.
\end{enumerate}
\end{definition}

\noindent The hypotheses $\mathrm{GD}(\mathbb K,l,e,\eta)$ are the
conditions under which the deformation argument in the proof of
Theorem~\ref{thm:plus-gk} applies: $l$ is odd and prime to the
auxiliary data, the twist $e$ is not divisible by $(l-1)/2$, and
$\mathbb K$ is not contained in the field cut out by the mod-$l$
cyclotomic character and $\chi'_l(\eta)$, so that these characters
stay distinct on $G_{\mathbb K}$. For an abelian CM field $E$ the
length of the plus part of
$\Sha_l(E,e)$ is a sum of eigen-lengths over the plus characters of
$\Gal(E/\Q)$, after the values of the characters are adjoined to the
coefficients; this is how the counts below bound the plus part.

\begin{theorem}[Alp\"oge]\label{thm:plus-gk}
Let $\mathbb K$ be a number field. There are
natural numbers $C_1$ and $C_2$ with the following
property. Let $e\ge1$ and $N\ge1$ be integers, $H'\le(\Z/N\Z)^\times$ a
subgroup of index $i$, and $p$ a prime with
$p\nmid30N|\operatorname{disc}\mathbb K|$. There is a nonzero integer
$\mathcal N$ with
\[
\log|\mathcal N|\le C_1\,\bigl|(\Z/i\Z)^\times\bigr|\,
(N(e+1))^{C_2}\log p
\]
such that, for every prime $l\ne p$ and every coefficient ring
$\mathcal O$ for
$l$: if $\mathrm{GD}(\mathbb K,l,e,\eta)$ holds for some
$\eta\in\mathcal X_{\mathcal O}(H')$, then
$h^1_{\mathcal O}(e,\eta)<\infty$ for every
$\eta\in\mathcal X_{\mathcal O}(H')$ and
\[
\sum_{\eta\in\mathcal X_{\mathcal O}(H')}\bar h^1_{\mathcal O}(e,
\eta)\le2v_l(\mathcal N).
\]
\end{theorem}

\noindent The integer $\mathcal N$ is chosen
before $l$ and $\mathcal O$, so one integer controls every prime
$l\ne p$ at once. The statement has content only for imaginary
quadratic $\mathbb K$, since $\mathrm{GD}$ requires it. The automorphy
input of the proof is Theorem~\ref{thm:galois-reps}\ref{item:gr-lgc}
of Chapter~\ref{sec:galois}.

\begin{proof}[Proof (outline)]
We give Alp\"oge's route at the level of its main steps.
Fix $\mathbb K$, $e$, $N$, $H'$ and $p$, and let $\sigma$ be a
Frobenius element at $p$. A nonzero class in the eigenspace of $\eta$,
of the wrong parity, gives a nontrivial tangent vector, of weight $-1$
for the torus action, at the Eisenstein point $\Psi\oplus1\oplus
\Psi^{-1}$ of an ordinary deformation ring of a residually reducible
three-dimensional representation of $G_{\mathbb K}$, where $\Psi$ is
built from $\eta$ and the cyclotomic character. The length of the
class is bounded by the length of a graded piece of the cotangent space
of this ring at the Eisenstein point. Every characteristic-zero point
of the ring is automorphic. This automorphy-lifting theorem of $R=T$
type for residually reducible representations is part of Alp\"oge's
argument, and is not taken from the literature: it is proved from a
patching statement for residually reducible determinant rings,
together with Hida theory and pro-modularity, and its automorphic
input is Theorem~\ref{thm:galois-reps}\ref{item:gr-lgc}. The
automorphic points are of two
kinds: irreducible three-dimensional representations whose cyclotomic
twist is attached to a regular algebraic conjugate self-dual cuspidal
representation of $\mathrm{GL}_3(\mathbb A_{\mathbb K})$, and
representations $W|_{G_{\mathbb K}}\oplus1$ with $W$ attached to a
newform of weight $2e+1$. The cotangent space is then bounded by the
congruences between the Eisenstein trace at $\sigma$ and the traces at
the automorphic points: $\mathcal N$ is a norm of the product of
the differences $\operatorname{tr}\rho(\sigma)-t_{\mathrm{Eis}}$ over
these points, where $t_{\mathrm{Eis}}$ is the trace of the Eisenstein
point at $\sigma$. It is nonzero because the automorphic traces are
smaller than the Eisenstein trace. In the unitary normalisation the
Jacquet--Shalika estimate \cite{JacquetShalika1981} gives
$p^{-1/2}<|\alpha|<p^{1/2}$ for each Satake parameter $\alpha$ at $p$
of a unitary cuspidal representation of $\mathrm{GL}_n$, and so, for
a three-dimensional point, only $|\operatorname{tr}\rho(\sigma)|<
3p^{1/2}$. The bound $|\operatorname{tr}\rho(\sigma)|\le\sqrt p+2$
used here follows when the eigenvalues of $\rho(\sigma)$ have the form
$\{\zeta,\alpha,\alpha^{-1}\}$ with $|\zeta|=1$, as for the points
$W|_{G_{\mathbb K}}\oplus1$ when $\det W(\sigma)=1$: then
$|\operatorname{tr}\rho(\sigma)|\le1+|\alpha|+|\alpha|^{-1}<1+p^{1/2}
+p^{-1/2}<\sqrt p+2$. That the eigenvalues have this form at all the
automorphic points is assumed here, not derived; it is where
conjugate self-duality enters. With the bound $\sqrt p+2$ the
Eisenstein trace is larger (Figure~\ref{fig:plus-route-b} shows the
case $e=1$, with $q$ in the role of $p$).
\end{proof}

\begin{theorem}[Alp\"oge]\label{thm:plus}
Let $\psi$ be a primitive Dirichlet
character. There are a natural number $C$ and primes $r_1,r_2$ with the
following property. Let $\mathfrak h$ be a host datum with character $\psi$
whose primes $q,d$ exceed $r_1$ and $r_2$, let $t$ be an integer for
which $\mathfrak h$ is large, and let $0<e<t$. Let $\overline\Q$ be an algebraic
closure of $\Q$, choose a coefficient ring $\mathcal O_l$ for every
prime $l$, and
let $L$ be a finite set of odd primes $l$ with $l\notin B_\psi$ and
$(l-1)\nmid2e$. Then $h^1_{\mathcal O_l}(e,\eta)<\infty$ for
every
$l\in L$ and $\eta\in \mathcal X^+_{\mathcal O_l}(\mathfrak h,e)$, and
\[
\sum_{l\in L}\log l\sum_{\eta\in \mathcal X^+_{\mathcal O_l}(\mathfrak h,e)}
\bar h^1_{\mathcal O_l}(e,\eta)\le(N_{\mathfrak h}t)^C.
\]
\end{theorem}

\noindent The condition
on $l$ in Theorem~\ref{thm:plus} says that every triple $(l,e,\eta)$ is
in Case~I; it depends only on $(l,e)$, which is why the sum runs over
all of $\mathcal X^+_{\mathcal O_l}(\mathfrak h,e)$.
The automorphy input of Theorem~\ref{thm:plus}, through
Theorem~\ref{thm:plus-gk}, is
Theorem~\ref{thm:galois-reps}\ref{item:gr-lgc}.

\begin{proof}[Proof (outline)]
We give Alp\"oge's route at the level of its main steps.
The proof applies a form of
Theorem~\ref{thm:plus-gk}
summed over orbits and over primes $l$, with
$\mathbb K=\Q(\sqrt{-p_K})$ for the auxiliary prime $p_K$ of
Definition~\ref{def:cases}, and with two fixed primes
$r_1,r_2>30N_\psi p_K$ in the role of $p$. The discriminant of
$\mathbb K$ is $-p_K$ and its class number is $h^*$, and both $p_K$ and
the prime divisors of $h^*$ lie in $B_\psi$; this is why $l\notin
B_\psi$ is imposed. The sum over $l$ costs nothing, since
$\sum_l v_l(\mathcal N)\log l=\log|\mathcal N|$, and the
bound on $\log|\mathcal N|$ is polynomial in $N_{\mathfrak h}$ and $t$.
\end{proof}

\begin{remark}[the plus part is bounded by automorphy lifting]
\label{rem:plus-mech}
Theorem~\ref{thm:plus-gk} is not proved as a converse to Ribet's
theorem \cite{Ribet1976} by a lattice argument on a unitary group in
three variables. The route is a cotangent bound on a deformation ring together
with the automorphy of
all its characteristic-zero points, and no unitary group occurs in the
statements. The group in the background is
$\mathrm{GL}_3\rtimes\{\pm1\}$, which encodes conjugate self-duality.
A lattice construction in the style of Ribet appears only in the
construction of one rank-two representation used along the way, and no
Ramanujan bound is used: the traces are bounded through the
Jacquet--Shalika estimate. Figure~\ref{fig:plus-route} shows the route.
\end{remark}

\begin{figure}[htbp]
\centering
\input{alpoge/figs/v2/plus-route.tex}
\caption{(a) The route to the plus-part bound (Theorems~\ref{thm:plus-gk}
and~\ref{thm:plus}),
schematically: a Sha class of the wrong parity is a tangent vector of
weight $-1$ at the Eisenstein point $x_0$ of an ordinary deformation
ring, and its length is bounded by the congruences between the trace
at one Frobenius and the traces at the automorphic points. (b) For
$e=1$, $|t_{\mathrm{Eis}}|\ge q-1-q^{-1}>\sqrt q+2$ once $q\ge6$, so a
trace with $|t|\le\sqrt q+2$ differs from the Eisenstein trace. This is
the bound of the proof of Theorem~\ref{thm:plus-gk}, which follows from
the Jacquet--Shalika estimate when the eigenvalues at the point pair
off as $\alpha,\alpha^{-1}$; the general Jacquet--Shalika bound
$3\sqrt q$ separates the traces only for $q\ge13$. The argument takes
$q\ge7$.}
\label{fig:plus-route}
\end{figure}

\begin{definition}[Kato configurations]\label{def:kato-config}
Let $l$ be a prime, $e\ge1$ and $N\ge1$ integers, $\mathbb K$ a number
field with a character $\theta$ of $\Gal(\overline\Q/\Q)$, $\eta$ a
character of $\Gal(\overline\Q/\Q)$ with values in
$\overline\Q_l^\times$, and
$W\colon\Gal(\overline\Q/\Q)\to\mathrm{GL}_2(\overline\Q_l)$ a
continuous representation. Let $\omega_l$ be the $l$-adic cyclotomic
character and put $\Psi=\omega_l^e\eta$. The tuple
$(l,e,N,\mathbb K,\eta,W)$ is a \emph{Kato configuration} if:
\begin{enumerate}[label=(\arabic*)]
\item $l$ is odd, $(l-1)\nmid2e$, $N$ is even, $\mathbb K$ is
imaginary quadratic, $l\nmid|\operatorname{disc}\mathbb K|h_{\mathbb
K}$, $\gcd(|\operatorname{disc}\mathbb K|,N)=1$, and $\theta$ is the
quadratic character of $\mathbb K$;
\item $\eta$ has finite order prime to $l$, $\eta^2$ is trivial on the
inertia groups at $l$, and $\Psi(c)=1$ for some complex conjugation
$c$;
\item $W$ has level dividing $N|\operatorname{disc}\mathbb K|$ and
Hodge--Tate weights $(e,-e)$, $\det W=\theta$,
$\operatorname{tr}W=\Psi+\Psi^{-1}$ on the inertia groups at $l$, and
$\operatorname{tr}W\equiv\theta\Psi+\Psi^{-1}$ modulo the maximal
ideal;
\item $W$ is ordinary at $l$, with a stable flag at a place above $l$
on whose graded pieces inertia acts by $\omega_l^{\pm e}$ times a
common character;
\item\label{item:kc-local}
$\eta$ and $W$ satisfy further local conditions: a condition on
$\eta$ relative to the roots of unity of order $N$, and, over every
imaginary quadratic field over which $W$ is irreducible, a condition
at the places above $l$.
\end{enumerate}
\end{definition}

\noindent The conditions in~\ref{item:kc-local} are part of the definition and
are not
written out here. The data mirror those of the hypotheses $\mathrm{GD}$ of
Definition~\ref{def:gd}, with $W$ a two-dimensional representation of
the kind that occurs among the automorphic points in the proof of
Theorem~\ref{thm:plus-gk}.

\begin{theorem}[Alp\"oge]\label{thm:bloch-kato}
Let $l$ be an odd prime, $e\ge1$, $N\ge1$,
$\mathbb K$ an imaginary quadratic field with quadratic character
$\theta$, $\eta$ a character of $\Gal(\overline\Q/\Q)$ and
$W\colon\Gal(\overline\Q/\Q)\to\mathrm{GL}_2(\overline\Q_l)$ a
continuous representation, such that $(l,e,N,\mathbb K,\eta,W)$ is a
Kato configuration. Suppose that $W(-e)$ is attached to a newform of
weight $2e+1$. Then the Bloch--Kato Selmer groups of $W$ and of
$W\otimes\theta$ vanish.
\end{theorem}

\noindent The statement is taken, in a weaker form, from
\cite[\nolinkurl{KatoTwoW.lean}, lines~10--14]{AlpogeArchive2026},
where it is the proposition named \texttt{KatoTwoW}. No proof of
Theorem~\ref{thm:bloch-kato} is given in this book, and no
published source for it is known to us. It is a vanishing theorem
of Kato type \cite{Kato2004}; Kato treats elliptic modular newforms
over $\Q$ and their Dirichlet twists, and we do not settle here
whether the twist by $\theta$ lies in that range. The vanishing is
of the kind that Kato's Euler system gives for modular forms.

\begin{theorem}[Alp\"oge]\label{thm:h2-bound}
There is a constant $C_K$ with the following
property. Let $l$ be a prime, $L$ a totally complex number field
containing the $l$-th roots of unity (the fourth roots of unity if
$l=2$), and $e\ge1$. Then $H^2(G_{L,S_l},\Z_l(e+1))$ is finite and
\[
\log\bigl|H^2(G_{L,S_l},\Z_l(e+1))\bigr|\le
\exp\bigl((2e[L:\Q])^{C_K}\log(2|\operatorname{disc}L|)\bigr).
\]
\end{theorem}

\noindent The statement is that of
\cite[\nolinkurl{LemmaK.lean}, lines~9--17]{AlpogeArchive2026},
where it is the proposition named \texttt{LemmaK}. No proof of
Theorem~\ref{thm:h2-bound} is given in this book. It is a
finiteness statement of Soul\'e type: Soul\'e \cite{Soule1979}
proves the finiteness of these groups, and the doubly exponential
bound is an addition, for which no published source is known to
us. The bound is weak, and it is used only for a bounded
number of exceptional twists and in the sub-case $V_s$ of Case~V
(Theorem~\ref{thm:remaining-counts}).

\begin{theorem}[Alp\"oge]\label{thm:windows-pigeonhole}
There are functions $R_0$ and $R_l$ (for primes $l$) from abelian data
to the natural numbers and a constant $C_K>0$ with the following
property. Let $C_t\ge1$. Let $M$ be an abelian datum, $B'$ a finite set
of odd primes, and for each $l\in B'$ let $M_l$ be an abelian datum
whose fields contain the $l$-th roots of unity. Let $\eps\in(0,1/52)$
and $A_1>0$, and put $R=R_0(M)+\sum_{l\in B'}R_l(M_l)$ and
$J=2^{R+2}(2C_KC_t/\eps)^R$. There is $X_0$ such that for all functions
$\nu,t$ on $[2,\infty)$ with
$2\le\nu(x)$, $2^{\nu(x)}\le x^{A_1}$ and $2\le t(x)\le x^{C_t}$, every
$X\ge X_0$ and every prime pair $(q,d)$ for $M$ with
$\log(qd)\le X^{\eps/(4J)}$, there is $Y$ with $X^{1/J}\le Y$ and
$Y^2\le X$ such that for every $x\in[Y,Y^2]$:
\begin{enumerate}[label=(\alph*)]
\item\label{item:pig-two} every field $E$ for $(M,\nu(x),q,d)$
satisfies $\shalen_2(e)<\infty$ for $0<e<t(x)$ and
$\sum_{1\le e<t(x)}\shalen_2(e)\le\exp(x^{4\eps/5})$;
\item\label{item:pig-bad} for every $l\in B'$, every field $L_0$ cut out by
the datum $M_l$ and every field $L_1=L_0(\mu_{l^k})$, $k\ge0$, with
$[L_1:L_0]=l^{j_l}$, where
$j_l=\lfloor\log_lR_l(M_l)\rfloor+3$, the field $L_1$ satisfies
$\shalen_l(e)<\infty$ for $0<e<t(x)$ and
$\sum_{1\le e<t(x)}\shalen_l(e)\le\exp(x^{4\eps/5})$.
\end{enumerate}
\end{theorem}

\noindent Theorem~\ref{thm:windows-pigeonhole} is unconditional, but its
proof uses other results: Theorem~\ref{thm:h2-bound}, Poitou--Tate
duality, the finiteness of local Galois cohomology and local Tate
duality, and the theorem
of Ferrero and Washington \cite{FerreroWashington1979}; the exceptional
twists are treated by Theorem~\ref{thm:h2-bound} alone.

\begin{proof}[Proof (outline)]
We give Alp\"oge's route at the level of its main steps.
Ferrero--Washington
\cite{FerreroWashington1979}, in the form used, bounds the $l$-ranks
of the class groups uniformly up a cyclotomic tower. A tower bound then
bounds the Sha lengths at a layer of the tower linearly in the layer,
for every twist $e$ outside a fixed $l$-adic neighbourhood of each of
at most $R$ centres. This leaves at most $R$ exceptional twists in a
given range, and for those only the doubly exponential bound of
Theorem~\ref{thm:h2-bound} is available. An
exceptional twist $e_i$ does harm at the scale $x$ only when it lies
between a small power $x^{\kappa}$ and $x^{C_t}$, that is, for $\log x$
in an interval of fixed ratio. Among at most $R$ such intervals a
pigeonhole argument finds, in a range of length depending only on $R$
and the constants, a window $[Y,Y^2]$ free of all of them. This is the
origin of the exponent $J$.
\end{proof}

Case~V is divided into two sub-cases. Attached to an odd prime $l$ and
an eigen-character $\eta$ are two natural numbers $R_1(l,\eta)$ and
$c_2(\eta)$, explicit in terms of $l$, $\eta$ and the host datum, whose
definitions we omit. A triple $(l,1,\eta)$ in Case~V is
in the \emph{main sub-case} if $R_1(l,\eta)\le2^{c_2(\eta)}$, and in
the \emph{sub-case} $V_s$ otherwise.

\begin{theorem}[Alp\"oge]\label{thm:remaining-counts}
Let $\psi$ be a primitive Dirichlet character. There is a natural
number $C$ with the following property. Let $\mathfrak h$ be a host datum with
character $\psi$, $t$ an integer for which $\mathfrak h$ is large, $\overline\Q$
an
algebraic closure of $\Q$, $l$ an odd prime, $\mathcal O$ a
coefficient ring for
$l$, $0<e<t$, and $\eta\in \mathcal X^+_{\mathcal O}(\mathfrak h,e)$. Suppose
that $(l,e,\eta)$ is in
Case~II, III or~Z, or in the main sub-case of Case~V. Then
$h^1_{\mathcal O}(e,\eta)<\infty$ and
\[
\bar h^1_{\mathcal O}(e,\eta)\log l\le(l\,N_{\mathfrak h}\,e)^C.
\]
\end{theorem}

\noindent Cases~IV and
$V_s$ are not covered by the theorem: they are bounded without
automorphy, Case~IV by restriction to a layer of a cyclotomic tower and
Theorem~\ref{thm:windows-pigeonhole}, and Case~$V_s$ by
Theorem~\ref{thm:h2-bound} at $e=1$.
The automorphy statements used for Cases~II, III and the main
sub-case of~V are, like that of Theorem~\ref{thm:plus}, consequences
of Theorem~\ref{thm:galois-reps}\ref{item:gr-lgc}.

\begin{proof}[On the proof]
We describe the route by case; this account is ours.
Cases~II and~III follow the pattern of the proof of
Theorem~\ref{thm:plus-gk} in rank four: ordinary deformations, with values
in $\mathrm{GSp}_4$, of a sum of four characters; an automorphy
statement in which the automorphic points are attached, over an
imaginary quadratic field, to cuspidal representations of
$\mathrm{GL}_4$ (no Siegel modular forms occur); a rigidity statement;
a congruence number; and an input on the residual nonvanishing of
Bernoulli numbers. The main sub-case of Case~V has the same shape over
$\Q(i)$. Case~Z concerns the trivial character, and its bound uses the
arithmetic of Bernoulli numbers only, with no automorphy; for $l\ge5$,
$e\ge2$ even and $(l-1)\mid2e$ the relevant eigen-length of the
trivial character vanishes.
\end{proof}

\begin{remark}[comparison with Vandiver and Greenberg]
\label{rem:vandiver}
Theorems~\ref{thm:plus} and~\ref{thm:remaining-counts} bound the
eigenspaces of the wrong parity of $\Sha^1(\Q_l/\Z_l(-e))$ over abelian
fields, polynomially in the conductor and uniformly in $l$. For $e=0$
statements of this kind concern the plus parts of class groups of
abelian fields, and are of the type of Vandiver's conjecture and of
Greenberg's conjecture on the Iwasawa invariants of totally real fields
\cite{Washington1997,Greenberg1976}. We know of no published result of
this strength.
\end{remark}

\section{The Sha bound on windows}

\begin{theorem}[Alp\"oge]\label{thm:sha-windows}
Every primitive Dirichlet character has Sha bounds on windows with
constant $401$ (Definition~\ref{def:sha-windows}).
\end{theorem}

\noindent So the Sha bound on windows holds. The automorphy used in
the case bounds, Theorems~\ref{thm:plus} and~\ref{thm:remaining-counts}, is
Theorem~\ref{thm:galois-reps}\ref{item:gr-lgc}; this dependence is
recorded in Remark~\ref{thm:lgc-to-sha}.
\begin{figure}[tbp]
\centering
\input{book/figs/ch10-sha-ledger.tex}
\end{figure}

\begin{proof}[On the proof]
We describe how the bounds combine.
Fix a
primitive $\psi$, $\eps\in(0,1/52)$ and an adapted $\nu$, and let $E$
be a $\psi$-field for $(\psi,\nu(x),q,d)$ with $x$ in a window. The
Sha budget of $E$ splits into three parts: the $2$-part, the minus part
at odd primes and the plus part at odd primes, so a bound $B_2$ for
the $2$-lengths and bounds $B_-$, $B_+$ for the other two give the
budget $(t,B_2\log2+B_-+B_+)$. The window and the $2$-part come
from Theorem~\ref{thm:windows-pigeonhole}\ref{item:pig-two}, with
$C_t=401$, applied to the datum
of $\psi$ with $t(x)=\lfloor x^{401}\rfloor$; the same window serves
the bad primes $l\in B_\psi$, with $M_l$ the datum of
$\Q(\zeta_{lN_\psi})$, through clause~\ref{item:pig-bad}. The
minus part is bounded by Theorem~\ref{thm:minus}. For the plus part,
the length at an odd prime is the sum of the eigen-lengths over the
plus characters, and Proposition~\ref{prop:cases-exhaustive} sorts
every eigenspace into one of the six cases: Case~I is bounded by
Theorem~\ref{thm:plus}; Cases~II, III, Z and the
main sub-case of~V by Theorem~\ref{thm:remaining-counts}; Case~IV by
restriction to a layer of a cyclotomic tower and the window bound of
clause~\ref{item:pig-bad}; and Case~$V_s$ by
Theorem~\ref{thm:h2-bound}. Poitou--Tate duality, the finiteness of
local Galois cohomology and local Tate duality
\cite{Tate1963,Poitou1967} are used to pass between the groups of the
field $E$ and the eigen-lengths over $\Q$, and between $\Sha^1$ and
$H^2$.

On a window, $\log q$ and $\log d$ are at most $x^{\eps/5}\log^2x$
and $2^{\nu(x)}\le2x^{2+8\eps}$, so $\log N_{\mathfrak h}=O(x^{\eps/5}\log^2x)$.
This is why bounds polynomial in $N_{\mathfrak h}$, $t$ and $l$ are small
against
the budget $\exp(x^\eps)^{1/2}$ (Figure~\ref{fig:sha-ledger}).
\end{proof}

\begin{remark}[where the windows and constants come from]
\label{rem:windows-origin}
The windows do not come from an Iwasawa growth formula
$\lambda_{\mathrm{Iw}}\nu+\nu_{\mathrm{Iw}}$ for the Sha lengths in a
cyclotomic tower, with invariants controlled only doubly
exponentially. No growth formula is used. The windows come from
the pigeonhole argument in the proof of
Theorem~\ref{thm:windows-pigeonhole}, applied to a bounded number of
exceptional twists, and the constants are only doubly exponential
because Theorem~\ref{thm:h2-bound} is. The arithmetic input exists
only on windows, so the analytic criterion of
Chapter~\ref{sec:landau} is windowed
(Theorem~\ref{thm:windowed-landau}). The window exponent of the Sha bound is
$2$, and the
supply of hosts in Chapter~\ref{sec:assembly} is stated with the
exponent $3/2$. How the one becomes the other is
Question~\ref{q:window-exponent}, which follows the proof of
Theorem~\ref{thm:host-supply}. One further constant has the
same origin as the windows. The constant $401$ makes the
budget cover $\lfloor x^{401}\rfloor$ Tate twists, while the number of
twists needed is the rank parameter $t$ of
Definition~\ref{def:ranks}, which satisfies $t\le x^{400}$. So $401$
leaves a margin of one in the exponent; we do not identify a step
that uses the margin.
\end{remark}

\begin{remark}[named inputs that the argument does not use]
\label{rem:unused-inputs}
The argument for Theorem~\ref{thm:sha-windows} outlined here does not
use Clozel's purity lemma \cite{Clozel1990}, the main theorem of
complex multiplication or the theorem of Deligne and Ribet
\cite{DeligneRibet1980}. The bound on traces that makes the
congruence number nonzero in Theorem~\ref{thm:plus-gk} is the
Jacquet--Shalika bound. No bound on Satake parameters from purity is
used, and the purity lemma in question concerns infinity types.
\end{remark}
\clearpage
\clearpage
\chapter{Galois representations and local--global compatibility}
\label{ch:galois}\label{sec:galois}

This chapter states Theorem~\ref{thm:galois-reps}: weak local
Langlands maps exist, and truncated local--global compatibility holds
for regular algebraic polarized cuspidal representations over CM
fields. Part~\ref{item:gr-llc} is stated without an outline of proof,
and part~\ref{item:gr-lgc} with an outline of the named steps of the
construction. The chapter also states the elementary
Proposition~\ref{prop:compact-group} on compact groups, which is the
step of the construction that recovers an $a$-dimensional
representation from an annihilating polynomial. In
Chapter~\ref{ch:assembly}, Theorem~\ref{thm:galois-reps} is the first
step of the proof of the main theorem. It is used as a black box
twice: as the automorphy input of the Sha bound on windows of
Chapter~\ref{ch:sha} (Remark~\ref{thm:lgc-to-sha}), and, through its
weak local Langlands maps, in the existence of admissible-family data
(Theorem~\ref{thm:packages}\ref{item:packages-data},
Chapter~\ref{ch:families}).

The automorphy statements of Chapter~\ref{sec:sha} need Galois
representations attached to regular algebraic polarized cuspidal
representations over CM fields, with a truncated form of local--global
compatibility. This chapter states what is needed, compares it with
the literature, and outlines how the representations are constructed.
What is needed is weaker than the published theorems in two ways: the
local Langlands maps are required to satisfy only a list of laws, and
the compatibility is asserted only at the places where the monodromy
vanishes.

\section{Weak local Langlands maps and truncated compatibility}

\begin{definition}[weak local Langlands maps]\label{def:weak-llc}
Let $E$ be a number field and $w$ a finite place of $E$, with residue
characteristic $p$.
\begin{enumerate}[label=(\arabic*)]
\item An irreducible admissible representation of some
$\mathrm{GL}_n(E_w)$ is \emph{good} if its cuspidal support consists of
characters, or, when $E_w\cong\Q_p$, of twists of depth-zero cuspidal
representations.
\item\label{item:wl-laws} A map $\mathrm{rc}$ from the irreducible admissible
representations of the groups $\mathrm{GL}_n(E_w)$, $n\ge1$, to
Weil--Deligne representations of the Weil group of $E_w$
\emph{satisfies the laws} if:
\begin{enumerate}[label=(\alph*)]
\item\label{item:wl-a} it preserves dimension and is local class field theory on
characters;
\item\label{item:wl-b}
it is compatible with the Zelevinsky classification by segments,
sends cuspidal representations to irreducible ones, and respects
unramified twists and the absolute value of the central character;
\item on good representations it respects twists by characters, sends
the central character to the determinant and the contragredient to
the dual, is injective, and sends depth zero to tame;
\item\label{item:wl-d}
it satisfies a compatibility of types over finite extensions.
\end{enumerate}
\item A \emph{weak local Langlands map} for $(E,w)$ is a map satisfying
the laws.
\end{enumerate}
\end{definition}

\noindent Clauses~\ref{item:wl-laws}\ref{item:wl-b}
and~\ref{item:wl-laws}\ref{item:wl-d} are summarised; their precise form
is a list of compatibilities that we do not need. The uniqueness that
goes with this definition holds only on good representations: two weak
local Langlands maps for $(E,w)$ agree on good representations.

\begin{remark}[the local Langlands input is weaker than the
correspondence]\label{rem:llc}
The argument does not use the local Langlands correspondence for
$\mathrm{GL}_n$ \cite{HarrisTaylor2001,Henniart2000,Scholze2013}. It
uses only the existence of weak local Langlands maps. These are not required to
be bijective, and no condition on $L$- or $\epsilon$-factors is
imposed. On representations that are not good, a weak local Langlands
map is constrained only by laws~\ref{item:wl-a}, \ref{item:wl-b}
and~\ref{item:wl-d}, and the existence statement makes an arbitrary
choice there. The correspondence is the stronger statement that for
every number field $E$ and finite place $w$ there is exactly one map
from the irreducible admissible representations of the groups
$\mathrm{GL}_n(E_w)$ to
Frobenius-semisimple Weil--Deligne representations that is a
bijection, satisfies the laws of Definition~\ref{def:weak-llc} on all
representations and preserves the $L$- and $\epsilon$-factors of
pairs, with these maps compatible for different fields.
\end{remark}

\begin{definition}[truncated local--global compatibility]
\label{def:truncated-lgc}
Let $E$ be a CM field with maximal totally real subfield $E^+$ and
complex conjugation $c$, let $l$ be a prime, $\iota\colon
\overline\Q_l\cong\C$ an isomorphism and $a\ge1$. Let $\pi$ be a
cuspidal automorphic representation of $\mathrm{GL}_a(\mathbb A_E)$.
\begin{enumerate}[label=(\arabic*)]
\item $\pi$ is \emph{polarized} if $\pi^c\cong\pi^\vee\otimes
(\chi\circ N_{E/E^+})$ for a character $\chi$ of
$\mathbb A_{E^+}^\times/(E^+)^\times$ with the sign condition at the
real places of $E^+$.
\item\label{item:tl-system}
A \emph{compatible system for} $\pi$ consists of a continuous
semisimple representation $r\colon G_E\to\mathrm{GL}_a(\overline\Q_l)$,
an integer $\wt$ and, for each embedding $\tau\colon
E\to\overline\Q_l$, a multiset $\mathrm{HT}_\tau$ of $a$ integers, such
that:
\begin{enumerate}[label=(\alph*)]
\item (weights) for every $\tau$, $\mathrm{HT}_\tau$, viewed in $\C$
by $\iota$, is the multiset of the numbers $(a-1)/2-z$ for $z$ in the
infinitesimal character of $\pi$ at $\iota\circ\tau$, and
$\mathrm{HT}_{\tau c}=\wt-\mathrm{HT}_\tau$;
\item\label{item:tl-off}
(places not above $l$) at every finite place $w\nmid l$ and for
every weak local Langlands map $\mathrm{rc}$ for $(E,w)$: the
dimensions of $\mathrm{rc}(\pi_w)$ and of $r|_{G_{E_w}}$ agree; if
$\mathrm{rc}(\pi_w)$ is a sum of characters with monodromy $N=0$, then
the Weil--Deligne representation of $r|_{G_{E_w}}$ is
$\mathrm{rc}(\pi_w)$; and if $E_w=\Q_p$ with $a=1$ or $p>16a^2$, and
$\mathrm{rc}(\pi_w)$ is tame, multiplicity free and has $N=0$, then the
Weil--Deligne representation of $r|_{G_{E_w}}$ has $N=0$ and the
characteristic polynomials of inertia on the two agree;
\item\label{item:tl-above}
(places above $l$) at every place $w\mid l$: if $\pi_w$ is
spherical, then $r|_{G_{E_w}}$ is crystalline with Hodge--Tate weights
$(\mathrm{HT}_\tau)_\tau$; and for every weak local Langlands map $\mathrm{rc}$
such that $\mathrm{rc}(\pi_w)$ has $N=0$ and is either a sum of
characters or satisfies the three conditions of~\ref{item:tl-off}, the
restriction
$r|_{G_{E_w}}$ is potentially semistable with these Hodge--Tate weights,
and its Weil--Deligne representation is $\mathrm{rc}(\pi_w)$ in the
first case and, in the second, has $N=0$ and the same characteristic
polynomials of inertia as $\mathrm{rc}(\pi_w)$.
\end{enumerate}
\item \emph{Truncated local--global compatibility} is the statement
that for every CM field $E$, every prime $l$, every $\iota$, every
$a\ge1$ and every regular algebraic polarized cuspidal automorphic
representation $\pi$ of $\mathrm{GL}_a(\mathbb A_E)$ there is a
compatible system for $\pi$.
\end{enumerate}
\end{definition}

\noindent Figure~\ref{fig:truncated-lgc} lists the local clauses case by case. Compared with local--global compatibility in the
literature, nothing is asserted about the monodromy operator except
where it vanishes, nothing about purity, and at places above $l$ only
the cases listed in~\ref{item:tl-system}\ref{item:tl-above}.
\begin{figure}[tbp]
\centering
\input{book/figs/ch11-truncated-lgc.tex}
\end{figure}

\begin{theorem}[Alp\"oge]\label{thm:galois-reps}
\begin{enumerate}[label=(\alph*)]
\item\label{item:gr-llc} For every number field $E$ and every finite
place $w$ of $E$, a weak local Langlands map for $(E,w)$ exists.
\item\label{item:gr-lgc} Truncated local--global compatibility
(Definition~\ref{def:truncated-lgc}) holds.
\end{enumerate}
\end{theorem}

\noindent In the literature the corresponding results are due to
Kottwitz, Clozel, Harris--Taylor, Shin, Chenevier--Harris, Caraiani and
others \cite{Kottwitz1992,HarrisTaylor2001,TaylorYoshida2007,Shin2011,
ChenevierHarris2013,Caraiani2012,Caraiani2014,
BarnetLambGeeGeraghtyTaylor2014}. The construction below agrees with
that literature in its frame and, as far as we know, differs from it at
three inner steps.

\begin{proof}[Proof (outline)]
\noindent\emph{Proof of \ref{item:gr-llc}.} We do not outline this
part. On representations that are not good, the map it provides is
an arbitrary choice compatible with laws~\ref{item:wl-a},
\ref{item:wl-b} and~\ref{item:wl-d}.

\noindent\emph{Proof of \ref{item:gr-lgc}.}
We describe only the named steps of the construction.
\begin{enumerate}[label=(\arabic*)]
\item\label{item:gr-s1}
The Shimura varieties are those of the unitary similitude group
of a hermitian space of odd dimension $a$ over $E$, of signature
$(1,a-1)$ at one infinite place and definite at the others. The order
of the PEL datum has rank one, so no division algebra occurs. The
cohomology used is that of the Kuga--Sato powers of the universal
abelian scheme, in all degrees.
\item\label{item:gr-s2}
The Hecke eigensystem of a form on the definite unitary group is
put into this cohomology by a holomorphic Lefschetz formula for
correspondences, compared with a Hecke trace on the definite group.
This gives occurrence of the eigensystem, not its multiplicity.
\item\label{item:gr-s3}
The Eichler--Shimura relation is used only as an annihilation
statement: a polynomial whose roots are the Satake parameters kills
the Frobenius.
\item\label{item:gr-s4}
An $a$-dimensional piece is cut out of the cohomology by a
statement on the Sen operator together with
Proposition~\ref{prop:compact-group}, using a set of Hodge--Tate
weights that is lacunary (Definition~\ref{def:lacunary}).
\item Crystallinity comes from a crystalline comparison theorem for all
cohomology degrees of smooth projective models.
\item\label{item:gr-s6}
General forms on the definite group are reached by $l$-adic
approximation and limits of pseudo-characters. Even rank is reached
from odd rank by theta lifts, which occur in the construction only for
this step. The passage from $\mathrm{GL}_a$ to the definite unitary
group is a descent by a twisted trace formula, and general CM fields
are reached by base-change patching.
\end{enumerate}
\end{proof}

\noindent The frame (descent by the trace formula, limits of
pseudo-characters, base-change patching) is that of the literature.
The departures are steps~\ref{item:gr-s2}, \ref{item:gr-s3}--\ref{item:gr-s4}
and the use of theta lifts in
step~\ref{item:gr-s6}. As far as we know, the published constructions compute
the
trace of Frobenius on the cohomology by counting points, in the manner
of Langlands and Kottwitz \cite{Kottwitz1992,HarrisTaylor2001,Shin2011},
which gives the whole characteristic polynomial of Frobenius. Here only
an
annihilating polynomial is available, and Step~\ref{item:gr-s4} supplies what is
missing: the Eichler--Shimura relation gives the eigenvalues without
their multiplicities, and the Sen operator with
Proposition~\ref{prop:compact-group} recovers an $a$-dimensional
representation from them. Figure~\ref{fig:galois-route} shows the
steps.

\begin{figure}[htbp]
\centering
\input{alpoge/figs/v2/galois-route.tex}
\caption{The construction behind Theorem~\ref{thm:galois-reps}. Plain
boxes follow the published constructions, as far as we know; dashed
boxes mark the three departures from them (step~\ref{item:gr-s2},
steps~\ref{item:gr-s3}--\ref{item:gr-s4}, and the theta lifts in
step~\ref{item:gr-s6}; Remark~\ref{rem:galois-inputs}). Step~\ref{item:gr-s4}
uses
Proposition~\ref{prop:compact-group}.}
\label{fig:galois-route}
\end{figure}

\begin{remark}[the Galois inputs are proved in the argument]
\label{rem:galois-inputs}
The argument does not take two separate Galois-theoretic inputs from
the literature as assumptions, namely Caraiani's local--global
compatibility and the existence of Galois representations attached to
polarized cuspidal representations. Both places need only truncated
local--global compatibility, the same statement, and
Theorem~\ref{thm:galois-reps} supplies it,
so no further hypothesis enters (Alp\"oge's claim; see
Remark~\ref{rem:status}). The third input of this kind, the existence of
weak local Langlands maps, is also part of
Theorem~\ref{thm:galois-reps}.

Truncated compatibility is weaker than the known theorems. For a
regular algebraic polarized cuspidal $\pi$ over a CM field, the Galois
representations constructed in
\cite{HarrisTaylor2001,Shin2011,ChenevierHarris2013} satisfy, under
the hypotheses of the papers cited, local--global compatibility
including the monodromy operator up to Frobenius semisimplification,
at the places not above $l$ \cite{TaylorYoshida2007,Caraiani2012} and
at the places above $l$
\cite{Caraiani2014}, \cite[Thm.~2.1.1]{BarnetLambGeeGeraghtyTaylor2014}.
Truncated
compatibility asks for less: no monodromy beyond the case $N=0$, no
purity, and at each place only the clauses of
Definition~\ref{def:truncated-lgc}. Since
clauses~\ref{item:tl-system}\ref{item:tl-off}
and~\ref{item:tl-system}\ref{item:tl-above}
are asserted for every weak local Langlands map, they say something
only where the laws determine the map. For example, if
$\mathrm{rc}(\pi_w)$ is a sum of characters, then by laws~\ref{item:wl-a}
and~\ref{item:wl-b}
the cuspidal support of $\pi_w$ consists of characters, so $\pi_w$ is
good, and on good representations two weak local Langlands maps agree
(laws~\ref{item:wl-a} and~\ref{item:wl-b} are used here in the summarised form
of
Definition~\ref{def:weak-llc}).
Provided the local Langlands correspondence itself satisfies the laws,
the clauses of Definition~\ref{def:truncated-lgc} are therefore
consequences of the theorems just cited. Under that proviso,
truncated compatibility, as a hypothesis, is implied by the
literature. It is also
what the application needs, for the abelian CM host fields of
Chapter~\ref{sec:hosts}.
\end{remark}

\section{A proposition on compact groups}

\begin{definition}[lacunary sets and product-generic polynomials]
\label{def:lacunary}
Let $K$ be a field of characteristic zero.
\begin{enumerate}[label=(\arabic*)]
\item A finite family $(e_1,\dots,e_n)$ of elements of $K$ is
\emph{lacunary} if no two disjoint nonempty sets $A,B\subset
\{1,\dots,n\}$ satisfy $|B|\sum_{i\in A}e_i=|A|\sum_{i\in B}e_i$.
\item A monic polynomial $P\in K[X]$ of degree $n$ is
\emph{product-generic} if it is separable and the product of its roots
is not the product of any other multiset of $n$ of its roots, taken
with repetition.
\end{enumerate}
\end{definition}

\noindent Equivalently, a family is lacunary if no two disjoint
nonempty subfamilies have equal averages.
\begin{figure}[tbp]
\centering
\input{book/figs/ch11-subset-averages.tex}
\end{figure}

\begin{definition}[tangent to the Zariski closure]
\label{def:zariski-tangent}
Let $K$ be a field of characteristic zero, $\Gamma$ a group,
$\rho\colon\Gamma\to\mathrm{GL}_N(K)$ a representation and $K'$ an
extension of $K$. A matrix $X\in M_N(K')$ is \emph{tangent to the
Zariski closure} of $\rho$ if for every polynomial function $f$ on
$M_N$ over $K$ that vanishes on $\rho(\Gamma)$, the differential of
$f$ at the identity vanishes on $X$.
\end{definition}

\begin{proposition}[Alp\"oge]\label{prop:compact-group}
Let $\Gamma$ be a compact topological group, $K$ an algebraically
closed field of characteristic zero with a Hausdorff ring topology,
$K'$ an algebraically closed extension of $K$, and $N,n\ge1$ integers.
Let $\rho\colon\Gamma\to\mathrm{GL}_N(K)$ be a continuous semisimple
representation. Let $S\subset\Gamma$ be dense, and for each $s\in S$
let $P_s\in K[X]$ be monic of degree $n$ with $P_s(\rho(s))=0$. Let
$X\in M_N(K')$ be tangent to the Zariski closure of $\rho$, and suppose
that $X$ has $n$ distinct eigenvalues that form a lacunary family. Let
$\psi\colon\Gamma\to K^\times$ be a continuous character with
$\psi(s)=(-1)^nP_s(0)$ for $s\in S$. Then there is a continuous
semisimple representation $r\colon\Gamma\to\mathrm{GL}_n(K)$ such that:
\begin{enumerate}[label=(\alph*)]
\item\label{item:cg-sub} $r$ is a subrepresentation of $\rho$, and
$\rho$ is isomorphic to a subrepresentation of $r^{\oplus k}$ for some
$k\ge1$;
\item\label{item:cg-kill} for every $g\in\Gamma$ the characteristic
polynomial of $r(g)$ kills $\rho(g)$, and $P_s(r(s))=0$ for every
$s\in S$;
\item\label{item:cg-det} $\det r=\psi$;
\item\label{item:cg-sep} the set of $g\in\Gamma$ at which the
characteristic polynomial of $r(g)$ is separable is open and dense;
\item\label{item:cg-char} for $s\in S$, the characteristic polynomial
of $r(s)$ equals $P_s$ whenever the characteristic polynomial of
$r(s)$ is separable, and whenever $P_s$ is product-generic.
\end{enumerate}
\end{proposition}

\noindent The statement is elementary: it concerns compact groups,
finite-dimensional representations and Zariski closures only, and it
can be checked independently of the rest of the argument. In the
application $\Gamma$ is the image of a Galois group, $\rho$ is a piece
of the cohomology of step~\ref{item:gr-s1}, the polynomials $P_s$ come from the
Eichler--Shimura relation of step~\ref{item:gr-s3}, and $X$ is the Sen operator,
whose eigenvalues are the Hodge--Tate weights.

\begin{proof}[On the proof]
We do not outline the proof. Part~\ref{item:cg-sub} is the main content of
the statement: every $\rho(s)$ has at most $n$ distinct eigenvalues,
and the conclusion is that $\rho$ is built, up to multiplicity, from a
single representation of dimension $n$. The lacunary eigenvalues of
$X$ are the hypothesis that excludes the other ways in which the
eigenvalues could be constrained. Parts~\ref{item:cg-kill}
to~\ref{item:cg-char} record the properties of $r$ that step~\ref{item:gr-s4} of
the proof of Theorem~\ref{thm:galois-reps} uses.
\end{proof}

\begin{example}[lacunary sets]\label{ex:lacunary}
The set $\{1,2,3\}$ is not lacunary: $\{1,3\}$ and $\{2\}$ both have
average $2$. More generally an arithmetic progression
$a,a+\delta,a+2\delta$ is never lacunary, since $\{a,a+2\delta\}$ and
$\{a+\delta\}$ have the same average. The sets $\{1,2,4\}$,
$\{1,10,100\}$ and $\{1,3,9,27\}$ are lacunary. Table~\ref{tab:lacunary}
records an exhaustive check for these and a few more small sets, and Figure~\ref{fig:subset-averages} shows the subset averages of the first two.
\end{example}

\begin{table}[htbp]
\centering
\input{alpoge/figs/v2/tab-lacunary.tex}
\caption{A finite set $J$ is lacunary if no two disjoint nonempty
subsets of $J$ have equal averages (Definition~\ref{def:lacunary}),
the condition on the Hodge--Tate weights in
Proposition~\ref{prop:compact-group}. Exhaustive check for small sets.
Arithmetic progressions of length at least $3$ are never lacunary.}
\label{tab:lacunary}
\end{table}

\section{What the Sha bound uses}

\begin{remark}[the Sha bound rests on Theorem~\ref{thm:galois-reps}]
\label{thm:lgc-to-sha}
The only input of Chapter~\ref{sec:sha} beyond classical Galois
cohomology and the theorems stated there is
Theorem~\ref{thm:galois-reps}\ref{item:gr-lgc}. Truncated
local--global compatibility (Definition~\ref{def:truncated-lgc})
supplies the Galois representations attached to regular algebraic
polarized cuspidal representations of $\mathrm{GL}_3$ and
$\mathrm{GL}_4$ over imaginary quadratic fields, together with the
local properties that make them points of the ordinary deformation
rings in the proofs of Theorems~\ref{thm:plus-gk}, \ref{thm:plus}
and~\ref{thm:remaining-counts}. Read the other way, the argument of
Chapter~\ref{sec:sha} shows that truncated local--global
compatibility implies the Sha bound on windows; a base change
statement used along the way is also part of Alp\"oge's argument.
\end{remark}

\clearpage
\clearpage
\chapter{Existence of families}
\label{ch:families}\label{sec:assembly}

This chapter states the existence results that the Chebyshev
inequality on windows of Chapter~\ref{ch:chebyshev} needs, and deduces
from them the zero-freeness of $L(s,\psi)$ for primitive $\psi$ in
$\operatorname{Re}s>1/2$ (Theorem~\ref{thm:zero-free-primitive}). The
existence results are the supply of hosts
(Theorem~\ref{thm:host-supply}), Green's parametrization
(Theorem~\ref{thm:green}), the existence of packages and of
admissible-family data (Theorem~\ref{thm:packages}) and the existence
of tied families (Theorem~\ref{thm:tied-families}); their proofs are
given in outline or by a list of inputs, and the deduction of
Theorem~\ref{thm:zero-free-primitive} is proved in full. All five
theorems are black boxes in the proof of the main theorem in
Chapter~\ref{ch:assembly}; Section~\ref{sec:families-to-assembly} says
which step uses which.

Chapters~\ref{sec:balance}--\ref{sec:chebwin} reduce the zero-freeness
of the $L(s,\psi)$ to two kinds of existence. First, for each
primitive $\psi$ there must be host fields at every scale of a system
of windows, meeting the explicit-formula condition and carrying a
small Sha budget. Second, there must be families of packages
(Definition~\ref{def:family}) tied at every large scale of their
windows (Definition~\ref{def:tied}), so that the Chebyshev inequality
on windows (Theorem~\ref{thm:chebyshev-windows}) applies. The chain
of implications that proves Theorem~\ref{thm:main} is assembled in
Section~\ref{sec:proof-main} of Chapter~\ref{ch:assembly}. The
mathematical content is in the results quoted; the assembly itself is short
(Remark~\ref{rem:content}).

\section{The supply of hosts}

\begin{definition}[host supply; schematic]\label{def:host-supply}
Let $\lambda>1$.
\begin{enumerate}[label=(\arabic*)]
\item A \emph{host requirement} $\mathscr R$ is a property of a scale
$(x,\eps,J)$ (Definition~\ref{def:scales}), a frame
(Definition~\ref{def:frame}), a host $\mathcal H$ for that frame
(Definition~\ref{def:host}), a finite set $\Top\subset A$, where $A$
is the group of the frame, and a positive integer $t$.
\item\label{item:hs-good}
Let $(x,\eps,J)$ be a scale, with $b=\exp(x^\eps)$, let
$\mathbf k=(k,m,t)$ be ranks (Definition~\ref{def:ranks}), let
$C_1,C_2,C_3,C_4$ be real numbers and $n_0$ a natural number. A host
$\mathcal H$ for $\mathbf k$, with auxiliary primes $q,d$ and field
$E\supset K=\Q(\sqrt{-d})$, is \emph{good at} $(x,\eps,J)$ for these
constants if it satisfies size conditions at this scale; it has norm primes
$\varpi_0,\varpi_1$ (Definition~\ref{def:norm-primes}); it has primes
that split completely in $E$; the
characters $\chi_q=(\cdot/q)$, $\chi_d=(\cdot/d)$ and $\chi_q\chi_d$
satisfy pointwise
bounds at scale $x$ and have zero-free boxes;
$|\operatorname{Re}L'(1,\chi_{-d})/L(1,\chi_{-d})|\le x^\eps$; the
thin-scale bound of Definition~\ref{def:thick} holds; and the number
of bad Frobenius classes is bounded.
\item The host requirement $\mathscr R$ is \emph{supplied on windows of
exponent} $\lambda$
if for every $\eps\in(0,1/52)$ there are real numbers
$C_1,\dots,C_5,x_0$ with $C_1\ge1$ and $C_2,C_3,C_4>0$, and a
system of scales of exponent $\lambda$ (Definition~\ref{def:scales}),
such that for every scale $(x,\eps,J)$ of the system with $x\ge x_0$
there are ranks $\mathbf k=(k,m,t)$ valid at this scale
(Definition~\ref{def:ranks}), a host $\mathcal H$ for $\mathbf k$ with
constant $C_5$, with top places $\Top_{\mathcal H}$, and a natural
number $n_0$, such that $\mathcal H$ is good at $(x,\eps,J)$ for
$C_1,C_2,C_3,C_4,n_0$ and $\mathscr R$ holds at the scale, the frame
of $\mathbf k$, $\mathcal H$, $\Top_{\mathcal H}$ and $t$.
\item For a primitive Dirichlet character $\psi$, the
\emph{requirement of} $\psi$ is the host requirement $\mathscr R_\psi$
that holds at a scale $(x,\eps,J)$, a frame, a host $\mathcal H$ with
field $E$, a set $\Top$ and an integer $t$ if $\mathcal H$ satisfies
the explicit-formula condition $\mathrm{EF}(\lev_\psi,\Top,x,\eps,4)$
of Definition~\ref{def:shares-ef}, for the cyclic level $\lev_\psi$
(Definition~\ref{def:level}), and $E$ has Sha budget $(t,b^{1/2})$
(Definition~\ref{def:sha-windows}), with $b=\exp(x^\eps)$.
\end{enumerate}
\end{definition}

\noindent Clause~\ref{item:hs-good} is descriptive. The size conditions of a
good
host, the finer size conditions on $\varpi_0,\varpi_1$, the bounds and boxes
for $\chi_q$, $\chi_d$ and
$\chi_q\chi_d$, the bound on bad Frobenius classes, the validity of
ranks at a scale and the constant $C_5$ of a host for given ranks
depend on the constants and on $n_0$; clause~\ref{item:hs-good} is
schematic in the sense of Section~\ref{sec:conventions}, and these
conditions
are not written out here. In $\mathscr R_\psi$ the Sha budget is part
of the requirement, and this is how the Sha bound of
Chapter~\ref{sec:sha} reaches the hosts.

\begin{theorem}[Alp\"oge]\label{thm:host-supply}
Let $\psi$ be a primitive Dirichlet character. If $\psi$ has Sha
bounds on windows with constant $401$
(Definition~\ref{def:sha-windows}), then
$\mathscr R_\psi$ is supplied on windows of exponent $3/2$.
\end{theorem}

\begin{proof}[Proof (outline)]
We give Alp\"oge's route at the level of its main steps.
Fix $\eps\in(0,1/52)$. The host field is the field of
Theorem~\ref{thm:fields-exist}, with a frame of large exponent and
twist form for $\psi$. In each window the auxiliary primes $(q,d)$
are a good pair, which Theorem~\ref{thm:good-pairs} supplies. A good
pair bounds the bounded twists
(Proposition~\ref{prop:bounded-twist-bounds}), and this gives the
explicit-formula condition $\mathrm{EF}(\lev_\psi,\Top,x,\eps,4)$ in
$\mathscr R_\psi$. The Sha budget in
$\mathscr R_\psi$ is the hypothesis on $\psi$.

The two windowings are compatible, because one pair $(q,d)$ serves
both. Theorem~\ref{thm:good-pairs} with $M=N_\psi$, $\Delta=J$ and
$Y=X^{1/J}$, for large $X$, gives a window pair $(q,d)$ for
$(\eps,X^{1/J})$ that is good for $(N_\psi,\eps,X^{1/J},X)$. For this
pair the Sha bound on windows gives a window $[Y',Y'^2]\subset
[X^{1/J},X]$. At every $x\in[Y',Y'^2]$ the pair is good for
$(N_\psi,\eps,X^{1/J},x)$, since $X^{1/J}\le x\le X$, so
Proposition~\ref{prop:bounded-twist-bounds} applies at $x$; the frame
exponent of Theorem~\ref{thm:fields-exist} is taken to be $\nu(x)$,
adapted to $(\eps,x)$ in the sense of Definition~\ref{def:sha-windows}. The Sha
bound gives
windows $[Y,Y^2]$, and the theorem is stated with the exponent
$3/2$; the step between the two is not reconstructed here
(Question~\ref{q:window-exponent}, Figure~\ref{fig:windows}).
\end{proof}

\begin{question}[the window exponent $3/2$]\label{q:window-exponent}
By what step do the windows $[Y,Y^2]$ of the Sha bound
(Theorem~\ref{thm:sha-windows}) become the windows $[W,W^{3/2}]$ of
the supply of hosts (Theorem~\ref{thm:host-supply})?
\end{question}

\noindent One candidate is a rescaling by the rank: the
analytic scale $x$ of a host would be $k$ times the scale $x'$ at
which the Sha bound is applied, where $k$ is the first entry of the
ranks datum, and a window $[W,W^2]$ for $x'$ would give the window
$[kW,(kW)^{3/2}]$ for $x$, since $(kW)^{3/2}\le W^2$ once
$W\ge k^3$. This candidate is inconsistent with the definitions.
\begin{figure}[tbp]
\centering
\input{book/figs/ch12-window-exponent.tex}
\end{figure}

\begin{remark}[rescaling by the rank does not give $3/2$]
\label{rem:window-exponent}
Let $(k,m,t)$ be the ranks datum of a host at the scale $x$. The
group $A$ of the frame surjects onto a group of order $2^\nu$
(Definition~\ref{def:frame}), with $\nu$ adapted to the scale, so
its order $|A|$ is at least $2^\nu\ge x^{2+8\eps}$; and
Definition~\ref{def:ranks} requires $2|A|\le k$.
Hence $k\ge2x^{2+8\eps}$ (Figure~\ref{fig:window-exponent}). Independently, the hypotheses
$xm\le k$ and $m\ge x^{1/2-\eps}$ of Proposition~\ref{prop:sqrtx}
give $k\ge x^{3/2-\eps}$. So $x=kx'$ would force
$x'\le x^{-1-8\eps}/2<1$, whereas the Sha bound is applied at
scales in a window $[Y,Y^2]$ with $Y$ large; and $W\ge k^3$ would
force $W\ge8x^{6+24\eps}$, which is incompatible with $kW\le x$.
The rank $k$ is not bounded independently of $x$, so it is not a
fixed factor either. One observation is elementary: a set with
windows of exponent $2$ has windows of exponent $3/2$
(Definition~\ref{def:windows}), because
$[Y,Y^{3/2}]\subset[Y,Y^2]$. Whether this is all that happens, or
whether the scale changes between the two statements, is not
determined here. Remark~\ref{rem:bb-open} lists it among the open
points.
\end{remark}

\section{Green's parametrization}

\begin{theorem}[Green \cite{Green1955}]\label{thm:green}
Let $p$ be a prime and $n\ge1$. There is a map $P\mapsto\chi_P$ from
the monic irreducible polynomials $P\in\F_p[X]$ of degree $n$ with
$P(0)\ne0$ to the cuspidal irreducible characters of
$\mathrm{GL}_n(\F_p)$ such that
\begin{enumerate}[label=(\alph*)]
\item $\chi_P(1)=\prod_{i=1}^{n-1}(p^i-1)$ for every such $P$;
\item $P\mapsto\chi_P$ is injective;
\item every cuspidal irreducible character of $\mathrm{GL}_n(\F_p)$ is
$\chi_P$ for some such $P$.
\end{enumerate}
\end{theorem}

\noindent The map depends on a choice of labels: an identification of
the characters of $\F_{p^n}^\times$ with $\F_{p^n}^\times$, under which
the Frobenius orbits of regular characters correspond to the roots of
the polynomials $P$.

\begin{example}\label{ex:green}
For $p=3$ and $n=2$ the monic irreducible quadratics over $\F_3$ are
$X^2+1$, $X^2+X+2$ and $X^2+2X+2$. Hence $\mathrm{GL}_2(\F_3)$ has
three cuspidal irreducible characters, each of degree $p-1=2$.
Table~\ref{tab:green} gives the counts for small $p$ and $n$.
\end{example}

\begin{table}[htbp]
\centering
\input{alpoge/figs/v2/tab-green.tex}
\caption{Cuspidal irreducible characters of $\mathrm{GL}_n(\F_p)$
correspond to monic irreducible polynomials $P$ of degree $n$ over
$\F_p$ with $P(0)\ne0$ (Theorem~\ref{thm:green}). The table gives
their number, by the M\"obius formula, and the degree
$\chi_P(1)=\prod_{i=1}^{n-1}(p^i-1)$.}
\label{tab:green}
\end{table}

\begin{remark}[where Green's theorem and the exponent $3/2$ enter]
\label{rem:green}
Theorem~\ref{thm:green} and the exponent $3/2$ appear in adjacent
steps of the proof of Theorem~\ref{thm:main}. The exponent does not
come from Green's theorem, and Green's theorem is not an analytic
input. Green's parametrization enters only through local types at the norm
primes $\varpi_0,\varpi_1$: it is used
to show that the rank defect of each member is small (the irreducibility
link, clause~\ref{item:lk-irr} of Definition~\ref{def:links}) and to label the
local
types in the ties. The exponent $3/2$ is the window exponent of
Theorem~\ref{thm:host-supply}; its relation to the exponent $2$ of
the Sha bound on windows is Question~\ref{q:window-exponent}.
\end{remark}

\section{Admissible families}

\begin{definition}[the links; schematic]\label{def:links}
Let $\Fam$ be a family at $\eps$. For a member $\mathcal P$ at $x$
with size datum $Z$, coefficient field $L$, ranks $(k,m,t)$ with
derived quantity $f$ (Definition~\ref{def:ranks}) and lift index $j$,
put $b=\exp(x^\eps)$, let $B_0$ and $N_c$ be the size quantities of
$Z$ (Definition~\ref{def:size-datum}), and let
\[
\Sigma_j=\sum_{\mathfrak p}e_{\mathfrak p}\,
\frac{\log N\mathfrak p}{[L:\Q]}\,\eta_j(\mathfrak p),
\]
the sum over the places $\mathfrak p$ of the support of $Z$ at which
the local type of $\mathcal P$ is twisted, where
$\eta_j(\mathfrak p)\ge0$ is the congruence invariant of $\mathcal P$
at $\mathfrak p$ and lift index $j$ (Definition~\ref{def:package}).
For a prime $l$ put $v(l)=\log l\sum_{1\le e<t}\shalen(l,e)$, with
$\shalen(l,e)$ the lengths of Definition~\ref{def:package}. Let
$D_0\ge1$ be the rank count of $\mathcal P$. The \emph{links} of
$\Fam$ are the following statements, each for all large $x$ and
every member at $x$.
\begin{enumerate}[label=(\arabic*)]
\item\label{item:lk-orth}
(Near-orthogonality.) The theta kernels of distinct words are nearly
orthogonal for the Hilbert--Schmidt pairing, and the kernel of each
word is nearly an isometry, with relative defects at most
$\exp(-x^{\eps/8})$.
\item\label{item:lk-irr}
(Irreducibility.) The rank defect of each member is nonnegative and
at most $\exp(-x^{e_z/7})$, where $e_z>0$ is an exponent attached to
the family.
\item\label{item:lk-cong} (Congruence.) For every lift index $j$,
\[
\Sigma_j\le D_0\Bigl(C\bigl(fmb\,x^{2\eps}+N_c\log^2x\bigr)
+\sum_lv(l)\Bigr),
\]
for a constant $C$ of the family.
\item\label{item:lk-div} (Divisibility.) For every lift index $j$,
$\Sigma_j\le D_0\,Cx^\eps(B_0+N_c)$.
\item\label{item:lk-glob} (Global.) The chosen data and the inputs of
clauses~\ref{item:lk-orth}--\ref{item:lk-div}, the upper link
(Definition~\ref{def:upper-link}) and the lower link
(Definition~\ref{def:lower-link}) imply that every large $x\in\Win$
is tied (Definition~\ref{def:tied}).
\end{enumerate}
\end{definition}

\noindent The rank count $D_0$ is an integer attached to the
package. The precise form of the near-orthogonality
in~\ref{item:lk-orth}, and the further hypotheses
of~\ref{item:lk-glob}, are schematic (Section~\ref{sec:conventions}).
The divisibility bound~\ref{item:lk-div} is a hypothesis of the upper link.

\begin{definition}[admissible-family data and the top existence
statement; schematic]\label{def:admissible}
Let $\lambda>1$, let $\mathscr R$ be a host requirement
(Definition~\ref{def:host-supply}) and let $\mathcal C$ and
$\mathcal A$ be classes of families (Definition~\ref{def:family}).
\begin{enumerate}[label=(\arabic*)]
\item\label{item:ad-sha}
A family at $\eps$ \emph{satisfies the Sha hypothesis} if there is $x_0$
such that for every $x\ge x_0$ the host field of every member at $x$
has Sha budget $(t,b^{1/2})$ (Definition~\ref{def:sha-windows}), with
$t$ the rank parameter and $b=\exp(x^\eps)$.
\item The \emph{class of $\mathscr R$ at $\lambda$},
$\mathcal C_{\mathscr R}(\lambda)$, consists of the families with
windows of exponent $\lambda$ all of whose members beyond the
threshold meet $\mathscr R$.
\item The class $\mathcal C$ \emph{supplies tied families} if there are real
numbers $e_D,e_L$ and $\eps_1>0$ such that for every
$\eps\in(0,\eps_1)$ some family at $\eps$ in $\mathcal C$ is tied
(Definition~\ref{def:tied}) at every large $x\in\Win$, with these
$e_D,e_L$ and some constants.
\item The \emph{top existence statement} is: for every $\lambda>1$ and
every host requirement $\mathscr R$, if
$\mathscr R$ is supplied on windows of exponent $\lambda$ and every family of
$\mathcal C_{\mathscr R}(\lambda)$ satisfies the Sha hypothesis, then
$\mathcal C_{\mathscr R}(\lambda)$ supplies tied families.
\item\label{item:ad-data}
\emph{Admissible-family data} for $\mathcal A$ consist of
realisation data, constants and a list of propositions about the
families of $\mathcal A$, in four groups:
\begin{enumerate}[label=(\alph*)]
\item automorphic forms, theta lifts and mass formulas: class number
equals mass times index; volume ratios by the Siegel--Weil formula
\cite{Weil1965,KudlaRallis1988,Ichino2004}; decay of matrix
coefficients; integrality of row maps; Hecke compatibility; isotypic
decomposition; integrality of periods up to the order of the
stabiliser;
\item Galois representations and local--global compatibility: twist
places, and the congruence link, clause~\ref{item:lk-cong} of
Definition~\ref{def:links};
\item\label{item:ad-c}
Arakelov geometry and covolume: propositions obtained from the
covolume condition (Definition~\ref{def:covolume-condition});
\item\label{item:ad-d}
counting and design: the existence of a large family in the
class, using Minkowski's bound on finite subgroups of
$\mathrm{GL}_n(\Z)$ \cite{Minkowski1887}, applied to
$\mathrm{GL}_n(\Q)$; sizes of level primes; thin-set statements;
near-orthogonality; good local types; level projections.
\end{enumerate}
The propositions are required to hold as closed properties
(Definition~\ref{def:closed}) on $\mathcal A$ or on subclasses of
$\mathcal A$ attached to the data; among them, the congruence link is
closed on such a subclass.
\end{enumerate}
\end{definition}

\noindent Item~\ref{item:ad-data} is descriptive: the realisation data, the
constants
and the statements of the individual propositions are schematic
(Section~\ref{sec:conventions}); the list in~\ref{item:ad-data} names
the topics they cover. Minkowski's bound applies to
$\mathrm{GL}_n(\Q)$ because a finite subgroup of $\mathrm{GL}_n(\Q)$
fixes a lattice. The top existence statement does not mention
$\mathcal A$.

The class $\mathcal C_{\mathrm{cov}}$ of
Definition~\ref{def:covolume-condition} is not one of the classes
$\mathcal C$, $\mathcal A$ of this definition. It is a fixed class, cut
out by conditions on the packages of a family, and the covolume
condition is asserted on it. The admissible-family data of a class
$\mathcal A$ are built with the covolume condition as one input
(Theorem~\ref{thm:packages}\ref{item:packages-data}), and their
propositions in group~\ref{item:ad-c} are then required on $\mathcal A$ or on
its
subclasses.

\begin{remark}[where the Sha bound is spent]\label{rem:sha-upper-link}
The Sha bound enters on the archimedean side. It does not enter as a
separate correction on the finite side at the twisted places. The divisibility
bound, clause~\ref{item:lk-div} of
Definition~\ref{def:links}, follows from the Sha hypothesis
(Definition~\ref{def:admissible}\ref{item:ad-sha}), the congruence link and a
rank
statement, and it is a hypothesis of the upper link
(Definition~\ref{def:upper-link}). The Sha budget $b^{1/2}$ is
absorbed into the bound $D_0Cx^\eps(B_0+N_c)$ because $b\le B_0$.
\end{remark}

\begin{theorem}[Alp\"oge]\label{thm:packages}
Let $\eps\in(0,1/52)$.
\begin{enumerate}[label=(\alph*)]
\item\label{item:packages-exist} For every scale $(x,\eps,J)$ with
$x$ sufficiently large there are packages
(Definition~\ref{def:package}) whose automorphic forms have the
weights, level and local types prescribed at that scale, on the tower
of unitary groups of Definition~\ref{def:spaces} over a CM field of
degree about $x^{2+8\eps}$.
\item\label{item:packages-points} The members of the families built
from these packages have points, and for each lift index $j$ the
lattice $M_j$ of the covolume condition
(Definition~\ref{def:covolume-condition}) is the lattice of integral
cusp forms on $U(W_j)$.
\item\label{item:packages-data} There is a class of families of
these packages that carries admissible-family data
(Definition~\ref{def:admissible}).
\end{enumerate}
\end{theorem}

\noindent \emph{Informal statement.} At every large scale there are
automorphic forms on the tower of unitary groups over a CM field of
degree about $x^{2+8\eps}$ with the prescribed weights, level and
local types. Their integral cusp forms give the lattices $M_j$ of the
covolume bound, and families of these packages satisfy the
propositions listed in Definition~\ref{def:admissible}\ref{item:ad-data}.

\begin{proof}[Idea of proof of \ref{item:packages-exist}
and~\ref{item:packages-points}]
We indicate the standard method. On the definite unitary groups of the
tower the forms are algebraic automorphic forms \cite{Gross1999}, so
forms with prescribed local components can be sought among functions
on finite class sets. The chain of theta kernels
(Definition~\ref{def:column}) links the levels of the tower, and its
analysis rests on the Siegel--Weil formula
\cite{Weil1965,KudlaRallis1988,Ichino2004} and Howe duality
\cite{Howe1989,Waldspurger1990,GanTakeda2016}.
\end{proof}

\begin{proof}[On the proof of \ref{item:packages-data}]
The data are assembled from the following inputs: the separation of
functionals by unramified Howe duality
\cite{Howe1989,Waldspurger1990,GanTakeda2016} and the Siegel--Weil
formula; the rationality of theta lifts; a fit of Fourier
coefficients; theta integrals, by M\'inguez's results
\cite{Minguez2008}; row functions, through discrete series and the
Casselman--Wallach theory \cite{Casselman1989,Wallach1992}; the
covolume bound (Theorem~\ref{thm:covolume}); the weak local Langlands
maps of Theorem~\ref{thm:galois-reps}; and a bound at the places $2$
and $3$.
\end{proof}

\noindent \emph{Informal statement of
Theorem~\ref{thm:tied-families}.} Suppose some class of families
carries admissible-family data. Let $\mathscr R$ be a host requirement
that is supplied on windows of exponent $\lambda$, and suppose the
host fields of the families built on these hosts satisfy the Sha
bound. Then for every small $\eps$ there is a family of packages, built
on the supplied hosts, that is tied (Definition~\ref{def:tied}) at
every large scale of its windows, with exponents $e_D=2$ and $e_L=0$.

\begin{theorem}[Alp\"oge]\label{thm:tied-families}
Suppose that some class $\mathcal A$ of families carries
admissible-family data (Definition~\ref{def:admissible}). Let
$\lambda>1$, and let $\mathscr R$ be a host requirement
(Definition~\ref{def:host-supply}) that is supplied on windows of
exponent $\lambda$ and such that every family of
$\mathcal C_{\mathscr R}(\lambda)$ satisfies the Sha hypothesis. Then
$\mathcal C_{\mathscr R}(\lambda)$ supplies tied families.
\end{theorem}

\noindent In other words, the top existence statement of
Definition~\ref{def:admissible} holds.

\begin{proof}[Proof (outline)]
We give Alp\"oge's route at the level of its main steps.
Let $\lambda>1$ and let $\mathscr R$ be a host requirement such that
$\mathscr R$ is supplied on windows of exponent $\lambda$ and every family of
$\mathcal C_{\mathscr R}(\lambda)$ satisfies the Sha hypothesis.
Consider the families of $\mathcal A$ that lie in
$\mathcal C_{\mathscr R}(\lambda)$.

The links of Definition~\ref{def:links}, the upper and lower links
and the remaining propositions of the data are closed on the subclass
attached to the data, so for every small $\eps$ one threshold serves
for all of them. Green's
parametrization (Theorem~\ref{thm:green}) enters here, in the closed
form of the irreducibility link. The design statement of the data
(item~\ref{item:ad-data}\ref{item:ad-d} of
Definition~\ref{def:admissible}),
fed by the hosts with which $\mathscr R$ is supplied,
gives for every small $\eps$ and every threshold a family at $\eps$ of
the class that is large for it. For that family the global link shows
that every large $x\in\Win$ is tied, with $e_D=2$ and $e_L=0$. Hence
the class supplies tied families.

In terms of the objects of Chapters~\ref{sec:slopes}--\ref{sec:chebwin},
the data encode nine steps:
\begin{enumerate}[label=(\arabic*)]
\item build the packages and their lattices (Theorem~\ref{thm:packages});
\item form the period map (Remark~\ref{rem:periodmap});
\item apply the slope inequality (Theorem~\ref{thm:frame});
\item bound the archimedean side (the upper link);
\item bound the covolume (Theorem~\ref{thm:covolume});
\item bound the finite side (the lower link);
\item balance the two sides (Theorem~\ref{thm:gap});
\item bound the defects at the twisted places by Sha lengths and the
Sha hypothesis (Remark~\ref{rem:sha-upper-link});
\item supply near-orthogonality, good local types, masses and class
numbers.
\end{enumerate}
\end{proof}

\begin{remark}[quantifiers]\label{rem:quantifiers}
Supplying tied families is an existence statement that cannot hold
vacuously. The closed properties
choose the threshold before the family, and the family produced is
large for that threshold. An empty class cannot supply tied
families, since the definition asks for a tied family at every small $\eps$.
Since the
top existence statement does not mention $\mathcal A$,
Theorem~\ref{thm:tied-families} says that one class carrying
admissible-family data suffices.
\end{remark}

\section{Zero-freeness for primitive characters}

\begin{theorem}[Alp\"oge]\label{thm:zero-free-primitive}
Let $\lambda>1$. Assume that the top existence statement
(Definition~\ref{def:admissible}) holds and that
$\mathscr R_\psi$ is supplied on windows of exponent $\lambda$ for every
primitive
Dirichlet character $\psi$. Then $L(s,\psi)\ne0$ for every primitive
$\psi$ and every $s$ with $\operatorname{Re}s>1/2$, except $s=1$ when
$\psi$ is trivial.
\end{theorem}

\begin{proof}
Fix a primitive $\psi$, and let $\mathcal C=\mathcal C
_{\mathscr R_\psi}(\lambda)$; Figure~\ref{fig:choice-order} shows the order in which the proof fixes its objects. Every member of a family in
$\mathcal C$ beyond the threshold meets $\mathscr R_\psi$. Hence every
family in $\mathcal C$ satisfies the Sha hypothesis, because the Sha
budget is part of $\mathscr R_\psi$, and the uniform explicit-formula
condition for $\lev_\psi$ (Definition~\ref{def:uniform-ef}), because
the explicit-formula condition is. By assumption $\mathscr R_\psi$ is
supplied on windows of exponent $\lambda$. The top existence statement
therefore applies to $\lambda$ and $\mathscr R_\psi$, and $\mathcal C$
supplies tied families: there are $e_D,e_L$ and $\eps_1>0$ such that
for every $\eps\in(0,\eps_1)$ some family at $\eps$ in $\mathcal C$ is
tied at every large $x\in\Win$, with these $e_D,e_L$ and some
constants.

Let $s$ satisfy $\operatorname{Re}s>1/2$, with $s\ne1$ when $\psi$ is
trivial, put $\mu=\min(\operatorname{Re}s-1/2,1/2)>0$, and choose
$\eps\in(0,\eps_1)$ with $(e_D+e_L+7)\eps\le\mu/2$. By the previous
paragraph there is a family $\Fam$ at this $\eps$ in $\mathcal C$ and
there are real numbers $c_W,c_S,e_K,C_L$ and $x_0$ such that $\Fam$ is
tied with constants $e_D,e_L,c_W,c_S,e_K,C_L$ at every $x\in\Win$ with
$x\ge x_0$. By the Chebyshev inequality on windows
(Theorem~\ref{thm:chebyshev-windows}), applied to $\lev_\psi$ and
$\Fam$, there is $x_1$ such that
\[
G_{\lev_\psi}(x)\le x^{1/2+(e_D+e_L+7)\eps}\le x^{\sigma_0},
\qquad \sigma_0=\tfrac12+\tfrac\mu2,
\]
for every $x\in\Win$ with $x\ge x_1$. Since $\Fam$ lies in
$\mathcal C$, the set $\Win$ has windows of exponent $\lambda$. Hence
the set $\{x:G_{\lev_\psi}(x)\le x^{\sigma_0}\}$, which contains every
large element of $\Win$, has windows of exponent $\lambda$, and
$\sigma_0\in(0,1)$. By the windowed Landau criterion for levels
(Theorem~\ref{thm:level-criterion}), $\lev_\psi$ satisfies the window
criterion at $\lambda$, so $L(s',\chi)\ne0$ for every
$\chi\in\Kept(\lev_\psi)$ and every $s'$ with
$\operatorname{Re}s'>\sigma_0$, where $s'\ne1$ when $\chi$ is the
trivial character. The character $\psi$ is
kept by $\lev_\psi$, and $\operatorname{Re}s>\sigma_0$ because
$\mu\le\operatorname{Re}s-1/2$. Hence $L(s,\psi)\ne0$.
\end{proof}

\begin{figure}[htbp]
\centering
\input{alpoge/figs/v2/eps.tex}
\caption{The final step (Theorem~\ref{thm:zero-free-primitive}). A
hypothetical zero $\rho$ with $\operatorname{Re}\rho>1/2$ is excluded
by choosing $\eps$ with $9\eps\le\mu/2$ in the Chebyshev inequality on
windows (Theorem~\ref{thm:chebyshev-windows}) and applying the
windowed Landau criterion for levels
(Theorem~\ref{thm:level-criterion}). Schematic.}
\label{fig:eps}
\end{figure}

\noindent For the families of Theorem~\ref{thm:tied-families} one has
$e_D=2$ and $e_L=0$, so $e_D+e_L+7=9$; this is the constant in
Figure~\ref{fig:eps}.
\begin{figure}[tbp]
\centering
\input{book/figs/ch12-choice-order.tex}
\end{figure}

\section{What this chapter supplies to the main proof}
\label{sec:families-to-assembly}

The results of this chapter, together with those of
Chapters~\ref{ch:covolume}--\ref{ch:galois}, complete the inputs to
the proof of Theorem~\ref{thm:main} given in
Section~\ref{sec:proof-main} of Chapter~\ref{ch:assembly}. That proof
is a chain of implications, and it uses the following statements as
black boxes, in this order.
\begin{enumerate}[label=(\arabic*)]
\item Theorem~\ref{thm:galois-reps} (Chapter~\ref{ch:galois}): weak
local Langlands maps exist, and truncated local--global compatibility
(Definition~\ref{def:truncated-lgc}) holds.
\item Theorem~\ref{thm:sha-windows} (Chapter~\ref{ch:sha}): every
primitive Dirichlet character has Sha bounds on windows with constant
$401$; its automorphy input is the previous item
(Remark~\ref{thm:lgc-to-sha}).
\item Theorem~\ref{thm:host-supply}: given the previous item,
$\mathscr R_\psi$ is supplied on windows of exponent $3/2$ for every
primitive $\psi$.
\item Theorem~\ref{thm:packages}\ref{item:packages-data}: some class
of families carries admissible-family data. Its proof uses the
covolume bound (Theorem~\ref{thm:covolume}, Chapter~\ref{ch:covolume})
and the weak local Langlands maps of the first item.
With Green's parametrization (Theorem~\ref{thm:green}),
Theorem~\ref{thm:tied-families} then gives the top existence
statement of Definition~\ref{def:admissible}.
\item Theorem~\ref{thm:zero-free-primitive} with $\lambda=3/2$, whose
proof applies the Chebyshev inequality on windows
(Theorem~\ref{thm:chebyshev-windows}, Chapter~\ref{ch:chebyshev}) and
the windowed Landau criterion for levels
(Theorem~\ref{thm:level-criterion}): $L(s,\psi)\ne0$ for every
primitive $\psi$ and every $s$ with $\operatorname{Re}s>1/2$, except
$s=1$ when $\psi$ is trivial.
\end{enumerate}
The one remaining step of the proof in Chapter~\ref{ch:assembly}, the
passage to all Dirichlet characters and to the critical line, uses
Proposition~\ref{prop:imprimitive} and Theorem~\ref{thm:fe} of
Chapter~\ref{ch:landau}, not the results of this chapter. The
analytic part of the chain is the gap bound (Theorem~\ref{thm:gap})
with the core bound (Theorem~\ref{thm:core-bound}). The constant
$e_D+e_L+7=9$ of Figure~\ref{fig:eps} comes from
Theorem~\ref{thm:tied-families}, which gives $e_D=2$ and $e_L=0$.
\clearpage

\appendix
\clearpage
\chapter{Software and AI methods}
\label{ch:ai}
\label{app:ai-methods}
\sloppy

At Taylor Dupuy's direction (2026-11-21), all material on the use of
software and AI in this book is collected in this appendix, apart
from the short discussion that closes the Introduction
(page~\pageref{intro:ai}). The heading ``Software and AI methods'' is
part of the same instruction. The book has no abstract.

\subsection*{Summary}
AI agents assembled and rewrote every chapter of this book on
2026-11-21, in Claude Code
\cite[\href{https://claude.com/claude-code}{product page}]{ClaudeCode}
(desktop app, version 2.1.328, with subagents), with the models
Claude Opus~5.5 (\texttt{claude-opus-5-5})
\cite[\href{https://platform.claude.com/docs/en/models/opus-5-5/overview}%
{model page}]{ClaudeOpus55} and Claude Fable~5.2
(\texttt{claude-fable-5-2})
\cite[\href{https://platform.claude.com/docs/en/models/fable-5-2/overview}%
{model page}]{ClaudeFable52}; the subagents ran at effort setting
\emph{medium}, and the coordinator's effort was not recorded. They
worked from four sources, all drafted by AI agents on 2026-11-20
and 2026-11-21 and none reviewed by a human: a short account of the
idea of the argument (now Chapter~\ref{ch:idea}); an account of the
Stepanov--Bombieri proof and an account of the sign-only argument
(now Chapter~\ref{ch:function-fields}); and a survey of Alp\"oge's
Lean archive (now Chapters~\ref{ch:assembly}--\ref{ch:families}). The
statement of Theorem~\ref{thm:main} was drafted for that survey by an
agent in the Codex desktop app
\cite[\href{https://openai.com/index/introducing-the-codex-app/}%
{product page}]{OpenAI2026codex} with GPT-6
\cite[\href{https://help.openai.com/en/articles/%
20001354-gpt-6-and-other-models-in-chatgpt}%
{model guide}]{OpenAI2026gpt} on 2026-11-21; its exact variant and
reasoning setting were not exposed. All mathematical prose of the book
is therefore AI-drafted. The results that originated with the agents
are listed in Section~\ref{sec:ai-proofreading}, each with the place
where the book states it. Chapters~\ref{ch:idea}--\ref{ch:assembly}
had one review by an AI agent; Chapters~\ref{ch:landau}--%
\ref{ch:families} had no review as chapters of this book (Figure~\ref{fig:ai-review-coverage}). No human has
reviewed the book, and no agent review is a verification.
This appendix was drafted by an AI agent (Section~\ref{sec:ai-statement}).

\subsection*{Starting material}
The following were available before the agent work on the book.
\begin{enumerate}[label=(\alph*)]
\item Alp\"oge's note \cite{AlpogeRemark2026}, reproduced after the
preface (human-authored, public), and his Lean archive
\cite[\href{https://alpo.ge/}{home page}]{AlpogeArchive2026}, imported
unmodified on 2026-11-20 and described by its author as AI-assisted;
that description was not checked.
\item The requests of Taylor Dupuy (human): the title, the remark in
the frontmatter, the order of the chapters (the idea first, then the
Stepanov--Bombieri proof as an analogy, then the whole proof with
black boxes, then the tools), the sign-only argument woven into
Chapter~\ref{ch:function-fields}, the style of lecture notes for
professional mathematicians and graduate students, the exclusion of
the project's material on the Deninger program, and one appendix for
all AI material.
\item The four sources (AI, not human-reviewed). The account of the
idea was drafted on 2026-11-20 by \texttt{claude-opus-5-5} in Claude
Code from the agent's earlier notes on the argument (AI), the remark,
and Lean statements read as text. The Stepanov--Bombieri account was
written on 2026-11-20 by \texttt{claude-opus-5-5} from a research report
of \texttt{claude-fable-5-2} (see Section~\ref{sec:ai-nlp}). The
sign-only account was drafted on 2026-11-21 by \texttt{claude-opus-5-5}
from the text of the remark and the project's notes. The survey was
rewritten on 2026-11-21 by \texttt{claude-opus-5-5} agents from an
earlier, longer report of 2026-11-20, which rested on readings of the
archive by \texttt{claude-fable-5-2} and on two earlier research notes
of the project (AI, \texttt{claude-fable-5-2}).
\item A plan for the book, written on 2026-11-21 by a
\texttt{claude-fable-5-2} agent and reviewed for wording by a
\texttt{claude-opus-5-5} agent (AI), and two transcriptions of the
remark into \LaTeX{} (AI, one by \texttt{claude-fable-5-2} in the plan
and one by \texttt{claude-opus-5-5}).
\item Vakil's book \cite{Vakil2025} (human, published), assumed
throughout, and the bibliographic data of the sources' bibliographies.
\item The project's written conventions for style, \LaTeX{}, layout
and AI disclosure (drafts; see Section~\ref{sec:ai-agents}).
\end{enumerate}

\section{Typesetting}
\label{sec:ai-typesetting}
The book is typeset with the class \texttt{amsbook}~2.20.6 of the
AMS document classes
\cite[\href{https://ctan.org/pkg/amscls}{CTAN page}]{AMScls} and built
with pdf\TeX{}~1.40.27
\cite[\href{https://ctan.org/pkg/pdftex}{CTAN page}]{pdfTeX} (\TeX{}
Live 2025) through \texttt{latexmk}~4.86a
\cite[\href{https://ctan.org/pkg/latexmk}{CTAN page}]{Latexmk}. The
figures are those of the sources, unchanged, drawn with TikZ/PGF
3.1.10 \cite[\href{https://ctan.org/pkg/pgf}{CTAN page}]{Tantau2023pgf},
pgfplots 1.18.1
\cite[\href{https://ctan.org/pkg/pgfplots}{CTAN page}]%
{Feuersanger2021pgfplots} and tikz-cd 1.0
\cite[\href{https://ctan.org/pkg/tikz-cd}{CTAN page}]{Stoffel2021tikzcd}.
Some builds of the survey used Xe\TeX{}
\cite[\href{https://tug.org/xetex/}{home page}]{XeTeX}.

Skills are named in \texttt{typewriter} type. Where a skill is in a
public snapshot of the skill repository \cite[\href{\vtwoskillurl/%
README.md}{README}]{Dupuy2026agentskills}, its name links to its copy
there. The copies executed were loaded from a local checkout of the
workflow repository \cite[\href{https://github.com/tdupu/mathcity}%
{repository}]{Dupuy2026mathcity}; no pinned public version of it is
recorded, and the executed copies are not shown to equal the snapshot.
The skills \texttt{finish-writing}, \texttt{using-writingpowers},
\texttt{check-ste100}, \texttt{take-notes}, \texttt{create-figure},
\texttt{check-result-statement}, \texttt{check-definition},
\texttt{check-notation} and \texttt{check-claims} are not in the public
snapshot.

The book was made in three steps on 2026-11-21. First, a plan, with
\vtwoskill{frontier-dump}: the planning agent read the four sources,
compared their preambles, labels, figure paths and bibliographies by
script, and built two throwaway prototypes of the assembly. One
prototype dropped spaces between words and was rejected. The plan
listed eighteen problems of assembly, among them colliding macros and
labels, two different figure files of the same name, and a shift by
two chapters in every number of the survey, so that a citation of the
survey by number does not match this book. Second, a skeleton:
one preamble for the four sources, a merged bibliography, label
renames for the collisions, and the remark, merged from the two
transcriptions by a comparison of their words with page images made by
the Poppler tools
\cite[\href{https://poppler.freedesktop.org/}{home page}]{Poppler}.
Third, four writers ported the sources in parallel under
\vtwoskill{using-latexpowers}, which passed the work to
\vtwoskill{math-workflow}: Chapters~\ref{ch:idea}
and~\ref{ch:function-fields} (\texttt{claude-fable-5-2}),
Chapter~\ref{ch:assembly} (\texttt{claude-fable-5-2}),
Chapters~\ref{ch:landau}--\ref{ch:theta} (\texttt{claude-opus-5-5}) and
Chapters~\ref{ch:covolume}--\ref{ch:families}
(\texttt{claude-opus-5-5}). The writers invoked no further skill.
Chapter~\ref{ch:assembly} restates twenty results of the later
chapters as black boxes; a script found each restatement identical to
its target, before and after revision.

The mechanical checks were \vtwoskill{check-labels-and-refs},
\vtwoskill{check-citations}, \vtwoskill{check-style},
\vtwoskill{check-latex}, \vtwoskill{check-latex-hygiene} and
\vtwoskill{check-layout}, run by a \texttt{claude-opus-5-5} agent. The
citation check failed: most citations of the book lack a pinpoint
locator, and many bibliography entries are not cited. The hygiene check
asked for revision, the layout check failed in an advisory mode, and the
style check passed in an advisory mode; these findings are open.
Two prose passes (\texttt{claude-opus-5-5}) with
\texttt{using-writingpowers}, \texttt{finish-writing} and
\vtwoskill{purge-claudisms} changed wording only, in the openings of
the chapters, in Chapters~\ref{ch:idea}--\ref{ch:assembly} and in the
preface.

The sources had been made as follows. The account of the idea was
written in one pass, revised after a review, and given figures by
\texttt{claude-opus-5-5} agents with \vtwoskill{generate-graphics} and
\vtwoskill{add-figure}. The Stepanov--Bombieri account was harvested
from its research report with \vtwoskill{harvest-dumps}, revised so
that each corrected expectation stands where the claim is met,
given figures by \texttt{claude-opus-5-5} agents with
\texttt{create-figure} and \vtwoskill{add-figure}, and shortened from
$29$ to $10$ pages at Taylor Dupuy's request; the book also uses the
commented-out longer version. The sign-only account was written in
three iterations with \texttt{take-notes}. The survey was rewritten
from the earlier report by a coordinating agent and its subagents: an
outline (\vtwoskill{rapid-prototype}, \vtwoskill{create-exposition})
reviewed with \vtwoskill{fp-finder-latex}, five writers, an assembling
agent that checked every statement, and a hygiene pass.

\section{Agent management}
\label{sec:ai-agents}
The project's written conventions are six contracts, each amended
only through its own writer skill: the style contract
(\vtwoskill{new-repo-style-policy}), the \LaTeX{} contract
(\vtwoskill{new-repo-latex-policy}), the layout contract
(\vtwoskill{new-repo-layout-policy}), the AI-disclosure contract
(\vtwoskill{new-repo-ai-policy}), the agents file
(\vtwoskill{new-repo-agents-policy}) and the decision record
(\vtwoskill{new-repo-adr-policy}). The AI-disclosure, style, layout and
\LaTeX{} contracts are drafts that have not been adopted.

The plan was one task: a coordinating session (\texttt{claude-opus-5-5},
Claude Code 2.1.330, desktop app) briefed a planning agent
(\texttt{claude-fable-5-2}, effort \emph{medium}) and a wording
reviewer (\texttt{claude-opus-5-5}, effort \emph{medium}). The book was
a second task: a coordinating session (\texttt{claude-opus-5-5}, Claude
Code 2.1.328, desktop app) and twelve subagents, all at effort
\emph{medium} except the record-keeping agent, whose effort was not
recorded. They were the skeleton agent, the four writers, a subagent
of the first writer that checked its proofs adversarially, a referee, the
agent of the mechanical checks, a revision agent, two prose agents and
a record-keeping agent. The writers of
Chapters~\ref{ch:idea}--\ref{ch:assembly}, the referee and the
revision agent were \texttt{claude-fable-5-2}; the others were
\texttt{claude-opus-5-5}. The models are those recorded in every
billed message of each agent's transcript. The coordinator set the
rulings for the revision and edited two sentences of
Chapters~\ref{ch:sha} and~\ref{ch:galois} on the status of their
results.

The sources were made by other agent teams, described in
Section~\ref{sec:ai-typesetting}. For the survey: the reading of the
archive was one task, in which a coordinating session
(\texttt{claude-opus-5-5}) briefed a worker (\texttt{claude-fable-5-2})
that split the archive into four branches read by four subagents; the
earlier report was a second task with $44$ agents
(\texttt{claude-opus-5-5}); the five further readings were five tasks
with five \texttt{claude-fable-5-2} workers and fourteen subagents;
the rewrite was done by a \texttt{claude-opus-5-5} coordinator and
its subagents.

\section{Natural language proofs}
\label{sec:ai-nlp}
The skill \vtwoskill{frontier-dump} sends a prompt to one model and
saves its report; here the model was always \texttt{claude-fable-5-2}.
It was used four times for the material of this book.
\begin{enumerate}[label=(\arabic*)]
\item On 2026-11-20, in Claude Code 2.1.326 at effort \emph{xhigh},
for the proof of the Riemann hypothesis for curves in the style of
Stepanov and Bombieri, with an analogy to Alp\"oge's argument. Its
result was a research report with a full proof, two variants of the
argument, a list of corrected expectations and a script of numerical
checks; the Stepanov--Bombieri sections of
Chapter~\ref{ch:function-fields} are built from it.
\item On 2026-11-20, for a reading of the Lean archive. Its result was
a reconstruction of the proof architecture (a refined Chebyshev method
giving a one-sided bound on windows of scales, fed into a windowed
Landau criterion), an assumption ledger, a check of the final
statements and an outline; the earlier report was built from it with
\vtwoskill{harvest-dumps}, \vtwoskill{create-exposition},
\vtwoskill{rapid-prototype} and \vtwoskill{fill-in-prototype}.
\item On 2026-11-20, five further readings, each following one part of
the argument below the top level. They gave the deduction from the
character of modulus~$1$ (Remark~\ref{rem:zeta-suffices}), the frame
theorem behind the slope clause (Theorem~\ref{thm:frame}), the routes to
the plus and minus parts of the Sha bound
(Remarks~\ref{rem:plus-mech} and~\ref{rem:minus-mech}), the
construction of the Galois representations
(Theorem~\ref{thm:galois-reps}) and the role of arithmetic geometry
(Remarks~\ref{rem:covol-supply}, \ref{rem:smooth-metrics}
and~\ref{rem:one-point}).
\item On 2026-11-21, at effort \emph{medium}, for the plan of this
book (Section~\ref{sec:ai-typesetting}).
\end{enumerate}

The proofs of the book were written as follows. In
Chapter~\ref{ch:idea}, the proofs of
Propositions~\ref{prop:stepanov}, \ref{prop:chebyshev},
\ref{prop:mellin}, \ref{prop:idea-landau} and~\ref{prop:slope} were
written by \texttt{claude-opus-5-5} on 2026-11-20 following the
standard arguments of the cited sources, which were not reopened; the
writer of the chapter (\texttt{claude-fable-5-2}) replaced the proof of
Proposition~\ref{prop:stepanov} by a sketch citing
Chapter~\ref{ch:function-fields}. In Chapter~\ref{ch:function-fields},
the proofs of the Stepanov--Bombieri sections were written by
\texttt{claude-opus-5-5} on 2026-11-20 from the research report of
item~(1), and the writer of the chapter restored eleven statements from
the longer version of that account. The proofs of
Propositions~\ref{lem:parseval} and~\ref{lem:growth} and of
Theorem~\ref{thm:sign} are reconstructions by \texttt{claude-opus-5-5}
(2026-11-21) of the sketch in Alp\"oge's note; hypothesis~\ref{H2} of
Theorem~\ref{thm:sign} is not proved (Question~\ref{q:sign-h2}). The
proofs and outlines of Chapters~\ref{ch:assembly}--\ref{ch:families}
are those of the survey, written by \texttt{claude-opus-5-5} with
\texttt{fill-in-prototype} and revised in the rewrite; the proof of
Theorem~\ref{thm:main} in Section~\ref{sec:proof-main} was expanded by
the writer of Chapter~\ref{ch:assembly} (\texttt{claude-fable-5-2}),
each step with the black box it calls.

Outside the scope of this book, and not withdrawn: six earlier
findings of the project on the conjecture of Birch and
Swinnerton-Dyer, and a block of $250$ refutations concerning the
Deninger program. None is a premise of the argument of this book.

\section{Adversarial reviews}
\label{sec:ai-reviews}
For the book, one referee review with \vtwoskill{referee-report}, by
a \texttt{claude-fable-5-2} agent that had not drafted the text,
covered Chapters~\ref{ch:idea}--\ref{ch:assembly} and asked for
revision. A \texttt{claude-fable-5-2} subagent of the writer of
Chapters~\ref{ch:idea} and~\ref{ch:function-fields} checked the
sign-only sections of Chapter~\ref{ch:function-fields} adversarially,
following the prompt of the skill \vtwoskill{doubt} without invoking
it. Chapters~\ref{ch:landau}--\ref{ch:families} had no referee review
as chapters of this book.

The sources had these reviews. The account of the idea had one review
with \vtwoskill{critical-review} by a \texttt{claude-opus-5-5}
subagent. The Stepanov--Bombieri account had one review with
\vtwoskill{fp-finder-latex} by a \texttt{claude-opus-5-5} agent, the
same model as the writer, which asked for minor revision and made the
revisions itself, with no return to the referee. The sign-only account
had no review. The earlier report behind the survey had seven
simulated referee reviews with \vtwoskill{fp-finder-latex}, two
reviews of its figures, a review of its introduction with
\vtwoskill{critical-review} and checks of its proved propositions with
\vtwoskill{doubt}; the rewrite had a review of its outline and an
integrated review of its body, which asked for a major revision and
reported no mathematical error in a proved statement. All these
reviews were by agents of the same model family as the drafters.
\begin{figure}[tbp]
\centering
\input{book/figs/ch00-ai-review-coverage.tex}
\end{figure}

\section{AI assisted computations}
\label{sec:ai-computations}
No computer algebra system was run for the book, and no Lean code was
run. The book's computations are Python scripts written and run by
the agents in 2026-11-21 (Python version not recorded): the comparison
of the black-box restatements of Chapter~\ref{ch:assembly} with their
targets; the values of hypothesis~\ref{H2} in characteristic~$2$ and
at $k=4,6,7$ for odd $q\le29$ (Remark~\ref{rem:h2-fails}); the factor
formula of Remark~\ref{rem:hyperbolic}; the orbit count of
Remark~\ref{rem:torus-orbit}; and the ratios of
Remark~\ref{rem:J-balance}. The adversarial subagent of
Section~\ref{sec:ai-reviews} checked the identity $qN=M^k+R$ and the
bound of
Proposition~\ref{lem:parseval} numerically for $q=5,7,9$ and
$k=2,\dots,5$. These are outputs of the scripts; they test small cases
and are not proofs.

For the Stepanov--Bombieri account, a Python~3.13.7 script
\cite[\href{https://www.python.org/downloads/release/python-3137/}%
{release page}]{Python} using only the standard library checked the
parameter inequalities, the constructions on small explicit curves,
point counts and $L$-polynomials by brute force, and some identities on
random inputs; the reviewer checked one identity numerically over
$\F_3$ and $\F_9$. For the survey, the counts of the static scans
(Remark~\ref{rem:ai-scan}) are outputs of a ripgrep~15.1.0
script \cite[\href{https://github.com/BurntSushi/ripgrep/releases/%
tag/15.1.0}{release 15.1.0}]{Ripgrep} and of Python~3.13.7 scripts,
written and run by the reading agents (\texttt{claude-fable-5-2}) and
not recomputed. The data of the figures and example tables of
Chapters~\ref{ch:assembly}--\ref{ch:families} come from Python scripts
written and run by the figure agents (\texttt{claude-opus-5-5}), some
with exact rational arithmetic and two with mpmath~1.3.0
\cite[\href{https://mpmath.org/}{home page}]{Mpmath}. The zeros
of~$\zeta$ in Figures~\ref{fig:gz} and~\ref{fig:landau} come from a
script run in SageMath~10.8
\cite[\href{https://www.sagemath.org/}{home page}]{SageMath} (NumPy
2.3.2 \cite[\href{https://doi.org/10.1038/s41586-020-2649-2}{DOI}]%
{Harris2020numpy} as Sage reported it), with byte-identical output on
a rerun. On 2026-11-21 an agent of the rewrite reran the scripts
behind Examples~\ref{ex:legendre}, \ref{ex:crt}, \ref{ex:lacunary}
and~\ref{ex:green} and Figure~\ref{fig:inert}, with byte-identical
outputs, and the value $L'(1,\chi_{-7})/L(1,\chi_{-7})=
0.015635689993720378957$ was reproduced independently twice with
mpmath. The inequalities of Proposition~\ref{prop:sqrtx} were derived
by hand by an agent and re-derived by a reviewer. For the account of
the idea and the sign-only account no machine computation was made;
their parameter balances and the bound at $k=3$ were checked by hand
by the agents. These computations illustrate the text; no result of
the book rests on them.

\section{AI statement generation}
\label{sec:ai-statement}
This appendix was drafted on 2026-11-21 by a \texttt{claude-opus-5-5}
subagent in Claude Code (harness version and effort not recorded) with
\vtwoskill{write-materials-and-methods}, from the project's
provenance and token records, the logs of the book's run and the
disclosures of the four sources. The record-writing skills
\vtwoskill{update-ai-usage} and \vtwoskill{update-tokens} were only
consulted. The same agent then audited it with the read-only skill
\vtwoskill{latex-ai-statement}; that audit is by the drafting model
and is not independent. The AI notice of the preface was drafted by
the skeleton agent (\texttt{claude-opus-5-5}) and revised by a prose
agent; the AI and Software Discussion of the Introduction was drafted
by the introduction agent (\texttt{claude-fable-5-2}). The
disclosures of the four sources were not ported into the
chapters; their content is in this appendix. No human has reviewed or
accepted this appendix. The author's statement of responsibility for
the book has not yet been written.

\section{Proofreading and adversarial review}
\label{sec:ai-proofreading}
The findings of the referee of the book were applied in a revision
pass with \vtwoskill{revise} (\texttt{claude-fable-5-2}), under the
coordinator's rulings, with \vtwoskill{track-down-reference} for the
references. The following results of the book originated with the
agents in the book's run. None has been checked by a human; each entry
names its origin and how far it was checked.
\begin{enumerate}[label=(\roman*)]
\item Hypothesis~\ref{H2} of Theorem~\ref{thm:sign} fails for the path
quadric at $k=3$ (Remark~\ref{rem:h2-fails}). Found by the adversarial
subagent and written by the writer of Chapter~\ref{ch:function-fields}.
The criterion through the K\"unneth formula was added in the revision.
Proved in the text. The computed failures at $k=4,6,7$ and in
characteristic~$2$ in the same remark are numerical evidence only.
\item The hyperbolic reading of the form, with its bound of order
$q^{3k/4}$ (Remark~\ref{rem:hyperbolic}), and the torus-orbit remark
(Remark~\ref{rem:torus-orbit}); from the referee, written in the
revision, proved in the text and checked by script.
\item The exclusion of a real zero by the sign of its residue in the
proof of Proposition~\ref{prop:idea-landau}; from the referee, written
in the revision.
\item Rescaling the window exponent by the rank is inconsistent with
Definitions~\ref{def:frame} and~\ref{def:ranks} and with
Proposition~\ref{prop:sqrtx} (Remark~\ref{rem:window-exponent}); from the
referee, proved in the revision from the definitions. The step from
the exponent $2$ to $3/2$ is left open
(Question~\ref{q:window-exponent}); the archive's formal statement of
host supply, which would decide it, was not consulted.
\item The remarks~\ref{rem:h2-half}, \ref{rem:origin},
\ref{rem:J-balance} and~\ref{rem:special-admissible}; from the writer
of Chapter~\ref{ch:function-fields} and the revision, proved or
computed in the text. Remark~\ref{rem:special-admissible} rests on
Clifford's theorem, used without proof or locator.
\item The $y(c)$ of the remark are theta coefficients only for $n=3$,
not for every~$n$ (Section~\ref{sec:role}); from the referee, written
in the revision after the cited sources were traced.
\item Corrections to the sources, each made in place: the theta
exponent $a/p+a/p^3$; Proposition~\ref{prop:formal-cex} at $n=1$; the
Parseval saving $q^{(k-1)/2}$; the hypothesis $g\ge1$ of
Theorem~\ref{thm:dichotomy}; and the restriction of an attribution to
Kazhdan and Patterson to $n=3$; and the observation that the
hypothesis $0\notin U$ of Theorem~\ref{thm:sign} is not used. From the
writers and the revision, and confirmed by the referee.
\end{enumerate}

The following results came from the sources, with their origin.
\begin{enumerate}[label=(\roman*),resume]
\item Stepanov--Bombieri account: the corrected range of
\eqref{eq:which-bound}, from its writer; and
Remark~\ref{rem:special-admissible} in its first form, from its
reviewer. The negative results of Chapter~\ref{ch:function-fields}
come from the research report of
Section~\ref{sec:ai-nlp}, item~(1), and its harvest.
\item Account of the idea: the agent's first claim that Landau's
theorem does the job of Bombieri's lower bound was incomplete; the
one-sided criterion works for $\zeta$ because $\zeta$ has no real zero
in $(0,1)$, and fails as stated over a finite field
(Remark~\ref{rem:one-sided-curves}).
\item Sign-only account: the hypothesis $0\notin U$ of
Theorem~\ref{thm:sign}, added in its second iteration (and found
unused in the book, item~(vii)); the link
between the path quadric and Gabber's hypersurface in Katz's argument
\cite{Katz2015} is the agent's inference.
\item Survey (\texttt{claude-opus-5-5} unless stated, from the reviews
and readings of Sections~\ref{sec:ai-nlp} and~\ref{sec:ai-reviews}):
the proof of the balance identity (Proposition~\ref{prop:balance}), by
orthogonality of characters, checked line by line by two reviewers;
the bookkeeping of Proposition~\ref{prop:sqrtx}; the slope clause is
not Bost's inequality but has the form of the conclusion of the frame
theorem (Remark~\ref{rem:not-bost}, Theorem~\ref{thm:frame}), weakened
after a review; the plus-part route
(Theorems~\ref{thm:plus-gk} and~\ref{thm:plus},
Remark~\ref{rem:plus-mech}); the spread of a slope record is set by
no tie (Definition~\ref{def:tied}); the covolume condition and its
one consumer of arithmetic geometry
(Definition~\ref{def:covolume-condition},
Remark~\ref{rem:member-hyp}, Theorem~\ref{thm:covolume},
Remark~\ref{rem:covol-supply}); the cancellations of the upper link
(Remarks~\ref{rem:cancel} and~\ref{rem:sha-upper-link}); the minus
part and its windows (Remark~\ref{rem:minus-mech},
Theorem~\ref{thm:minus}, Theorem~\ref{thm:windows-pigeonhole},
Remark~\ref{rem:windows-origin}); the route of the Galois construction
and its three departures from the published constructions
(Theorem~\ref{thm:galois-reps}, Proposition~\ref{prop:compact-group},
Remarks~\ref{rem:llc} and~\ref{rem:galois-inputs}); the divisor data
(Remarks~\ref{rem:one-point} and~\ref{rem:smooth-metrics}); the
standing hypotheses and parity character of the case split
(Definition~\ref{def:cases}, Proposition~\ref{prop:cases-exhaustive},
Remark~\ref{rem:cases-parity}), from a review; the deduction from the
character of modulus~$1$ (Remark~\ref{rem:zeta-suffices}); the
strength of the arithmetic-geometry inputs and three inputs not used
(Remarks~\ref{rem:ag-strength} and~\ref{rem:unused-inputs}), from the
agents' memory and a textual search, not a term-level dependency
check; Remark~\ref{rem:green} and the exponents of
Theorem~\ref{thm:tied-families}; corrections to
Definitions~\ref{def:ranks} and~\ref{def:spaces} and
Remark~\ref{rem:topspace}, made while drawing the figures; the
interpretation of the exponent $1/2$ (Remark~\ref{rem:exponent-half});
the hypothesis on the trivial character in
Proposition~\ref{prop:averaged}; the size of the modulus in the idea
of proof of Proposition~\ref{prop:meansq}; the dependence on~$\Delta$
of the constant in the route sketched for
Proposition~\ref{prop:box-from-headsum}; and
part~\ref{item:packages-data} of Theorem~\ref{thm:packages}. The
readings that produced these points were by \texttt{claude-fable-5-2}
agents; the derivations were checked only as stated.
\end{enumerate}

\section{Literature review}
\label{sec:ai-literature}
For the book, the revision agent located references with
\vtwoskill{track-down-reference}: the locators in Vakil's book
\cite{Vakil2025} were checked in the public PDF of its version of
October 2025; the pages of \cite{Chen2010} against the publisher's
page; the publisher of \cite{Waldspurger1990} from citing papers, not
from the volume; and \cite{Patterson1977,KazhdanPatterson1984,%
ChintaFriedbergHoffstein2012} only through a secondary source, without
opening them. No search for prior occurrences was run for the results
(i)--(vi) of Section~\ref{sec:ai-proofreading}; their novelty is not
claimed, and whether they are known is unresolved.

For the sources: the Stepanov--Bombieri report read Bombieri's
Bourbaki expos\'e \cite{Bombieri1974} (pp.~234--240) and Tao's blog
post \cite{Tao2014} directly, through the web tools of Claude Code
\cite[\href{https://code.claude.com/docs/en/tools-reference}%
{tools reference}]{ClaudeCodeTools}, \texttt{curl}
\cite[\href{https://curl.se/}{home page}]{curl} and a Semantic
Scholar search connector, which serves the Semantic Scholar API
\cite[\href{https://www.semanticscholar.org/product/api}%
{API page}]{SemanticScholarAPI} and refused queries for rate
limiting; other references were taken from publisher or index pages
or cited from memory. No source was found for
Theorems~\ref{thm:route2-input} and~\ref{thm:route2prime} or
Propositions~\ref{prop:explicit-G} and~\ref{prop:which-bound} in this
form, and whether Proposition~\ref{prop:all-q} is known is unresolved.
The sign-only account cited \cite{Katz2015,Poonen2004,Deligne1974}
from memory and earlier records, without rereading them. For the
survey, two \texttt{claude-opus-5-5} agents checked the references
with \vtwoskill{track-down-reference} against zbMATH Open
\cite[\href{https://zbmath.org/}{home page}]{zbMATHOpen}, Crossref
\cite[\href{https://www.crossref.org/documentation/retrieve-metadata/%
rest-api/}{REST API}]{Crossref}, Numdam
\cite[\href{https://www.numdam.org/}{home page}]{Numdam} and
publishers, and ran bounded searches for its distinctive steps through
the arXiv API \cite[\href{https://info.arxiv.org/help/api/index.html}%
{API page}]{arXivAPI} in the manner of \vtwoskill{search-arxiv}, web
search, and the Semantic Scholar API in the manner of
\vtwoskill{search-scholar}. A search without hits is not a claim of
novelty. For the windowed Landau criterion no prior occurrence was
found; the closest work \cite[Thm.~3.1]{BrouckeRevesz2026landau}
starts from a pointwise bound (Remark~\ref{rem:landau-compare},
Question~\ref{q:landau-new}). For a Chebyshev or slope inequality
through theta lifts, none was found. For a uniform Vandiver-type bound
on Sha groups of Tate twists none was found; the related
\cite[\S5, p.~232]{Kurihara1992} works over $\Q$, one prime at a time
(Remark~\ref{rem:vandiver}). For the balance identity, the
orthogonality it uses is standard
\cite[Cor.~4.2, Ex.~4.3]{Conrad2026characters}, and no prior occurrence
of the identity was found. The comparisons of the plus part, the minus
part and the Galois construction with published work were made from
the reading agents' memory, except the check of
Remark~\ref{rem:galois-inputs} against
\cite[Thm.~2.1.1]{BarnetLambGeeGeraghtyTaylor2014} in full text and
\cite{Caraiani2012,Caraiani2014} in abstract; one point is open, the
twist by a power of $|\det|$ in the normalisation of the local
Langlands correspondence. Six entries added in the rewrite
\cite{Kawamata1982,Viehweg1982,Nadel1990,Messing1972,%
AnschutzLeBras2023,Walfisz1936} were checked as bibliographic records,
not for what they say. No human has checked the references of the
book.

\section{Formalization}
\label{sec:ai-formalization}
This book formalizes nothing, and no Lean skill was run for it. Its
object is Alp\"oge's Lean archive \cite{AlpogeArchive2026}, on Lean~4
v4.33.1
\cite[\href{https://github.com/leanprover/lean4/releases/tag/v4.33.1}%
{release v4.33.1}]{Lean4}
\cite[\href{https://doi.org/10.1007/978-3-030-79876-5_37}{DOI}]%
{Moura2021} and Mathlib revision \texttt{0df444a3}
\cite[\href{\vtwomathlib}{commit}]{Mathlib}. It was read as text and
never built. The formal names of the statements are kept only in
source comments next to each statement. The readings covered
definitions and statements, and proofs only at the top two to four
levels of each branch, eight to twelve levels for the construction of
the Galois representations; about $175{,}000$ of its $175{,}187$ files
were not opened. A build of the main file, its axiom report and an
independent replay of the kernel check
\cite[\href{https://github.com/leanprover/lean4checker}%
{repository}]{lean4checker} would decide whether the final statements
are proved relative to the kernel; static reading cannot. Elaboration
runs the archive's meta code, so such a check belongs in a disposable
sandbox.

The readings established facts about the archive. Each fact that
bears on a statement of the book is stated next to that statement.
The final statements of the archive are compared with
Theorem~\ref{thm:main} in Remark~\ref{rem:ai-final}. The scan for
assumptions follows Remark~\ref{rem:bb-unconditional}, as
Remark~\ref{rem:ai-scan}. The facts on imports and on transcribed
theorems follow Remark~\ref{rem:bb-imported}, as
Remarks~\ref{rem:ai-imports} and~\ref{rem:ai-transcribed}. One fact
concerns the method of reading.

\begin{remark}[names]\label{rem:ai-names}
At every level, the names of the files and declarations of the archive
do not describe their content.
\end{remark}

Further facts found by the readings: the archive contains meta code
that adds declarations and defers type checking to the kernel, so that
``it elaborates'' is no evidence, and it shows signs of mechanical
assembly (duplicated declarations, $1727$ one-line wrappers); none of
this is evidence of an error. It claims in-archive proofs of
resolution of singularities in characteristic zero, the arithmetic
Hilbert--Samuel theorem, the Siegel--Weil formula, Howe duality and
about $190$ further named theorems. A Weil--Deligne parameter has
default values when none exists, whose branches were not checked to be
unreached. No counterexample to a statement of the archive was found,
and no ill-posed definition at the level of the final statements.

\section{Record keeping and token accounting}
\label{sec:ai-records}
The project keeps one provenance record and one token record. They
are working records and are not distributed with this book. The
provenance record, written only with \vtwoskill{update-ai-usage},
gives for each task the agents with their models and harness, the
skill chain, the inputs and their origin, the results with their
derivation, the searches, the reviews and the limits. The token
record, written only with \vtwoskill{update-tokens}, gives per agent
the token counts by class where available, the rates with their date,
and the cost, labelled as measured or estimated. ``Not recorded''
means not measured, not zero.

Table~\ref{tab:ai-tokens} gives the token use of the tasks that wrote
the four sources and the book. Measured rows sum the per-message
counts of the agents' transcripts by token class, cache reads
included, and price each class at its rate; most of them are lower
bounds, because a coordinating session or the last turns of a
record-keeping agent were not counted. The estimated rows are carried
from the survey's own table: per-agent totals without a split by
class, as the harness reported them; they probably exclude cache
reads. Their costs price \texttt{claude-opus-5-5} totals at the input
rate and \texttt{claude-fable-5-2} totals with an assumed split of
$25\%$ output and $75\%$ cache write, with bounds at a low and a high
rate. Costs are token charges in US dollars at the rates per million
tokens recorded on 2026-11-20 and checked again on 2026-11-21 on the
provider's pricing page, which shows no effective date: for
\texttt{claude-opus-5-5} input $4.00$, output $20.00$, cache write
$5.00$ (5 minutes) and $8.00$ (1 hour), cache read $0.20$; for
\texttt{claude-fable-5-2} input $8.00$, output $40.00$, cache write
$10.00$ and $16.00$, cache read $0.20$. Not recorded: the import of
the archive, the introduction and Theorem~\ref{thm:main} drafted in
the Codex app, the rewrite's coordinator after a point and the last
turns of its body revision, the coordinator's last turns in the book's
run, and the drafting and audit of this appendix.

\begin{table}[htbp]
\centering\footnotesize\setlength{\tabcolsep}{4pt}
\begin{tabular}{@{}>{\raggedright\arraybackslash}p{0.42\linewidth}%
lrrl@{}}
\toprule
Task & Model & Tokens & USD & Source\\
\midrule
\multicolumn{5}{@{}l}{\emph{Chapter~\ref{ch:idea}: account of the idea
(2026-11-20)}}\\
Draft & \texttt{opus} & 3{,}961{,}002 & 2.04 & measured\\
Review and revision & \texttt{opus} & 9{,}656{,}142 & 3.96 & measured\\
Figures & \texttt{opus} & 14{,}149{,}338 & 4.63 & measured\\
Figure insertion & \texttt{opus} & 11{,}032{,}067 & 2.87 & measured\\
\midrule
\multicolumn{5}{@{}l}{\emph{Chapter~\ref{ch:function-fields}:
Stepanov--Bombieri (2026-11-20) and sign-only (2026-11-21)}}\\
Research report & both & 7{,}283{,}238 & 18.50 & measured\\
Harvest and review & \texttt{opus} & 18{,}264{,}997 & 9.35
& measured\\
Expectations placed & \texttt{opus} & 16{,}291{,}335 & 5.70
& measured\\
Figures & \texttt{opus} & 50{,}892{,}392 & 24.92 & measured\\
Shortening & \texttt{opus} & 26{,}162{,}340 & 8.31 & measured\\
Figure insertion & \texttt{opus} & 15{,}492{,}962 & 4.54 & measured\\
Sign-only, iteration 1 & \texttt{opus} & 6{,}464{,}552 & 4.32
& measured\\
Sign-only, iterations 2--3 & \texttt{opus} & 3{,}945{,}668 & 2.21
& measured\\
\midrule
\multicolumn{5}{@{}l}{\emph{Chapters~\ref{ch:assembly}--%
\ref{ch:families}: the survey (2026-11-20 unless stated)}}\\
Import of the archive & \texttt{opus} & not recorded & -- & --\\
Note on Birch--Swinnerton-Dyer & both & 24{,}829{,}392 & 28.15
& measured\\
Note on the Deninger program & both & 66{,}519{,}676 & 55.97
& measured\\
Reading of the archive & \texttt{fable} & 338{,}616 & 5.93
& estimated\\
Further readings 1--5 & \texttt{fable} & 4{,}851{,}181 & 84.90
& estimated\\
Earlier report & \texttt{opus} & 7{,}407{,}733 & 29.63 & estimated\\
Figures of the earlier report & \texttt{opus} & 4.36--11.88\,M
& 17.42--47.53 & estimated\\
Rewrite into the survey (2026-11-21) & \texttt{opus}
& 222{,}670{,}060 & 84.01 & measured\\
\midrule
\multicolumn{5}{@{}l}{\emph{The book (2026-11-21)}}\\
Transcription of the remark & \texttt{opus} & 2{,}320{,}296 & 1.25
& measured\\
Plan & both & 12{,}720{,}908 & 10.58 & measured\\
Skeleton, chapters, review, revision, prose & both
& 104{,}511{,}334 & 70.73 & measured\\
\bottomrule
\end{tabular}
\caption{Token use of the tasks behind this book. Model names omit
the prefix \texttt{claude-} and the version; ``both'' means agents of
the two models, priced per model by token class. The measured rows sum
to $617{,}167{,}699$ tokens and $342.04$ US dollars; most are lower
bounds. The estimated rows of the survey sum to $12{,}597{,}530$ tokens
and $120.45$ US dollars (bounds $43.00$ to $355.75$), and its figure
pass is a range from $4{,}355{,}591$ to $11{,}881{,}382$ tokens,
because its forks re-read the coordinator's context. In the book's
run, \texttt{claude-fable-5-2} agents used $50{,}581{,}916$ tokens
($47.77$ US dollars) and \texttt{claude-opus-5-5} agents
$53{,}929{,}418$ ($22.97$ US dollars).}
\label{tab:ai-tokens}
\end{table}
\fussy
\clearpage

\backmatter
\bibliographystyle{amsalpha}
\bibliography{book/riemann-book}

\end{document}
